% Generated mechanically from frozen input canon/ko-kr/integration-20260908-r10-p03-register/inputs/p03/ko/sites-cohomology.tex.
% Do not edit this generated file; edit the bound locale source instead.

% OK, start here.
%

\setcounter{chapter}{20}
\chapter{사이트 위의 코호몰로지}



\phantomsection
\label{sites-cohomology-section-phantom}


\section{도입}
\label{sites-cohomology-section-introduction}

\noindent
이 문서에서는 층의 코호몰로지에 관한 몇 가지 주제를 전개한다.
사이트 위의 층에서 어떤 일이 일어나는지 다루지만,
흔히 위상적 경우의 논의, 구성, 증명을
이 경우에 그대로 되풀이할 뿐이다.
기본 참고문헌은 \cite{SGA4}, \cite{Godement}, \cite{Iversen}이다.




\section{층의 코호몰로지}
\label{sites-cohomology-section-cohomology-sheaves}

\noindent
$\mathcal{C}$를 사이트라 하자. Sites, Definition
\ref{sites-definition-site}를 보라.
$\mathcal{F}$를 $\mathcal{C}$ 위의 아벨 층이라 하자.
$\mathcal{C}$ 위 아벨 층의 범주에는 단사 대상이 충분히 많음을 알고 있다.
Injectives, Theorem \ref{injectives-theorem-sheaves-injectives}를 보라.
따라서 단사 분해
$\mathcal{F}[0] \to \mathcal{I}^\bullet$을 택할 수 있다.
사이트 $\mathcal{C}$의 임의의 대상 $U$에 대해
\begin{equation}
\label{sites-cohomology-equation-cohomology-object-site}
H^i(U, \mathcal{F}) = H^i(\Gamma(U, \mathcal{I}^\bullet))
\end{equation}
로 정의하고 이를 {\it 대상 $U$ 위에서의 아벨 층
$\mathcal{F}$의 제$i$ 코호몰로지군}이라 한다. 다시 말해 이들은 함자
$\mathcal{F} \mapsto \mathcal{F}(U)$의 오른쪽 유도함자들이다.
함자족 $H^i(U, -)$는
$\textit{Ab}(\mathcal{C}) \to \textit{Ab}$인 보편 $\delta$-함자를 이룬다.

\medskip\noindent
사이트 $\mathcal{C}$에 끝 대상이 없는 경우가 있다. 이때
$\mathcal{C}$ 위 집합의 준층 $\mathcal{F}$의 {\it 대역 절단}을 집합
\begin{equation}
\label{sites-cohomology-equation-global-sections}
\Gamma(\mathcal{C}, \mathcal{F}) =
\Mor_{\textit{PSh}(\mathcal{C})}(e, \mathcal{F})
\end{equation}
으로 정의한다. 여기서 $e$는 $\mathcal{C}$ 위 준층의 범주에서 끝 대상이다.
이 경우 $\mathcal{C}$ 위의 아벨 층 $\mathcal{F}$가 주어지면
{\it $\mathcal{C}$ 위에서의 $\mathcal{F}$의 제$i$ 코호몰로지군}을
다음과 같이 정의한다.
\begin{equation}
\label{sites-cohomology-equation-cohomology}
H^i(\mathcal{C}, \mathcal{F}) = H^i(\Gamma(\mathcal{C}, \mathcal{I}^\bullet))
\end{equation}
다시 말해 이는 대역 절단 함자의 제$i$ 오른쪽 유도함자이다.
함자족 $H^i(\mathcal{C}, -)$는
$\textit{Ab}(\mathcal{C}) \to \textit{Ab}$인 보편 $\delta$-함자를 이룬다.

\medskip\noindent
$f : \Sh(\mathcal{C}) \to \Sh(\mathcal{D})$를 토포스의 사상이라 하자.
Sites, Definition \ref{sites-definition-topos}를 보라.
위와 같은 $\mathcal{F}[0] \to \mathcal{I}^\bullet$을 사용하여
\begin{equation}
\label{sites-cohomology-equation-higher-direct-image}
R^if_*\mathcal{F} = H^i(f_*\mathcal{I}^\bullet)
\end{equation}
로 정의하고 이를 {\it $\mathcal{F}$의 제$i$ 고차 직상}이라 한다.
이들은 $f_*$의 오른쪽 유도함자들이다.
함자족 $R^if_*$는
$\textit{Ab}(\mathcal{C}) \to \textit{Ab}(\mathcal{D})$인 보편
$\delta$-함자를 이룬다.

\medskip\noindent
$(\mathcal{C}, \mathcal{O})$를 환 달린 사이트라 하자. Modules on Sites,
Definition \ref{sites-modules-definition-ringed-site}를 보라.
$\mathcal{F}$를 $\mathcal{O}$-가군이라 하자.
$\mathcal{O}$-가군의 범주에는 단사 대상이 충분히 많음을 알고 있다.
Injectives, Theorem
\ref{injectives-theorem-sheaves-modules-injectives}를 보라.
따라서 단사 분해 $\mathcal{F}[0] \to \mathcal{I}^\bullet$을 택할 수 있다.
사이트 $\mathcal{C}$의 임의의 대상 $U$에 대해
\begin{equation}
\label{sites-cohomology-equation-cohomology-object-site-modules}
H^i(U, \mathcal{F}) = H^i(\Gamma(U, \mathcal{I}^\bullet))
\end{equation}
로 정의하고 이를 {\it $U$ 위에서의 $\mathcal{F}$의 제$i$ 코호몰로지군}이라
한다. 함자족 $H^i(U, -)$는
$\textit{Mod}(\mathcal{O}) \to \text{Mod}_{\mathcal{O}(U)}$인 보편
$\delta$-함자를 이룬다. 마찬가지로
\begin{equation}
\label{sites-cohomology-equation-cohomology-modules}
H^i(\mathcal{C}, \mathcal{F}) = H^i(\Gamma(\mathcal{C}, \mathcal{I}^\bullet))
\end{equation}
로 정의한 것은 {\it $\mathcal{C}$ 위에서의 $\mathcal{F}$의 제$i$
코호몰로지군}이다. 함자족 $H^i(\mathcal{C}, -)$는
$\textit{Mod}(\mathcal{C}) \to \text{Mod}_{\Gamma(\mathcal{C}, \mathcal{O})}$인
보편 $\delta$-함자를 이룬다.

\medskip\noindent
$f : (\Sh(\mathcal{C}), \mathcal{O}) \to (\Sh(\mathcal{D}), \mathcal{O}')$를
환 달린 토포스의 사상이라 하자. Modules on Sites, Definition
\ref{sites-modules-definition-ringed-topos}를 보라.
위와 같은 $\mathcal{F}[0] \to \mathcal{I}^\bullet$을 사용하여
\begin{equation}
\label{sites-cohomology-equation-higher-direct-image-modules}
R^if_*\mathcal{F} = H^i(f_*\mathcal{I}^\bullet)
\end{equation}
로 정의하고 이를 {\it $\mathcal{F}$의 제$i$ 고차 직상}이라 한다.
이들은 $f_*$의 오른쪽 유도함자들이다.
함자족 $R^if_*$는
$\textit{Mod}(\mathcal{O}) \to \textit{Mod}(\mathcal{O}')$인 보편
$\delta$-함자를 이룬다.








\section{유도함자}
\label{sites-cohomology-section-derived-functors}

\noindent
분해 함자를 이용하여 오른쪽 유도함자를 얻는 접근법을 간단히 설명한다.
즉, $(\mathcal{C}, \mathcal{O})$를 환 달린 사이트라 하자. 이 장에서는
$$
K(\mathcal{O}) = K(\textit{Mod}(\mathcal{O}))
\quad
\text{and}
\quad
D(\mathcal{O}) = D(\textit{Mod}(\mathcal{O}))
$$
로 쓰며, Derived Categories, Definition
\ref{derived-definition-complexes-notation}과 Definition
\ref{derived-definition-unbounded-derived-category}에서 도입한
삼각범주의 유계 형태에도 같은 표기를 사용한다. Derived Categories,
Remark \ref{derived-remark-big-abelian-category}에 의해 분해 함자
$$
j = j_{(\mathcal{C}, \mathcal{O})} :
K^{+}(\textit{Mod}(\mathcal{O}))
\longrightarrow
K^{+}(\mathcal{I})
$$
가 존재한다. 여기서 $\mathcal{I}$는 단사 $\mathcal{O}$-가군들로 이루어진
$\textit{Mod}(\mathcal{O})$의 엄밀 충만 가법 부분범주이다. 임의의 아벨
범주 $\mathcal{B}$로 가는 왼쪽 완전 함자
$F : \textit{Mod}(\mathcal{O}) \to \mathcal{B}$에 대해, Derived
Categories, Section \ref{derived-section-right-derived-functor}에
설명되어 있고 방금의 분해 함자 $j$를 사용하여 구성되는 오른쪽 유도함자를
$RF$로 나타낸다.
\begin{equation}
\label{sites-cohomology-equation-RF}
RF = F \circ j' : D^{+}(\mathcal{O}) \longrightarrow D^{+}(\mathcal{B})
\end{equation}
표기는 Derived Categories, Lemma
\ref{derived-lemma-right-derived-functor}를 보라. 상황에 따라 $RF$는
$\textit{Mod}(\mathcal{O})$, $\text{Comp}^{+}(\textit{Mod}(\mathcal{O}))$
또는 $K^{+}(\mathcal{O})$ 위에 정의된 것으로 생각할 수 있다. Derived
Categories, Definition \ref{derived-definition-higher-derived-functors}에
따라 제$i$ 오른쪽 유도함자
\begin{equation}
\label{sites-cohomology-equation-RFi}
R^iF = H^i \circ RF : \textit{Mod}(\mathcal{O}) \longrightarrow \mathcal{B}
\end{equation}
를 얻는다. 따라서 $R^0F = F$이고
$\{R^iF, \delta\}_{i \geq 0}$는 보편 $\delta$-함자이다. Derived
Categories, Lemma \ref{derived-lemma-higher-derived-functors}를 보라.

\medskip\noindent
이 구성의 두 가지 특수한 경우를 들자. 환 $R$이 주어지면
$K(R) = K(\text{Mod}_R)$과 $D(R) = D(\text{Mod}_R)$로 쓰고 유계
형태에도 같은 표기를 사용한다. $\mathcal{C}$의 임의의 대상 $U$에 대해
왼쪽 완전 함자
$
\Gamma(U, -) :
\textit{Mod}(\mathcal{O})
\longrightarrow
\text{Mod}_{\mathcal{O}(U)}
$
가 있고, 위의 논의로부터
$$
R\Gamma(U, -) :
D^{+}(\mathcal{O})
\longrightarrow
D^{+}(\mathcal{O}(U))
$$
을 얻는다. $H^i(U, -) = R^i\Gamma(U, -)$는 위의
(\ref{sites-cohomology-equation-cohomology-object-site-modules})와 양립한다. 마찬가지로
$$
R\Gamma(\mathcal{C}, -) :
D^{+}(\mathcal{O})
\longrightarrow
D^{+}(\Gamma(\mathcal{C}, \mathcal{O}))
$$
를 얻으며, 이는 (\ref{sites-cohomology-equation-cohomology-modules})와 양립한다.
$f : (\Sh(\mathcal{C}), \mathcal{O}) \to (\Sh(\mathcal{D}), \mathcal{O}')$
가 환 달린 토포스의 사상이면 왼쪽 완전 함자
$f_* : \textit{Mod}(\mathcal{O}) \to \textit{Mod}(\mathcal{O}')$
가 있고, 이로부터 {\it 유도 직상}
$$
Rf_* : D^{+}(\mathcal{O}) \to D^+(\mathcal{O}')
$$
을 얻는다. $Rf_*\mathcal{F}^\bullet$의 제$i$ 코호몰로지 층은
$R^if_*\mathcal{F}^\bullet$로 나타내며,
(\ref{sites-cohomology-equation-higher-direct-image-modules})에 따라 제$i$
{\it 고차 직상}이라 한다. 위에 표시한 함자들은 유도범주 사이의 완전함자이다.







\section{제1 코호몰로지와 토서}
\label{sites-cohomology-section-h1-torsors}

\begin{definition}
\label{sites-cohomology-definition-torsor}
$\mathcal{C}$를 사이트라 하고, $\mathcal{G}$를 $\mathcal{C}$ 위의
(가환하지 않을 수도 있는) 군의 층이라 하자. {\it 유사 토서}, 더 정확히는
{\it 유사 $\mathcal{G}$-토서}란 작용
$\mathcal{G} \times \mathcal{F} \to \mathcal{F}$가 주어진
$\mathcal{C}$ 위의 집합의 층 $\mathcal{F}$로서 다음을 만족하는 것을 말한다.
\begin{enumerate}
\item $\mathcal{F}(U)$가 공집합이 아닐 때마다 작용
$\mathcal{G}(U) \times \mathcal{F}(U) \to \mathcal{F}(U)$
은 단순 추이적이다.
\end{enumerate}
{\it 유사 $\mathcal{G}$-토서의 사상} $\mathcal{F} \to \mathcal{F}'$은
$\mathcal{G}$-작용과 양립하는 집합의 층의 사상이다. {\it 토서}, 더 정확히는
{\it $\mathcal{G}$-토서}란 다음 조건도 만족하는 유사 $\mathcal{G}$-토서이다.
\begin{enumerate}
\item[(2)] 모든 $U \in \Ob(\mathcal{C})$에 대해 $U$의 피복
$\{U_i \to U\}_{i \in I}$가 존재하여 모든 $i \in I$에 대해
$\mathcal{F}(U_i)$가 공집합이 아니다.
\end{enumerate}
{\it $\mathcal{G}$-토서의 사상}은 유사 $\mathcal{G}$-토서의 사상이다.
{\it 자명한 $\mathcal{G}$-토서}는 자명한 왼쪽 $\mathcal{G}$-작용이 주어진
층 $\mathcal{G}$이다.
\end{definition}

\noindent
토서의 사상은 자동으로 동형사상임이 분명하다.

\begin{lemma}
\label{sites-cohomology-lemma-trivial-torsor}
$\mathcal{C}$를 사이트라 하고, $\mathcal{G}$를 $\mathcal{C}$ 위의
(가환하지 않을 수도 있는) 군의 층이라 하자. $\mathcal{G}$-토서
$\mathcal{F}$가 자명할 필요충분조건은
$\Gamma(\mathcal{C}, \mathcal{F}) \not = \emptyset$인 것이다.
\end{lemma}

\begin{proof}
생략한다.
\end{proof}

\begin{lemma}
\label{sites-cohomology-lemma-torsors-h1}
$\mathcal{C}$를 사이트라 하고, $\mathcal{H}$를 $\mathcal{C}$ 위의
아벨 층이라 하자. $\mathcal{H}$-토서의 동형류 집합과
$H^1(\mathcal{C}, \mathcal{H})$ 사이에는 표준 전단사가 있다.
\end{lemma}

\begin{proof}
$\mathcal{F}$를 $\mathcal{H}$-토서라 하자. $\mathcal{F}$ 위의 자유 아벨 층
$\mathbf{Z}[\mathcal{F}]$을 생각하자. 이것은
$U \in \Ob(\mathcal{C})$에 $n_i \in \mathbf{Z}$이고
$s_i \in \mathcal{F}(U)$인 유한 형식합 $\sum n_i[s_i]$들의 모임을
대응시키는 규칙의 층화이다. 자연스러운 사상
$$
\sigma : \mathbf{Z}[\mathcal{F}] \longrightarrow \underline{\mathbf{Z}}
$$
이 있으며, 국소 단면 $\sum n_i[s_i]$를 $\sum n_i$로 보낸다.
$\sigma$의 핵은 $[s] - [s']$ 꼴의 단면들로 생성된다. 표준 사상
$a : \Ker(\sigma) \to \mathcal{H}$가 있으며,
$[s] - [s'] \mapsto h$로 보내는데, 여기서 $h$는
$h \cdot s' = s$를 만족하는 $\mathcal{H}$의 국소 단면이다. 밀어내기 도식
$$
\xymatrix{
0 \ar[r] &
\Ker(\sigma) \ar[r] \ar[d]^a &
\mathbf{Z}[\mathcal{F}] \ar[r] \ar[d] &
\underline{\mathbf{Z}} \ar[r] \ar[d] &
0 \\
0 \ar[r] &
\mathcal{H} \ar[r] &
\mathcal{E} \ar[r] &
\underline{\mathbf{Z}} \ar[r] &
0
}
$$
을 생각하자. 여기서 $\mathcal{E}$는 밀어내기로 얻은 확대이다. 아래의
짧은 완전열에 딸린 코호몰로지 긴 완전열에서
$\xi = \xi_\mathcal{F} \in H^1(\mathcal{C}, \mathcal{H})$를 얻는데,
이는 경계 사상을
$1 \in H^0(\mathcal{C}, \underline{\mathbf{Z}})$에 적용하여 얻는다.

\medskip\noindent
거꾸로 $\xi \in H^1(\mathcal{C}, \mathcal{H})$가 주어지면 다음과 같이
$\xi$에 토서를 대응시킬 수 있다. $\mathcal{H}$를 단사 아벨 층
$\mathcal{I}$에 넣는 매장 $\mathcal{H} \to \mathcal{I}$를 택하자.
$\mathcal{Q} = \mathcal{I}/\mathcal{H}$로 두면 다음 짧은 완전열을 얻는다.
$$
\xymatrix{
0 \ar[r] &
\mathcal{H} \ar[r] &
\mathcal{I} \ar[r] &
\mathcal{Q} \ar[r] &
0
}
$$
$H^1(\mathcal{C}, \mathcal{I}) = 0$이므로(Derived Categories, Lemma
\ref{derived-lemma-higher-derived-functors} 참조), 원소 $\xi$는 전역 단면
$q \in H^0(\mathcal{C}, \mathcal{Q})$의 상이다. 층 $\mathcal{Q}$에서
$q$로 가는 단면들로 이루어진 부분층(집합의 층으로서)을
$\mathcal{F} \subset \mathcal{I}$라 하자. $\mathcal{F}$가
$\mathcal{H}$-토서임은 쉽게 확인할 수 있다.

\medskip\noindent
위에서 준 두 구성이 서로 역이라는 확인은 생략한다.
\end{proof}







\section{제1 코호몰로지와 확대}
\label{sites-cohomology-section-h1-extensions}

\begin{lemma}
\label{sites-cohomology-lemma-h1-extensions}
$(\mathcal{C}, \mathcal{O})$를 환 달린 사이트라 하고, $\mathcal{F}$를
$\mathcal{C}$ 위의 $\mathcal{O}$-가군층이라 하자. 표준 전단사
$$
\Ext^1_{\textit{Mod}(\mathcal{O})}(\mathcal{O}, \mathcal{F})
\longrightarrow
H^1(\mathcal{C}, \mathcal{F})
$$
가 있으며, 이 전단사는 확대
$$
0 \to \mathcal{F} \to \mathcal{E} \to \mathcal{O} \to 0
$$
에 $1 \in \Gamma(\mathcal{C}, \mathcal{O})$의
$H^1(\mathcal{C}, \mathcal{F})$에서의 상을 대응시킨다.
\end{lemma}

\begin{proof}
보조정리에서 준 사상의 역을 구성하자.
$\xi \in H^1(\mathcal{C}, \mathcal{F})$라 하자.
$\textit{Mod}(\mathcal{O})$에서 $\mathcal{I}$가 단사 대상이 되도록
단사 사상 $\mathcal{F} \subset \mathcal{I}$를 택한다.
$\mathcal{Q} = \mathcal{I}/\mathcal{F}$로 두자. 코호몰로지 긴 완전열에 의해
$\xi$는 단면
$\tilde \xi \in \Gamma(\mathcal{C}, \mathcal{Q}) =
\Hom_\mathcal{O}(\mathcal{O}, \mathcal{Q})$
의 상이다. 이제 다음 당김을 만들면 된다.
$$
\xymatrix{
0 \ar[r] &
\mathcal{F} \ar[r] \ar@{=}[d] &
\mathcal{E} \ar[r] \ar[d] &
\mathcal{O} \ar[r] \ar[d]^{\tilde \xi} &
0 \\
0 \ar[r] &
\mathcal{F} \ar[r] &
\mathcal{I} \ar[r] &
\mathcal{Q} \ar[r] &
0
}
$$
Homology, Section \ref{homology-section-extensions}을 보라.
\end{proof}

\noindent
다음 보조정리는 더 일반적인 보조정리
\ref{sites-cohomology-lemma-cohomology-modules-abelian-agree}로 대체될 것이다.

\begin{lemma}
\label{sites-cohomology-lemma-h1-mod-ab-agree}
$(\mathcal{C}, \mathcal{O})$를 환 달린 사이트라 하고, $\mathcal{F}$를
$\mathcal{C}$ 위의 $\mathcal{O}$-가군층이라 하자. $\mathcal{F}_{ab}$로
그 밑의 아벨군층을 나타내자. 그러면 함자적 동형사상
$$
H^1(\mathcal{C}, \mathcal{F}_{ab})
=
H^1(\mathcal{C}, \mathcal{F})
$$
이 있다. 여기서 왼쪽은 $\textit{Ab}(\mathcal{C})$에서 계산한
코호몰로지이고, 오른쪽은 $\textit{Mod}(\mathcal{O})$에서 계산한
코호몰로지이다.
\end{lemma}

\begin{proof}
$\underline{\mathbf{Z}}$로 상수층 $\mathbf{Z}$를 나타내자.
$\textit{Ab}(\mathcal{C}) = \textit{Mod}(\underline{\mathbf{Z}})$이므로
보조정리 \ref{sites-cohomology-lemma-h1-extensions}를 두 번 적용할 수 있다. 따라서
다음을 보이면 된다.
$$
\Ext^1_{\textit{Mod}(\mathcal{O})}(\mathcal{O}, \mathcal{F})
=
\Ext^1_{\textit{Mod}(\underline{\mathbf{Z}})}(
\underline{\mathbf{Z}}, \mathcal{F}_{ab}).
$$
$0 \to \mathcal{F} \to \mathcal{E} \to \mathcal{O} \to 0$을
$\textit{Mod}(\mathcal{O})$의 확대라 하자. 그러면 아벨 층의 자명한 사상
$1 : \underline{\mathbf{Z}} \to \mathcal{O}$을 사용하여 당김을 취함으로써
다음과 같이 확대 $\mathcal{E}_{ab}$를 얻을 수 있다.
$$
\xymatrix{
0 \ar[r] &
\mathcal{F}_{ab} \ar[r] \ar@{=}[d] &
\mathcal{E}_{ab} \ar[r] \ar[d] &
\underline{\mathbf{Z}} \ar[r] \ar[d]^{1} &
0 \\
0 \ar[r] &
\mathcal{F} \ar[r] &
\mathcal{E} \ar[r] &
\mathcal{O} \ar[r] &
0
}
$$
거꾸로 하는 쪽이 조금 더 흥미롭다.
$0 \to \mathcal{F}_{ab} \to \mathcal{E}_{ab} \to \underline{\mathbf{Z}} \to 0$
을 $\textit{Mod}(\underline{\mathbf{Z}})$의 확대라 하자.
$\underline{\mathbf{Z}}$는 평탄한 $\underline{\mathbf{Z}}$-가군이므로 완전열
$$
0 \to \mathcal{F}_{ab} \otimes_{\underline{\mathbf{Z}}} \mathcal{O}
\to \mathcal{E}_{ab} \otimes_{\underline{\mathbf{Z}}} \mathcal{O}
\to \underline{\mathbf{Z}} \otimes_{\underline{\mathbf{Z}}} \mathcal{O}
\to 0
$$
을 얻는다. Modules on Sites, Lemma
\ref{sites-modules-lemma-flat-tor-zero}를 보라. 물론
$\underline{\mathbf{Z}} \otimes_{\underline{\mathbf{Z}}} \mathcal{O}
= \mathcal{O}$이다. 따라서 ($\mathcal{O}$-선형인) 곱셈 사상
$\mu : \mathcal{F} \otimes_{\underline{\mathbf{Z}}} \mathcal{O}
\to \mathcal{F}$를 따라 밀어내기를 취하여 다음과 같이
$\mathcal{F}$에 의한 $\mathcal{O}$의 확대를 얻을 수 있다.
$$
\xymatrix{
0 \ar[r] &
\mathcal{F}_{ab} \otimes_{\underline{\mathbf{Z}}} \mathcal{O}
\ar[r] \ar[d]^\mu &
\mathcal{E}_{ab} \otimes_{\underline{\mathbf{Z}}} \mathcal{O}
\ar[r] \ar[d] &
\mathcal{O} \ar[r] \ar@{=}[d] &
0 \\
0 \ar[r] &
\mathcal{F} \ar[r] &
\mathcal{E} \ar[r] &
\mathcal{O} \ar[r] &
0
}
$$
이것이 원하는 확대이다. 이 두 구성이 서로 역이라는 확인은 생략한다.
\end{proof}





\section{제1 코호몰로지와 가역층}
\label{sites-cohomology-section-invertible-sheaves}

\noindent
환 달린 사이트의 피카르 군은 Modules on Sites, Section
\ref{sites-modules-section-invertible}에서 정의하였다.

\begin{lemma}
\label{sites-cohomology-lemma-h1-invertible}
$(\mathcal{C}, \mathcal{O})$를 국소 환 달린 사이트라 하자.
아벨군의 표준 동형사상
$$
H^1(\mathcal{C}, \mathcal{O}^*) = \Pic(\mathcal{O}).
$$
이 있다.
\end{lemma}

\begin{proof}
$\mathcal{L}$을 가역 $\mathcal{O}$-가군이라 하자. 다음 규칙으로 정의된
준층 $\mathcal{L}^*$를 생각하자.
$$
U \longmapsto \{s \in \mathcal{L}(U)
\text{ 중 } \mathcal{O}_U \xrightarrow{s \cdot -} \mathcal{L}_U
\text{ 가 동형사상인 것}\}
$$
이 준층은 층 조건을 만족한다. 또한
$f \in \mathcal{O}^*(U)$이고 $s \in \mathcal{L}^*(U)$이면 분명히
$fs \in \mathcal{L}^*(U)$이다. 마찬가지로
$s, s' \in \mathcal{L}^*(U)$이면 $fs = s'$를 만족하는
$f \in \mathcal{O}^*(U)$가 유일하게 존재한다. 더욱이 Modules on Sites,
Lemma \ref{sites-modules-lemma-invertible-is-locally-free-rank-1}에 의해
층 $\mathcal{L}^*$는 국소적으로 단면을 가진다. 다시 말해
$\mathcal{L}^*$는 $\mathcal{O}^*$-토서이다. 따라서 다음 사상을 얻는다.
$$
\begin{matrix}
\text{환 달린 사이트 }(\mathcal{C}, \mathcal{O})\text{ 위의 가역층} \\
\text{의 동형류}
\end{matrix}
\longrightarrow
\begin{matrix}
\mathcal{O}^*\text{-토서} \\
\text{의 동형류}
\end{matrix}
$$
이 사상이 아벨군의 준동형이라는 확인은 생략한다. 보조정리
\ref{sites-cohomology-lemma-torsors-h1}에 의해 오른쪽은
$H^1(\mathcal{C}, \mathcal{O}^*)$와 표준적으로 전단사이다.
따라서 이 사상이 단사이고 전사임을 보이면 된다.

\medskip\noindent
단사성. 토서 $\mathcal{L}^*$가 자명하면 보조정리
\ref{sites-cohomology-lemma-trivial-torsor}에 의해 $\mathcal{L}^*$는 전역 단면을 가진다.
따라서 이는 정확히 $\mathcal{L} \cong \mathcal{O}$가
$\Pic(\mathcal{O})$의 항등원이라는 뜻이다.

\medskip\noindent
전사성. $\mathcal{F}$를 $\mathcal{O}^*$-토서라 하자. 집합의 준층
$$
\mathcal{L}_1 : U \longmapsto
(\mathcal{F}(U) \times \mathcal{O}(U))/\mathcal{O}^*(U)
$$
을 생각하자. 여기서 $f \in \mathcal{O}^*(U)$가
$(s, g)$에 작용한 결과는 $(fs, f^{-1}g)$이다. $s'/s$를
$fs = s'$를 만족하는 $\mathcal{O}^*$의 국소 단면 $f$라 하고,
$(s, g) + (s', g') = (s, g + (s'/s)g')$로 두며, $\mathcal{O}$의
국소 단면 $h$에 대해 $h(s, g) = (s, hg)$로 두면 $\mathcal{L}_1$은
$\mathcal{O}$-가군의 준층이 된다. 그 층화
$\mathcal{L} = \mathcal{L}_1^\#$가 가역 $\mathcal{O}$-가군이고,
이에 딸린 $\mathcal{O}^*$-토서 $\mathcal{L}^*$가
$\mathcal{F}$와 동형이라는 확인은 생략한다.
\end{proof}










\section{코호몰로지의 국소성}
\label{sites-cohomology-section-locality}

\noindent
다음 보조정리는 사이트의 대상 위에서 층 $\mathcal{F}$의 코호몰로지를
정의하는 데 모호함이 없음을 말한다.

\begin{lemma}
\label{sites-cohomology-lemma-cohomology-of-open}
$(\mathcal{C}, \mathcal{O})$를 환 달린 사이트라 하고, $U$를
$\mathcal{C}$의 대상이라 하자.
\begin{enumerate}
\item $\mathcal{I}$가 단사 $\mathcal{O}$-가군이면
$\mathcal{I}|_U$는 단사 $\mathcal{O}_U$-가군이다.
\item 임의의 $\mathcal{O}$-가군층 $\mathcal{F}$에 대해
$H^p(U, \mathcal{F}) = H^p(\mathcal{C}/U, \mathcal{F}|_U)$이다.
\end{enumerate}
\end{lemma}

\begin{proof}
$U$로 제한하는 함자 $j_U^{-1}$는 $0$에 의한 연장 함자 $j_{U!}$의
오른쪽 수반임을 상기하라. Modules on Sites, Section
\ref{sites-modules-section-localize}를 보라. 또한 $j_{U!}$는 완전하다.
따라서 (1)은 Homology, Lemma
\ref{homology-lemma-adjoint-preserve-injectives}에서 따른다.

\medskip\noindent
정의에 따라 $H^p(U, \mathcal{F}) = H^p(\mathcal{I}^\bullet(U))$이다.
여기서 $\mathcal{F} \to \mathcal{I}^\bullet$은
$\textit{Mod}(\mathcal{O})$에서의 단사 분해이다. 위에서 본 바에 따라
$\mathcal{F}|_U \to \mathcal{I}^\bullet|_U$는
$\textit{Mod}(\mathcal{O}_U)$에서의 단사 분해이다. 따라서
$H^p(U, \mathcal{F}|_U)$는
$H^p(\mathcal{I}^\bullet|_U(U))$와 같다. 물론 $\mathcal{C}$ 위의
임의의 층 $\mathcal{F}$에 대해
$\mathcal{F}(U) = \mathcal{F}|_U(U)$이다. 따라서 (2)의 등식이 따른다.
\end{proof}

\noindent
다음 보조정리는 부분 우주를 바꿀 때 어떤 일이 일어나는지를 알아보거나,
작은 에탈 사이트와 큰 에탈 사이트의 코호몰로지를 비교하는 데 쓰일 것이다.

\begin{lemma}
\label{sites-cohomology-lemma-cohomology-bigger-site}
$\mathcal{C}$와 $\mathcal{D}$를 사이트라 하고,
$u : \mathcal{C} \to \mathcal{D}$를 함자라 하자. $u$가 Sites, Lemma
\ref{sites-lemma-bigger-site}의 가정을 만족한다고 하자.
$g : \Sh(\mathcal{C}) \to \Sh(\mathcal{D})$를 이에 딸린
토포스의 사상이라 하자. $\mathcal{D}$ 위의 임의의 아벨 층
$\mathcal{F}$에 대해 동형사상
$$
R\Gamma(\mathcal{C}, g^{-1}\mathcal{F}) = R\Gamma(\mathcal{D}, \mathcal{F}),
$$
이 있고, 특히
$H^p(\mathcal{C}, g^{-1}\mathcal{F}) = H^p(\mathcal{D}, \mathcal{F})$이다.
또 임의의 $U \in \Ob(\mathcal{C})$에 대해 동형사상
$$
R\Gamma(U, g^{-1}\mathcal{F}) = R\Gamma(u(U), \mathcal{F}),
$$
이 있고, 특히
$H^p(U, g^{-1}\mathcal{F}) = H^p(u(U), \mathcal{F})$이다.
이 동형사상들은 모두 $\mathcal{F}$에 관해 함자적이다.
\end{lemma}

\begin{proof}
가정 (e)에 의해
$\Gamma(\mathcal{C}, g^{-1}\mathcal{F}) = \Gamma(\mathcal{D}, \mathcal{F})$
임은 분명하므로, $g^{-1}$가 단사 아벨 층을 단사 아벨 층으로 옮김을 보이면
충분하다. 통상대로 Homology, Lemma
\ref{homology-lemma-adjoint-preserve-injectives}를 사용한다.
Sites, Lemma \ref{sites-lemma-bigger-site}의 표기로 $g^{-1}$의 왼쪽 수반은
$g_! = f^{-1}$이고 이는 완전함자이다. 따라서 그 보조정리를 적용할 수 있다.
\end{proof}

\noindent
$(\mathcal{C}, \mathcal{O})$를 환 달린 사이트라 하고, $\mathcal{F}$를
$\mathcal{O}$-가군층이라 하자. $\varphi : U \to V$를
$\mathcal{O}$의 사상이라 하자. 그러면 표준적인 {\it 제한 사상}
\begin{equation}
\label{sites-cohomology-equation-restriction-mapping}
H^n(V, \mathcal{F})
\longrightarrow
H^n(U, \mathcal{F}), \quad
\xi \longmapsto \xi|_U
\end{equation}
이 있으며 이는 $\mathcal{F}$에 관해 함자적이다. 실제로 임의의 단사 분해
$\mathcal{F} \to \mathcal{I}^\bullet$을 택하자. 층 $\mathcal{I}^p$의
제한 사상들은 복합체의 사상
$$
\Gamma(V, \mathcal{I}^\bullet)
\longrightarrow
\Gamma(U, \mathcal{I}^\bullet)
$$
을 준다. 왼쪽은 $R\Gamma(V, \mathcal{F})$를 나타내는 복합체이고,
오른쪽은 $R\Gamma(U, \mathcal{F})$를 나타내는 복합체이다. 함자 $H^n$을
적용하여 코호몰로지군 사이의 사상을 얻는다. 앞서 표시했듯이 이 사상을
$\xi \mapsto \xi|_U$로 나타낸다. 따라서 규칙
$U \mapsto H^n(U, \mathcal{F})$는 $\mathcal{O}$-가군의 준층이다.
이 준층은 관례상 $\underline{H}^n(\mathcal{F})$로 나타낸다. 보조정리
\ref{sites-cohomology-lemma-include}에서 이 준층을 달리 해석할 것이다.

\medskip\noindent
다음 보조정리는 피복으로 넘어가면 고차 코호몰로지류를 소멸시킬 수 있음을
말한다.

\begin{lemma}
\label{sites-cohomology-lemma-kill-cohomology-class-on-covering}
$(\mathcal{C}, \mathcal{O})$를 환 달린 사이트라 하고, $\mathcal{F}$를
$\mathcal{O}$-가군층이라 하자. $U$를 $\mathcal{C}$의 대상이라 하자.
$n > 0$이고 $\xi \in H^n(U, \mathcal{F})$라 하자. 그러면 모든
$i \in I$에 대해 $\xi|_{U_i} = 0$이 되는 $\mathcal{C}$의 피복
$\{U_i \to U\}$가 존재한다.
\end{lemma}

\begin{proof}
$\mathcal{F} \to \mathcal{I}^\bullet$을 단사 분해라 하자. 그러면
$$
H^n(U, \mathcal{F}) =
\frac{\Ker(\mathcal{I}^n(U) \to \mathcal{I}^{n + 1}(U))}
{\Im(\mathcal{I}^{n - 1}(U) \to \mathcal{I}^n(U))}.
$$
위의 표시에 나오는 코호몰로지류를 나타내는 원소
$\tilde \xi \in \mathcal{I}^n(U)$를 택하자. $\mathcal{I}^\bullet$은
$\mathcal{F}$의 단사 분해이고 $n > 0$이므로, 복합체
$\mathcal{I}^\bullet$은 차수 $n$에서 완전하다. 따라서 층으로서
$\Im(\mathcal{I}^{n - 1} \to \mathcal{I}^n) =
\Ker(\mathcal{I}^n \to \mathcal{I}^{n + 1})$이다. $\tilde \xi$는
$U$ 위에서 핵층의 단면이므로, $\tilde \xi|_{U_i}$가 어떤 단면
$\xi_i \in \mathcal{I}^{n - 1}(U_i)$의 $d$에 의한 상이 되는 사이트의
피복 $\{U_i \to U\}$가 존재한다. 제한 $\xi|_{U_i}$를
$\tilde \xi|_{U_i}$의 류에 대응하는 것으로 정의했으므로 결론이 따른다.
\end{proof}

\begin{lemma}
\label{sites-cohomology-lemma-higher-direct-images}
$f : (\mathcal{C}, \mathcal{O}_\mathcal{C}) \to
(\mathcal{D}, \mathcal{O}_\mathcal{D})$를 연속함자
$u : \mathcal{D} \to \mathcal{C}$에 대응하는 환 달린 사이트의 사상이라
하자. 임의의
$\mathcal{F} \in \Ob(\textit{Mod}(\mathcal{O}_\mathcal{C}))$에 대해 층
$R^if_*\mathcal{F}$는 준층
$$
V \longmapsto H^i(u(V), \mathcal{F})
$$
에 결부된 층이다.
\end{lemma}

\begin{proof}
$\mathcal{F} \to \mathcal{I}^\bullet$을 단사 분해라 하자. 정의에 따라
$R^if_*\mathcal{F}$는 복합체
$$
f_*\mathcal{I}^0 \to f_*\mathcal{I}^1 \to f_*\mathcal{I}^2 \to \ldots
$$
의 제$i$ 코호몰로지층이다. $\mathcal{O}_\mathcal{D}$-가군의 아벨 범주
구조의 정의에 따라 이 코호몰로지층은 준층
$$
V
\longmapsto
\frac{\Ker(f_*\mathcal{I}^i(V) \to f_*\mathcal{I}^{i + 1}(V))}
{\Im(f_*\mathcal{I}^{i - 1}(V) \to f_*\mathcal{I}^i(V))}
$$
에 결부된 층이다. 이것은 분명히
$$
\frac{\Ker(\mathcal{I}^i(u(V)) \to \mathcal{I}^{i + 1}(u(V)))}
{\Im(\mathcal{I}^{i - 1}(u(V)) \to \mathcal{I}^i(u(V)))}
$$
와 같고, 이는 $H^i(u(V), \mathcal{F})$와 같다. 이로써 끝난다.
\end{proof}








\section{체흐 복합체와 체흐 코호몰로지}
\label{sites-cohomology-section-cech}

\noindent
$\mathcal{C}$를 범주라 하자.
$\mathcal{U} = \{U_i \to U\}_{i \in I}$를 공역이 고정된 사상족이라 하자.
Sites, Definition \ref{sites-definition-family-morphisms-fixed-target}를 보라.
모든 섬유곱 $U_{i_0} \times_U \ldots \times_U U_{i_p}$가
$\mathcal{C}$에 존재한다고 하자. $\mathcal{F}$를 $\mathcal{C}$ 위의
아벨 준층이라 하고 다음과 같이 놓는다.
$$
\check{\mathcal{C}}^p(\mathcal{U}, \mathcal{F})
=
\prod\nolimits_{(i_0, \ldots, i_p) \in I^{p + 1}}
\mathcal{F}(U_{i_0} \times_U \ldots \times_U U_{i_p}).
$$
이는 아벨군이다.
$s \in \check{\mathcal{C}}^p(\mathcal{U}, \mathcal{F})$에 대해 인자
$\mathcal{F}(U_{i_0} \times_U \ldots \times_U U_{i_p})$에서의 그 값을
$s_{i_0\ldots i_p}$로 나타낸다. 다음을 정의한다.
$$
d : \check{\mathcal{C}}^p(\mathcal{U}, \mathcal{F})
\longrightarrow
\check{\mathcal{C}}^{p + 1}(\mathcal{U}, \mathcal{F})
$$
그 공식은
\begin{equation}
\label{sites-cohomology-equation-d-cech}
d(s)_{i_0\ldots i_{p + 1}} =
\sum\nolimits_{j = 0}^{p + 1}
(-1)^j s_{i_0\ldots \hat i_j \ldots i_{p + 1}}
|_{U_{i_0} \times_U \ldots \times_U U_{i_{p + 1}}}
\end{equation}
이다. 여기서 제한은 사영사상
$$
U_{i_0} \times_U \ldots \times_U U_{i_{p + 1}} \longrightarrow
U_{i_0} \times_U \ldots \times_U \widehat{U_{i_j}} \times_U
\ldots \times_U U_{i_{p + 1}}.
$$
을 따른다. $d \circ d = 0$임은 곧바로 알 수 있다. 다시 말해
$\check{\mathcal{C}}^\bullet(\mathcal{U}, \mathcal{F})$는 복합체이다.

\begin{definition}
\label{sites-cohomology-definition-cech-complex}
$\mathcal{C}$를 범주라 하자.
$\mathcal{U} = \{U_i \to U\}_{i \in I}$를 공역이 고정된 사상족으로서
모든 섬유곱 $U_{i_0} \times_U \ldots \times_U U_{i_p}$가
$\mathcal{C}$에 존재하는 것이라 하자. $\mathcal{F}$를 $\mathcal{C}$
위의 아벨 준층이라 하자. 복합체
$\check{\mathcal{C}}^\bullet(\mathcal{U}, \mathcal{F})$를
$\mathcal{F}$와 사상족 $\mathcal{U}$에 결부된 {\it 체흐 복합체}라고 한다.
그 코호몰로지군
$H^i(\check{\mathcal{C}}^\bullet(\mathcal{U}, \mathcal{F}))$를
$\mathcal{U}$에 대한 $\mathcal{F}$의 {\it 체흐 코호몰로지군}이라고 하며,
$\check H^i(\mathcal{U}, \mathcal{F})$로 나타낸다.
\end{definition}

\noindent
사이트 $\mathcal{C}$의 임의의 피복 $\{U_i \to U\}$는 이 정의를 적용할 수
있는 공역이 고정된 사상족임을 관찰하자.

\begin{lemma}
\label{sites-cohomology-lemma-cech-h0}
$\mathcal{C}$를 사이트라 하고, $\mathcal{F}$를 $\mathcal{C}$ 위의
아벨 준층이라 하자. 다음은 서로 동치이다.
\begin{enumerate}
\item $\mathcal{F}$는 $\mathcal{C}$ 위의 아벨 층이다.
\item 사이트 $\mathcal{C}$의 모든 피복
$\mathcal{U} = \{U_i \to U\}_{i \in I}$에 대해 자연스러운 사상
$$
\mathcal{F}(U) \to \check{H}^0(\mathcal{U}, \mathcal{F})
$$
은 전단사이다(Sites, Section \ref{sites-section-sheafification} 참조).
\end{enumerate}
\end{lemma}

\begin{proof}
층 조건은 정확히 $\mathcal{C}$의 모든 피복에 대해
$\mathcal{F}(U) \to \check{H}^0(\mathcal{U}, \mathcal{F})$가
전단사라는 조건이므로 성립한다.
\end{proof}

\noindent
$\mathcal{C}$를 범주라 하자.
$\mathcal{U} = \{U_i \to U\}_{i\in I}$를 공역이 고정된
$\mathcal{C}$의 사상족으로서 모든 섬유곱
$U_{i_0} \times_U \ldots \times_U U_{i_p}$가 $\mathcal{C}$에
존재하는 것이라 하자.
$\mathcal{V} = \{V_j \to V\}_{j\in J}$를 또 하나의 그러한 사상족이라 하자.
$f : U \to V$, $\alpha : I \to J$ 및 $f_i : U_i \to V_{\alpha(i)}$를
공역이 고정된 사상족의 사상이라 하자. Sites, Section
\ref{sites-section-refinements}를 보라. 이 경우 체흐 복합체의 사상
\begin{equation}
\label{sites-cohomology-equation-map-cech-complexes}
\varphi : \check{\mathcal{C}}^\bullet(\mathcal{V}, \mathcal{F})
\longrightarrow
\check{\mathcal{C}}^\bullet(\mathcal{U}, \mathcal{F})
\end{equation}
을 얻으며, 차수 $p$에서 이는
$$
\varphi(s)_{i_0 \ldots i_p} =
(f_{i_0} \times \ldots \times f_{i_p})^*s_{\alpha(i_0) \ldots \alpha(i_p)}
$$
로 주어진다.




\section{준층 위의 함자로서의 체흐 코호몰로지}
\label{sites-cohomology-section-cech-functor}

\noindent
경고: 이 절에서는 오직 범주 위의 아벨 준층만을 다룬다. 여기의 결과들은
층과 층의 범주라는 설정에서는 완전히 틀린다!

\medskip\noindent
$\mathcal{C}$를 범주라 하자.
$\mathcal{U} = \{U_i \to U\}_{i \in I}$를 공역이 고정된 사상족으로서
모든 섬유곱 $U_{i_0} \times_U \ldots \times_U U_{i_p}$가
$\mathcal{C}$에 존재하는 것이라 하자. $\mathcal{F}$를 $\mathcal{C}$
위의 아벨 준층이라 하자. 구성
$$
\mathcal{F} \longmapsto \check{\mathcal{C}}^\bullet(\mathcal{U}, \mathcal{F})
$$
은 $\mathcal{F}$에 대해 함자적이다. 실제로 이는 함자
\begin{equation}
\label{sites-cohomology-equation-cech-functor}
\check{\mathcal{C}}^\bullet(\mathcal{U}, -) :
\textit{PAb}(\mathcal{C})
\longrightarrow
\text{Comp}^{+}(\textit{Ab})
\end{equation}
이다. 표기는 Derived Categories, Definition
\ref{derived-definition-complexes-notation}을 보라. 아벨 범주에서 아래로
유계인 복합체들의 범주는 아벨 범주임을 상기하라. Homology, Lemma
\ref{homology-lemma-cat-cochain-abelian}을 보라.

\begin{lemma}
\label{sites-cohomology-lemma-cech-exact-presheaves}
Equation (\ref{sites-cohomology-equation-cech-functor})으로 주어진 함자는 완전함자이다
(Homology, Lemma \ref{homology-lemma-exact-functor} 참조).
\end{lemma}

\begin{proof}
$\mathcal{C}$의 임의의 대상 $W$에 대해 함자
$\mathcal{F} \mapsto \mathcal{F}(W)$는
$\textit{PAb}(\mathcal{C})$에서 $\textit{Ab}$로 가는 가법 완전함자이다.
복합체의 항 $\check{\mathcal{C}}^p(\mathcal{U}, \mathcal{F})$는 이러한
완전함자들의 곱이므로 완전하다. 또한 복합체의 열이 완전일 필요충분조건은
각 차수에서 항들의 열이 완전인 것이다. 따라서 보조정리가 따른다.
\end{proof}

\begin{lemma}
\label{sites-cohomology-lemma-cech-cohomology-delta-functor-presheaves}
$\mathcal{C}$를 범주라 하자.
$\mathcal{U} = \{U_i \to U\}_{i \in I}$를 공역이 고정된 사상족으로서
모든 섬유곱 $U_{i_0} \times_U \ldots \times_U U_{i_p}$가
$\mathcal{C}$에 존재하는 것이라 하자. 함자들
$\mathcal{F} \mapsto \check{H}^n(\mathcal{U}, \mathcal{F})$는
아벨 범주 $\textit{PAb}(\mathcal{C})$에서 $\mathbf{Z}$-가군의 범주로
가는 $\delta$-함자를 이룬다(Homology, Definition
\ref{homology-definition-cohomological-delta-functor} 참조).
\end{lemma}

\begin{proof}
보조정리 \ref{sites-cohomology-lemma-cech-exact-presheaves}에 의해 아벨 준층들의 짧은 완전열
$0 \to \mathcal{F}_1 \to \mathcal{F}_2 \to \mathcal{F}_3 \to 0$은
$\mathbf{Z}$-가군 복합체들의 짧은 완전열로 보내진다. 따라서 Homology,
Lemma \ref{homology-lemma-long-exact-sequence-cochain}을 사용하여 경계 사상
$\delta_{\mathcal{F}_1 \to \mathcal{F}_2 \to \mathcal{F}_3} :
\check{H}^n(\mathcal{U}, \mathcal{F}_3) \to
\check{H}^{n + 1}(\mathcal{U}, \mathcal{F}_1)$
과 이에 대응하는 긴 완전열을 얻는다. 이 사상들이 준층들의 짧은 완전열
사이의 사상들과 양립한다는 확인은 생략한다.
\end{proof}

\begin{lemma}
\label{sites-cohomology-lemma-cech-map-into}
$\mathcal{C}$를 범주라 하자.
$\mathcal{U} = \{U_i \to U\}_{i \in I}$를 공역이 고정된 사상족으로서
모든 섬유곱 $U_{i_0} \times_U \ldots \times_U U_{i_p}$가
$\mathcal{C}$에 존재하는 것이라 하자. 아벨 준층의 사슬 복합체
$\mathbf{Z}_{\mathcal{U}, \bullet}$
$$
\ldots
\to
\bigoplus_{i_0i_1i_2} \mathbf{Z}_{U_{i_0} \times_U U_{i_1} \times_U U_{i_2}}
\to
\bigoplus_{i_0i_1} \mathbf{Z}_{U_{i_0} \times_U U_{i_1}}
\to
\bigoplus_{i_0} \mathbf{Z}_{U_{i_0}}
\to 0 \to \ldots
$$
를 생각하자. 여기서 마지막 영이 아닌 항은 차수 $0$에 놓이고, 사상
$$
\mathbf{Z}_{U_{i_0} \times_U \ldots \times_u U_{i_{p + 1}}}
\longrightarrow
\mathbf{Z}_{U_{i_0} \times_U
\ldots \widehat{U_{i_j}} \ldots \times_U U_{i_{p + 1}}}
$$
은 표준 사상의 $(-1)^j$배로 주어진다. 그러면 동형사상
$$
\Hom_{\textit{PAb}(\mathcal{C})}(\mathbf{Z}_{\mathcal{U}, \bullet}, \mathcal{F})
=
\check{\mathcal{C}}^\bullet(\mathcal{U}, \mathcal{F})
$$
은 $\mathcal{F} \in \Ob(\textit{PAb}(\mathcal{C}))$에 대해 함자적이다.
\end{lemma}

\begin{proof}
이는 다음 사실에 바탕을 둔 항진명제이다.
\begin{align*}
\Hom_{\textit{PAb}(\mathcal{C})}(
\bigoplus_{i_0 \ldots i_p}
\mathbf{Z}_{U_{i_0} \times_U \ldots \times_U U_{i_p}},
\mathcal{F})
& =
\prod_{i_0 \ldots i_p}
\Hom_{\textit{PAb}(\mathcal{C})}(
\mathbf{Z}_{U_{i_0} \times_U \ldots \times_U U_{i_p}},
\mathcal{F}) \\
& =
\prod_{i_0 \ldots i_p}
\mathcal{F}(U_{i_0} \times_U \ldots \times_U U_{i_p})
\end{align*}
Modules on Sites, Lemma \ref{sites-modules-lemma-obvious-adjointness}를 보라.
\end{proof}

\begin{lemma}
\label{sites-cohomology-lemma-homology-complex}
$\mathcal{C}$를 범주라 하고,
$\mathcal{U} = \{f_i : U_i \to U\}_{i \in I}$를 공역이 고정된
사상족으로서 모든 섬유곱
$U_{i_0} \times_U \ldots \times_U U_{i_p}$가 $\mathcal{C}$에
존재하는 것이라 하자. 위 보조정리 \ref{sites-cohomology-lemma-cech-map-into}의 준층
사슬 복합체 $\mathbf{Z}_{\mathcal{U}, \bullet}$는 양의 차수에서 완전하다.
즉, $i > 0$이면 호몰로지 준층
$H_i(\mathbf{Z}_{\mathcal{U}, \bullet})$는 영이다.
\end{lemma}

\begin{proof}
$V$를 $\mathcal{C}$의 대상이라 하자. 아벨군의 사슬 복합체
$\mathbf{Z}_{\mathcal{U}, \bullet}(V)$가 차수 $> 0$에서 완전함을 보이면
된다. 이 복합체는
$$
\xymatrix{
\ldots \ar[d] \\
\bigoplus_{i_0i_1i_2}
\mathbf{Z}[
\Mor_\mathcal{C}(V, U_{i_0} \times_U U_{i_1} \times_U U_{i_2})
]
\ar[d] \\
\bigoplus_{i_0i_1}
\mathbf{Z}[
\Mor_\mathcal{C}(V, U_{i_0} \times_U U_{i_1})
]
\ar[d] \\
\bigoplus_{i_0}
\mathbf{Z}[
\Mor_\mathcal{C}(V, U_{i_0})
] \ar[d] \\
0
}
$$
이다. 임의의 사상 $\varphi : V \to U$에 대해
$\Mor_\varphi(V, U_i) = \{\varphi_i : V \to U_i \mid
f_i \circ \varphi_i = \varphi\}$로 나타내자.
$\Mor_\varphi(V, U_{i_0} \times_U \ldots \times_U U_{i_p})$에도
비슷한 표기를 사용한다. 섬유곱
$U_{i_0} \times_U \ldots \times_U U_{i_p}$ 사이의 여러 사영사상과
합성해도 이 사상 집합들이 보존됨에 유의하자. 따라서 위의 복합체는
다음 복합체와 같다.
$$
\xymatrix{
\ldots \ar[d] \\
\bigoplus_\varphi
\bigoplus_{i_0i_1i_2}
\mathbf{Z}[
\Mor_\varphi(V, U_{i_0} \times_U U_{i_1} \times_U U_{i_2})
]
\ar[d] \\
\bigoplus_\varphi
\bigoplus_{i_0i_1}
\mathbf{Z}[
\Mor_\varphi(V, U_{i_0} \times_U U_{i_1})
]
\ar[d] \\
\bigoplus_\varphi
\bigoplus_{i_0}
\mathbf{Z}[
\Mor_\varphi(V, U_{i_0})
] \ar[d] \\
0
}
$$
다음으로
$$
\Mor_\varphi(V, U_{i_0} \times_U \ldots \times_U U_{i_p})
=
\Mor_\varphi(V, U_{i_0}) \times \ldots
\times \Mor_\varphi(V, U_{i_p})
$$
임을 관찰하자. 이 등식과 $\mathbf{Z}[A] \oplus \mathbf{Z}[B] =
\mathbf{Z}[A \amalg B]$라는 사실을 사용하면 복합체는 다음과 같이 된다.
$$
\xymatrix{
\ldots \ar[d] \\
\bigoplus_\varphi
\mathbf{Z}\left[
\coprod_{i_0i_1i_2}
\Mor_\varphi(V, U_{i_0}) \times \Mor_\varphi(V, U_{i_1}) \times
\Mor_\varphi(V, U_{i_2})
\right]
\ar[d] \\
\bigoplus_\varphi
\mathbf{Z}\left[
\coprod_{i_0i_1}
\Mor_\varphi(V, U_{i_0}) \times \Mor_\varphi(V, U_{i_1})
\right]
\ar[d] \\
\bigoplus_\varphi
\mathbf{Z}\left[
\coprod_{i_0}
\Mor_\varphi(V, U_{i_0})
\right] \ar[d] \\
0
}
$$
마지막으로 $S_\varphi = \coprod_{i \in I} \Mor_\varphi(V, U_i)$로 두면
다음을 얻는다.
$$
\bigoplus\nolimits_\varphi \left(\ldots \to
\mathbf{Z}[S_\varphi \times S_\varphi \times S_\varphi] \to
\mathbf{Z}[S_\varphi \times S_\varphi] \to
\mathbf{Z}[S_\varphi] \to 0 \to \ldots
\right)
$$
이로써 과제가 단순해졌다. 즉, 임의의 공집합이 아닌 집합 $S$에 대해
자유 아벨군의 확대 복합체
$$
\ldots \to
\mathbf{Z}[S \times S \times S] \to
\mathbf{Z}[S \times S] \to
\mathbf{Z}[S] \xrightarrow{\Sigma} \mathbf{Z} \to 0 \to \ldots
$$
가 모든 차수에서 완전함을 보이면 충분하다. 이를 위해 원소 $s \in S$를
고정하고, 자명한 표기 아래 호모토피
$$
n_{(s_0, \ldots, s_p)} \longmapsto n_{(s, s_0, \ldots, s_p)}
$$
를 사용하면 된다.
\end{proof}

\begin{lemma}
\label{sites-cohomology-lemma-complex-tensored-still-exact}
\begin{slogan}
정수 준층 체흐 복합체는 정수 상수 준층의 평탄 분해이다.
\end{slogan}
$\mathcal{C}$를 범주라 하고,
$\mathcal{U} = \{f_i : U_i \to U\}_{i \in I}$를 공역이 고정된
사상족으로서 모든 섬유곱
$U_{i_0} \times_U \ldots \times_U U_{i_p}$가 $\mathcal{C}$에
존재하는 것이라 하자. $\mathcal{O}$를 $\mathcal{C}$ 위의 환의 준층이라
하자. 사슬 복합체
$$
\mathbf{Z}_{\mathcal{U}, \bullet}
\otimes_{p, \mathbf{Z}}
\mathcal{O}
$$
는 양의 차수에서 완전하다. 여기서 $\mathbf{Z}_{\mathcal{U}, \bullet}$는
보조정리 \ref{sites-cohomology-lemma-cech-map-into}의 사슬 복합체이고, 텐서곱은 값이
$\mathbf{Z}$인 환의 상수 준층 위에서 취한다.
\end{lemma}

\begin{proof}
$V$를 $\mathcal{C}$의 대상이라 하자. 보조정리
\ref{sites-cohomology-lemma-homology-complex}의 증명에서
$\mathbf{Z}_{\mathcal{U}, \bullet}(V)$는 복합체로서 차수 영에 놓인
$\mathbf{Z}$와 호모토픽인 복합체들의 직합과 동형임을 보았다. 따라서
$\mathbf{Z}_{\mathcal{U}, \bullet}(V) \otimes_\mathbf{Z} \mathcal{O}(V)$
도 복합체로서 차수 영에 놓인 $\mathcal{O}(V)$와 호모토픽인 복합체들의
직합과 동형이다. 또는 Modules on Sites, Lemma
\ref{sites-modules-lemma-flat-resolution-of-flat}을 사용할 수도 있다.
실제로 준층 $\mathbf{Z}_{\mathcal{U}, i}$들은 평탄하고, 보조정리
\ref{sites-cohomology-lemma-homology-complex}의 증명은
$H_0(\mathbf{Z}_{\mathcal{U}, \bullet})$도 평탄한 준층임을 보인다.
\end{proof}

\begin{lemma}
\label{sites-cohomology-lemma-cech-cohomology-derived-presheaves}
$\mathcal{C}$를 범주라 하고,
$\mathcal{U} = \{f_i : U_i \to U\}_{i \in I}$를 공역이 고정된
사상족으로서 모든 섬유곱
$U_{i_0} \times_U \ldots \times_U U_{i_p}$가 $\mathcal{C}$에
존재하는 것이라 하자. 체흐 코호몰로지 함자
$\check{H}^p(\mathcal{U}, -)$들은 하나의 $\delta$-함자로서 다음 함자의
오른쪽 유도함자들과 표준적으로 동형이다.
$$
\check{H}^0(\mathcal{U}, -) :
\textit{PAb}(\mathcal{C})
\longrightarrow
\textit{Ab}.
$$
또한 함자적 유사동형
$$
\check{\mathcal{C}}^\bullet(\mathcal{U}, \mathcal{F})
\longrightarrow
R\check{H}^0(\mathcal{U}, \mathcal{F})
$$
이 있다. 여기서 오른쪽은 왼쪽 완전함자
$\check{H}^0(\mathcal{U}, -)$의 유도함자
$$
R\check{H}^0(\mathcal{U}, -) :
D^{+}(\textit{PAb}(\mathcal{C}))
\longrightarrow
D^{+}(\mathbf{Z})
$$
를 나타낸다.
\end{lemma}

\begin{proof}
아벨 준층의 범주에는 단사가 충분히 있다. Injectives, Proposition
\ref{injectives-proposition-presheaves-injectives}를 보라. 또한
$\check{H}^0(\mathcal{U}, -)$는 아벨 준층의 범주에서
$\mathbf{Z}$-가군의 범주로 가는 왼쪽 완전함자이다. 따라서 유도함자와
오른쪽 유도함자가 존재한다. Derived Categories, Section
\ref{derived-section-right-derived-functor}를 보라.

\medskip\noindent
$\mathcal{I}$를 단사 아벨 준층이라 하자. 이 경우 함자
$\Hom_{\textit{PAb}(\mathcal{C})}(-, \mathcal{I})$
는 $\textit{PAb}(\mathcal{C})$에서 완전하다. 보조정리
\ref{sites-cohomology-lemma-cech-map-into}에 의해
$$
\Hom_{\textit{PAb}(\mathcal{C})}(
\mathbf{Z}_{\mathcal{U}, \bullet}, \mathcal{I})
=
\check{\mathcal{C}}^\bullet(\mathcal{U}, \mathcal{I}).
$$
이다. 보조정리 \ref{sites-cohomology-lemma-homology-complex}에 의해
$\mathbf{Z}_{\mathcal{U}, \bullet}$는 양의 차수에서 완전하다. 따라서
위에서 언급한 $\mathcal{I}$로 들어가는 Hom의 완전성으로부터 모든
$i > 0$에 대해 $\check{H}^i(\mathcal{U}, \mathcal{I}) = 0$임을 알 수
있다. 그러므로 $\delta$-함자 $(\check{H}^n, \delta)$(보조정리
\ref{sites-cohomology-lemma-cech-cohomology-delta-functor-presheaves} 참조)는 Homology,
Lemma \ref{homology-lemma-efface-implies-universal}의 가정을 만족하며,
따라서 보편 $\delta$-함자이다.

\medskip\noindent
Derived Categories, Lemma \ref{derived-lemma-higher-derived-functors}에 의해
열 $R^i\check{H}^0(\mathcal{U}, -)$도 보편 $\delta$-함자를 이룬다.
보편 $\delta$-함자의 유일성(Homology, Lemma
\ref{homology-lemma-uniqueness-universal-delta-functor} 참조)에 의해
$R^i\check{H}^0(\mathcal{U}, -) = \check{H}^i(\mathcal{U}, -)$를 얻는다.
이는 대부분의 응용에 충분하므로 독자는 증명의 나머지를 건너뛰어도 좋다.

\medskip\noindent
$\mathcal{F}$를 $\mathcal{C}$ 위의 임의의 아벨 준층이라 하자. 범주
$\textit{PAb}(\mathcal{C})$에서 단사 분해
$\mathcal{F} \to \mathcal{I}^\bullet$을 택하자. 항들이
$\check{\mathcal{C}}^p(\mathcal{U}, \mathcal{I}^q)$인 이중 복합체
$\check{\mathcal{C}}^\bullet(\mathcal{U}, \mathcal{I}^\bullet)$
와 결부된 전체 복합체
$\text{Tot}(\check{\mathcal{C}}^\bullet(\mathcal{U}, \mathcal{I}^\bullet))$
를 생각하자. Homology, Section
\ref{homology-section-double-complexes}를 보라. 사상들
$\check{\mathcal{C}}^p(\mathcal{U}, \mathcal{F})
\to \check{\mathcal{C}}^p(\mathcal{U}, \mathcal{I}^0)$에서 오는 복합체의 사상
$$
\check{\mathcal{C}}^\bullet(\mathcal{U}, \mathcal{F})
\longrightarrow
\text{Tot}(\check{\mathcal{C}}^\bullet(\mathcal{U}, \mathcal{I}^\bullet))
$$
이 있고, 사상들
$\check{H}^0(\mathcal{U}, \mathcal{I}^q) \to
\check{\mathcal{C}}^0(\mathcal{U}, \mathcal{I}^q)$에서 오는 복합체의 사상
$$
\check{H}^0(\mathcal{U}, \mathcal{I}^\bullet)
\longrightarrow
\text{Tot}(\check{\mathcal{C}}^\bullet(\mathcal{U}, \mathcal{I}^\bullet))
$$
도 있다. Homology, Lemma
\ref{homology-lemma-double-complex-gives-resolution}을 적용하면 이 두 사상은
모두 유사동형이다. 실제로 함자로서의 체흐 복합체가 완전하므로(보조정리
\ref{sites-cohomology-lemma-cech-exact-presheaves}) 이중 복합체의 열들은 양의 차수에서
완전하고, 방금 본 바와 같이 단사 준층 $\mathcal{I}^q$의 고차 체흐
코호몰로지군들이 영이므로 이중 복합체의 행들도 양의 차수에서 완전하다.
유사동형은 $D^{+}(\mathbf{Z})$에서 가역이 되므로 보조정리의 마지막에
표시한 사상을 얻는다. 이 사상이 함자적이라는 확인은 생략한다.
\end{proof}





\section{체흐 코호몰로지와 코호몰로지}
\label{sites-cohomology-section-cech-cohomology-cohomology}

\noindent
코호몰로지와 체흐 코호몰로지의 관계는 단사 아벨 층의 체흐 코호몰로지가
영이라는 사실에서 비롯된다. 이를 보기 위해 단사 아벨 층이 단사 아벨
준층임을 관찰한 뒤 앞 절의 체흐 코호몰로지에 관한 결과들을 적용한다.

\begin{lemma}
\label{sites-cohomology-lemma-injective-abelian-sheaf-injective-presheaf}
$\mathcal{C}$를 사이트라 하자. 단사 아벨 층은 범주
$\textit{PAb}(\mathcal{C})$의 대상으로서도 단사이다.
\end{lemma}

\begin{proof}
범주 $\mathcal{A} = \textit{Ab}(\mathcal{C})$와
$\mathcal{B} = \textit{PAb}(\mathcal{C})$, 포함 함자와 층화에 Homology,
Lemma \ref{homology-lemma-adjoint-preserve-injectives}를 적용한다.
보조정리의 모든 가정이 만족됨은 Modules on Sites, Section
\ref{sites-modules-section-abelian-sheaves}를 보라.
\end{proof}

\begin{lemma}
\label{sites-cohomology-lemma-injective-trivial-cech}
$\mathcal{C}$를 사이트라 하고,
$\mathcal{U} = \{U_i \to U\}_{i \in I}$를 $\mathcal{C}$의 피복이라 하자.
$\mathcal{I}$를 단사 아벨 층, 즉 $\textit{Ab}(\mathcal{C})$의 단사
대상이라 하자. 그러면
$$
\check{H}^p(\mathcal{U}, \mathcal{I}) =
\left\{
\begin{matrix}
\mathcal{I}(U) & \text{if} & p = 0 \\
0 & \text{if} & p > 0
\end{matrix}
\right.
$$
\end{lemma}

\begin{proof}
보조정리 \ref{sites-cohomology-lemma-injective-abelian-sheaf-injective-presheaf}에 의해
$\mathcal{I}$는 $\textit{PAb}(\mathcal{C})$의 단사 대상이다. 따라서
보조정리 \ref{sites-cohomology-lemma-cech-cohomology-derived-presheaves}(또는 그 증명)를
적용하면 고차 체흐 코호몰로지군의 소멸을 알 수 있다. 제영차에 대해서는
보조정리 \ref{sites-cohomology-lemma-cech-h0}을 보라.
\end{proof}

\begin{lemma}
\label{sites-cohomology-lemma-cech-cohomology}
$\mathcal{C}$를 사이트라 하고,
$\mathcal{U} = \{U_i \to U\}_{i \in I}$를 $\mathcal{C}$의 피복이라 하자.
함자들 $\textit{Ab}(\mathcal{C}) \to D^{+}(\mathbf{Z})$ 사이의 변환
$$
\check{\mathcal{C}}^\bullet(\mathcal{U}, -)
\longrightarrow
R\Gamma(U, -)
$$
이 있다. 특히 $\mathcal{F}$가 $\textit{Ab}(\mathcal{C})$를 달릴 때
이는 함자들의 변환
$\check{H}^p(\mathcal{U}, \mathcal{F}) \to H^p(U, \mathcal{F})$를 준다.
\end{lemma}

\begin{proof}
$\mathcal{F}$를 아벨 층이라 하고 단사 분해
$\mathcal{F} \to \mathcal{I}^\bullet$을 택하자. 항들이
$\check{\mathcal{C}}^p(\mathcal{U}, \mathcal{I}^q)$인 이중 복합체
$\check{\mathcal{C}}^\bullet(\mathcal{U}, \mathcal{I}^\bullet)$와 이에
결부된 전체 복합체
$\text{Tot}(\check{\mathcal{C}}^\bullet(\mathcal{U}, \mathcal{I}^\bullet))$
를 생각하자. Homology, Definition
\ref{homology-definition-associated-simple-complex}을 보라. 사상들
$\mathcal{I}^q(U) \to \check{H}^0(\mathcal{U}, \mathcal{I}^q)$에서 오는
복합체의 사상
$$
\alpha :
\Gamma(U, \mathcal{I}^\bullet)
\longrightarrow
\text{Tot}(\check{\mathcal{C}}^\bullet(\mathcal{U}, \mathcal{I}^\bullet))
$$
과 사상 $\mathcal{F} \to \mathcal{I}^0$에서 오는 복합체의 사상
$$
\beta :
\check{\mathcal{C}}^\bullet(\mathcal{U}, \mathcal{F})
\longrightarrow
\text{Tot}(\check{\mathcal{C}}^\bullet(\mathcal{U}, \mathcal{I}^\bullet))
$$
이 있다. Homology, Lemma
\ref{homology-lemma-double-complex-gives-resolution}을 적용하면 $\alpha$가
유사동형임을 알 수 있다. 실제로 보조정리
\ref{sites-cohomology-lemma-injective-trivial-cech}에 의해 이중 복합체
$\check{\mathcal{C}}^\bullet(\mathcal{U}, \mathcal{I}^\bullet)$의
$q$번째 행은 $\Gamma(U, \mathcal{I}^q)$의 분해이다. 따라서 $\alpha$는
$D^{+}(\mathbf{Z})$에서 가역이 되고, 보조정리의 변환은 $\beta$에
$\alpha$의 역을 뒤이어 합성한 것이다. 이것이 함자적이라는 확인은 생략한다.
\end{proof}

\begin{lemma}
\label{sites-cohomology-lemma-cech-h1}
$\mathcal{C}$를 사이트, $\mathcal{G}$를 $\mathcal{C}$ 위의 아벨 층,
$\mathcal{U} = \{U_i \to U\}_{i \in I}$를 $\mathcal{C}$의 피복이라 하자.
사상
$$
\check{H}^1(\mathcal{U}, \mathcal{G})
\longrightarrow
H^1(U, \mathcal{G})
$$
은 단사이며, 보조정리 \ref{sites-cohomology-lemma-torsors-h1}의 전단사를 통해
$\check{H}^1(\mathcal{U}, \mathcal{G})$를 각 $U_i$ 위로 제한했을 때
자명한 토서가 되는 $\mathcal{G}|_U$-토서들의 동형류 집합과 식별한다.
\end{lemma}

\begin{proof}
이를 보이기 위해 역 사상을 구성한다. 즉, $\mathcal{F}$를
$\mathcal{C}/U$ 위의 $\mathcal{G}|_U$-토서로서
$\mathcal{C}/U_i$로의 제한이 자명한 것이라 하자. 보조정리
\ref{sites-cohomology-lemma-trivial-torsor}에 의해 이는 단면
$s_i \in \mathcal{F}(U_i)$가 존재한다는 뜻이다.
$U_{i_0} \times_U U_{i_1}$ 위에는 다음을 만족하는 $\mathcal{G}$의 유일한
단면 $s_{i_0i_1}$이 있다.
$s_{i_0i_1} \cdot s_{i_0}|_{U_{i_0} \times_U U_{i_1}} =
s_{i_1}|_{U_{i_0} \times_U U_{i_1}}$. 계산하면 $s_{i_0i_1}$이 체흐
코사이클이고 그 류가 잘 정의됨을 알 수 있다(즉, 단면 $s_i$의 선택에
의존하지 않는다). 역 사상은 $\mathcal{F}$의 동형류를 코사이클
$(s_{i_0i_1})$의 코호몰로지류로 보낸다. 이 사상이 실제로 역이라는 확인은
생략한다.
\end{proof}

\begin{lemma}
\label{sites-cohomology-lemma-include}
$\mathcal{C}$를 사이트라 하자. 함자
$i : \textit{Ab}(\mathcal{C}) \to \textit{PAb}(\mathcal{C})$를
생각하자. 이는 왼쪽 완전함자이고 그 오른쪽 유도함자는 다음과 같다.
$$
R^pi(\mathcal{F}) = \underline{H}^p(\mathcal{F}) :
U \longmapsto H^p(U, \mathcal{F})
$$
Section \ref{sites-cohomology-section-locality}의 논의를 보라.
\end{lemma}

\begin{proof}
$i$가 왼쪽 완전임은 분명하다. 단사 분해
$\mathcal{F} \to \mathcal{I}^\bullet$을 택하자. 정의에 따라 $R^pi$는
복합체 $\mathcal{I}^\bullet$의 $p$번째 코호몰로지 {\it 준층}이다.
다시 말해 $\mathcal{C}$의 대상 $U$ 위에서 $R^pi(\mathcal{F})$의 단면은
다음과 같다.
$$
\frac{\Ker(\mathcal{I}^n(U) \to \mathcal{I}^{n + 1}(U))}
{\Im(\mathcal{I}^{n - 1}(U) \to \mathcal{I}^n(U))}.
$$
이는 $H^p(U, \mathcal{F})$의 정의이다.
\end{proof}

\begin{lemma}
\label{sites-cohomology-lemma-cech-spectral-sequence}
$\mathcal{C}$를 사이트라 하고,
$\mathcal{U} = \{U_i \to U\}_{i \in I}$를 $\mathcal{C}$의 피복이라 하자.
임의의 아벨 층 $\mathcal{F}$에 대해
$$
E_2^{p, q} = \check{H}^p(\mathcal{U}, \underline{H}^q(\mathcal{F}))
$$
를 가지며 $H^{p + q}(U, \mathcal{F})$로 수렴하는 스펙트럼열
$(E_r, d_r)_{r \geq 0}$이 있다. 이 스펙트럼열은 $\mathcal{F}$에 관해
함자적이다.
\end{lemma}

\begin{proof}
이는 함자들
$$
i :  \textit{Ab}(\mathcal{C}) \to \textit{PAb}(\mathcal{C})
\quad\text{and}\quad
\check{H}^0(\mathcal{U}, - ) : \textit{PAb}(\mathcal{C})
\to \textit{Ab}.
$$
에 대한 Grothendieck 스펙트럼열이다(Derived Categories, Lemma
\ref{derived-lemma-grothendieck-spectral-sequence} 참조). 실제로 보조정리
\ref{sites-cohomology-lemma-cech-h0}에 의해
$\check{H}^0(\mathcal{U}, i(\mathcal{F})) = \mathcal{F}(U)$이다. 보조정리
\ref{sites-cohomology-lemma-injective-trivial-cech}에 의해 $i(\mathcal{I})$는 체흐
비순환이다. 그리고 보조정리
\ref{sites-cohomology-lemma-cech-cohomology-derived-presheaves}에 의해
$\textit{PAb}(\mathcal{C})$ 위의 함자로서
$\check{H}^p(\mathcal{U}, -) = R^p\check{H}^0(\mathcal{U}, -)$이다.
이들을 모두 합치면 보조정리를 얻는다.
\end{proof}

\begin{lemma}
\label{sites-cohomology-lemma-cech-spectral-sequence-application}
$\mathcal{C}$를 사이트,
$\mathcal{U} = \{U_i \to U\}_{i \in I}$를 피복,
$\mathcal{F} \in \Ob(\textit{Ab}(\mathcal{C}))$라 하자. 모든 $i > 0$,
모든 $p \geq 0$, 모든 $i_0, \ldots, i_p \in I$에 대해
$H^i(U_{i_0} \times_U \ldots \times_U U_{i_p}, \mathcal{F}) = 0$이라고
가정하자. 그러면
$\check{H}^p(\mathcal{U}, \mathcal{F}) = H^p(U, \mathcal{F})$이다.
\end{lemma}

\begin{proof}
보조정리 \ref{sites-cohomology-lemma-cech-spectral-sequence}의 스펙트럼열을 사용한다.
가정은 $q \not = 0$인 모든 $(p, q)$에 대해
$E_2^{p, q} = 0$이라는 뜻이다. 따라서 스펙트럼열은 $E_2$에서 퇴화하고
결과가 따른다.
\end{proof}

\begin{lemma}
\label{sites-cohomology-lemma-ses-cech-h1}
$\mathcal{C}$를 사이트라 하자.
$$
0 \to \mathcal{F} \to \mathcal{G} \to \mathcal{H} \to 0
$$
을 $\mathcal{C}$ 위의 아벨 층들의 짧은 완전열이라 하자. $U$를
$\mathcal{C}$의 대상이라 하자.
$\check{H}^1(\mathcal{U}, \mathcal{F}) = 0$을 만족하는 $U$의 피복
$\mathcal{U}$들의 공종계가 존재하면 사상
$\mathcal{G}(U) \to \mathcal{H}(U)$는 전사이다.
\end{lemma}

\begin{proof}
원소 $s \in \mathcal{H}(U)$를 택하자. (a)
$\check{H}^1(\mathcal{U}, \mathcal{F}) = 0$이고 (b) $s|_{U_i}$가 어떤
단면 $s_i \in \mathcal{G}(U_i)$의 상이 되도록 피복
$\mathcal{U} = \{U_i \to U\}_{i \in I}$를 택하자. (b)를 만족하는 피복은
확실히 찾을 수 있으므로, 보조정리의 가정에 의해 (a)와 (b)를 모두 만족하는
피복을 찾을 수 있다. 단면들
$$
s_{i_0i_1} =
s_{i_1}|_{U_{i_0} \times_U U_{i_1}} - s_{i_0}|_{U_{i_0} \times_U U_{i_1}}.
$$
을 생각하자. $s_i$가 $s$의 올림이므로
$s_{i_0i_1} \in \mathcal{F}(U_{i_0} \times_U U_{i_1})$임을 알 수 있다.
$\check{H}^1(\mathcal{U}, \mathcal{F})$가 영이므로 다음을 만족하는 단면
$t_i \in \mathcal{F}(U_i)$들을 찾을 수 있다.
$$
s_{i_0i_1} =
t_{i_1}|_{U_{i_0} \times_U U_{i_1}} - t_{i_0}|_{U_{i_0} \times_U U_{i_1}}.
$$
그러면 단면들 $s_i - t_i$는 분명 층 조건을 만족하며 서로 붙어서 $s$로
가는 $U$ 위의 $\mathcal{G}$의 단면이 된다. 따라서 결론을 얻는다.
\end{proof}

\begin{lemma}
\label{sites-cohomology-lemma-cech-vanish-collection}
(Cohomology, Lemma \ref{cohomology-lemma-cech-vanish}의 변형.)
$\mathcal{C}$를 사이트라 하고, $\text{Cov}_\mathcal{C}$를
$\mathcal{C}$의 피복들의 집합이라 하자(Sites, Definition
\ref{sites-definition-site} 참조).
$\mathcal{B} \subset \Ob(\mathcal{C})$와
$\text{Cov} \subset \text{Cov}_\mathcal{C}$를 부분집합들이라 하자.
$\mathcal{F}$를 $\mathcal{C}$ 위의 아벨 층이라 하고 다음을 가정하자.
\begin{enumerate}
\item $\mathcal{U} = \{U_i \to U\}_{i \in I}$인 모든
$\mathcal{U} \in \text{Cov}$에 대해 $U, U_i \in \mathcal{B}$이고 모든
$U_{i_0} \times_U \ldots \times_U U_{i_p} \in \mathcal{B}$이다.
\item 모든 $U \in \mathcal{B}$에 대해 $\text{Cov}$에 나타나는 $U$의
피복들은 $U$의 피복들의 공종계를 이룬다.
\item 모든 $\mathcal{U} \in \text{Cov}$와 모든 $p > 0$에 대해
$\check{H}^p(\mathcal{U}, \mathcal{F}) = 0$이다.
\end{enumerate}
그러면 모든 $p > 0$과 모든 $U \in \mathcal{B}$에 대해
$H^p(U, \mathcal{F}) = 0$이다.
\end{lemma}

\begin{proof}
$\mathcal{F}$와 $\text{Cov}$가 보조정리에서와 같다고 하자. 이를
“$\mathcal{F}$의 고차 체흐 코호몰로지는 임의의
$\mathcal{U} \in \text{Cov}$에 대해 영이다”라고 말하겠다.
$\mathcal{F} \to \mathcal{I}$를 단사 아벨 층으로 가는 매입으로 택하자.
보조정리 \ref{sites-cohomology-lemma-injective-trivial-cech}에 의해 $\mathcal{I}$의 고차
체흐 코호몰로지는 임의의 $\mathcal{U} \in \text{Cov}$에 대해 영이다.
$\mathcal{Q} = \mathcal{I}/\mathcal{F}$로 놓으면 짧은 완전열
$$
0 \to \mathcal{F} \to \mathcal{I} \to \mathcal{Q} \to 0.
$$
을 얻는다. 보조정리 \ref{sites-cohomology-lemma-ses-cech-h1}과 가정 (2)에 의해 이 열은 모든
$U \in \mathcal{B}$에 대해 완전열
$$
0 \to \mathcal{F}(U) \to \mathcal{I}(U) \to \mathcal{Q}(U) \to 0.
$$
을 낳는다. 따라서 임의의 $\mathcal{U} \in \text{Cov}$에 대해 체흐
복합체들의 짧은 완전열
$$
0 \to
\check{\mathcal{C}}^\bullet(\mathcal{U}, \mathcal{F}) \to
\check{\mathcal{C}}^\bullet(\mathcal{U}, \mathcal{I}) \to
\check{\mathcal{C}}^\bullet(\mathcal{U}, \mathcal{Q}) \to 0
$$
을 얻는다. 가정 (1)에 의해 체흐 복합체의 각 항은 $\mathcal{B}$의
원소들 위에서 취한 값들의 곱으로 이루어지기 때문이다. 특히 임의의 피복
$\mathcal{U} \in \text{Cov}$에 대해 체흐 코호몰로지군들의 긴 완전열을
얻는다. 이로부터 $\mathcal{Q}$도 모든 $\mathcal{U} \in \text{Cov}$에
대해 고차 체흐 코호몰로지가 영인 아벨 층임을 알 수 있다.

\medskip\noindent
이제 임의의 $U \in \mathcal{B}$에 대한 코호몰로지 긴 완전열
$$
\xymatrix{
0 \ar[r] &
H^0(U, \mathcal{F}) \ar[r] &
H^0(U, \mathcal{I}) \ar[r] &
H^0(U, \mathcal{Q}) \ar[lld] \\
&
H^1(U, \mathcal{F}) \ar[r] &
H^1(U, \mathcal{I}) \ar[r] &
H^1(U, \mathcal{Q}) \ar[lld] \\
&
\ldots & \ldots & \ldots \\
}
$$
을 살펴보자. $\mathcal{I}$가 단사이므로 $n > 0$에 대해
$H^n(U, \mathcal{I}) = 0$이다(Derived Categories, Lemma
\ref{derived-lemma-higher-derived-functors} 참조). 위에서
$H^0(U, \mathcal{I}) \to H^0(U, \mathcal{Q})$가 전사임을 보았으므로
$H^1(U, \mathcal{F}) = 0$이다. $\mathcal{F}$는 모든
$\mathcal{U} \in \text{Cov}$에 대해 고차 체흐 코호몰로지가 영인 임의의
아벨 층이었고 $\mathcal{Q}$도 그러한 층이므로(위 참조)
$H^1(U, \mathcal{Q}) = 0$도 얻는다. 그러면 긴 완전열에 의해
$H^2(U, \mathcal{F}) = 0$이 따른다. 같은 과정을 계속한다.
\end{proof}














\section{제2 코호몰로지와 제르브}
\label{sites-cohomology-section-gerbes}

\noindent
$p : \mathcal{S} \to \mathcal{C}$를 모든 자기동형군이 가환인 사이트 위의
제르브라 하자. 이 경우 Stacks, Lemma
\ref{stacks-lemma-gerbe-abelian-auts}의 자기동형층의 제1 및 제2
코호몰로지군은 대상의 존재를 제어한다.

\medskip\noindent
다음 보조정리는 앞으로 이 관계에 관한 더 완전한 논의를 추가하면
그 논의로 대체될 것이다.

\begin{lemma}
\label{sites-cohomology-lemma-existence}
$\mathcal{C}$를 사이트라 하고,
$p : \mathcal{S} \to \mathcal{C}$를 자기동형층이 아벨인 사이트 위의
제르브라 하자. $\mathcal{G}$를 Stacks, Lemma
\ref{stacks-lemma-gerbe-abelian-auts}에서 구성한 아벨군층이라 하자.
$U$를 $\mathcal{C}$의 대상으로서 다음을 만족하는 것이라 하자.
\begin{enumerate}
\item 모든 $i, j$에 대해 $H^1(U_i, \mathcal{G}) = 0$이고
$H^1(U_i \times_U U_j, \mathcal{G}) = 0$인 $\mathcal{C}$에서의
$U$의 피복 $\{U_i \to U\}$들의 공종계가 존재한다.
\item $H^2(U, \mathcal{G}) = 0$.
\end{enumerate}
그러면 $U$ 위에 놓이는 $\mathcal{S}$의 대상이 존재한다.
\end{lemma}

\begin{proof}
Stacks, Definition \ref{stacks-definition-gerbe}에 의해 피복
$\mathcal{U} = \{U_i \to U\}$와 $U_i$ 위에 놓이는 $\mathcal{S}$의 대상
$x_i$가 존재한다. $U_{ij} = U_i \times_U U_j$로 쓰자. (1)에 의해 피복을
세분한 뒤 $H^1(U_i, \mathcal{G}) = 0$이고
$H^1(U_{ij}, \mathcal{G}) = 0$이라고 가정할 수 있다. $\mathcal{C}/U_{ij}$
위의 층
$$
\mathcal{F}_{ij} =
\mathit{Isom}(x_i|_{U_{ij}}, x_j|_{U_{ij}})
$$
을 생각하자.
$\mathcal{G}|_{U_{ij}} = \mathit{Aut}(x_i|_{U_{ij}})$이므로 앞합성에 의한
작용
$$
\mathcal{G}|_{U_{ij}} \times \mathcal{F}_{ij} \to \mathcal{F}_{ij}
$$
이 있다. $\mathcal{F}_{ij}$는 유사 $\mathcal{G}|_{U_{ij}}$-토서임이
분명하며, 제르브의 임의의 두 대상은 국소적으로 동형이므로 실제로 토서이다.
피복의 선택과 보조정리 \ref{sites-cohomology-lemma-torsors-h1}에 의해 이 토서들은
자명하다(따라서 보조정리 \ref{sites-cohomology-lemma-trivial-torsor}에 의해 전역 단면을
가진다). 다시 말해 동형사상
$$
\varphi_{ij} : x_i|_{U_{ij}} \longrightarrow x_j|_{U_{ij}}
$$
들을 택할 수 있다. $U$ 위의 대상 $x$를 찾기 위해 이 $\varphi_{ij}$들의
선택을 조정하여 내림 자료를 얻을 것이다
($p : \mathcal{S} \to \mathcal{C}$가 스택이므로 이 내림 자료는 반드시
유효하다). 내림 자료가 되는 데 대한 장애는 코사이클 조건이 성립하지 않을
수 있다는 것이다. $U_{ijk} = U_i \times_U U_j \times_U U_k$로 놓고
$$
g_{ijk} = \varphi_{ik}^{-1}|_{U_{ijk}} \circ \varphi_{jk}|_{U_{ijk}} \circ
\varphi_{ij}|_{U_{ijk}}
$$
를 생각하자. 이는 $U_{ijk}$ 위에서 $x_i$의 자기동형이다. 따라서
$g_{ijk}$를 $U_{ijk}$ 위의 $\mathcal{G}$의 단면으로 본다. 계산(생략)은
$(g_{i_0i_1i_2})$가 피복 $\mathcal{U}$에 대한 $\mathcal{G}$의 체흐 복합체
${\check C}^\bullet(\mathcal{U}, \mathcal{G})$의 $2$-코사이클임을 보인다.
보조정리 \ref{sites-cohomology-lemma-cech-spectral-sequence}의 스펙트럼열과 모든 $i$에
대한 $H^1(U_i, \mathcal{G}) = 0$으로부터 사상
${\check H}^2(\mathcal{U}, \mathcal{G}) \to H^2(U, \mathcal{G})$가
단사임을 알 수 있다. 따라서 $H^2(U, \mathcal{G}) = 0$이라는 가정에 의해
$(g_{i_0i_1i_2})$는 코경계이다. 그러므로 모든 $i, j, k$에 대해
$g_{ik}^{-1}|_{U_{ijk}} g_{jk}|_{U_{ijk}}  g_{ij}|_{U_{ijk}} = g_{ijk}$
를 만족하는 단면 $g_{ij} \in \mathcal{G}(U_{ij})$들을 찾을 수 있다.
$\varphi_{ij}$를 $\varphi_{ij}g_{ij}^{-1}$로 바꾸면
$\varphi_{ij}$는 $U_i$ 위의 대상 $x_i$들에 대한 내림 자료를 주며,
이로써 증명이 끝난다.
\end{proof}












\section{가군의 코호몰로지}
\label{sites-cohomology-section-cohomology-modules}

\noindent
아벨 군층의 코호몰로지에 관해 말한 모든 내용은 가군의
코호몰로지에도 그대로 성립한다. 실제로 이 둘은 일치한다.

\begin{lemma}
\label{sites-cohomology-lemma-injective-module-injective-presheaf}
$(\mathcal{C}, \mathcal{O})$를 환 달린 사이트라 하자.
단사 가군층은 범주 $\textit{PMod}(\mathcal{O})$의
대상으로서도 단사이다.
\end{lemma}

\begin{proof}
범주 $\mathcal{A} = \textit{Mod}(\mathcal{O})$와
$\mathcal{B} = \textit{PMod}(\mathcal{O})$, 포함 함자 및 층화에
Homology, Lemma \ref{homology-lemma-adjoint-preserve-injectives}를
적용한다. (이 보조정리의 모든 가정이 성립함은
Modules on Sites,
Section \ref{sites-modules-section-sheafification-presheaves-modules}을
참조하라.)
\end{proof}

\begin{lemma}
\label{sites-cohomology-lemma-include-modules}
$(\mathcal{C}, \mathcal{O})$를 환 달린 사이트라 하자.
함자
$i : \textit{Mod}(\mathcal{C}) \to \textit{PMod}(\mathcal{C})$.
를 생각하자. 이는 왼쪽 완전 함자이고 그 오른쪽 유도함자는
다음으로 주어진다.
$$
R^pi(\mathcal{F}) = \underline{H}^p(\mathcal{F}) :
U \longmapsto H^p(U, \mathcal{F})
$$
이에 관한 논의는 Section \ref{sites-cohomology-section-locality}를 참조하라.
\end{lemma}

\begin{proof}
$i$가 왼쪽 완전임은 분명하다.
$\textit{Mod}(\mathcal{O})$에서 단사 분해
$\mathcal{F} \to \mathcal{I}^\bullet$을 택하자.
정의상 $R^pi$는 복합체 $\mathcal{I}^\bullet$의 제$p$
코호몰로지 {\it 준층}이다. 다시 말해 $\mathcal{C}$의 대상 $U$ 위에서
$R^pi(\mathcal{F})$의 단면은
$$
\frac{\Ker(\mathcal{I}^n(U) \to \mathcal{I}^{n + 1}(U))}
{\Im(\mathcal{I}^{n - 1}(U) \to \mathcal{I}^n(U))}.
$$
로 주어지며, 이것이 $H^p(U, \mathcal{F})$의 정의이다.
\end{proof}

\begin{lemma}
\label{sites-cohomology-lemma-injective-module-trivial-cech}
$(\mathcal{C}, \mathcal{O})$를 환 달린 사이트라 하자.
$\mathcal{U} = \{U_i \to U\}_{i \in I}$를 $\mathcal{C}$의 덮개라 하자.
$\mathcal{I}$를 단사 $\mathcal{O}$-가군, 즉
$\textit{Mod}(\mathcal{O})$의 단사 대상이라 하자. 그러면
$$
\check{H}^p(\mathcal{U}, \mathcal{I}) =
\left\{
\begin{matrix}
\mathcal{I}(U) & \text{일 때} & p = 0 \\
0 & \text{일 때} & p > 0
\end{matrix}
\right.
$$
\end{lemma}

\begin{proof}
다음 등식의 연쇄에서 첫째 등식은
Lemma \ref{sites-cohomology-lemma-cech-map-into}가 준다.
\begin{align*}
\check{\mathcal{C}}^\bullet(\mathcal{U}, \mathcal{I})
& =
\Mor_{\textit{PAb}(\mathcal{C})}(
\mathbf{Z}_{\mathcal{U}, \bullet}, \mathcal{I}) \\
& =
\Mor_{\textit{PMod}(\mathbf{Z})}(
\mathbf{Z}_{\mathcal{U}, \bullet}, \mathcal{I}) \\
& =
\Mor_{\textit{PMod}(\mathcal{O})}(
\mathbf{Z}_{\mathcal{U}, \bullet} \otimes_{p, \mathbf{Z}} \mathcal{O},
\mathcal{I})
\end{align*}
셋째 등식은 Modules on Sites,
Lemma \ref{sites-modules-lemma-adjointness-tensor-restrict-presheaves}에서
따른다. Lemma \ref{sites-cohomology-lemma-injective-module-injective-presheaf}에 의해
$\mathcal{I}$는 $\textit{PMod}(\mathcal{O})$의 단사 대상이다.
따라서 $\Hom_{\textit{PMod}(\mathcal{O})}(-, \mathcal{I})$는
완전 함자이다. Lemma \ref{sites-cohomology-lemma-complex-tensored-still-exact}에 의해
고차 체흐 코호몰로지 군들은 소멸한다.
제0차의 경우는 Lemma \ref{sites-cohomology-lemma-cech-h0}를 참조하라.
\end{proof}

\begin{lemma}
\label{sites-cohomology-lemma-cohomology-modules-abelian-agree}
$\mathcal{C}$를 사이트라 하고 $\mathcal{O}$를 $\mathcal{C}$ 위의
환층이라 하자. $\mathcal{F}$를 $\mathcal{O}$-가군이라 하고,
그 기저 아벨 군층을 $\mathcal{F}_{ab}$로 나타내자. 그러면
$$
H^i(\mathcal{C}, \mathcal{F}_{ab})
=
H^i(\mathcal{C}, \mathcal{F})
$$
이고, $\mathcal{C}$의 임의의 대상 $U$에 대해서도
$$
H^i(U, \mathcal{F}_{ab})
=
H^i(U, \mathcal{F}).
$$
여기서 왼쪽은 $\textit{Ab}(\mathcal{C})$에서 계산한 코호몰로지이고,
오른쪽은 $\textit{Mod}(\mathcal{O})$에서 계산한 코호몰로지이다.
\end{lemma}

\begin{proof}
Derived Categories, Lemma \ref{derived-lemma-higher-derived-functors}에
의해 $\delta$-함자
$(\mathcal{F} \mapsto H^p(U, \mathcal{F}))_{p \geq 0}$는
보편적이다. 함자
$\textit{Mod}(\mathcal{O}) \to \textit{Ab}(\mathcal{C})$,
$\mathcal{F} \mapsto \mathcal{F}_{ab}$는 완전하다. 따라서
$(\mathcal{F} \mapsto H^p(U, \mathcal{F}_{ab}))_{p \geq 0}$
도 $\delta$-함자이다. 이제
$(\mathcal{F} \mapsto H^p(U, \mathcal{F}_{ab}))_{p \geq 0}$
도 보편적임을 보였다고 하자. 보편 $\delta$-함자의 유일성에 의해
이는 보조정리의 둘째 명제를 함의한다. Homology,
Lemma \ref{homology-lemma-uniqueness-universal-delta-functor}를 참조하라.
$\textit{Mod}(\mathcal{O})$에는 단사 대상이 충분히 있으므로,
$\textit{Mod}(\mathcal{O})$의 임의의 단사 대상 $\mathcal{I}$에 대해
$H^i(U, \mathcal{I}_{ab}) = 0$임을 보이면 충분하다. Homology,
Lemma \ref{homology-lemma-efface-implies-universal}를 참조하라.

\medskip\noindent
$\mathcal{I}$를 $\textit{Mod}(\mathcal{O})$의 단사 대상이라 하자.
$\mathcal{F} = \mathcal{I}$, $\mathcal{B} = \mathcal{C}$ 및
$\text{Cov} = \text{Cov}_\mathcal{C}$로 놓고
Lemma \ref{sites-cohomology-lemma-cech-vanish-collection}을 적용한다.
그 보조정리의 가정 (3)은
Lemma \ref{sites-cohomology-lemma-injective-module-trivial-cech}에 의해 성립한다.
따라서 $\mathcal{C}$의 모든 대상 $U$에 대해
$H^i(U, \mathcal{I}_{ab}) = 0$이다.

\medskip\noindent
$\mathcal{C}$에 끝 대상이 있으면 이 사실은 첫째 등식도 함의한다.
끝 대상이 없으면 Sites, Lemma \ref{sites-lemma-topos-good-site}에 의해
환 달린 토포스 $(\Sh(\mathcal{C}), \mathcal{O})$는 기저 사이트에
끝 대상이 있는 환 달린 토포스와 동치이다. 따라서 보조정리가 따른다.
\end{proof}

\begin{lemma}
\label{sites-cohomology-lemma-cohomology-products}
$\mathcal{C}$를 사이트라 하고 $I$를 집합이라 하자. 각 $i \in I$에
대해 $\mathcal{F}_i$를 $\mathcal{C}$ 위의 아벨 군층이라 하자.
$U \in \Ob(\mathcal{C})$라 하자. 정준 사상
$$
H^p(U, \prod\nolimits_{i \in I} \mathcal{F}_i)
\longrightarrow
\prod\nolimits_{i \in I} H^p(U, \mathcal{F}_i)
$$
은 $p = 0$일 때 동형이고 $p = 1$일 때 단사이다.
\end{lemma}

\begin{proof}
$p = 0$에 관한 명제는 층들의 곱이 그 기저 준층들의 곱과 같으므로
성립한다. Sites, Lemma \ref{sites-lemma-limit-sheaf}를 참조하라.
$p = 1$인 경우를 증명하자. $\mathcal{F} = \prod \mathcal{F}_i$로 놓자.
$\xi \in H^1(U, \mathcal{F})$가
$\prod H^1(U, \mathcal{F}_i)$에서 영으로 사상된다고 하자.
코호몰로지의 국소성, 즉
Lemma \ref{sites-cohomology-lemma-kill-cohomology-class-on-covering}에 의해 모든 $j$에
대해 $\xi|_{U_j} = 0$이 되는 덮개
$\mathcal{U} = \{U_j \to U\}$가 존재한다.
Lemma \ref{sites-cohomology-lemma-cech-h1}에 의해 이는 $\xi$가 어떤 원소
$\check \xi \in \check H^1(\mathcal{U}, \mathcal{F})$에서
온다는 뜻이다. 모든 $i$에 대해 사상
$\check H^1(\mathcal{U}, \mathcal{F}_i) \to H^1(U, \mathcal{F}_i)$
은 단사이고(Lemma \ref{sites-cohomology-lemma-cech-h1}), $\xi$의 상은
$\prod H^1(U, \mathcal{F}_i)$에서 영이므로
$\check H^1(\mathcal{U}, \mathcal{F}_i)$에서 그 상
$\check \xi_i = 0$이다. 한편
$\mathcal{F} = \prod \mathcal{F}_i$이므로
$\check{\mathcal{C}}^\bullet(\mathcal{U}, \mathcal{F})$는 복합체
$\check{\mathcal{C}}^\bullet(\mathcal{U}, \mathcal{F}_i)$들의 곱이다.
따라서 Homology,
Lemma \ref{homology-lemma-product-abelian-groups-exact}에 의해
원하는 대로 $\check \xi = 0$이다.
\end{proof}

\begin{lemma}
\label{sites-cohomology-lemma-restriction-along-monomorphism-surjective}
$(\mathcal{C}, \mathcal{O})$를 환 달린 사이트라 하자.
$a : U' \to U$를 $\mathcal{C}$의 단사 사상이라 하자.
그러면 임의의 단사 $\mathcal{O}$-가군 $\mathcal{I}$에 대해 제한 사상
$\mathcal{I}(U) \to \mathcal{I}(U')$은 전사이다.
\end{lemma}

\begin{proof}
$j : \mathcal{C}/U \to \mathcal{C}$와
$j' : \mathcal{C}/U' \to \mathcal{C}$를 국소화 사상이라 하자
(Modules on Sites, Section \ref{sites-modules-section-localize}).
$j_!$는 제한의 왼쪽 수반이므로 임의의 $\mathcal{O}$-가군층
$\mathcal{F}$에 대해
$$
\Hom_\mathcal{O}(j_!\mathcal{O}_U, \mathcal{F})
=
\Hom_{\mathcal{O}_U}(\mathcal{O}_U, \mathcal{F}|_U)
=
\mathcal{F}(U)
$$
이다. 마찬가지로 층 $j'_!\mathcal{O}_{U'}$는 함자
$\mathcal{F} \mapsto \mathcal{F}(U')$를 표현한다.
더 나아가 아래에서 $\mathcal{O}$-가군의 정준 사상
$$
j'_!\mathcal{O}_{U'} \longrightarrow j_!\mathcal{O}_U
$$
을 기술한다. 이 사상은 요네다 보조정리
(Categories, Lemma \ref{categories-lemma-yoneda})를 통해 제한 사상
$\mathcal{F}(U) \to \mathcal{F}(U')$에 대응한다.
표시된 가군 사상이 단사임을 보이면 충분하다. Homology,
Lemma \ref{homology-lemma-characterize-injectives}를 참조하라.

\medskip\noindent
우리의 사상을 구성하려면 $\mathcal{C}$의 대상 $V$에 각각
다음 $\mathcal{O}(V)$-가군을 대응시키는 준층들 사이의 사상을
구성하면 충분하다.
$$
\bigoplus\nolimits_{\varphi' \in \Mor_\mathcal{C}(V, U')} \mathcal{O}(V)
\quad\text{and}\quad
\bigoplus\nolimits_{\varphi \in \Mor_\mathcal{C}(V, U)} \mathcal{O}(V)
$$
Modules on Sites, Lemma \ref{sites-modules-lemma-extension-by-zero}를
참조하라. $\varphi'$에 대응하는 직합 성분을
$\varphi = a \circ \varphi'$에 대응하는 직합 성분으로
$\mathcal{O}(V)$의 항등 사상을 통해 보내는 사상을 취한다.
$a$가 단사 사상이므로 이 사상은 단사이다.
층화가 완전하므로 결론이 따른다.
\end{proof}






\section{완전 비순환층}
\label{sites-cohomology-section-limp}

\noindent
$(\mathcal{C}, \mathcal{O})$를 환 달린 사이트라 하자.
$K$를 $\mathcal{C}$ 위의 집합 준층이라 하자(아벨 군층과 구별하기 위해
여기서는 의도적으로 로마체 대문자를 사용한다).
$\mathcal{O}$-가군층 $\mathcal{F}$가 주어졌을 때
$$
\mathcal{F}(K) =
\Mor_{\textit{PSh}(\mathcal{C})}(K, \mathcal{F}) =
\Mor_{\Sh(\mathcal{C})}(K^\#, \mathcal{F})
$$
로 놓는다. 함자 $\mathcal{F} \mapsto \mathcal{F}(K)$는
$\textit{Mod}(\mathcal{O}) \to \textit{Ab}$인 왼쪽 완전 함자이므로
그 오른쪽 유도함자들이 존재한다. 이들을 $H^p(K, \mathcal{F})$로
나타내어 $H^0(K, \mathcal{F}) = \mathcal{F}(K)$가 되게 하자.

\medskip\noindent
몇 가지 관찰을 적어 두자.
\begin{enumerate}
\item $\mathcal{F}(K) = \mathcal{F}(K^\#)$이므로
$H^p(K, \mathcal{F}) = H^p(K^\#, \mathcal{F})$.
위 정의에서 $K$를 준층으로 허용하는 것은 순전히 표기상의 편의이다.
\item $\mathcal{C}$의 어떤 대상 $U$에 대해 $K = h_U$ 또는
$K = h_U^\#$라고 하자. 그러면
$H^p(K, \mathcal{F}) = H^p(U, \mathcal{F})$이다. 실제로
$\Mor_{\Sh(\mathcal{C})}(h_U^\#, \mathcal{F}) = \mathcal{F}(U)$이다.
Sites, Section \ref{sites-section-representable-sheaves}를 참조하라.
\item $\mathcal{O} = \mathbf{Z}$가 상수층이면 코호몰로지 군들은 함자
$H^p(K, - ) : \textit{Ab}(\mathcal{C}) \to \textit{Ab}$
이다. 이 경우
$\textit{Mod}(\mathcal{O}) = \textit{Ab}(\mathcal{C})$이기 때문이다.
\end{enumerate}
이미 증명한 결과 몇 가지를 이 언어로 옮겨 쓸 수 있다.

\begin{lemma}
\label{sites-cohomology-lemma-compute-cohomology-on-sheaf-sets}
$(\mathcal{C}, \mathcal{O})$를 환 달린 사이트라 하자.
$K$를 $\mathcal{C}$ 위의 집합 준층이라 하자.
$\mathcal{F}$를 $\mathcal{O}$-가군이라 하고 그 기저 아벨 군층을
$\mathcal{F}_{ab}$로 나타내자. 그러면
$H^p(K, \mathcal{F}) = H^p(K, \mathcal{F}_{ab})$이다.
\end{lemma}

\begin{proof}
$K$를 그 층화로 바꾸어 $K$가 층이라고 가정해도 된다.
$H^p(K, \mathcal{F})$와 $H^p(K, \mathcal{F}_{ab})$는 모두 기저
사이트가 아니라 토포스에만 의존한다. 따라서 Sites,
Lemma \ref{sites-lemma-topos-good-site}에 의해 $\mathcal{C}$를
``더 큰'' 사이트로 바꾸어, $\mathcal{C}$의 어떤 대상 $U$에 대해
$K = h_U$가 되게 할 수 있다. 이 경우 결과는
Lemma \ref{sites-cohomology-lemma-cohomology-modules-abelian-agree}에서 따른다.
\end{proof}

\begin{lemma}
\label{sites-cohomology-lemma-cech-to-cohomology-sheaf-sets}
$\mathcal{C}$를 사이트라 하자. $K' \to K$를 그 층화가 전사인
$\mathcal{C}$ 위의 집합 준층 사상이라 하자.
$K'_p = K' \times_K \ldots \times_K K'$ ($p + 1$개 인자)로 놓자.
모든 아벨 군층 $\mathcal{F}$에 대해
$E_1^{p, q} = H^q(K'_p, \mathcal{F})$이고
$H^{p + q}(K, \mathcal{F})$로 수렴하는 스펙트럼열이 존재한다.
\end{lemma}

\begin{proof}
층화는 완전하므로 $(K_p')^\#$는
$(K')^\# \times_{K^\#} \ldots \times_{K^\#} (K')^\#$
($p + 1$개 인자)와 같다. 따라서 $K$와 $K'$을 각각의 층화로 바꾸어
$K \to K'$가 층들의 전사 사상이라고 가정해도 된다.
Sites, Lemma \ref{sites-lemma-topos-good-site}에서와 같이
$\mathcal{C}$를 ``더 큰'' 사이트로 바꾸면 $K, K'$이
$\mathcal{C}$의 대상이고 $\mathcal{U} = \{K' \to K\}$가
덮개라고 가정할 수 있다. 그러면
Lemma \ref{sites-cohomology-lemma-cech-spectral-sequence}의 체흐-코호몰로지
스펙트럼열을 얻으며, 그 $E_1$ 쪽은 보조정리의 명제에 적힌 것과 같다.
\end{proof}

\begin{lemma}
\label{sites-cohomology-lemma-cohomology-on-sheaf-sets}
$\mathcal{C}$를 사이트라 하고 $K$를 $\mathcal{C}$ 위의 집합층이라 하자.
토포스의 사상
$j : \Sh(\mathcal{C}/K) \to \Sh(\mathcal{C})$를 생각하자.
Sites, Lemma \ref{sites-lemma-localize-topos-site}를 참조하라.
그러면 $j^{-1}$은 단사 대상을 보존하며, $\mathcal{C}$ 위의
임의의 아벨 군층 $\mathcal{F}$에 대해
$H^p(K, \mathcal{F}) = H^p(\mathcal{C}/K, j^{-1}\mathcal{F})$
이다.
\end{lemma}

\begin{proof}
Sites, Lemmas \ref{sites-lemma-localize-topos} and
\ref{sites-lemma-localize-topos-site}에 의해 토포스의 사상 $j$는
국소화와 동치이다. 따라서 결과는
Lemma \ref{sites-cohomology-lemma-cohomology-of-open}에서 따른다.
\end{proof}

\noindent
Lemma \ref{sites-cohomology-lemma-compute-cohomology-on-sheaf-sets}을 염두에 두면
다음 정의는 환 달린 사이트 위의 가군층에 대해서도
``올바른 정의''임을 알 수 있다.

\begin{definition}
\label{sites-cohomology-definition-limp}
$\mathcal{C}$를 사이트라 하자.
아벨 군층 $\mathcal{F}$가 {\it 완전 비순환}이라는 것은\footnote{이
용어가 \cite[Vbis, Proposition 1.3.10]{SGA4}에서 사용되기는 하지만,
아마도 표준 용어는 아니다.
\cite[V, Definition 4.1]{SGA4}에서는 이 성질을 ``flasque''라 부르지만,
위상 공간 위의 flasque 층에 대한 우리의 정의와 충돌하므로 이 말을
사용할 수 없다. 더 나은 제안이 있으면
\href{mailto:stacks.project@gmail.com}{stacks.project@gmail.com}
로 이메일을 보내 주기 바란다.} 모든 집합층 $K$와 모든 $p \geq 1$에
대해 $H^p(K, \mathcal{F}) = 0$이라는 뜻이다.
\end{definition}

\noindent
완전 비순환이라는 것은 내재적 성질, 즉 토포스의 동치 아래 보존되는
성질임이 분명하다. 완전 비순환층은 사이트의 모든 대상 위에서
고차 코호몰로지가 소멸하지만, 일반적으로 완전 비순환이라는 조건은
이보다 진정으로 더 강하다. 때때로 유용한 완전 비순환층의
특징짓기는 다음과 같다.

\begin{lemma}
\label{sites-cohomology-lemma-characterize-limp}
$\mathcal{C}$를 사이트라 하고 $\mathcal{F}$를 아벨 군층이라 하자.
다음이 성립한다고 하자.
\begin{enumerate}
\item $p > 0$ 및 $U \in \Ob(\mathcal{C})$에 대해
$H^p(U, \mathcal{F}) = 0$이다.
\item 모든 집합층의 전사 사상 $K' \to K$에 대해 확대 체흐 복합체
$$
0 \to H^0(K, \mathcal{F}) \to H^0(K', \mathcal{F}) \to
H^0(K' \times_K K', \mathcal{F}) \to \ldots
$$
는 완전하다.
\end{enumerate}
그러면 $\mathcal{F}$는 완전 비순환이다(역도 성립한다).
\end{lemma}

\begin{proof}
가정 (1)에 의해 모든 $p > 0$과 $\mathcal{C}$의 모든 대상 $U$에 대해
$H^p(h_U^\#, g^{-1}\mathcal{I}) = 0$이다.
$K = \coprod K_i$가 $\mathcal{C}$ 위의 집합층들의 쌍대곱이면
$H^p(K, g^{-1}\mathcal{I}) = \prod H^p(K_i, g^{-1}\mathcal{I})$임을
주목하라. 임의의 집합층 $K$에 대해 전사 사상
$$
K' = \coprod h_{U_i}^\# \longrightarrow K
$$
가 존재한다. Sites,
Lemma \ref{sites-lemma-sheaf-coequalizer-representable}를 참조하라.
따라서 다음을 얻는다. (*) 모든 집합층 $K$에 대해, 모든 $p > 0$에서
$H^p(K', \mathcal{F}) = 0$이 되는 집합층의 전사 사상
$K' \to K$가 존재한다. (*)와 조건 (2)가 $\mathcal{F}$의
완전 비순환성을 함의한다고 주장한다.
조건 (*)와 (2)는 기저 사이트가 아니라 토포스
$\Sh(\mathcal{C})$의 대상으로서의 $\mathcal{F}$에만 의존함을
주목하라. (증명의 나머지 부분에서는 성질 (1)을 사용하지 않는다.)

\medskip\noindent
$n \geq 0$에 대한 귀납법으로 (*)와 (2)가 다음 귀납 가설 $IH_n$을
함의함을 보이겠다. 모든 $0 < p \leq n$과 모든 집합층 $K$에 대해
$H^p(K, \mathcal{F}) = 0$이다. $IH_0$는 성립한다.
$IH_n$을 가정하자. 집합층 $K$를 택하고, 모든 $p > 0$에 대해
$H^p(K', \mathcal{F}) = 0$이 되는 전사 사상 $K' \to K$를 택하자.
다음을 만족하는 스펙트럼열이 있다.
$$
E_1^{p, q} = H^q(K'_p, \mathcal{F})
$$
이는 $H^{p + q}(K, \mathcal{F})$로 수렴한다.
Lemma \ref{sites-cohomology-lemma-cech-to-cohomology-sheaf-sets}를 참조하라.
$IH_n$에 의해 $0 < q \leq n$이면 $E_1^{p, q} = 0$이고,
가정 (2)에 의해 $p > 0$이면 $E_2^{p, 0} = 0$이다.
끝으로 $K'$의 선택에 의해 $H^q(K', \mathcal{F}) = 0$이므로
$q > 0$이면 $E_1^{0, q} = 0$이다. 따라서 $p + q = n + 1$인
모든 항 $E_2^{p, q}$가 영이므로
$H^{n + 1}(K, \mathcal{F}) = 0$이라는 결론을 얻는다.
\end{proof}







\section{Leray 스펙트럼열}
\label{sites-cohomology-section-leray}

\noindent
Leray 스펙트럼열의 존재를 증명하는 핵심은 다음 보조정리이다.

\begin{lemma}
\label{sites-cohomology-lemma-direct-image-injective-sheaf}
$f : (\Sh(\mathcal{C}), \mathcal{O}_\mathcal{C}) \to
(\Sh(\mathcal{D}), \mathcal{O}_\mathcal{D})$를 환 달린 토포스의
사상이라 하자. 그러면 $\textit{Mod}(\mathcal{O}_\mathcal{C})$의
임의의 단사 대상 $\mathcal{I}$에 대해 직상
$f_*\mathcal{I}$는 완전 비순환이다.
\end{lemma}

\begin{proof}
$K$를 $\mathcal{D}$ 위의 집합층이라 하자.
Modules on Sites, Lemma
\ref{sites-modules-lemma-morphism-ringed-topoi-comes-from-morphism-ringed-sites}
에 의해 $\mathcal{C}$와 $\mathcal{D}$를 ``더 큰'' 사이트로
바꾸어도 된다. 그 결과 $f$는 연속 함자
$u : \mathcal{D} \to \mathcal{C}$가 유도하는 환 달린 사이트의
사상에서 오고, $\mathcal{D}$의 어떤 대상 $V$에 대해
$K = h_V$가 되게 할 수 있다.

\medskip\noindent
따라서 $f$가 환 달린 사이트의 사상으로 주어졌을 때, 모든
$q > 0$과 $\mathcal{D}$의 모든 대상 $V$에 대해
$H^q(V, f_*\mathcal{I})$가 영임을 보이면 된다.
$\mathcal{V} = \{V_j \to V\}$를 $\mathcal{D}$의 임의의 덮개라 하자.
$u$가 연속이므로
$\mathcal{U} = \{u(V_j) \to u(V)\}$는 $\mathcal{C}$의 덮개이다.
그러면 체흐 복합체의 등식
$$
\check{\mathcal{C}}^\bullet(\mathcal{V}, f_*\mathcal{I})
=
\check{\mathcal{C}}^\bullet(\mathcal{U}, \mathcal{I})
$$
을 얻는다. 이는 $f_*$의 정의에서 따른다.
Lemma \ref{sites-cohomology-lemma-injective-module-trivial-cech}에 의해 이 복합체의
코호몰로지는 양의 차수에서 영이다. 이제
Lemma \ref{sites-cohomology-lemma-cech-vanish-collection}이 결론을 준다.
\end{proof}

\noindent
평탄 사상에 대해서는 함자 $f_*$가 단사 가군을 보존한다.
특히 함자
$f_* : \textit{Ab}(\mathcal{C}) \to \textit{Ab}(\mathcal{D})$
는 언제나 단사 아벨 군층을 단사 아벨 군층으로 보낸다.

\begin{lemma}
\label{sites-cohomology-lemma-pushforward-injective-flat}
$f : (\Sh(\mathcal{C}), \mathcal{O}_\mathcal{C}) \to
(\Sh(\mathcal{D}), \mathcal{O}_\mathcal{D})$를 환 달린 토포스의
사상이라 하자. $f$가 평탄이면, 임의의 단사
$\mathcal{O}_\mathcal{C}$-가군 $\mathcal{I}$에 대해
$f_*\mathcal{I}$는 단사 $\mathcal{O}_\mathcal{D}$-가군이다.
\end{lemma}

\begin{proof}
이 경우 함자 $f^*$는 완전하다. Modules on Sites,
Lemma \ref{sites-modules-lemma-flat-pullback-exact}를 참조하라.
따라서 결과는 Homology,
Lemma \ref{homology-lemma-adjoint-preserve-injectives}에서 따른다.
\end{proof}

\begin{lemma}
\label{sites-cohomology-lemma-limp-acyclic}
$(\Sh(\mathcal{C}), \mathcal{O}_\mathcal{C})$를 환 달린 토포스라 하자.
완전 비순환층은 다음 함자들에 대해 오른쪽 비순환이다.
\begin{enumerate}
\item $\mathcal{C}$의 임의의 대상 $U$에 대한 함자 $H^0(U, -)$,
\item 임의의 집합 준층 $K$에 대한 함자
$\mathcal{F} \mapsto \mathcal{F}(K)$,
\item 대역 절단 함자 $\Gamma(\mathcal{C}, -)$,
\item 환 달린 토포스의 임의의 사상
$f : (\Sh(\mathcal{C}), \mathcal{O}_\mathcal{C}) \to
(\Sh(\mathcal{D}), \mathcal{O}_\mathcal{D})$에 대한 함자 $f_*$.
\end{enumerate}
\end{lemma}

\begin{proof}
(2)는 완전 비순환층의 정의이다.
완전 비순환층의 정의 뒤의 논의에서 지적했듯이 (1)은 (2)의 귀결이다.
(3)은 $K = e$가 $\Sh(\mathcal{C})$의 끝 대상인 (2)의 특수한 경우이다.

\medskip\noindent
(4)를 증명하자. Modules on Sites, Lemma
\ref{sites-modules-lemma-morphism-ringed-topoi-comes-from-morphism-ringed-sites}
에 의해 $f$가 사이트의 사상으로 주어진다고 가정해도 된다.
이 경우 Lemma \ref{sites-cohomology-lemma-higher-direct-images}의 고차 직상에 대한
기술에 의해 완전 비순환층의 $R^if_*$, $i > 0$은 영이다.
\end{proof}

\begin{remark}
\label{sites-cohomology-remark-before-Leray}
위 결과들의 귀결로 Derived Categories,
Lemma \ref{derived-lemma-compose-derived-functors}를 여러 상황에
적용할 수 있다. 예를 들어 환 달린 토포스의 사상
$f : (\Sh(\mathcal{C}), \mathcal{O}_\mathcal{C}) \to
(\Sh(\mathcal{D}), \mathcal{O}_\mathcal{D})$가 주어졌을 때,
임의의 $\mathcal{O}_\mathcal{C}$-가군층 $\mathcal{F}$에 대해
$$
R\Gamma(\mathcal{D}, Rf_*\mathcal{F}) = R\Gamma(\mathcal{C}, \mathcal{F})
$$
이다. 실제로 단사 $\mathcal{O}_\mathcal{X}$-가군 $\mathcal{I}$에
대해 $\mathcal{O}_\mathcal{D}$-가군 $f_*\mathcal{I}$는
Lemma \ref{sites-cohomology-lemma-direct-image-injective-sheaf}에 의해 완전 비순환이고,
완전 비순환층은 Lemma \ref{sites-cohomology-lemma-limp-acyclic}에 의해
$\Gamma(\mathcal{D}, -)$에 대해 비순환이다.
\end{remark}

\begin{lemma}[Leray 스펙트럼열]
\label{sites-cohomology-lemma-Leray}
$f : (\Sh(\mathcal{C}), \mathcal{O}_\mathcal{C}) \to
(\Sh(\mathcal{D}), \mathcal{O}_\mathcal{D})$를 환 달린 토포스의
사상이라 하자. $\mathcal{F}^\bullet$을 아래로 유계인
$\mathcal{O}_\mathcal{C}$-가군 복합체라 하자. 다음 스펙트럼열이
존재한다.
$$
E_2^{p, q} = H^p(\mathcal{D}, R^qf_*(\mathcal{F}^\bullet))
$$
이는 $H^{p + q}(\mathcal{C}, \mathcal{F}^\bullet)$로 수렴한다.
\end{lemma}

\begin{proof}
이는 함자의 합성
$\Gamma(\mathcal{C}, -) = \Gamma(\mathcal{D}, -) \circ f_*$에서
나오는 Grothendieck 스펙트럼열
Derived Categories, Lemma \ref{derived-lemma-grothendieck-spectral-sequence}
이다. 이 보조정리의 가정들이 성립함은
Derived Categories, Lemma \ref{derived-lemma-grothendieck-spectral-sequence}
와 Lemmas \ref{sites-cohomology-lemma-direct-image-injective-sheaf} 및
\ref{sites-cohomology-lemma-limp-acyclic}를 참조하라.
\end{proof}

\begin{lemma}
\label{sites-cohomology-lemma-apply-Leray}
$f : (\Sh(\mathcal{C}), \mathcal{O}_\mathcal{C}) \to
(\Sh(\mathcal{D}), \mathcal{O}_\mathcal{D})$를 환 달린 토포스의
사상이라 하자. $\mathcal{F}$를 $\mathcal{O}_\mathcal{C}$-가군이라 하자.
\begin{enumerate}
\item $q > 0$에서 $R^qf_*\mathcal{F} = 0$이면, 모든 $p$에 대해
$H^p(\mathcal{C}, \mathcal{F}) = H^p(\mathcal{D}, f_*\mathcal{F})$이다.
\item 모든 $q$와 $p > 0$에 대해
$H^p(\mathcal{D}, R^qf_*\mathcal{F}) = 0$이면, 모든 $q$에 대해
$H^q(\mathcal{C}, \mathcal{F}) = H^0(\mathcal{D}, R^qf_*\mathcal{F})$이다.
\end{enumerate}
\end{lemma}

\begin{proof}
이들은 Leray 스펙트럼열의 수렴으로 결론을 얻게 하는 두 가지 간단한
조건이다. 이 사실들은 스펙트럼열을 사용하지 않고 직접 증명할 수도
있으며, 이는 층의 코호몰로지에 관한 좋은 연습문제이다.
\end{proof}

\begin{lemma}[상대 Leray 스펙트럼열]
\label{sites-cohomology-lemma-relative-Leray}
다음 두 환 달린 토포스의 사상을 생각하자.
$f : (\Sh(\mathcal{C}), \mathcal{O}_\mathcal{C}) \to
(\Sh(\mathcal{D}), \mathcal{O}_\mathcal{D})$
및
$g : (\Sh(\mathcal{D}), \mathcal{O}_\mathcal{D}) \to
(\Sh(\mathcal{E}), \mathcal{O}_\mathcal{E})$.
$\mathcal{F}$를 $\mathcal{O}_\mathcal{C}$-가군이라 하자.
다음을 만족하는 스펙트럼열이 존재한다.
$$
E_2^{p, q} = R^pg_*(R^qf_*\mathcal{F})
$$
이는 $R^{p + q}(g \circ f)_*\mathcal{F}$로 수렴한다.
이 스펙트럼열은 $\mathcal{F}$에 대해 함자적이며, 아래로 유계인
$\mathcal{O}_\mathcal{C}$-가군 복합체에 대한 판도 있다.
\end{lemma}

\begin{proof}
이는 함자의 합성에 대한 Grothendieck 스펙트럼열이다. Derived Categories,
Lemma \ref{derived-lemma-grothendieck-spectral-sequence}와
Lemmas \ref{sites-cohomology-lemma-direct-image-injective-sheaf} 및
\ref{sites-cohomology-lemma-limp-acyclic}를 참조하라.
\end{proof}







\section{기저변환 사상}
\label{sites-cohomology-section-base-change-map}

\noindent
이 절에서는 몇몇 경우의 기저변환 사상을 구성한다.
일반적인 경우는 Remark \ref{sites-cohomology-remark-base-change}에서 다룬다.
여기서는 환 달린 사이트의 평탄 사상에 의한 기저변환으로 제한하여
유도 당김을 사용하지 않는다. 결과를 서술하기 전에 유도 범주 위의
평탄 당김을 논의하자.
$g : (\Sh(\mathcal{C}), \mathcal{O}_\mathcal{C})
\to (\Sh(\mathcal{D}), \mathcal{O}_\mathcal{D})$
를 환 달린 토포스의 평탄 사상이라 하자. Modules on Sites,
Lemma \ref{sites-modules-lemma-flat-pullback-exact}에 의해 함자
$g^* : \textit{Mod}(\mathcal{O}_\mathcal{D}) \to
\textit{Mod}(\mathcal{O}_\mathcal{C})$는 완전하다.
따라서 유도함자
$$
g^* : D(\mathcal{O}_\mathcal{D}) \to D(\mathcal{O}_\mathcal{C})
$$
를 가지며, 이는 $D(\mathcal{O}_\mathcal{D})$의 주어진 대상을 나타내는
복합체를 단순히 당겨 계산한다. Derived Categories,
Lemma \ref{derived-lemma-right-derived-exact-functor}를 참조하라.
이 함자는 위로 또는 아래로 유계인 부분범주들을 보존한다.
따라서 위와 같이 이 함자를 $Lg^*$가 아니라 $g^*$로 나타낸다.

\begin{lemma}
\label{sites-cohomology-lemma-base-change-map-flat-case}
다음 환 달린 토포스의 가환 그림을 생각하자.
$$
\xymatrix{
(\Sh(\mathcal{C}'), \mathcal{O}_{\mathcal{C}'})
\ar[r]_{g'} \ar[d]_{f'} &
(\Sh(\mathcal{C}), \mathcal{O}_\mathcal{C}) \ar[d]^f \\
(\Sh(\mathcal{D}'), \mathcal{O}_{\mathcal{D}'})
\ar[r]^g &
(\Sh(\mathcal{D}), \mathcal{O}_\mathcal{D})
}
$$
$\mathcal{F}^\bullet$을 아래로 유계인
$\mathcal{O}_\mathcal{C}$-가군 복합체라 하자.
$g$와 $g'$이 모두 평탄하다고 가정하자. 그러면
$$
g^*Rf_*\mathcal{F}^\bullet
\longrightarrow
R(f')_*(g')^*\mathcal{F}^\bullet
$$
인 $D^{+}(\mathcal{O}_{\mathcal{D}'})$의 정준 기저변환 사상이 존재한다.
\end{lemma}

\begin{proof}
단사 분해 $\mathcal{F}^\bullet \to \mathcal{I}^\bullet$과
$(g')^*\mathcal{F}^\bullet \to \mathcal{J}^\bullet$을 택하자.
Lemma \ref{sites-cohomology-lemma-pushforward-injective-flat}에 의해
$(g')_*\mathcal{J}^\bullet$은
$R(g')_*(g')^*\mathcal{F}^\bullet$을 나타내는 단사 대상의 복합체이다.
따라서 Derived Categories, Lemmas \ref{derived-lemma-morphisms-lift}
및 \ref{derived-lemma-morphisms-equal-up-to-homotopy}에 의해
다음 그림의 화살표 $\beta$는 존재하며 호모토피를 제외하고 유일하다.
$$
\xymatrix{
(g')_*(g')^*\mathcal{F}^\bullet \ar[r] &
(g')_*\mathcal{J}^\bullet \\
\mathcal{F}^\bullet \ar[u]^{adjunction} \ar[r] &
\mathcal{I}^\bullet \ar[u]_\beta
}
$$
$\mathcal{D}$로 직상시키면
$$
f_*\beta :
f_*\mathcal{I}^\bullet
\longrightarrow
f_*(g')_*\mathcal{J}^\bullet
=
g_*(f')_*\mathcal{J}^\bullet
$$
를 얻는다. $g^*$와 $g_*$의 수반성에 의해 복합체의 사상
$g^*f_*\mathcal{I}^\bullet \to (f')_*\mathcal{J}^\bullet$을 얻는다.
전체 과정에서 선택한 것은 사상 $\beta$뿐이고 모든 작업이
복합체 수준에서 이루어졌으므로, 이 사상은 호모토피를 제외하고
유일함을 주목하라.
\end{proof}








\section{코호몰로지와 여극한}
\label{sites-cohomology-section-limits}

\noindent
$(\mathcal{C}, \mathcal{O})$를 환 달린 사이트라 하자.
$\mathcal{I} \to \textit{Mod}(\mathcal{O})$,
$i \mapsto \mathcal{F}_i$를 지표 범주 $\mathcal{I}$ 위의 도식이라 하자.
Categories, Section \ref{categories-section-limits}를 참조하라.
각 $i$에 대한 정준 사상
$\mathcal{F}_i \to \colim_i \mathcal{F}_i$는 코호몰로지 위의 사상을
유도한다. 따라서 모든 $p \geq 0$과 $\mathcal{C}$의 모든 대상 $U$에
대해 정준 사상
$$
\colim_i H^p(U, \mathcal{F}_i)
\longrightarrow
H^p(U, \colim_i \mathcal{F}_i)
$$
을 얻는다. 이 사상들은 일반적으로 동형이 아니며 $p = 0$일 때조차 그렇다.

\medskip\noindent
다음 보조정리는 Sites,
Lemma \ref{sites-lemma-directed-colimits-sections}의
코호몰로지에 대한 대응물이다.

\begin{lemma}
\label{sites-cohomology-lemma-colim-works-over-collection}
$\mathcal{C}$를 사이트라 하자.
$\text{Cov}_\mathcal{C}$를 $\mathcal{C}$의 덮개들의 집합이라 하자
(Sites, Definition \ref{sites-definition-site} 참조).
$\mathcal{B} \subset \Ob(\mathcal{C})$ 및
$\text{Cov} \subset \text{Cov}_\mathcal{C}$를 부분집합이라 하자.
다음을 가정하자.
\begin{enumerate}
\item 모든 $\mathcal{U} \in \text{Cov}$는
$\mathcal{U} = \{U_i \to U\}_{i \in I}$의 꼴이며, $I$는 유한하고
$U, U_i \in \mathcal{B}$이며 모든
$U_{i_0} \times_U \ldots \times_U U_{i_p} \in \mathcal{B}$이다.
\item 모든 $U \in \mathcal{B}$에 대해 $\text{Cov}$에 나타나는
$U$의 덮개들은 $U$의 덮개들의 공종계를 이룬다.
\end{enumerate}
그러면 사상
$$
\colim_i H^p(U, \mathcal{F}_i)
\longrightarrow
H^p(U, \colim_i \mathcal{F}_i)
$$
은 모든 $p \geq 0$, 모든 $U \in \mathcal{B}$ 및 모든 여과 도식
$\mathcal{I} \to \textit{Ab}(\mathcal{C})$에 대해 동형이다.
\end{lemma}

\begin{proof}
$p$에 대한 귀납법으로 보조정리를 증명하자.
(1)에서 덮개 $\mathcal{U} \in \text{Cov}$가 유한하도록 요구했으므로
$\mathcal{B}$의 모든 원소는 준콤팩트이다. 따라서 (2)와 (1)에 의해
임의의 $U \in \mathcal{B}$는 Sites,
Lemma \ref{sites-lemma-directed-colimits-sections} (4)의 가정을 만족한다.
그러므로 $p = 0$일 때 결과가 성립한다. 이제 $p$에 대해 보조정리가
성립한다고 가정하고 $p + 1$에 대해 증명하자.

\medskip\noindent
여과 도식
$\mathcal{F} : \mathcal{I} \to \textit{Ab}(\mathcal{C})$,
$i \mapsto \mathcal{F}_i$를 택하자.
$\textit{Ab}(\mathcal{C})$에는 함자적 단사 매입들이 있으므로
(Injectives, Theorem \ref{injectives-theorem-sheaves-injectives} 참조),
여과 도식의 사상
$\mathcal{F} \to \mathcal{I}$
을 찾을 수 있다. 여기서 각 $\mathcal{F}_i \to \mathcal{I}_i$는
아벨 군층에서 단사 아벨 군층으로 가는 단사 사상이다.
그 여핵을 $\mathcal{Q}_i$로 나타내면 짧은 완전열
$$
0 \to
\mathcal{F}_i \to
\mathcal{I}_i \to
\mathcal{Q}_i \to 0.
$$
을 얻는다. 층의 여극한은 준층 수준에서 취한 여극한의 층화이고,
층화는 완전하며, 아벨 군의 여과 여극한은 완전하므로
(Algebra, Lemma \ref{algebra-lemma-directed-colimit-exact} 참조),
열
$$
0 \to
\colim_i \mathcal{F}_i \to
\colim_i \mathcal{I}_i \to
\colim_i \mathcal{Q}_i \to 0.
$$
도 짧은 완전열이다. 모든 $U \in \mathcal{B}$와 모든 $q \geq 1$에
대해 $H^q(U, \colim_i \mathcal{I}_i) = 0$이라고 주장한다.
우선 이 주장을 받아들이고 다음 그림을 생각하자.
$$
\xymatrix{
\colim_i H^p(U, \mathcal{I}_i) \ar[d] \ar[r] &
\colim_i H^p(U, \mathcal{Q}_i) \ar[d] \ar[r] &
\colim_i H^{p + 1}(U, \mathcal{F}_i) \ar[d] \ar[r] &
0 \ar[d] \\
H^p(U, \colim_i \mathcal{I}_i) \ar[r] &
H^p(U, \colim_i \mathcal{Q}_i) \ar[r] &
H^{p + 1}(U, \colim_i \mathcal{F}_i) \ar[r] &
0
}
$$
오른쪽 아래의 영은 위 주장으로부터, 오른쪽 위의 영은 층
$\mathcal{I}_i$들이 단사라는 사실로부터 나온다.
Algebra, Lemma \ref{algebra-lemma-directed-colimit-exact}를 적용하면
위 행은 완전하다. 따라서 뱀 보조정리에 의해 $p + 1$에 대한 결과를 얻는다.

\medskip\noindent
이제 주장을 증명하면 된다.
Lemma \ref{sites-cohomology-lemma-cech-vanish-collection}을 사용할 것이다.
$p = 0$인 경우의 결과에 의해 $\mathcal{U} \in \text{Cov}$에 대해
$$
\check{\mathcal{C}}^\bullet(\mathcal{U}, \colim_i \mathcal{I}_i)
=
\colim_i \check{\mathcal{C}}^\bullet(\mathcal{U}, \mathcal{I}_i)
$$
이다. 모든 $U_{j_0} \times_U \ldots \times_U U_{j_p}$가
$\mathcal{B}$에 속하기 때문이다.
Lemma \ref{sites-cohomology-lemma-injective-trivial-cech}에 의해 체흐 복합체의 여극한
안의 각 복합체는 차수 $\geq 1$에서 비순환이다. 따라서 Algebra,
Lemma \ref{algebra-lemma-directed-colimit-exact}에 의해 체흐 복합체
$\check{\mathcal{C}}^\bullet(\mathcal{U}, \colim_i \mathcal{I}_i)$도
차수 $\geq 1$에서 비순환이다. 다시 말해 모든 $p \geq 1$에 대해
$\check{H}^p(\mathcal{U}, \colim_i \mathcal{I}_i) = 0$이다.
그러므로 Lemma \ref{sites-cohomology-lemma-cech-vanish-collection}의 가정들이
만족되어 주장이 따른다.
\end{proof}

\begin{lemma}
\label{sites-cohomology-lemma-colim-global}
$\mathcal{C}$를 사이트라 하고
$S \subset \Ob(\Sh(\mathcal{C}))$를 부분집합이라 하자.
$\Sh(\mathcal{C})$의 끝 대상을 $*$로 나타내자. 다음을 가정하자.
\begin{enumerate}
\item 어떤 $K \in S$에 대해 사상 $K \to *$는 전사이다.
\item $K \in S$이고 층의 전사 사상 $\mathcal{F} \to K$가 주어지면,
$K' \in S$와 사상 $K' \to \mathcal{F}$가 존재하여 합성
$K' \to K$가 전사이다.
\item $K, K' \in S$가 주어지면, $K'' \in S$인 전사 사상
$K'' \to K \times K'$이 존재한다.
\item $K, K' \in S$이고 $a, b : K \to K'$가 주어지면,
$K'' \in S$인 전사 사상
$K'' \to \text{Equalizer}(a, b)$가 존재한다.
\item 모든 $K \in S$는 준콤팩트이다
(Sites, Definition \ref{sites-definition-quasi-compact-topos}).
\end{enumerate}
그러면 모든 $p \geq 0$에 대해 사상
$$
\colim_\lambda H^p(\mathcal{C}, \mathcal{F}_\lambda)
\longrightarrow
H^p(\mathcal{C}, \colim_\lambda \mathcal{F}_\lambda)
$$
은 모든 여과 도식
$\Lambda \to \textit{Ab}(\mathcal{C})$,
$\lambda \mapsto \mathcal{F}_\lambda$에 대해 동형이다.
\end{lemma}

\begin{proof}
$p$에 대한 귀납법으로 증명하자. 기초 단계 $p = 0$은
Sites, Lemma \ref{sites-lemma-directed-colimits-global-sections} (4)에서
따른다. 가정들이 성립함을 확인하겠지만, 독자는 이 부분을 건너뛰어도 좋다.
$\mathcal{F} \to *$가 전사라고 하자. (1)과 같이
$K \in S$이고 $K \to *$가 전사인 대상을 택하자. 그러면
$\mathcal{F} \times K \to K$는 전사이다. (2)와 같이
$K' \in S$이고 $K' \to K$가 전사인
$K' \to \mathcal{F} \times K$를 택하자.
그러면 사상 $K' \to \mathcal{F}$가 존재하고 $K' \to *$는 전사이다.
따라서 Sites, Lemma \ref{sites-lemma-directed-colimits-global-sections}의
가정 (4)(a)가 성립한다.
Sites, Lemma \ref{sites-lemma-image-of-quasi-compact}와 가정 (3), (5)에
의해 모든 $K \in S$에 대해 $K \times K$는 준콤팩트이다.
따라서 Sites, Lemma \ref{sites-lemma-directed-colimits-global-sections}의
가정 (4)(b)가 성립한다. 이로써 기초 단계의 증명이 끝난다.

\medskip\noindent
귀납 단계를 보이자. (1)--(5)를 만족하는 집합
$S \subset \Ob(\Sh(\mathcal{C}))$를 찾을 수 있는 모든 토포스
$\Sh(\mathcal{C})$에 대해, $p \leq p_0$인 $H^p$에 대해 결과가
성립한다고 가정하자.
Lemma \ref{sites-cohomology-lemma-colim-works-over-collection}의 증명과 똑같이 논증하면,
단사 아벨 군층 $\mathcal{I}_\lambda$들의 여과 여극한
$\mathcal{I} = \colim \mathcal{I}_\lambda$이 주어졌을 때
$H^{p_0 + 1}(\mathcal{C}, \mathcal{I}) = 0$임을 보이면 충분하다.
(1)과 같이 $K \in S$이고 전사인 $K \to *$를 택하자.
$K$의 $n$중 자기곱을 $K^n$으로 나타내자. 스펙트럼열
$$
E_1^{p, q} = H^q(K^{p + 1}, \mathcal{I}) \Rightarrow
H^{p + q}(*, \mathcal{I}) = H^{p + q}(\mathcal{C}, \mathcal{I})
$$
을 생각하자. 이는 Lemma
\ref{sites-cohomology-lemma-cech-to-cohomology-sheaf-sets}의 스펙트럼열이다.
임의의 $\mathcal{C}$ 위의 아벨 군층에 대해
$H^q(K^{p + 1}, \mathcal{F}) = H^q(\mathcal{C}/K^{p + 1}, j^{-1}\mathcal{F})$,
임을 상기하라. Lemma \ref{sites-cohomology-lemma-cohomology-on-sheaf-sets}를 참조하라.
$j^{-1}\mathcal{I} = \colim j^{-1}\mathcal{I}_\lambda$
이다. $j^{-1}$이 여극한과 가환하기 때문이다.
제한 $j^{-1}\mathcal{I}_\lambda$들은
Lemma \ref{sites-cohomology-lemma-cohomology-of-open}에 의해
$\mathcal{C}/K^{p + 1}$ 위의 단사 아벨 군층이다.
아래에서 귀납 가설을 $\mathcal{C}/K^{p + 1}$에 적용할 수 있음을
보일 것이다. 따라서
$H^q(K^{p + 1}, \mathcal{I}) = \colim H^q(K^{p + 1}, \mathcal{I}_\lambda) = 0$
이 $q < p_0 + 1$에 대해 성립한다
($\mathcal{I}_\lambda$가 단사이므로 소멸한다). 따라서
$$
H^{p_0 + 1}(\mathcal{C}, \mathcal{I}) =
H^{p_0 + 1}\left(\ldots \to
H^0(K^{p_0}, \mathcal{I}) \to
H^0(K^{p_0 + 1}, \mathcal{I}) \to
H^0(K^{p_0 + 2}, \mathcal{I}) \to \ldots\right)
$$
이다. 다시 귀납 가설을 사용하면 오른쪽에 표시된 복합체는 다음 복합체들의
$\Lambda$에 걸친 여극한이다.
$$
\ldots \to
H^0(K^{p_0}, \mathcal{I}_\lambda) \to
H^0(K^{p_0 + 1}, \mathcal{I}_\lambda) \to
H^0(K^{p_0 + 2}, \mathcal{I}_\lambda) \to \ldots
$$
$\mathcal{I}_\lambda$가 단사 아벨 군층이므로 이 복합체들은 완전하다
(예를 들어 스펙트럼열에서 이를 알 수 있다).
아벨 군의 범주에서 여과 여극한은 완전하므로 원하는 소멸을 얻는다.

\medskip\noindent
모든 $n \geq 1$에 대해 귀납 가설을 사이트
$\mathcal{C}/K^n$에 적용할 수 있음을 보여야 한다.
$\Sh(\mathcal{C}/K^n) = \Sh(\mathcal{C})/K^n$임을 상기하라.
Sites, Lemma \ref{sites-lemma-localize-topos-site}를 참조하라.
따라서 $\Sh(\mathcal{C})/K^n$에서 작업해도 된다.
$S_n \subset \Ob(\Sh(\mathcal{C}/K^n)$을 $K' \to K^n$ 꼴의
대상들의 집합이라 하자. 각 성질을 차례로 확인하자.
\begin{enumerate}
\item (3)과 귀납법에 의해 $K' \in S$인 전사 사상
$K' \to K^n$이 존재한다. 그러면
$(K' \to K^n) \to (K^n \to K^n)$은
$\Sh(\mathcal{C})/K^n$의 전사 사상이고,
$K^n \to K^n$은 $\Sh(\mathcal{C})/K^n$의 끝 대상이다.
따라서 (1)은 $S_n$에 대해 성립한다.
\item $S_n$에 대한 성질 (2)는 $S$에 대한 (2)의 직접적인 귀결이다.
\item $a : K_1 \to K^n$ 및 $b : K_2 \to K^n$이 $S_n$에 속한다고 하자.
그러면 $(K_1 \to K^n) \times (K_2 \to K^n)$은
$\Sh(\mathcal{C})/K^n$의 대상
$K_1 \times_{K^n} K_2 \to K^n$이다.
부분층 $K_1 \times_{K^n} K_2 \subset K_1 \times K_2$는
$a \circ \text{pr}_1$과 $b \circ \text{pr}_2$의 등화자이다.
$a = (a_1, \ldots, a_n)$ 및 $b = (b_1, \ldots, b_n)$로 쓰자.
$K_3 \in S$이고 전사인 $K_3 \to K_1 \times K_2$를 택하자.
이는 $\mathcal{C}$에 대한 가정 (3)에 의해 가능하다. 전사 사상
$$
K_4
\longrightarrow
\text{Equalizer}(K_3 \to K_1 \times K_2 \xrightarrow{a_1, b_1} K)
$$
을 $K_4 \in S$가 되도록 택하자.
이는 $\mathcal{C}$에 대한 가정 (4)에 의해 가능하다. 전사 사상
$$
K_5
\longrightarrow
\text{Equalizer}(K_4 \to K_1 \times K_2 \xrightarrow{a_2, b_2} K)
$$
을 $K_5 \in S$가 되도록 택하자. 역시 가능하다.
이 과정을 계속하여 $K_{3 + n} \in S$인 전사 사상
$$
K_{3 + n}
\longrightarrow
\text{Equalizer}(K_{2 + n} \to K_1 \times K_2 \xrightarrow{a_n, b_n} K)
$$
을 얻는다. 구성에 의해
$K_{3 + n} \to K_1 \times_{K^n} K_2$
은 전사이다. 따라서 $(K_{3 + n} \to K^n)$은 $S_n$에 속하고
곱 $(K_1 \to K^n) \times (K_2 \to K^n)$ 위로 전사한다.
그러므로 (3)은 $S_n$에 대해 성립한다.
\item $S_n$에 대한 성질 (4)는 $S$에 대한 성질 (4)의 직접적인 귀결이다.
\item $S_n$에 대한 성질 (5)는 $S$에 대한 성질 (5)의 귀결이다.
실제로 $\Sh(\mathcal{C})/K^n$의 대상
$\mathcal{F} \to K^n$은 $\mathcal{F}$가
$\Sh(\mathcal{C})$의 준콤팩트 대상일 때 그리고 그때에만
$\Sh(\mathcal{C}/K^n)$의 준콤팩트 대상에 대응한다.
\end{enumerate}
이로써 보조정리의 증명이 끝난다.
\end{proof}

\begin{remark}
\label{sites-cohomology-remark-colim-global}
$\mathcal{C}$를 사이트라 하고
$\mathcal{B} \subset \Ob(\mathcal{C})$를 부분집합이라 하자.
$S \subset \Ob(\Sh(\mathcal{C}))$를 다음 꼴의 층 $K$들의 집합이라 하자.
$$
K = \coprod\nolimits_{i = 1, \ldots, n} h_{U_i}^\#
$$
여기서 $U_1, \ldots, U_n \in \mathcal{B}$이다.
이 집합은 언제 Lemma \ref{sites-cohomology-lemma-colim-global}의 가정들을 만족하는가?
한 가지 충분조건은 다음과 같다.
\begin{enumerate}
\item 어떤 $n \geq 0$ 및 $U_1, \ldots, U_n \in \mathcal{B}$에 대해 사상
$\coprod_{i = 1, \ldots, n} h_{U_i}^\# \to *$가 전사이다.
\item 모든 $U \in \mathcal{B}$의 덮개는
$U_i \in \mathcal{B}$인 꼴
$\{U_i \to U\}_{i = 1, \ldots, n}$의 덮개로 세분될 수 있다.
\item $U, U' \in \mathcal{B}$가 주어지면 $n \geq 0$,
$U_1, \ldots, U_n \in \mathcal{B}$ 및 사상
$U_i \to U$, $U_i \to U'$가 존재하여 사상
$\coprod_{i = 1, \ldots, n} h_{U_i}^\# \to
h_U^\# \times h_{U'}^\#$가 전사이다.
\item $\mathcal{C}$에서 $U, U' \in \mathcal{B}$인 사상
$a, b : U \to U'$가 주어지면
$U_1, \ldots, U_n \in \mathcal{B}$와 $a, b$를 등화하는 사상
$U_i \to U$가 존재하여
$\coprod_{i = 1, \ldots, n} h_{U_i}^\# \to
\text{Equalizer}(h_a^\#, h_b^\# : h_U^\# \to h_{U'}^\#)$
가 전사이다.
\end{enumerate}
상세한 검증은 생략한다. 다만 위의 (2)는 $\mathcal{B}$의 모든 원소가
준콤팩트임을 보장하고, 따라서 Sites,
Lemma \ref{sites-lemma-coproduct-quasi-compact}에 의해 모든
$K \in S$도 준콤팩트임을 언급해 둔다.
\end{remark}

\begin{lemma}
\label{sites-cohomology-lemma-higher-direct-image-colimit}
$f : \mathcal{C} \to \mathcal{D}$를 연속 함자
$u : \mathcal{D} \to \mathcal{C}$에 대응하는 사이트의 사상이라 하자.
$(\mathcal{F}_i, \varphi_{ii'})$를 $\mathcal{C}$ 위의 아벨 군층의
계라 하고 $\mathcal{F} = \colim \mathcal{F}_i$로 놓자.
$p \geq 0$을 정수라 하자.
$\mathcal{B}$를 다음을 만족하는 $V \in \Ob(\mathcal{D})$들의 집합이라 하자.
$H^p(u(V), \mathcal{F}) = \colim H^p(u(V), \mathcal{F}_i)$.
$\mathcal{D}$의 모든 대상이 $\mathcal{B}$의 원소들에 의한 덮개를 가지면
$R^pf_*\mathcal{F} = \colim R^pf_*\mathcal{F}_i$이다.
\end{lemma}

\begin{proof}
$R^pf_*\mathcal{F}$는 $V$를 $H^p(u(V), \mathcal{F})$로 보내는 준층
$\mathcal{G}$의 층화임을 상기하라. Lemma
\ref{sites-cohomology-lemma-higher-direct-images}를 참조하라. 마찬가지로
$R^pf_*\mathcal{F}_i$는 $V$를
$H^p(u(V), \mathcal{F}_i)$로 보내는 준층 $\mathcal{G}_i$의 층화이다.
층화는 층에서 준층으로 가는 포함의 왼쪽 수반임을 상기하라.
Sites, Section \ref{sites-section-sheafification}을 참조하라.
따라서 층화는 여극한과 가환한다. Categories,
Lemma \ref{categories-lemma-adjoint-exact}를 참조하라.
그러므로 준층의 사상(여극한은 준층의 범주에서 취한다)
$$
\colim \mathcal{G}_i \longrightarrow \mathcal{G}
$$
이 층화 위에 동형을 유도함을 보이면 충분하다. 이는 Sites,
Lemma \ref{sites-lemma-isomorphism-sheafifications}와
$\mathcal{B}$에 대한 가정에서 따른다.
\end{proof}

\begin{lemma}
\label{sites-cohomology-lemma-colim-sites-injective}
$\mathcal{I}$를 공여과 지표 범주라 하고
$(\mathcal{C}_i, f_a)$를 Sites,
Situation \ref{sites-situation-inverse-limit-sites}에서와 같은
$\mathcal{I}$ 위의 사이트의 역계라 하자. Sites,
Lemmas \ref{sites-lemma-colimit-sites} 및
\ref{sites-lemma-compute-pullback-to-limit}에서와 같이
$\mathcal{C} = \colim \mathcal{C}_i$로 놓자. 또한 다음이 주어졌다고 하자.
\begin{enumerate}
\item 모든 $i \in \Ob(\mathcal{I})$에 대해
$\mathcal{C}_i$ 위의 아벨 군층 $\mathcal{F}_i$,
\item $a : j \to i$에 대해 $\mathcal{C}_j$ 위의 아벨 군층 사상
$\varphi_a : f_a^{-1}\mathcal{F}_i \to \mathcal{F}_j$
\end{enumerate}
이며, $c = a \circ b$일 때마다
$\varphi_c = \varphi_b \circ f_b^{-1}\varphi_a$라고 하자.
그러면 계의 사상
$(\mathcal{F}_i, \varphi_a) \to (\mathcal{G}_i, \psi_a)$
이 존재하여 $\mathcal{F}_i \to \mathcal{G}_i$는 단사이고
$\mathcal{G}_i$는 단사 아벨 군층이다.
\end{lemma}

\begin{proof}
각 $i$에 대해 단사 사상 $\mathcal{F}_i \to \mathcal{A}_i$를 택하자.
여기서 $\mathcal{A}_i$는 $\mathcal{C}_i$ 위의 단사 아벨 군층이다.
그러면 사상들의 족
$$
\gamma_i :
\mathcal{F}_i
\longrightarrow
\prod\nolimits_{b : k \to i} f_{b, *}\mathcal{A}_k = \mathcal{G}_i
$$
을 생각할 수 있다. 그 성분 사상들은 사상
$f_b^{-1}\mathcal{F}_i \to \mathcal{F}_k \to \mathcal{A}_k$에
수반인 사상들이다. $\mathcal{I}$의 $a : j \to i$에 대해 정준 사상
$$
\psi_a : f_a^{-1}\mathcal{G}_i \to \mathcal{G}_j
$$
이 존재하며, 그 성분들은 $b : k \to j$에 대한 정준 사상
$f_b^{-1}f_{a \circ b, *}\mathcal{A}_k \to f_{b, *}\mathcal{A}_k$
이다. 따라서 아벨 군층의 계들의 단사 사상
$(\gamma_i) : (\mathcal{F}_i, \varphi_a) \to (\mathcal{G}_i, \psi_a)$
을 얻는다. $\mathcal{G}_i$는 $\mathcal{C}_i$ 위의 단사 아벨 군층이다.
Lemma \ref{sites-cohomology-lemma-pushforward-injective-flat}과 Homology,
Lemma \ref{homology-lemma-product-injectives}를 참조하라.
이로써 구성이 끝난다.
\end{proof}

\begin{lemma}
\label{sites-cohomology-lemma-colimit}
Lemma \ref{sites-cohomology-lemma-colim-sites-injective}의 상황에서
$\mathcal{F} = \colim f_i^{-1}\mathcal{F}_i$로 놓자.
$i \in \Ob(\mathcal{I})$ 및
$X_i \in \text{Ob}(\mathcal{C}_i)$라 하자. 그러면 모든 $p \geq 0$에 대해
$$
\colim_{a : j \to i} H^p(u_a(X_i), \mathcal{F}_j) =
H^p(u_i(X_i), \mathcal{F})
$$
이다.
\end{lemma}

\begin{proof}
$p = 0$인 경우는 Sites, Lemma \ref{sites-lemma-colimit}이다.

\medskip\noindent
Lemma \ref{sites-cohomology-lemma-colim-sites-injective}에서와 같은
$(\mathcal{F}_i, \varphi_a) \to (\mathcal{G}_i, \psi_a)$를 택하자.
Lemma \ref{sites-cohomology-lemma-colim-works-over-collection}의 증명과 똑같이 논증하면
$p > 0$에 대해
$H^p(X, \colim f_i^{-1}\mathcal{G}_i) = 0$임을 보이는 것으로 충분하다.

\medskip\noindent
$\mathcal{G} = \colim f_i^{-1}\mathcal{G}_i$로 놓자.
$\mathcal{C}$의 모든 대상 위에서 $\mathcal{G}$의 코호몰로지가
소멸함을 보이기 위해, $\mathcal{C}$의 임의의 덮개 $\mathcal{U}$에
대한 $\mathcal{G}$의 체흐 코호몰로지가 영임을 보이자
(Lemma \ref{sites-cohomology-lemma-cech-vanish-collection}).
덮개 $\mathcal{U}$는 어떤 $i$에 대한
$\mathcal{C}_i$의 덮개 $\mathcal{U}_i$에서 온다. $p = 0$인 경우에 의해
$$
\check{\mathcal{C}}^\bullet(\mathcal{U}, \mathcal{G}) =
\colim_{a : j \to i}
\check{\mathcal{C}}^\bullet(u_a(\mathcal{U}_i), \mathcal{G}_j)
$$
이다. 오른쪽은 Lemma \ref{sites-cohomology-lemma-injective-trivial-cech}에 의한
비순환 복합체들의 여과 여극한이므로 양의 차수에서 비순환이다.
Algebra, Lemma \ref{algebra-lemma-directed-colimit-exact}를 참조하라.
\end{proof}















\section{평탄 분해}
\label{sites-cohomology-section-flat}

\noindent
이 절에서는 환 달린 사이트와 환 달린 토포스의 맥락에서
Cohomology, Section \ref{cohomology-section-flat}의
논의를 다시 전개한다.

\begin{lemma}
\label{sites-cohomology-lemma-derived-tor-exact}
$(\mathcal{C}, \mathcal{O})$를 환 달린 사이트라 하자.
$\mathcal{G}^\bullet$를 $\mathcal{O}$-가군 복합체라 하자. 함자
$$
K(\textit{Mod}(\mathcal{O}))
\longrightarrow
K(\textit{Mod}(\mathcal{O})),
\quad
\mathcal{F}^\bullet \longmapsto
\text{Tot}(\mathcal{G}^\bullet \otimes_\mathcal{O} \mathcal{F}^\bullet)
$$
와
$$
K(\textit{Mod}(\mathcal{O}))
\longrightarrow
K(\textit{Mod}(\mathcal{O})),
\quad
\mathcal{F}^\bullet \longmapsto
\text{Tot}(\mathcal{F}^\bullet \otimes_\mathcal{O} \mathcal{G}^\bullet)
$$
는 삼각범주의 완전 함자이다.
\end{lemma}

\begin{proof}
이는 Derived Categories, Remark
\ref{derived-remark-double-complex-as-tensor-product-of}에서 따른다.
\end{proof}

\begin{definition}
\label{sites-cohomology-definition-K-flat}
$(\mathcal{C}, \mathcal{O})$를 환 달린 사이트라 하자.
$\mathcal{O}$-가군 복합체 $\mathcal{K}^\bullet$가
{\it K-평탄}이라는 것은 모든 비순환 $\mathcal{O}$-가군 복합체
$\mathcal{F}^\bullet$에 대해 복합체
$$
\text{Tot}(\mathcal{F}^\bullet \otimes_\mathcal{O} \mathcal{K}^\bullet)
$$
가 비순환이라는 뜻이다.
\end{definition}

\begin{lemma}
\label{sites-cohomology-lemma-K-flat-quasi-isomorphism}
$(\mathcal{C}, \mathcal{O})$를 환 달린 사이트라 하자.
$\mathcal{K}^\bullet$를 K-평탄 복합체라 하자. 그러면 함자
$$
K(\textit{Mod}(\mathcal{O}))
\longrightarrow
K(\textit{Mod}(\mathcal{O})), \quad
\mathcal{F}^\bullet
\longmapsto
\text{Tot}(\mathcal{F}^\bullet \otimes_\mathcal{O} \mathcal{K}^\bullet)
$$
는 유사동형을 유사동형으로 보낸다.
\end{lemma}

\begin{proof}
Lemma \ref{sites-cohomology-lemma-derived-tor-exact}과 유사동형은 원뿔이 비순환인
사상으로 특징지어진다는 사실에서 따른다.
\end{proof}

\begin{lemma}
\label{sites-cohomology-lemma-restriction-K-flat}
$(\mathcal{C}, \mathcal{O})$를 환 달린 사이트라 하고,
$U$를 $\mathcal{C}$의 객체라 하자.
$\mathcal{K}^\bullet$가 K-평탄 $\mathcal{O}$-가군 복합체이면
$\mathcal{K}^\bullet|_U$는 K-평탄 $\mathcal{O}_U$-가군 복합체이다.
\end{lemma}

\begin{proof}
$\mathcal{G}^\bullet$를 완전 $\mathcal{O}_U$-가군 복합체라 하자.
$j_{U!}$가 완전하고
(Modules on Sites, Lemma \ref{sites-modules-lemma-extension-by-zero-exact})
$\mathcal{K}^\bullet$가 K-평탄 $\mathcal{O}$-가군 복합체이므로, 복합체
$$
j_{U!}(\text{Tot}(\mathcal{G}^\bullet \otimes_{\mathcal{O}_U}
\mathcal{K}^\bullet|_U)) =
\text{Tot}(j_{U!}\mathcal{G}^\bullet \otimes_\mathcal{O} \mathcal{K}^\bullet)
$$
는 완전이다. 여기서 등식은 Modules on Sites, Lemma
\ref{sites-modules-lemma-j-shriek-and-tensor}와 $j_{U!}$가
(왼쪽 수반이므로) 직합과 가환한다는 사실에서 나온다.
Modules on Sites, Lemma
\ref{sites-modules-lemma-j-shriek-reflects-exactness}에 의해
$j_{U!}$가 완전성을 반영하므로 결론을 얻는다.
\end{proof}

\begin{lemma}
\label{sites-cohomology-lemma-tensor-product-K-flat}
$(\mathcal{C}, \mathcal{O})$를 환 달린 사이트라 하자.
$\mathcal{K}^\bullet$, $\mathcal{L}^\bullet$가 K-평탄
$\mathcal{O}$-가군 복합체이면
$\text{Tot}(\mathcal{K}^\bullet \otimes_\mathcal{O} \mathcal{L}^\bullet)$
도 K-평탄 $\mathcal{O}$-가군 복합체이다.
\end{lemma}

\begin{proof}
다음 동형과 정의에서 따른다.
$$
\text{Tot}(\mathcal{M}^\bullet \otimes_\mathcal{O}
\text{Tot}(\mathcal{K}^\bullet \otimes_\mathcal{O} \mathcal{L}^\bullet))
=
\text{Tot}(\text{Tot}(\mathcal{M}^\bullet \otimes_\mathcal{O}
\mathcal{K}^\bullet) \otimes_\mathcal{O} \mathcal{L}^\bullet)
$$
\end{proof}

\begin{lemma}
\label{sites-cohomology-lemma-K-flat-two-out-of-three}
$(\mathcal{C}, \mathcal{O})$를 환 달린 사이트라 하자.
$(\mathcal{K}_1^\bullet, \mathcal{K}_2^\bullet, \mathcal{K}_3^\bullet)$를
$K(\textit{Mod}(\mathcal{O}))$ 안의 특수 삼각형이라 하자.
세 $\mathcal{K}_i^\bullet$ 가운데 둘이 K-평탄이면 나머지 하나도
K-평탄이다.
\end{lemma}

\begin{proof}
Lemma \ref{sites-cohomology-lemma-derived-tor-exact}과
$K(\textit{Mod}(\mathcal{O}))$ 안의 특수 삼각형에서 세 대상 가운데
둘이 비순환이면 나머지 하나도 비순환이라는 사실에서 따른다.
\end{proof}

\begin{lemma}
\label{sites-cohomology-lemma-K-flat-two-out-of-three-ses}
$(\mathcal{C}, \mathcal{O})$를 환 달린 사이트라 하자.
$0 \to \mathcal{K}_1^\bullet \to \mathcal{K}_2^\bullet \to
\mathcal{K}_3^\bullet \to 0$를 복합체들의 짧은 완전열이라 하고,
$\mathcal{K}_3^\bullet$의 각 항은 평탄 $\mathcal{O}$-가군이라고 하자.
세 $\mathcal{K}_i^\bullet$ 가운데 둘이 K-평탄이면 나머지 하나도
K-평탄이다.
\end{lemma}

\begin{proof}
Modules on Sites, Lemma \ref{sites-modules-lemma-flat-tor-zero}에 의해
모든 복합체 $\mathcal{L}^\bullet$에 대해 복합체들의 짧은 완전열
$$
0 \to
\text{Tot}(\mathcal{L}^\bullet \otimes_\mathcal{O} \mathcal{K}_1^\bullet) \to
\text{Tot}(\mathcal{L}^\bullet \otimes_\mathcal{O} \mathcal{K}_2^\bullet) \to
\text{Tot}(\mathcal{L}^\bullet \otimes_\mathcal{O} \mathcal{K}_3^\bullet) \to 0
$$
을 얻는다. 따라서 보조정리는 코호몰로지 층들의 긴 완전열과
K-평탄 복합체의 정의에서 따른다.
\end{proof}

\begin{lemma}
\label{sites-cohomology-lemma-bounded-flat-K-flat}
$(\mathcal{C}, \mathcal{O})$를 환 달린 사이트라 하자. 평탄
$\mathcal{O}$-가군들로 이루어진 위로 유계인 복합체는 K-평탄이다.
\end{lemma}

\begin{proof}
$\mathcal{K}^\bullet$를 평탄 $\mathcal{O}$-가군들로 이루어진 위로
유계인 복합체라 하고, $\mathcal{L}^\bullet$를 비순환
$\mathcal{O}$-가군 복합체라 하자. 항별 여극한을 취하면
$\mathcal{L}^\bullet = \colim_m \tau_{\leq m}\mathcal{L}^\bullet$
이고, 따라서 항별로
$$
\text{Tot}(\mathcal{K}^\bullet \otimes_\mathcal{O} \mathcal{L}^\bullet)
=
\colim_m \text{Tot}(
\mathcal{K}^\bullet \otimes_\mathcal{O} \tau_{\leq m}\mathcal{L}^\bullet)
$$
이다. 그러므로 왼쪽 복합체가 비순환임을 보이려면 오른쪽의 각
복합체가 비순환임을 보이면 충분하다.
$\tau_{\leq m}\mathcal{L}^\bullet$가 비순환이므로
$\mathcal{L}^\bullet$가 위로 유계인 경우로 환원된다.
이 경우 Homology, Lemma \ref{homology-lemma-first-quadrant-ss}의
스펙트럼 열은
$$
{}'E_1^{p, q} = H^p(\mathcal{L}^\bullet \otimes_R \mathcal{K}^q)
$$
을 가지며, $\mathcal{K}^q$가 평탄이고 $\mathcal{L}^\bullet$가
비순환이므로 이는 영이다. 이로써 결론을 얻는다.
\end{proof}

\begin{lemma}
\label{sites-cohomology-lemma-colimit-K-flat}
$(\mathcal{C}, \mathcal{O})$를 환 달린 사이트라 하자.
$\mathcal{K}_1^\bullet \to \mathcal{K}_2^\bullet \to \ldots$를
K-평탄 복합체들의 계라 하자. 그러면
$\colim_i \mathcal{K}_i^\bullet$는 K-평탄이다.
\end{lemma}

\begin{proof}
항별 여극한을 취하므로 다음 등식은 자명하다.
$$
\colim_i \text{Tot}(
\mathcal{F}^\bullet \otimes_\mathcal{O} \mathcal{K}_i^\bullet)
=
\text{Tot}(\mathcal{F}^\bullet \otimes_\mathcal{O}
\colim_i \mathcal{K}_i^\bullet)
$$
따라서 보조정리는 여과 여극한이 완전하다는 사실에서 따른다.
\end{proof}

\begin{lemma}
\label{sites-cohomology-lemma-resolution-by-direct-sums-extensions-by-zero}
$(\mathcal{C}, \mathcal{O})$를 환 달린 사이트라 하자. 임의의
$\mathcal{O}$-가군 복합체 $\mathcal{G}^\bullet$에 대해
$\mathcal{O}$-가군 복합체들의 가환도식
$$
\xymatrix{
\mathcal{K}_1^\bullet \ar[d] \ar[r] &
\mathcal{K}_2^\bullet \ar[d] \ar[r] & \ldots \\
\tau_{\leq 1}\mathcal{G}^\bullet \ar[r] &
\tau_{\leq 2}\mathcal{G}^\bullet \ar[r] & \ldots
}
$$
이 존재하며 다음 성질을 만족한다. (1) 수직 화살표들은 유사동형이고
항별 전사이다. (2) 각 $\mathcal{K}_n^\bullet$는 위로 유계인
복합체이며, 그 항들은 $j_{U!}\mathcal{O}_U$ 꼴의
$\mathcal{O}$-가군들의 직합이다. (3) 사상
$\mathcal{K}_n^\bullet \to \mathcal{K}_{n + 1}^\bullet$는 항별 분할
단사이고, 그 여핵들은 $j_{U!}\mathcal{O}_U$ 꼴의
$\mathcal{O}$-가군들의 직합이다. 더욱이 사상
$\colim \mathcal{K}_n^\bullet \to \mathcal{G}^\bullet$는 유사동형이다.
\end{lemma}

\begin{proof}
도식의 존재와 성질 (1), (2), (3)은 Modules on Sites, Lemma
\ref{sites-modules-lemma-module-quotient-flat}과 Derived Categories, Lemma
\ref{derived-lemma-special-direct-system}에서 곧바로 따른다. 유도된 사상
$\colim \mathcal{K}_n^\bullet \to \mathcal{G}^\bullet$는 여과 여극한이
완전하므로 유사동형이다.
\end{proof}

\begin{lemma}
\label{sites-cohomology-lemma-K-flat-resolution}
$(\mathcal{C}, \mathcal{O})$를 환 달린 사이트라 하자. 임의의 복합체
$\mathcal{G}^\bullet$에 대해, 각 항이 평탄 $\mathcal{O}$-가군인
$K$-평탄 복합체 $\mathcal{K}^\bullet$와 항별 전사인 유사동형
$\mathcal{K}^\bullet \to \mathcal{G}^\bullet$가 존재한다.
\end{lemma}

\begin{proof}
Lemma \ref{sites-cohomology-lemma-resolution-by-direct-sums-extensions-by-zero}의 도식을
고르자. 각 복합체 $\mathcal{K}_n^\bullet$는 평탄 가군들로 이루어진
위로 유계인 복합체이다. Modules on Sites, Lemma
\ref{sites-modules-lemma-j-shriek-flat}을 보라. 따라서 Lemma
\ref{sites-cohomology-lemma-bounded-flat-K-flat}에 의해 $\mathcal{K}_n^\bullet$는
K-평탄이다. 그러므로 Lemma \ref{sites-cohomology-lemma-colimit-K-flat}에 의해
$\colim \mathcal{K}_n^\bullet$는 K-평탄이다. 구성상 유도된 사상
$\colim \mathcal{K}_n^\bullet \to \mathcal{G}^\bullet$는
유사동형이며 항별 전사이다. Lemma
\ref{sites-cohomology-lemma-resolution-by-direct-sums-extensions-by-zero}의 성질 (3)에
의하면 $\colim \mathcal{K}_n^m$은 평탄 가군들의 직합이므로 평탄이다.
이로써 마지막 주장도 증명된다.
\end{proof}

\begin{lemma}
\label{sites-cohomology-lemma-derived-tor-quasi-isomorphism-other-side}
$(\mathcal{C}, \mathcal{O})$를 환 달린 사이트라 하자.
$\alpha : \mathcal{P}^\bullet \to \mathcal{Q}^\bullet$를 K-평탄
$\mathcal{O}$-가군 복합체들의 유사동형이라 하자. 모든
$\mathcal{O}$-가군 복합체 $\mathcal{F}^\bullet$에 대해 유도된 사상
$$
\text{Tot}(\text{id}_{\mathcal{F}^\bullet} \otimes \alpha) :
\text{Tot}(\mathcal{F}^\bullet \otimes_\mathcal{O} \mathcal{P}^\bullet)
\longrightarrow
\text{Tot}(\mathcal{F}^\bullet \otimes_\mathcal{O} \mathcal{Q}^\bullet)
$$
은 유사동형이다.
\end{lemma}

\begin{proof}
Lemma \ref{sites-cohomology-lemma-K-flat-resolution}에 따라 $\mathcal{K}^\bullet$가
K-평탄인 유사동형
$\mathcal{K}^\bullet \to \mathcal{F}^\bullet$를 고르자. 다음
가환도식을 생각하자.
$$
\xymatrix{
\text{Tot}(\mathcal{K}^\bullet
\otimes_\mathcal{O} \mathcal{P}^\bullet) \ar[r] \ar[d] &
\text{Tot}(\mathcal{K}^\bullet
\otimes_\mathcal{O} \mathcal{Q}^\bullet) \ar[d] \\
\text{Tot}(\mathcal{F}^\bullet
\otimes_\mathcal{O} \mathcal{P}^\bullet) \ar[r] &
\text{Tot}(\mathcal{F}^\bullet
\otimes_\mathcal{O} \mathcal{Q}^\bullet)
}
$$
Lemma \ref{sites-cohomology-lemma-K-flat-quasi-isomorphism}에 의해 수직 화살표들과 위쪽
수평 화살표가 유사동형이므로 결과가 따른다.
\end{proof}

\noindent
$(\mathcal{C}, \mathcal{O})$를 환 달린 사이트라 하자.
$\mathcal{F}^\bullet$를 $D(\mathcal{O})$의 대상이라 하자. Lemma
\ref{sites-cohomology-lemma-K-flat-resolution}에 따라 K-평탄 분해
$\mathcal{K}^\bullet \to \mathcal{F}^\bullet$를 고르자. Lemma
\ref{sites-cohomology-lemma-derived-tor-exact}에 의해 삼각범주의 완전 함자
$$
K(\mathcal{O})
\longrightarrow
K(\mathcal{O}),
\quad
\mathcal{G}^\bullet
\longmapsto
\text{Tot}(\mathcal{G}^\bullet \otimes_\mathcal{O} \mathcal{K}^\bullet)
$$
를 얻는다. Lemma \ref{sites-cohomology-lemma-K-flat-quasi-isomorphism}에 의해 이 함자는
함자 $D(\mathcal{O}) \to D(\mathcal{O})$를 유도한다. 이는
$D(\mathcal{O})$가 $K(\mathcal{O})$를 유사동형들에 대해
국소화한 것이기 때문이다. $\mathcal{F}^\bullet$의 $K$-평탄 분해들의
범주는 공여과되어 있고 Lemma
\ref{sites-cohomology-lemma-derived-tor-quasi-isomorphism-other-side}가 성립하므로,
그 결과 얻는 함자는 동형을 무시하면 $\mathcal{K}^\bullet$의 선택에
의존하지 않는다.

\begin{definition}
\label{sites-cohomology-definition-derived-tor}
$(\mathcal{C}, \mathcal{O})$를 환 달린 사이트라 하자.
$\mathcal{F}^\bullet$를 $D(\mathcal{O})$의 대상이라 하자.
{\it 유도 텐서곱}
$$
- \otimes_\mathcal{O}^{\mathbf{L}} \mathcal{F}^\bullet :
D(\mathcal{O})
\longrightarrow
D(\mathcal{O})
$$
은 위에서 설명한 삼각범주의 완전 함자이다.
\end{definition}

\noindent
명시적인 구성에서 표준 동형
$$
\mathcal{F}^\bullet \otimes_\mathcal{O}^{\mathbf{L}} \mathcal{G}^\bullet
\cong
\mathcal{G}^\bullet \otimes_\mathcal{O}^{\mathbf{L}} \mathcal{F}^\bullet
$$
이 $D(\mathcal{O})$ 안의 $\mathcal{G}^\bullet$와
$\mathcal{F}^\bullet$에 대해 존재함이 명백하다. 여기서는 More on
Algebra, Section \ref{more-algebra-section-sign-rules}의 부호 규칙을
사용한다. 따라서
$\mathcal{F}^\bullet \otimes_\mathcal{O}^{\mathbf{L}} \mathcal{G}^\bullet$
라고 쓸 때에는 보통 어느 변수를 사용해 유도 텐서곱을 정의하는지
구별하지 않을 것이다.

\begin{definition}
\label{sites-cohomology-definition-tor}
$(\mathcal{C}, \mathcal{O})$를 환 달린 사이트라 하자.
$\mathcal{F}$, $\mathcal{G}$를 $\mathcal{O}$-가군이라 하자.
$\mathcal{F}$와 $\mathcal{G}$의 {\it Tor}들은 위에서 정의한 유도
텐서곱을 사용하여 다음 공식으로 정의한다.
$$
\text{Tor}_p^\mathcal{O}(\mathcal{F}, \mathcal{G}) =
H^{-p}(\mathcal{F} \otimes_\mathcal{O}^\mathbf{L} \mathcal{G})
$$
\end{definition}

\noindent
이 정의에 따르면 $\mathcal{O}$-가군들의 모든 짧은 완전열
$0 \to \mathcal{F}_1 \to \mathcal{F}_2 \to \mathcal{F}_3 \to 0$과
모든 $\mathcal{O}$-가군 $\mathcal{G}$에 대해 코호몰로지 긴 완전열
$$
\xymatrix{
\mathcal{F}_1 \otimes_\mathcal{O} \mathcal{G} \ar[r] &
\mathcal{F}_2 \otimes_\mathcal{O} \mathcal{G} \ar[r] &
\mathcal{F}_3 \otimes_\mathcal{O} \mathcal{G} \ar[r] & 0 \\
\text{Tor}_1^\mathcal{O}(\mathcal{F}_1, \mathcal{G}) \ar[r] &
\text{Tor}_1^\mathcal{O}(\mathcal{F}_2, \mathcal{G}) \ar[r] &
\text{Tor}_1^\mathcal{O}(\mathcal{F}_3, \mathcal{G}) \ar[ull]
}
$$
을 얻는다. 이를 이 상황에 딸린 $\text{Tor}$의 긴 완전열이라 한다.

\begin{lemma}
\label{sites-cohomology-lemma-flat-tor-zero}
$(\mathcal{C}, \mathcal{O})$를 환 달린 사이트라 하자.
$\mathcal{F}$를 $\mathcal{O}$-가군이라 하자. 다음은 서로 동치이다.
\begin{enumerate}
\item $\mathcal{F}$는 평탄 $\mathcal{O}$-가군이다.
\item 모든 $\mathcal{O}$-가군 $\mathcal{G}$에 대해
$\text{Tor}_1^\mathcal{O}(\mathcal{F}, \mathcal{G}) = 0$이다.
\end{enumerate}
\end{lemma}

\begin{proof}
만일 $\mathcal{F}$가 평탄이면 $\mathcal{F} \otimes_\mathcal{O} -$는
완전 함자이고 위성함자들은 영이다. 역으로 (2)가 성립한다고 가정하자.
$\mathcal{G} \to \mathcal{H}$가 단사이고 그 여핵이
$\mathcal{Q}$이면, $\text{Tor}$의 긴 완전열에 의해
$\mathcal{F} \otimes_\mathcal{O} \mathcal{G} \to
\mathcal{F} \otimes_\mathcal{O} \mathcal{H}$
의 핵은 $\text{Tor}_1^\mathcal{O}(\mathcal{F}, \mathcal{Q})$의 몫이다.
이는 가정에 의해 영이다. 따라서 $\mathcal{F}$는 평탄이다.
\end{proof}

\begin{lemma}
\label{sites-cohomology-lemma-K-flat-flat-acyclic}
$(\mathcal{C}, \mathcal{O})$를 환 달린 사이트라 하자. 각 항이 평탄인
K-평탄 비순환 복합체를 $\mathcal{K}^\bullet$라 하자. 그러면
$\mathcal{F} = \Ker(\mathcal{K}^n \to \mathcal{K}^{n + 1})$
는 평탄 $\mathcal{O}$-가군이다.
\end{lemma}

\begin{proof}
다음 열은 가군 $\mathcal{F}$의 평탄 분해이다.
$$
\ldots \to \mathcal{K}^{n - 2} \to \mathcal{K}^{n - 1} \to
\mathcal{F} \to 0
$$
평탄 가군들로 이루어진 위로 유계인 복합체는 K-평탄이므로
(Lemma \ref{sites-cohomology-lemma-bounded-flat-K-flat}), 이 분해를 사용하여 임의의
$\mathcal{O}$-가군 $\mathcal{G}$에 대한
$\text{Tor}_i(\mathcal{F}, \mathcal{G})$를 계산할 수 있다. 한편
$\mathcal{K}^\bullet \otimes_\mathcal{O}^\mathbf{L} \mathcal{G}$
는 $\mathcal{K}^\bullet$가 비순환이므로 $D(\mathcal{O})$에서 영이고,
다른 한편으로는
$\mathcal{K}^\bullet \otimes_\mathcal{O} \mathcal{G}$로 표현된다.
따라서 다음 열은 완전이다.
$$
\mathcal{K}^{n - 3} \otimes_\mathcal{O} \mathcal{G} \to
\mathcal{K}^{n - 2} \otimes_\mathcal{O} \mathcal{G} \to
\mathcal{K}^{n - 1} \otimes_\mathcal{O} \mathcal{G}
$$
그러므로 $\text{Tor}_1(\mathcal{F}, \mathcal{G}) = 0$이고,
Lemma \ref{sites-cohomology-lemma-flat-tor-zero}에 의해 결론을 얻는다.
\end{proof}

\begin{lemma}
\label{sites-cohomology-lemma-factor-through-K-flat}
$(\mathcal{C}, \mathcal{O})$를 환 달린 사이트라 하자.
$a : \mathcal{K}^\bullet \to \mathcal{L}^\bullet$를
$\mathcal{O}$-가군 복합체들의 사상이라 하자.
$\mathcal{K}^\bullet$가 K-평탄이면 복합체 $\mathcal{N}^\bullet$와
복합체 사상 $b : \mathcal{K}^\bullet \to \mathcal{N}^\bullet$,
$c : \mathcal{N}^\bullet \to \mathcal{L}^\bullet$가 존재하여 다음을
만족한다.
\begin{enumerate}
\item $\mathcal{N}^\bullet$는 K-평탄이다.
\item $c$는 유사동형이다.
\item $a$는 $c \circ b$와 호모토픽이다.
\end{enumerate}
$\mathcal{K}^\bullet$의 항들이 평탄하면 $\mathcal{N}^\bullet$에도
같은 성질이 성립하도록 $\mathcal{N}^\bullet$, $b$, $c$를 고를 수 있다.
\end{lemma}

\begin{proof}
호모토피 범주 $K(\textit{Mod}(\mathcal{O}))$가 삼각범주라는 사실을
사용할 것이다. Derived Categories, Proposition
\ref{derived-proposition-homotopy-category-triangulated}을 보라.
특수 삼각형
$\mathcal{K}^\bullet \to \mathcal{L}^\bullet \to
\mathcal{C}^\bullet \to \mathcal{K}^\bullet[1]$.
을 고르자. Lemma \ref{sites-cohomology-lemma-K-flat-resolution}에 따라, 항들이 평탄이고
$\mathcal{M}^\bullet$가 K-평탄인 유사동형
$\mathcal{M}^\bullet \to \mathcal{C}^\bullet$를 고르자.
삼각범주의 공리들에 의해 합성
$\mathcal{M}^\bullet \to \mathcal{C}^\bullet \to \mathcal{K}^\bullet[1]$
을 특수 삼각형
$\mathcal{K}^\bullet \to \mathcal{N}^\bullet \to
\mathcal{M}^\bullet \to \mathcal{K}^\bullet[1]$.
에 넣을 수 있다. Lemma \ref{sites-cohomology-lemma-K-flat-two-out-of-three}에 의해
$\mathcal{N}^\bullet$가 K-평탄임을 알 수 있다. 다시 삼각범주의 공리들을
사용하여 다음 특수 삼각형의 사상에 들어가는 사상
$\mathcal{N}^\bullet \to \mathcal{L}^\bullet$를 고를 수 있다.
$$
\xymatrix{
\mathcal{K}^\bullet \ar[r] \ar[d] &
\mathcal{N}^\bullet \ar[r] \ar[d] &
\mathcal{M}^\bullet \ar[r] \ar[d] &
\mathcal{K}^\bullet[1] \ar[d] \\
\mathcal{K}^\bullet \ar[r] &
\mathcal{L}^\bullet \ar[r] &
\mathcal{C}^\bullet \ar[r] &
\mathcal{K}^\bullet[1]
}
$$
화살표 셋 가운데 둘이 유사동형이므로 세 번째 화살표
$\mathcal{N}^\bullet \to \mathcal{L}^\bullet$도 이 특수 삼각형들에
딸린 코호몰로지 긴 완전열들에 의해 유사동형이다(원한다면 이 도식의
$D(\mathcal{O})$ 안의 상을 보고 Derived Categories, Lemma
\ref{derived-lemma-third-isomorphism-triangle}를 사용해도 된다).
이로써 (1), (2), (3)의 증명이 끝난다. 마지막 주장을 증명하기 위해
Derived Categories, Lemma
\ref{derived-lemma-improve-distinguished-triangle-homotopy}에 따라
$\mathcal{N}^n \cong \mathcal{M}^n \oplus \mathcal{K}^n$이 되도록
$\mathcal{N}^\bullet$를 고를 수 있다. 따라서
$\mathcal{K}^\bullet$의 항들이 평탄하면 원하는 평탄성을 얻는다.
\end{proof}












\section{유도 당김}
\label{sites-cohomology-section-derived-pullback}

\noindent
$f : (\Sh(\mathcal{C}), \mathcal{O}) \to
(\Sh(\mathcal{C}'), \mathcal{O}')$를 환 달린 토포스의 사상이라 하자.
K-평탄 분해를 사용하여 유도 당김 함자
$$
Lf^* : D(\mathcal{O}') \to D(\mathcal{O})
$$
를 정의할 수 있다.

\begin{lemma}
\label{sites-cohomology-lemma-pullback-K-flat}
$f : (\Sh(\mathcal{C}'), \mathcal{O}') \to (\Sh(\mathcal{C}), \mathcal{O})$를
환 달린 토포스의 사상이라 하자. $\mathcal{K}^\bullet$를 각 항이 평탄
$\mathcal{O}$-가군인 K-평탄 $\mathcal{O}$-가군 복합체라 하자. 그러면
$f^*\mathcal{K}^\bullet$는 각 항이 평탄 $\mathcal{O}'$-가군인 K-평탄
$\mathcal{O}'$-가군 복합체이다.
\end{lemma}

\begin{proof}
Modules on Sites, Lemma \ref{sites-modules-lemma-pullback-flat}에 의해
각 항 $f^*\mathcal{K}^n$은 평탄 $\mathcal{O}'$-가군이다. Lemma
\ref{sites-cohomology-lemma-resolution-by-direct-sums-extensions-by-zero}와 같은 도식
$$
\xymatrix{
\mathcal{K}_1^\bullet \ar[d] \ar[r] &
\mathcal{K}_2^\bullet \ar[d] \ar[r] & \ldots \\
\tau_{\leq 1}\mathcal{K}^\bullet \ar[r] &
\tau_{\leq 2}\mathcal{K}^\bullet \ar[r] & \ldots
}
$$
을 고르자. 더 언급하지 않고 그 보조정리에 적힌 모든 성질을 사용할
것이다. 각 $\mathcal{K}_n^\bullet$는 평탄 가군들로 이루어진 위로
유계인 복합체이다. Modules on Sites, Lemma
\ref{sites-modules-lemma-j-shriek-flat}을 보라. $\mathcal{M}^\bullet$를
정의하는 복합체들의 짧은 완전열
$$
0 \to \mathcal{M}^\bullet \to
\colim \mathcal{K}_n^\bullet \to
\mathcal{K}^\bullet \to 0
$$
을 생각하자. Lemmas \ref{sites-cohomology-lemma-bounded-flat-K-flat}과
\ref{sites-cohomology-lemma-colimit-K-flat}에 의해 복합체
$\colim \mathcal{K}_n^\bullet$는 K-평탄이고, Modules on Sites, Lemma
\ref{sites-modules-lemma-colimits-flat}에 의해 그 항들은 평탄이다.
Modules on Sites, Lemma \ref{sites-modules-lemma-flat-ses}에 의해
$\mathcal{M}^\bullet$의 항들이 평탄이고, Lemma
\ref{sites-cohomology-lemma-K-flat-two-out-of-three-ses}에 의해
$\mathcal{M}^\bullet$가 K-평탄이며, 코호몰로지 긴 완전열에 의해
$\mathcal{M}^\bullet$가 비순환이다(두 번째 화살표가 유사동형이기
때문이다). 당김
$f^*(\colim \mathcal{K}_n^\bullet) = \colim f^*\mathcal{K}_n^\bullet$
은 평탄 $\mathcal{O}'$-가군들로 이루어진 아래로 유계인 복합체들의
여극한이므로 K-평탄이다(위와 같은 보조정리들을 사용한다).
우리 짧은 완전열의 당김
$$
0 \to f^*\mathcal{M}^\bullet \to
f^*(\colim \mathcal{K}_n^\bullet) \to
f^*\mathcal{K}^\bullet \to 0
$$
은 Modules on Sites, Lemma \ref{sites-modules-lemma-pullback-ses}에 의해
복합체들의 짧은 완전열이다. 따라서 Lemma
\ref{sites-cohomology-lemma-K-flat-two-out-of-three-ses}에 의해
$f^*\mathcal{M}^\bullet$가 K-평탄임을 보이면 충분하다. 이는 다음
문단에서 다루는 경우로 환원된다.

\medskip\noindent
이제 $\mathcal{K}^\bullet$가 K-평탄이고 비순환이며 각 항이 평탄이라고
가정하자. 그러면 Lemma \ref{sites-cohomology-lemma-K-flat-flat-acyclic}에 의해
$\tau_{\leq n}\mathcal{K}^\bullet$의 모든 항은 평탄
$\mathcal{O}$-가군이다. 위와 같은 도식을 고르고, 그 도식에 대해 위에서
증명한 모든 성질을 사용할 것이다. 복합체 사상
$\mathcal{K}_n^\bullet \to \tau_{\leq n}\mathcal{K}^\bullet$
의 핵을 $\mathcal{M}_n^\bullet$로 나타내자. 그러면 복합체들의 짧은 완전열
$$
0 \to \mathcal{M}_n^\bullet \to \mathcal{K}_n^\bullet \to
\tau_{\leq n}\mathcal{K}^\bullet \to 0
$$
을 얻는다. Modules on Sites, Lemma \ref{sites-modules-lemma-flat-ses}에
의해 복합체 $\mathcal{M}_n^\bullet$의 항들은 평탄이다. 따라서 이 경우
$\mathcal{M} = \colim \mathcal{M}_n^\bullet$은 평탄 가군들로 이루어진
아래로 유계인 복합체들의 여과 여극한이다. 그러므로
$f^*\mathcal{M}^\bullet$는 K-평탄이고(위와 같은 논증), 결론을 얻는다.
\end{proof}

\begin{lemma}
\label{sites-cohomology-lemma-derived-base-change}
$f : (\Sh(\mathcal{C}), \mathcal{O}) \to (\Sh(\mathcal{C}'), \mathcal{O}')$를
환 달린 토포스의 사상이라 하자. 삼각범주의 완전 함자
$$
Lf^* : D(\mathcal{O}') \longrightarrow D(\mathcal{O})
$$
가 존재하여, 각 항이 평탄인 모든 K-평탄 복합체
$\mathcal{K}^\bullet$에 대해
$Lf^*\mathcal{K}^\bullet = f^*\mathcal{K}^\bullet$이고, 특히 평탄
$\mathcal{O}'$-가군들로 이루어진 모든 위로 유계인 복합체에 대해
이 등식이 성립한다.
\end{lemma}

\begin{proof}
Derived Categories, Section \ref{derived-section-derived-functors}에서
전개한 일반 이론을 사용한다.
$\mathcal{D} = K(\mathcal{O}')$, $\mathcal{D}' = D(\mathcal{O})$로 두자.
규칙 $F(\mathcal{G}^\bullet) = f^*\mathcal{G}^\bullet$로 정의되는
삼각범주의 완전 함자를 $F : \mathcal{D} \to \mathcal{D}'$로 쓰자.
또한 $\mathcal{D} = K(\mathcal{O}')$ 안의 유사동형들의 집합을 $S$라 하자.
그러면 Derived Categories, Situation
\ref{derived-situation-derived-functor}의 상황이므로 Derived Categories,
Definition \ref{derived-definition-right-derived-functor-defined}을 적용할
수 있다. $LF$가 모든 곳에서 정의됨을 주장한다. 이는 Derived
Categories, Lemma \ref{derived-lemma-find-existence-computes}에서
$\mathcal{P} \subset \Ob(\mathcal{D})$를 각 항이 평탄인 K-평탄 복합체
$\mathcal{K}^\bullet$들의 모임으로 두면 따른다. 실제로 (1)은 Lemma
\ref{sites-cohomology-lemma-K-flat-resolution}에서 따른다. (2)를 보이려면
$\mathcal{P}$의 원소 사이의 유사동형
$\mathcal{K}_1^\bullet  \to \mathcal{K}_2^\bullet$에 대해 사상
$f^*\mathcal{K}_1^\bullet  \to f^*\mathcal{K}_2^\bullet$가
유사동형임을 보이면 된다. 이 사상을
$$
f^{-1}\mathcal{K}_1^\bullet \otimes_{f^{-1}\mathcal{O}'} \mathcal{O}
\longrightarrow
f^{-1}\mathcal{K}_2^\bullet \otimes_{f^{-1}\mathcal{O}'} \mathcal{O}
$$
로 쓰자. 함자 $f^{-1}$은 완전하므로 사상
$f^{-1}\mathcal{K}_1^\bullet  \to f^{-1}\mathcal{K}_2^\bullet$는
유사동형이다. 환 달린 토포스의 사상
$(\Sh(\mathcal{C}), f^{-1}\mathcal{O}') \to
(\Sh(\mathcal{C}'), \mathcal{O}')$을 생각할 수 있으므로, Lemma
\ref{sites-cohomology-lemma-pullback-K-flat}에 의해 복합체
$f^{-1}\mathcal{K}_1^\bullet$와 $f^{-1}\mathcal{K}_2^\bullet$는 K-평탄
$f^{-1}\mathcal{O}'$-가군 복합체이다. 따라서 Lemma
\ref{sites-cohomology-lemma-derived-tor-quasi-isomorphism-other-side}에 의해 표시된 사상은
유사동형이다. 이로써 유도 함자
$$
LF :
D(\mathcal{O}') = S^{-1}\mathcal{D}
\longrightarrow
\mathcal{D}' = D(\mathcal{O})
$$
를 얻는다. Derived Categories, Equation
(\ref{derived-equation-everywhere})을 보라. 마지막으로 Derived Categories,
Lemma \ref{derived-lemma-find-existence-computes}에 의해
$LF(\mathcal{K}^\bullet) = F(\mathcal{K}^\bullet) = f^*\mathcal{K}^\bullet$
는 $\mathcal{K}^\bullet$가 $\mathcal{P}$에 속할 때 성립한다.
Lemma \ref{sites-cohomology-lemma-bounded-flat-K-flat}에 의해 평탄 가군들로 이루어진
위로 유계인 복합체들이 $\mathcal{P}$에 속함을 관찰하면 증명이 끝난다.
\end{proof}

\begin{lemma}
\label{sites-cohomology-lemma-derived-pullback-composition}
환 달린 토포스의 사상
$f : (\Sh(\mathcal{C}), \mathcal{O}_\mathcal{C}) \to
(\Sh(\mathcal{D}), \mathcal{O}_\mathcal{D})$
과
$g : (\Sh(\mathcal{D}), \mathcal{O}_\mathcal{D}) \to
(\Sh(\mathcal{E}), \mathcal{O}_\mathcal{E})$.
을 생각하자. 그러면 $D(\mathcal{O}_\mathcal{E}) \to
D(\mathcal{O}_\mathcal{C})$인 함자로서
$Lf^* \circ Lg^* = L(g \circ f)^*$이다.
\end{lemma}

\begin{proof}
$E$를 $D(\mathcal{O}_\mathcal{E})$의 대상이라 하자. Lemma
\ref{sites-cohomology-lemma-K-flat-resolution}에 따라 각 항이 평탄인 K-평탄 복합체
$\mathcal{K}^\bullet$로 $E$를 나타낼 수 있다. 구성에 의해 $Lg^*E$는
$g^*\mathcal{K}^\bullet$로 계산된다. Lemma
\ref{sites-cohomology-lemma-derived-base-change}를 보라. Lemma
\ref{sites-cohomology-lemma-pullback-K-flat}에 의해 복합체
$g^*\mathcal{K}^\bullet$는 각 항이 평탄인 K-평탄 복합체이다. 따라서
$Lf^*Lg^*E$는 $f^*g^*\mathcal{K}^\bullet$로 표현된다. 또한
$L(g \circ f)^*E$도
$(g \circ f)^*\mathcal{K}^\bullet = f^*g^*\mathcal{K}^\bullet$로
표현되므로 결론을 얻는다.
\end{proof}

\begin{lemma}
\label{sites-cohomology-lemma-pullback-tensor-product}
$f : (\Sh(\mathcal{C}), \mathcal{O}) \to (\Sh(\mathcal{D}), \mathcal{O}')$를
환 달린 토포스의 사상이라 하자. 표준 이함자적 동형
$$
Lf^*(
\mathcal{F}^\bullet \otimes_{\mathcal{O}'}^{\mathbf{L}} \mathcal{G}^\bullet
) =
Lf^*\mathcal{F}^\bullet 
\otimes_{\mathcal{O}}^{\mathbf{L}}
Lf^*\mathcal{G}^\bullet 
$$
이 모든
$\mathcal{F}^\bullet, \mathcal{G}^\bullet \in \Ob(D(\mathcal{O}'))$에
대해 존재한다.
\end{lemma}

\begin{proof}
Lemma \ref{sites-cohomology-lemma-derived-base-change}의 유도 당김 구성과 Lemma
\ref{sites-cohomology-lemma-K-flat-resolution}의 분해 존재성에 의해
$\mathcal{F}^\bullet$와 $\mathcal{G}^\bullet$를 각 항이 평탄인 K-평탄
$\mathcal{O}'$-가군 복합체로 바꾸어도 된다. 이 경우
$\mathcal{F}^\bullet \otimes_{\mathcal{O}'}^{\mathbf{L}} \mathcal{G}^\bullet$
는 이중 복합체
$\mathcal{F}^\bullet \otimes_{\mathcal{O}'} \mathcal{G}^\bullet$에
딸린 전체 복합체일 뿐이다. 복합체
$\text{Tot}(\mathcal{F}^\bullet \otimes_{\mathcal{O}'} \mathcal{G}^\bullet)$
는 Lemma \ref{sites-cohomology-lemma-tensor-product-K-flat}과 Modules on Sites, Lemma
\ref{sites-modules-lemma-tensor-flats}에 의해 각 항이 평탄인 K-평탄
복합체이다. 따라서 보조정리의 동형은 동형
$$
\text{Tot}(f^*\mathcal{F}^\bullet \otimes_{\mathcal{O}}
f^*\mathcal{G}^\bullet)
\longrightarrow
f^*\text{Tot}(\mathcal{F}^\bullet \otimes_{\mathcal{O}'} \mathcal{G}^\bullet)
$$
에서 오며, 그 각 성분은 Modules on Sites, Lemma
\ref{sites-modules-lemma-tensor-product-pullback}의 동형
$f^*\mathcal{F}^p \otimes_{\mathcal{O}} f^*\mathcal{G}^q \to
f^*(\mathcal{F}^p \otimes_{\mathcal{O}'} \mathcal{G}^q)$이다.
\end{proof}

\begin{lemma}
\label{sites-cohomology-lemma-variant-derived-pullback}
$f : (\Sh(\mathcal{C}), \mathcal{O}) \to (\Sh(\mathcal{C}'), \mathcal{O}')$를
환 달린 토포스의 사상이라 하자. 표준 이함자적 동형
$$
\mathcal{F}^\bullet
\otimes_\mathcal{O}^{\mathbf{L}}
Lf^*\mathcal{G}^\bullet
=
\mathcal{F}^\bullet 
\otimes_{f^{-1}\mathcal{O}'}^{\mathbf{L}}
f^{-1}\mathcal{G}^\bullet 
$$
이 $D(\mathcal{O})$의 $\mathcal{F}^\bullet$와
$D(\mathcal{O}')$의 $\mathcal{G}^\bullet$에 대해 존재한다.
\end{lemma}

\begin{proof}
$\mathcal{F}$를 $\mathcal{O}$-가군, $\mathcal{G}$를
$\mathcal{O}'$-가군이라 하자. 그러면
$\mathcal{F} \otimes_{\mathcal{O}} f^*\mathcal{G} =
\mathcal{F} \otimes_{f^{-1}\mathcal{O}'} f^{-1}\mathcal{G}$
이다. 이는
$f^*\mathcal{G} =
\mathcal{O} \otimes_{f^{-1}\mathcal{O}'} f^{-1}\mathcal{G}$.
이기 때문이다. 이 사실과 정의들에서 보조정리가 따른다.
\end{proof}











\begin{lemma}
\label{sites-cohomology-lemma-check-K-flat-stalks}
$(\mathcal{C}, \mathcal{O})$를 환 달린 사이트라 하고,
$\mathcal{K}^\bullet$를 $\mathcal{O}$-가군 복합체라 하자.
\begin{enumerate}
\item $\mathcal{K}^\bullet$가 K-평탄이면, 사이트 $\mathcal{C}$의 모든 점
$p$에 대해 $\mathcal{O}_p$-가군 복합체 $\mathcal{K}_p^\bullet$는 More
on Algebra, Definition \ref{more-algebra-definition-K-flat}의 의미에서
K-평탄이다.
\item $\mathcal{C}$가 충분한 점을 가지면 그 역도 참이다.
\end{enumerate}
\end{lemma}

\begin{proof}
(2)의 증명. $\mathcal{C}$가 충분한 점을 가지고 $\mathcal{C}$의 모든 점
$p$에 대해 $\mathcal{K}_p^\bullet$가 K-평탄이라고 하자. $\otimes$와
직합은 줄기를 취하는 것과 가환하고 완전성은 줄기에서 확인할 수
있으므로 $\mathcal{K}^\bullet$가 K-평탄임을 알 수 있다. Modules on
Sites, Lemma \ref{sites-modules-lemma-check-exactness-stalks}를 보라.

\medskip\noindent
(1)의 증명. $\mathcal{K}^\bullet$가 K-평탄이라고 가정하자. Lemma
\ref{sites-cohomology-lemma-K-flat-resolution}에 따라, 각 항이 평탄이고
$\mathcal{L}^\bullet$가 K-평탄인 유사동형
$a : \mathcal{L}^\bullet \to \mathcal{K}^\bullet$를 고르자.
$\mathcal{L}^\bullet$의 모든 당김은 K-평탄이다. Lemma
\ref{sites-cohomology-lemma-pullback-K-flat}을 보라. 특히 줄기
$\mathcal{L}_p^\bullet$는 K-평탄 $\mathcal{O}_p$-가군 복합체이다.
따라서 $a$의 원뿔 $C(a)$는 K-평탄인
(Lemma \ref{sites-cohomology-lemma-K-flat-two-out-of-three}) 비순환
$\mathcal{O}$-가군 복합체이고, $C(a)$의 줄기가 K-평탄임을 보이면
충분하다(More on Algebra, Lemma
\ref{more-algebra-lemma-K-flat-two-out-of-three}). 그러므로
$\mathcal{K}^\bullet$가 K-평탄이고 비순환인 경우를 가정해도 된다.

\medskip\noindent
$\mathcal{K}^\bullet$가 비순환이고 K-평탄이라고 가정하자. 계속하기
전에 Sites, Lemma \ref{sites-lemma-topos-good-site}에서처럼 사이트
$\mathcal{C}$를 다른 사이트로 바꾸어 $\mathcal{C}$가 모든 유한 극한을
갖도록 한다. 그러면 $p$의 근방들의 범주는 여과 범주이고
(Sites, Lemma \ref{sites-lemma-neighbourhoods-directed}), 층의 줄기를
정의하는 여극한도 여과 여극한이다. $M$을 유한 표시
$\mathcal{O}_p$-가군이라 하자. More on Algebra, Lemma
\ref{more-algebra-lemma-universally-acyclic-K-flat}에 의해
$\mathcal{K}^\bullet \otimes_{\mathcal{O}_p} M$이 비순환임을 보이면
충분하다. $U$가 $p$의 근방들을 달릴 때 $\mathcal{O}_p$는
$\mathcal{O}(U)$들의 여과 여극한이므로, $p$의 근방 $(U, x)$와 밑변환이
$\mathcal{O}_p$ 위에서 $M$인 유한 표시 $\mathcal{O}(U)$-가군 $M'$을
찾을 수 있다. Algebra, Lemma
\ref{algebra-lemma-colimit-category-fp-modules}를 보라. Lemma
\ref{sites-cohomology-lemma-restriction-K-flat}에 의해
$\mathcal{C}, \mathcal{O}, \mathcal{K}^\bullet$를
$\mathcal{C}/U, \mathcal{O}_U, \mathcal{K}^\bullet|_U$로 바꾸어도 된다.
따라서 $M \cong \mathcal{F}_p$인 $\mathcal{O}$-가군
$\mathcal{F}$가 존재한다고 가정해도 된다. $\mathcal{K}^\bullet$가
K-평탄이고 비순환이므로
$\mathcal{K}^\bullet \otimes_\mathcal{O} \mathcal{F}$는 비순환이다
(정의에 따라 유도 텐서곱을 계산하기 때문이다). 줄기를 취하는 함자는
완전이므로, 원하는 대로
$\mathcal{K}^\bullet \otimes_{\mathcal{O}_p} M$이 비순환임을 얻는다.
\end{proof}

\begin{lemma}
\label{sites-cohomology-lemma-pullback-K-flat-points}
$f : (\Sh(\mathcal{C}), \mathcal{O}) \to (\Sh(\mathcal{C}'), \mathcal{O}')$를
환 달린 토포스의 사상이라 하자. $\mathcal{C}$가 충분한 점을 가지면,
K-평탄 $\mathcal{O}'$-가군 복합체의 당김은 K-평탄
$\mathcal{O}$-가군 복합체이다.
\end{lemma}

\begin{proof}
이는 Lemma \ref{sites-cohomology-lemma-check-K-flat-stalks}, Modules on Sites, Lemma
\ref{sites-modules-lemma-pullback-stalk}, More on Algebra, Lemma
\ref{more-algebra-lemma-base-change-K-flat}에서 따른다.
\end{proof}

\begin{lemma}
\label{sites-cohomology-lemma-tensor-pull-compatibility}
$f : (\Sh(\mathcal{C}), \mathcal{O}_\mathcal{C}) \to
(\Sh(\mathcal{D}), \mathcal{O}_\mathcal{D})$를 환 달린 토포스의 사상이라
하자. $\mathcal{K}^\bullet$와 $\mathcal{M}^\bullet$를
$\mathcal{O}_\mathcal{D}$-가군 복합체라 하자. 도식
$$
\xymatrix{
Lf^*(\mathcal{K}^\bullet
\otimes_{\mathcal{O}_\mathcal{D}}^\mathbf{L}
\mathcal{M}^\bullet) \ar[r] \ar[d] &
Lf^*\text{Tot}(\mathcal{K}^\bullet
\otimes_{\mathcal{O}_\mathcal{D}}
\mathcal{M}^\bullet) \ar[d] \\
Lf^*\mathcal{K}^\bullet \otimes_{\mathcal{O}_\mathcal{C}}^\mathbf{L}
Lf^*\mathcal{M}^\bullet \ar[d] &
f^*\text{Tot}(\mathcal{K}^\bullet
\otimes_{\mathcal{O}_\mathcal{D}}
\mathcal{M}^\bullet) \ar[d] \\
f^*\mathcal{K}^\bullet \otimes_{\mathcal{O}_\mathcal{C}}^\mathbf{L}
f^*\mathcal{M}^\bullet \ar[r] &
\text{Tot}(f^*\mathcal{K}^\bullet \otimes_{\mathcal{O}_\mathcal{C}}
f^*\mathcal{M}^\bullet)
}
$$
은 가환한다.
\end{lemma}

\begin{proof}
각 항이 평탄인 K-평탄 분해가 존재한다는 사실
(Lemma \ref{sites-cohomology-lemma-K-flat-resolution}), 유도 당김이 그러한 복합체로
계산된다는 사실(Lemma \ref{sites-cohomology-lemma-derived-base-change}), 그리고 당김이
이 성질들을 보존한다는 사실(Lemma \ref{sites-cohomology-lemma-pullback-K-flat})을
사용할 것이다. 그러한 분해
$\mathcal{P}^\bullet \to \mathcal{K}^\bullet$와
$\mathcal{Q}^\bullet \to \mathcal{M}^\bullet$를 고르면 도식
$$
\xymatrix{
Lf^*\text{Tot}(\mathcal{P}^\bullet
\otimes_{\mathcal{O}_\mathcal{D}}
\mathcal{Q}^\bullet) \ar[r] \ar[d] &
Lf^*\text{Tot}(\mathcal{K}^\bullet
\otimes_{\mathcal{O}_\mathcal{D}}
\mathcal{M}^\bullet) \ar[d] \\
f^*\text{Tot}(\mathcal{P}^\bullet
\otimes_{\mathcal{O}_\mathcal{D}}
\mathcal{Q}^\bullet) \ar[d] \ar[r] &
f^*\text{Tot}(\mathcal{K}^\bullet
\otimes_{\mathcal{O}_\mathcal{D}}
\mathcal{M}^\bullet) \ar[d] \\
\text{Tot}(f^*\mathcal{P}^\bullet \otimes_{\mathcal{O}_\mathcal{C}}
f^*\mathcal{Q}^\bullet) \ar[r] &
\text{Tot}(f^*\mathcal{K}^\bullet \otimes_{\mathcal{O}_\mathcal{C}}
f^*\mathcal{M}^\bullet)
}
$$
이 가환함을 알 수 있다. 그러나 이제
$\mathcal{P}^\bullet$와 $\mathcal{Q}^\bullet$의 선택과 Lemma
\ref{sites-cohomology-lemma-tensor-product-K-flat}에 의해 이 도식의 왼쪽 변은
이 도식의 왼쪽 변이다.
\end{proof}









\section{비유계 복합체의 코호몰로지}
\label{sites-cohomology-section-unbounded}

\noindent
$(\mathcal{C}, \mathcal{O})$를 환 달린 사이트라 하자. 범주
$\textit{Mod}(\mathcal{O})$는 Grothendieck 아벨 범주이다. 즉 모든
여극한을 가지며, 여과 여극한은 완전하고, 다음 생성자를 갖는다.
$$
\bigoplus\nolimits_{U \in \Ob(\mathcal{C})} j_{U!}\mathcal{O}_U,
$$
Modules on Sites, Section \ref{sites-modules-section-kernels}과 Lemmas
\ref{sites-modules-lemma-j-shriek-flat},
\ref{sites-modules-lemma-module-quotient-flat}을 보라. Injectives, Theorem
\ref{injectives-theorem-K-injective-embedding-grothendieck}에 의해 모든
$\mathcal{O}$-가군 복합체 $\mathcal{F}^\bullet$에는 K-단사
$\mathcal{O}$-가군 복합체로 가는 단사 유사동형
$\mathcal{F}^\bullet \to \mathcal{I}^\bullet$가 존재하며, 이 매장은
$\mathcal{F}^\bullet$에 대해 함자적으로 고를 수 있다. Derived
Categories, Lemma \ref{derived-lemma-enough-K-injectives-implies}에 의해
다음이 따른다.
\begin{enumerate}
\item 삼각범주 $\mathcal{D}$로 가는 임의의 완전 함자
$F : K(\textit{Mod}(\mathcal{O})) \to \mathcal{D}$는 오른쪽 유도 함자
$RF : D(\mathcal{O}) \to \mathcal{D}$를 갖는다.
\item 아벨 범주 $\mathcal{A}$로 가는 임의의 가법 함자
$F : \textit{Mod}(\mathcal{O}) \to \mathcal{A}$에 대해, $F$가 유도하는
완전 함자 $F : K(\textit{Mod}(\mathcal{O})) \to K(\mathcal{A})$를
생각하면 오른쪽 유도 함자
$RF : D(\mathcal{O}) \to D(\mathcal{A})$를 얻는다.
\end{enumerate}
구성에 의해 $RF(\mathcal{F}^\bullet) = F(\mathcal{I}^\bullet)$이다.
여기서 $\mathcal{F}^\bullet \to \mathcal{I}^\bullet$는 위와 같다.

\medskip\noindent
위 사실의 몇 가지 예를 들면 다음과 같다.
\begin{enumerate}
\item 함자 $\Gamma(\mathcal{C}, -) : \textit{Mod}(\mathcal{O}) \to
\text{Mod}_{\Gamma(\mathcal{C}, \mathcal{O})}$는
$$
R\Gamma(\mathcal{C}, -) :
D(\mathcal{O})
\longrightarrow
D(\Gamma(\mathcal{C}, \mathcal{O}))
$$
를 준다. 코호몰로지에는 표기
$H^i(\mathcal{C}, K) = H^i(R\Gamma(\mathcal{C}, K))$를 사용할 것이다.
\item $\mathcal{C}$의 객체 $U$에 대해 함자
$\Gamma(U, -) : \textit{Mod}(\mathcal{O}) \to
\text{Mod}_{\Gamma(U, \mathcal{O})}$를 생각하자. 이는
$$
R\Gamma(U, -) : D(\mathcal{O}) \to D(\Gamma(U, \mathcal{O}))
$$
를 준다. 코호몰로지에는 표기
$H^i(U, K) = H^i(R\Gamma(U, K))$를 사용할 것이다.
\item 환 달린 토포스의 사상
$f : (\Sh(\mathcal{C}), \mathcal{O}) \to (\Sh(\mathcal{D}), \mathcal{O}')$에
대해 함자
$f_* : \textit{Mod}(\mathcal{O}) \to \textit{Mod}(\mathcal{O}')$를 생각하자.
이는 비유계 유도범주 위의 전체 직접상
$$
Rf_* : D(\mathcal{O}) \longrightarrow D(\mathcal{O}')
$$
을 준다.
\end{enumerate}

\begin{lemma}
\label{sites-cohomology-lemma-adjoint}
$f : (\Sh(\mathcal{C}), \mathcal{O}) \to (\Sh(\mathcal{D}), \mathcal{O}')$를
환 달린 토포스의 사상이라 하자. 위에서 정의한 함자 $Rf_*$와 Lemma
\ref{sites-cohomology-lemma-derived-base-change}에서 정의한 함자 $Lf^*$는 서로 수반이다.
즉
$$
\Hom_{D(\mathcal{O})}(Lf^*\mathcal{G}^\bullet, \mathcal{F}^\bullet)
=
\Hom_{D(\mathcal{O}')}(\mathcal{G}^\bullet, Rf_*\mathcal{F}^\bullet)
$$
이며, 이는 $\mathcal{F}^\bullet \in \Ob(D(\mathcal{O}))$와
$\mathcal{G}^\bullet \in \Ob(D(\mathcal{O}'))$에 대해 이함자적이다.
\end{lemma}

\begin{proof}
이는 $Rf_*$와 $Lf^*$가 존재한다는 사실에서 형식적으로 따른다.
Derived Categories, Lemma
\ref{derived-lemma-derived-adjoint-functors}를 보라.
\end{proof}

\begin{lemma}
\label{sites-cohomology-lemma-derived-pushforward-composition}
환 달린 토포스의 사상
$f : (\Sh(\mathcal{C}), \mathcal{O}_\mathcal{C}) \to
(\Sh(\mathcal{D}), \mathcal{O}_\mathcal{D})$
과
$g : (\Sh(\mathcal{D}), \mathcal{O}_\mathcal{D}) \to
(\Sh(\mathcal{E}), \mathcal{O}_\mathcal{E})$
을 생각하자. 그러면 $D(\mathcal{O}_\mathcal{C}) \to
D(\mathcal{O}_\mathcal{E})$인 함자로서
$Rg_* \circ Rf_* = R(g \circ f)_*$이다.
\end{lemma}

\begin{proof}
Lemma \ref{sites-cohomology-lemma-adjoint}에 의해 $Rg_* \circ Rf_*$는
$Lf^* \circ Lg^*$의 수반이다. Lemma
\ref{sites-cohomology-lemma-derived-pullback-composition}에 의해
$Lf^* \circ Lg^* = L(g \circ f)^*$이고, 따라서 수반함자의 유일성에
의해 $Rg_* \circ Rf_* = R(g \circ f)_*$이다.
\end{proof}

\begin{remark}
\label{sites-cohomology-remark-base-change}
비유계 유도 함자 $Lf^*$와 $Rf_*$의 구성으로 밑변환 사상을 가장
일반적인 상황에서 구성할 수 있다. 구체적으로 환 달린 토포스의 가환도식
$$
\xymatrix{
(\Sh(\mathcal{C}'), \mathcal{O}_{\mathcal{C}'})
\ar[r]_{g'} \ar[d]_{f'} &
(\Sh(\mathcal{C}), \mathcal{O}_\mathcal{C}) \ar[d]^f \\
(\Sh(\mathcal{D}'), \mathcal{O}_{\mathcal{D}'})
\ar[r]^g &
(\Sh(\mathcal{D}), \mathcal{O}_\mathcal{D})
}
$$
이 주어졌다고 하자. $K$를 $D(\mathcal{O}_\mathcal{C})$의 대상이라
하자. 그러면 $D(\mathcal{O}_{\mathcal{D}'})$ 안에 표준 밑변환 사상
$$
Lg^*Rf_*K \longrightarrow R(f')_*L(g')^*K
$$
이 존재한다. 이 사상은 사상
$L(f')^*Lg^*Rf_*K \to L(g')^*K$와 수반인 사상이다.
$L(f')^* \circ Lg^* = L(g')^* \circ Lf^*$이므로 이는 사상
$L(g')^*Lf^*Rf_*K \to L(g')^*K$와 같다. 이 사상으로는 수반 사상
$Lf^*Rf_*K \to K$에 $L(g')^*$를 적용한 것을 취할 수 있다.
\end{remark}

\begin{remark}
\label{sites-cohomology-remark-compose-base-change}
환 달린 토포스의 가환도식
$$
\xymatrix{
(\Sh(\mathcal{B}'), \mathcal{O}_{\mathcal{B}'})
\ar[r]_k \ar[d]_{f'} &
(\Sh(\mathcal{B}), \mathcal{O}_\mathcal{B}) \ar[d]^f \\
(\Sh(\mathcal{C}'), \mathcal{O}_{\mathcal{C}'})
\ar[r]^l \ar[d]_{g'} &
(\Sh(\mathcal{C}), \mathcal{O}_\mathcal{C}) \ar[d]^g \\
(\Sh(\mathcal{D}'), \mathcal{O}_{\mathcal{D}'})
\ar[r]^m &
(\Sh(\mathcal{D}), \mathcal{O}_\mathcal{D}) \\
}
$$
을 생각하자. 두 정사각형에 대한 Remark \ref{sites-cohomology-remark-base-change}의
밑변환 사상들을 합성하면 바깥 직사각형에 대한 밑변환 사상을 얻는다.
더 정확히 말해, 합성
\begin{align*}
Lm^* \circ R(g \circ f)_*
& =
Lm^* \circ Rg_* \circ Rf_* \\
& \to Rg'_* \circ Ll^* \circ Rf_* \\
& \to Rg'_* \circ Rf'_* \circ Lk^* \\
& = R(g' \circ f')_* \circ Lk^*
\end{align*}
은 그 직사각형에 대한 밑변환 사상이다. 확인은 생략한다.
\end{remark}

\begin{remark}
\label{sites-cohomology-remark-compose-base-change-horizontal}
환 달린 토포스의 가환도식
$$
\xymatrix{
(\Sh(\mathcal{C}''), \mathcal{O}_{\mathcal{C}''})
\ar[r]_{g'} \ar[d]_{f''} &
(\Sh(\mathcal{C}'), \mathcal{O}_{\mathcal{C}'})
\ar[r]_g \ar[d]_{f'} &
(\Sh(\mathcal{C}), \mathcal{O}_\mathcal{C}) \ar[d]^f \\
(\Sh(\mathcal{D}''), \mathcal{O}_{\mathcal{D}''})
\ar[r]^{h'} &
(\Sh(\mathcal{D}'), \mathcal{O}_{\mathcal{D}'})
\ar[r]^h &
(\Sh(\mathcal{D}), \mathcal{O}_\mathcal{D})
}
$$
을 생각하자. 두 정사각형에 대한 Remark \ref{sites-cohomology-remark-base-change}의
밑변환 사상들을 합성하면 바깥 직사각형에 대한 밑변환 사상을 얻는다.
더 정확히 말해, 합성
\begin{align*}
L(h \circ h')^* \circ Rf_*
& =
L(h')^* \circ Lh^* \circ Rf_* \\
& \to L(h')^* \circ Rf'_* \circ Lg^* \\
& \to Rf''_* \circ L(g')^* \circ Lg^* \\
& = Rf''_* \circ L(g \circ g')^*
\end{align*}
은 그 직사각형에 대한 밑변환 사상이다. 확인은 생략한다.
\end{remark}



\begin{lemma}
\label{sites-cohomology-lemma-adjoints-push-pull-compatibility}
$f : (\Sh(\mathcal{C}), \mathcal{O}_\mathcal{C}) \to
(\Sh(\mathcal{D}), \mathcal{O}_\mathcal{D})$를 환 달린 토포스의 사상이라
하고, $\mathcal{K}^\bullet$를 $\mathcal{O}_\mathcal{C}$-가군 복합체라
하자. 도식
$$
\xymatrix{
Lf^*f_*\mathcal{K}^\bullet \ar[r] \ar[d] &
f^*f_*\mathcal{K}^\bullet \ar[d] \\
Lf^*Rf_*\mathcal{K}^\bullet \ar[r] &
\mathcal{K}^\bullet
}
$$
은 복합체 위의 $Lf^* \to f^*$, 복합체 위의 $f_* \to Rf_*$, 그리고
수반 $Lf^* \circ Rf_* \to \text{id}$에서 오며
$D(\mathcal{O}_\mathcal{C})$ 안에서 가환한다.
\end{lemma}

\begin{proof}
K-평탄 분해와 K-단사 분해의 존재성을 사용할 것이다. Lemmas
\ref{sites-cohomology-lemma-K-flat-resolution}, \ref{sites-cohomology-lemma-derived-base-change},
\ref{sites-cohomology-lemma-pullback-K-flat}과 위의 논의를 보라. $\mathcal{I}^\bullet$가
$\mathcal{O}_\mathcal{C}$-가군 복합체로서 K-단사가 되는 유사동형
$\mathcal{K}^\bullet \to \mathcal{I}^\bullet$를 고르자. 각 항이 평탄이고
$\mathcal{Q}^\bullet$가 K-평탄 $\mathcal{O}_\mathcal{D}$-가군 복합체인
유사동형
$\mathcal{Q}^\bullet \to f_*\mathcal{I}^\bullet$를 고르자. 각 항이
평탄인 K-평탄 $\mathcal{O}_\mathcal{D}$-가군 복합체
$\mathcal{P}^\bullet$와 복합체 사상들의 도식
$$
\xymatrix{
\mathcal{P}^\bullet \ar[r] \ar[d] &
f_*\mathcal{K}^\bullet \ar[d] \\
\mathcal{Q}^\bullet \ar[r] & f_*\mathcal{I}^\bullet
}
$$
을 고를 수 있는데, 이는 호모토피까지 가환하고 위쪽 수평 화살표는
유사동형이다. 실제로 복합체들의 호모토피 범주에서 유사동형들이
곱셈계를 이루므로 먼저 어떤 복합체 $\mathcal{P}^\bullet$에 대해
그러한 도식을 고른 뒤, 각 항이 평탄인 K-평탄 복합체에 의한
$\mathcal{P}^\bullet$의 분해를 고를 수 있다. 당김을 취하면 복합체
사상들의 도식
$$
\xymatrix{
f^*\mathcal{P}^\bullet \ar[r] \ar[d] &
f^*f_*\mathcal{K}^\bullet \ar[d] \ar[r] &
\mathcal{K}^\bullet \ar[d] \\
f^*\mathcal{Q}^\bullet \ar[r] &
f^*f_*\mathcal{I}^\bullet \ar[r] &
\mathcal{I}^\bullet
}
$$
을 얻으며 이는 호모토피까지 가환한다. 바깥 직사각형이 보조정리의
진술이 참임을 보여 준다.
\end{proof}



\begin{remark}
\label{sites-cohomology-remark-cup-product}
$f : (\Sh(\mathcal{C}), \mathcal{O}_\mathcal{C}) \to
(\Sh(\mathcal{D}), \mathcal{O}_\mathcal{D})$를 환 달린 토포스의
사상이라 하자. $Lf^*$와 $Rf_*$의 수반성으로 상대 컵곱
$$
Rf_*K \otimes_{\mathcal{O}_\mathcal{D}}^\mathbf{L} Rf_*L
\longrightarrow
Rf_*(K \otimes_{\mathcal{O}_\mathcal{C}}^\mathbf{L} L)
$$
을 $D(\mathcal{O}_\mathcal{C})$의 모든 $K, L$에 대해
$D(\mathcal{O}_\mathcal{D})$ 안에서 구성할 수 있다. 이 사상은 사상
$Lf^*(Rf_*K \otimes_{\mathcal{O}_\mathcal{D}}^\mathbf{L} Rf_*L) \to
K \otimes_{\mathcal{O}_\mathcal{C}}^\mathbf{L} L$와 수반인 사상이다.
이 사상으로는 동형
$Lf^*(Rf_*K \otimes_{\mathcal{O}_\mathcal{D}}^\mathbf{L} Rf_*L) =
Lf^*Rf_*K \otimes_{\mathcal{O}_\mathcal{C}}^\mathbf{L} Lf^*Rf_*L$
(Lemma \ref{sites-cohomology-lemma-pullback-tensor-product})과 여단위
$Lf^* \circ Rf_* \to \text{id}$에서 오는 사상
$Lf^*Rf_*K \otimes_{\mathcal{O}_\mathcal{C}}^\mathbf{L} Lf^*Rf_*L
\to K \otimes_{\mathcal{O}_\mathcal{C}}^\mathbf{L} L$
의 합성을 취할 수 있다.
\end{remark}

\begin{lemma}
\label{sites-cohomology-lemma-torsion}
$\mathcal{C}$를 사이트라 하자. 꼬임 아벨층들로 이루어진 Serre 부분범주를
$\mathcal{A} \subset \textit{Ab}(\mathcal{C})$로 나타내자. 그러면 함자
$D(\mathcal{A}) \to D_\mathcal{A}(\mathcal{C})$는 동치이다.
\end{lemma}

\begin{proof}
핵심 관찰은 단사 아벨층 $\mathcal{I}$가 가분이라는 것이다. 실제로
$s \in \mathcal{I}(U)$가 국소 단면이면 $s$를 사상
$s : j_{U!}\mathbf{Z} \to \mathcal{I}$로 해석하고, 층들의 단사 사상
$n : j_{U!}\mathbf{Z} \to j_{U!}\mathbf{Z}$에 단사 대상의 정의 성질을
적용한다. 그러면 $ns' = s$인 $s' \in \mathcal{I}(U)$가 존재한다.

\medskip\noindent
층 $\mathcal{F}$의 꼬임 부분층을 $\mathcal{F}_{tor}$로 나타내자.
$\mathcal{I}^\bullet$가 단사 아벨층들의 복합체이고 그 코호몰로지 층들이
꼬임층이면
$$
\mathcal{I}^\bullet_{tor} \to \mathcal{I}^\bullet
$$
가 유사동형이라고 주장한다. 실제로 $\mathbf{Q}$가 $\mathbf{Z}$ 위에서
평탄하므로
$$
H^p(\mathcal{I}^\bullet) \otimes_\mathbf{Z} \mathbf{Q} =
H^p(\mathcal{I}^\bullet \otimes_\mathbf{Z} \mathbf{Q})
$$
이고, 코호몰로지 층들이 꼬임층이므로 이는 영이다. 위에서 보인
가분성에 의해
$\mathcal{I}^\bullet \to \mathcal{I}^\bullet \otimes_\mathbf{Z} \mathbf{Q}$
는 전사이고 그 핵은 $\mathcal{I}^\bullet_{tor}$이다. 이때 얻는 짧은
완전열에 딸린 코호몰로지 층들의 긴 완전열에서 주장이 따른다.

\medskip\noindent
보조정리를 증명하기 위해 오른쪽 수반
$T : D(\mathcal{C}) \to D(\mathcal{A})$를 구성할 것이다. 즉,
$D(\mathcal{C})$의 $K$가 주어지면 Injectives, Theorem
\ref{injectives-theorem-K-injective-embedding-grothendieck}에 따라
코호몰로지 층들이 단사인 K-단사 복합체 $\mathcal{I}^\bullet$로
$K$를 나타낼 수 있다. 이제 $T(K) = \mathcal{I}^\bullet_{tor}$로 둔다.
다시 말해 $T$는 꼬임 부분을 취하는 함자의 오른쪽 유도 함자이다.
함자 $T$는 $i : D(\mathcal{A}) \to D_\mathcal{A}(\mathcal{C})$의 오른쪽
수반이다. 이는 $\mathcal{F}^\bullet$가 꼬임층들의 복합체이면
$$
\Hom_{K(\mathcal{A})}(\mathcal{F}^\bullet, I^\bullet_{tor}) =
\Hom_{K(\textit{Ab}(\mathcal{C}))}(\mathcal{F}^\bullet, I^\bullet)
$$
임을 관찰하면 곧바로 따른다. 특히 $\mathcal{I}^\bullet_{tor}$는
$\mathcal{A}$의 K-단사 복합체이다. 몇 가지 세부는 생략한다. 의문이
있다면 더 일반적인 Derived Categories, Lemma
\ref{derived-lemma-derived-adjoint-functors}에서도 따른다. 위 주장은
$D(\mathcal{A})$의 $L$에 대해 $L = T(i(L))$이고,
$K$가 $D_\mathcal{A}(\mathcal{C})$에 속하면 $i(T(K)) = K$임을 준다.
Categories, Lemma \ref{categories-lemma-adjoint-fully-faithful}를 적용하면
결과가 따른다.
\end{proof}




\section{K-단사 복합체의 몇 가지 성질}
\label{sites-cohomology-section-properties-K-injective}

\noindent
$(\mathcal{C}, \mathcal{O})$를 환 달린 사이트라 하고, $U$를
$\mathcal{C}$의 객체라 하자. 이에 대응하는 국소화 사상을
$j : (\Sh(\mathcal{C}/U), \mathcal{O}_U) \to (\Sh(\mathcal{C}), \mathcal{O})$
로 나타내자. 당김 함자 $j^*$는 제한 함자에 불과하므로 완전하다.
따라서 유도 당김 $Lj^*$는 임의의 복합체를 단순히 제한하여 계산된다.
우리는 흔히 이에 대응하는 함자를 간단히
$$
D(\mathcal{O}) \to D(\mathcal{O}_U),
\quad
E \mapsto j^*E = E|_U
$$
로 나타낸다. 마찬가지로 영에 의한 연장
$j_! : \textit{Mod}(\mathcal{O}_U) \to \textit{Mod}(\mathcal{O})$ (see
Modules on Sites, Definition
\ref{sites-modules-definition-localize-ringed-site})
은 완전 함자이다
(Modules on Sites, Lemma \ref{sites-modules-lemma-extension-by-zero-exact}).
그러므로 이는 함자
$$
j_! : D(\mathcal{O}_U) \to D(\mathcal{O}), \quad
F \mapsto j_!F
$$
를 유도하며, 이는 대상 $F$를 나타내는 임의의 복합체에 $j_!$를 단순히
적용하여 계산한다.

\begin{lemma}
\label{sites-cohomology-lemma-restrict-K-injective-to-open}
$(\mathcal{C}, \mathcal{O})$를 환 달린 사이트라 하고, $U$를
$\mathcal{C}$의 객체라 하자. K-단사 $\mathcal{O}$-가군 복합체를
$\mathcal{C}/U$로 제한하면 K-단사 $\mathcal{O}_U$-가군 복합체가 된다.
\end{lemma}

\begin{proof}
이는 Derived Categories, Lemma
\ref{derived-lemma-adjoint-preserve-K-injectives}와 제한 함자가 완전한
왼쪽 수반 $j_!$를 갖는다는 사실에서 곧바로 따른다. 위의 논의를 보라.
\end{proof}

\begin{lemma}
\label{sites-cohomology-lemma-unbounded-cohomology-of-open}
$(\mathcal{C}, \mathcal{O})$를 환 달린 사이트라 하고,
$U \in \Ob(\mathcal{C})$라 하자. $D(\mathcal{O})$의 $K$에 대해
$H^p(U, K) = H^p(\mathcal{C}/U, K|_{\mathcal{C}/U})$.
\end{lemma}

\begin{proof}
$K$를 나타내는 K-단사 $\mathcal{O}$-가군 복합체를
$\mathcal{I}^\bullet$라 하자. 그러면 코호몰로지의 구성에 의해
$$
H^q(U, K) = H^q(\Gamma(U, \mathcal{I}^\bullet)) =
H^q(\Gamma(\mathcal{C}/U, \mathcal{I}^\bullet|_{\mathcal{C}/U}))
$$
이다. Lemma \ref{sites-cohomology-lemma-restrict-K-injective-to-open}에 의해 복합체
$\mathcal{I}^\bullet|_{\mathcal{C}/U}$는
$K|_{\mathcal{C}/U}$를 나타내는 K-단사 복합체이므로 결론이 따른다.
\end{proof}

\begin{lemma}
\label{sites-cohomology-lemma-sheafification-cohomology}
$(\mathcal{C}, \mathcal{O})$를 환 달린 사이트라 하고, $K$를
$D(\mathcal{O})$의 대상이라 하자. 준층
$$
U \mapsto H^q(U, K) = H^q(\mathcal{C}/U, K|_{\mathcal{C}/U})
$$
의 층화는 $K$의 $q$차 코호몰로지 층 $H^q(K)$이다.
\end{lemma}

\begin{proof}
등식 $H^q(U, K) = H^q(\mathcal{C}/U, K|_{\mathcal{C}/U})$는 Lemma
\ref{sites-cohomology-lemma-unbounded-cohomology-of-open}에서 따른다. $K$를 나타내는
K-단사 복합체 $\mathcal{I}^\bullet$를 고르자. 그러면 코호몰로지의
구성에 의해
$$
H^q(U, K) =
\frac{\Ker(\mathcal{I}^q(U) \to \mathcal{I}^{q + 1}(U))}
{\Im(\mathcal{I}^{q - 1}(U) \to \mathcal{I}^q(U))}.
$$
한편 $H^q(K) = \Ker(\mathcal{I}^q \to \mathcal{I}^{q + 1})/
\Im(\mathcal{I}^{q - 1} \to \mathcal{I}^q)$이므로 결과는 명백하다.
\end{proof}

\begin{lemma}
\label{sites-cohomology-lemma-restrict-direct-image-open}
연속 함자 $u : \mathcal{D} \to \mathcal{C}$에 대응하는 환 달린 사이트의
사상을 $f : (\mathcal{C}, \mathcal{O}_\mathcal{C}) \to
(\mathcal{D}, \mathcal{O}_\mathcal{D})$라 하자. $V \in \mathcal{D}$가
주어지면 $U = u(V)$로 두고 유도된 환 달린 사이트의 사상을
$g : (\mathcal{C}/U, \mathcal{O}_U) \to (\mathcal{D}/V, \mathcal{O}_V)$
로 나타내자
(Modules on Sites, Lemma
\ref{sites-modules-lemma-localize-morphism-ringed-sites}).
그러면 $D(\mathcal{O}_\mathcal{C})$의 $E$에 대해
$(Rf_*E)|_{\mathcal{D}/V} = Rg_*(E|_{\mathcal{C}/U})$이다.
\end{lemma}

\begin{proof}
$E$를 K-단사 $\mathcal{O}_\mathcal{C}$-가군 복합체
$\mathcal{I}^\bullet$로 나타내자. 그러면
$Rf_*(E) = f_*\mathcal{I}^\bullet$
이고, Lemma \ref{sites-cohomology-lemma-restrict-K-injective-to-open}에 의해
$Rg_*(E|_{\mathcal{C}/U}) = g_*(\mathcal{I}^\bullet|_{\mathcal{C}/U})$이다.
사이트 $\mathcal{C}$ 위의 임의의 층 $\mathcal{F}$에 대해
$(f_*\mathcal{F})|_{\mathcal{D}/V} = g_*(\mathcal{F}|_{\mathcal{C}/U})$
임은 명백하므로 결과가 따른다(Modules on Sites, Lemma
\ref{sites-modules-lemma-localize-morphism-ringed-sites} 또는 더 기본적인
Sites, Lemma \ref{sites-lemma-localize-morphism}를 보라).
\end{proof}

\begin{lemma}
\label{sites-cohomology-lemma-Leray-unbounded}
연속 함자 $u : \mathcal{D} \to \mathcal{C}$에 대응하는 환 달린 사이트의
사상을 $f : (\mathcal{C}, \mathcal{O}_\mathcal{C}) \to
(\mathcal{D}, \mathcal{O}_\mathcal{D})$라 하자. 그러면 함자
$D(\mathcal{O}_\mathcal{C}) \to D(\Gamma(\mathcal{O}_\mathcal{D}))$로서
$R\Gamma(\mathcal{D}, -) \circ Rf_* = R\Gamma(\mathcal{C}, -)$이다.
더 일반적으로 $V \in \mathcal{D}$와 $U = u(V)$에 대해
$R\Gamma(U, -) = R\Gamma(V, -) \circ Rf_*$이다.
\end{lemma}

\begin{proof}
환 $\Gamma(\mathcal{O}_\mathcal{D})$가 주는 $\mathcal{O}_{pt}$를 갖춘
점 토포스 $pt$를 생각하자. 구조층의 대역 단면 위에서 주어진 식별을
유도하는 환 달린 토포스의 표준 사상
$(\mathcal{D}, \mathcal{O}_\mathcal{D}) \to (pt, \mathcal{O}_{pt})$
이 있다. 그러면
$D(\mathcal{O}_{pt}) = D(\Gamma(\mathcal{O}_\mathcal{D}))$.
등식
$R\Gamma(\mathcal{D}, -) \circ Rf_* = R\Gamma(\mathcal{C}, -)$
은 다음 합성에 Lemma \ref{sites-cohomology-lemma-derived-pushforward-composition}을
적용하면 따른다.
$$
(\mathcal{C}, \mathcal{O}_\mathcal{C}) \to
(\mathcal{D}, \mathcal{O}_\mathcal{D}) \to (pt, \mathcal{O}_{pt})
$$
둘째인 더 일반적인 명제는 Lemma
\ref{sites-cohomology-lemma-restrict-direct-image-open}을 적용한 뒤 첫째 명제에서 따른다.
\end{proof}

\begin{lemma}
\label{sites-cohomology-lemma-unbounded-describe-higher-direct-images}
연속 함자 $u : \mathcal{D} \to \mathcal{C}$에 대응하는 환 달린 사이트의
사상을 $f : (\mathcal{C}, \mathcal{O}_\mathcal{C}) \to
(\mathcal{D}, \mathcal{O}_\mathcal{D})$라 하고, $K$를
$D(\mathcal{O}_\mathcal{C})$의 대상이라 하자. 그러면 $H^i(Rf_*K)$는 준층
$$
V \mapsto H^i(u(V), K) = H^i(V, Rf_*K)
$$
에 연관된 층이다.
\end{lemma}

\begin{proof}
등식 $H^i(u(V), K) = H^i(V, Rf_*K)$는 Lemma
\ref{sites-cohomology-lemma-Leray-unbounded}의 둘째 명제에서 코호몰로지를 취하면 따른다.
층화에 관한 명제는 Lemma \ref{sites-cohomology-lemma-sheafification-cohomology}에서
따른다.
\end{proof}

\begin{lemma}
\label{sites-cohomology-lemma-modules-abelian-unbounded}
$(\mathcal{C}, \mathcal{O}_\mathcal{C})$를 환 달린 사이트라 하고,
$K$를 $D(\mathcal{O}_\mathcal{C})$의 대상이라 하자.
$D(\underline{\mathbf{Z}}_\mathcal{C})$에서의 그 상을 $K_{ab}$로 나타내자.
\begin{enumerate}
\item 표준 사상
$R\Gamma(\mathcal{C}, K) \to R\Gamma(\mathcal{C}, K_{ab})$
이 있으며, 이는 $D(\textit{Ab})$ 안의 동형이다.
\item 모든 $U \in \mathcal{C}$에 대해 표준 사상
$R\Gamma(U, K) \to R\Gamma(U, K_{ab})$
이 있으며, 이는 $D(\textit{Ab})$ 안의 동형이다.
\item $f : (\mathcal{C}, \mathcal{O}_\mathcal{C}) \to
(\mathcal{D}, \mathcal{O}_\mathcal{D})$를 환 달린 사이트의 사상이라 하자.
표준 사상 $Rf_*K \to Rf_*(K_{ab})$가 있으며, 이는
$D(\underline{\mathbf{Z}}_\mathcal{D})$ 안의 동형이다.
\end{enumerate}
\end{lemma}

\begin{proof}
사상은 다음과 같이 구성한다. $K$를 나타내는 K-단사 복합체
$\mathcal{I}^\bullet$를 고르자. $\mathcal{J}^\bullet$가 아벨군의 K-단사
복합체인 유사동형
$\mathcal{I}^\bullet \to \mathcal{J}^\bullet$를 고르자. 그러면 (1)의
사상은
$\Gamma(\mathcal{C}, \mathcal{I}^\bullet) \to
\Gamma(\mathcal{C}, \mathcal{J}^\bullet)$
로, (2)의 사상은
$\Gamma(U, \mathcal{I}^\bullet) \to \Gamma(U, \mathcal{J}^\bullet)$
로, (3)의 사상은
$f_*\mathcal{I}^\bullet \to f_*\mathcal{J}^\bullet$로 주어진다.
이 사상들이 동형임을 보이려면 코호몰로지 군과 코호몰로지 층에서
동형을 유도함을 보이면 충분하다. Lemmas
\ref{sites-cohomology-lemma-unbounded-cohomology-of-open}와
\ref{sites-cohomology-lemma-unbounded-describe-higher-direct-images}에 의해 사상
$$
H^0(\mathcal{C}, K) \longrightarrow H^0(\mathcal{C}, K_{ab})
$$
이 동형임을 보이면 충분하다. 다음을 관찰하자.
$$
H^0(\mathcal{C}, K) =
\Hom_{D(\mathcal{O}_\mathcal{C})}(\mathcal{O}_\mathcal{C}, K)
$$
다른 군도 마찬가지이다. $K$를 나타내는 임의의
$\mathcal{O}_\mathcal{C}$-가군 복합체 $\mathcal{K}^\bullet$를 고르자.
유도범주를 국소화로 구성한 방식에 의해
$$
\Hom_{D(\mathcal{O}_\mathcal{C})}(\mathcal{O}_\mathcal{C}, K) =
\colim_{s : \mathcal{F}^\bullet \to \mathcal{O}_\mathcal{C}}
\Hom_{K(\mathcal{O}_\mathcal{C})}(\mathcal{F}^\bullet, \mathcal{K}^\bullet)
$$
이며, 여기서 여극한은 $\mathcal{O}_\mathcal{C}$-가군 복합체의
유사동형 $s$ 전체에 걸쳐 취한다. 마찬가지로
$$
\Hom_{D(\underline{\mathbf{Z}}_\mathcal{C})}
(\underline{\mathbf{Z}}_\mathcal{C}, K) =
\colim_{s : \mathcal{G}^\bullet \to \underline{\mathbf{Z}}_\mathcal{C}}
\Hom_{K(\underline{\mathbf{Z}}_\mathcal{C})}
(\mathcal{G}^\bullet, \mathcal{K}^\bullet)
$$
이다. 다음으로 $\mathcal{G}^\bullet$가 평탄
$\underline{\mathbf{Z}}_\mathcal{C}$-가군들로 이루어진 위로 유계인
복합체인 유사동형
$s : \mathcal{G}^\bullet \to \underline{\mathbf{Z}}_\mathcal{C}$들이 이
계에서 공종임을 관찰하자. (이는 Modules on Sites, Lemma
\ref{sites-modules-lemma-module-quotient-flat} and
Derived Categories, Lemma \ref{derived-lemma-subcategory-left-resolution};
Section \ref{sites-cohomology-section-flat}의 논의를 보라.) 따라서 원소
$\xi \in H^0(\mathcal{C}, K_{ab})$를 쌍
$$
(s : \mathcal{G}^\bullet \to \underline{\mathbf{Z}}_\mathcal{C},
a : \mathcal{G}^\bullet \to \mathcal{K}^\bullet)
$$
으로 나타내자. 여기서 $\mathcal{G}^\bullet$는 평탄
$\underline{\mathbf{Z}}_\mathcal{C}$-가군들로 이루어진 위로 유계인
복합체이다. 이 쌍을
$$
(\mathcal{G}^\bullet \otimes_{\underline{\mathbf{Z}}_\mathcal{C}}
\mathcal{O}_\mathcal{C}
\to \mathcal{O}_\mathcal{C},
\mathcal{G}^\bullet  \otimes_{\underline{\mathbf{Z}}_\mathcal{C}}
\mathcal{O}_\mathcal{C}
\to \mathcal{K}^\bullet)
$$
로 보내면 사상
$H^0(\mathcal{C}, K) \longrightarrow H^0(\mathcal{C}, K_{ab})$의 역을
구성할 수 있다. 여기서 유의할 것은 첫째 화살표가 Lemmas
\ref{sites-cohomology-lemma-derived-tor-quasi-isomorphism-other-side}와
\ref{sites-cohomology-lemma-bounded-flat-K-flat}에 의해 유사동형이라는 점뿐이다.
이 구성이 실제로 역임을 확인하는 자세한 과정은 생략한다.
\end{proof}

\begin{lemma}
\label{sites-cohomology-lemma-adjoint-lower-shriek-restrict}
$(\mathcal{C}, \mathcal{O})$를 환 달린 사이트라 하고, $U$를
$\mathcal{C}$의 객체라 하자. 이에 대응하는 국소화 사상을
$j : (\Sh(\mathcal{C}/U), \mathcal{O}_U) \to (\Sh(\mathcal{C}), \mathcal{O})$
로 나타내자. 제한 함자 $D(\mathcal{O}) \to D(\mathcal{O}_U)$는 영에 의한
연장 $j_! : D(\mathcal{O}_U) \to D(\mathcal{O})$의 오른쪽 수반이다.
\end{lemma}

\begin{proof}
다음을 보이면 된다.
$$
\Hom_{D(\mathcal{O})}(j_!E, F) = \Hom_{D(\mathcal{O}_U)}(E, F|_U)
$$
$E$를 나타내는 $\mathcal{O}_U$-가군 복합체 $\mathcal{E}^\bullet$와
$F$를 나타내는 K-단사 복합체 $\mathcal{I}^\bullet$를 고르자. Lemma
\ref{sites-cohomology-lemma-restrict-K-injective-to-open}에 의해 복합체
$\mathcal{I}^\bullet|_U$도 K-단사이다. 따라서 위 공식은
$$
\Hom_{D(\mathcal{O})}(j_!\mathcal{E}^\bullet, \mathcal{I}^\bullet) =
\Hom_{D(\mathcal{O}_U)}(\mathcal{E}^\bullet, \mathcal{I}^\bullet|_U)
$$
가 된다. 이는 $|_U$와 $j_!$가 수반함자이기 때문에 성립한다(Modules on
Sites, Lemma \ref{sites-modules-lemma-extension-by-zero}와 Derived
Categories, Lemma \ref{derived-lemma-K-injective}).
\end{proof}

\begin{lemma}
\label{sites-cohomology-lemma-j-shriek-and-tensor}
$(\mathcal{C}, \mathcal{O})$를 환 달린 사이트라 하고,
$U \in \Ob(\mathcal{C})$라 하자. $D(\mathcal{O}_U)$의 $L$과
$D(\mathcal{O})$의 $K$에 대해
$j_!L \otimes_\mathcal{O}^\mathbf{L} K =
j_!(L \otimes_{\mathcal{O}_U}^\mathbf{L} K|_U)$.
\end{lemma}

\begin{proof}
$L$을 $\mathcal{O}_U$-가군 복합체로 나타내고 $K$를 K-평탄
$\mathcal{O}$-가군 복합체로 나타낸 뒤 Modules on Sites, Lemma
\ref{sites-modules-lemma-j-shriek-and-tensor}를 적용한다. 세부는 생략한다.
\end{proof}

\begin{lemma}
\label{sites-cohomology-lemma-K-injective-flat}
$f : (\Sh(\mathcal{C}), \mathcal{O}_\mathcal{C}) \to
(\Sh(\mathcal{D}), \mathcal{O}_\mathcal{D})$를 환 달린 토포스의 평탄
사상이라 하자. $\mathcal{I}^\bullet$가 K-단사
$\mathcal{O}_\mathcal{C}$-가군 복합체이면
$f_*\mathcal{I}^\bullet$는 $\mathcal{O}_\mathcal{D}$-가군 복합체로서
K-단사이다.
\end{lemma}

\begin{proof}
이는
$$
\Hom_{K(\mathcal{O}_\mathcal{D})}(\mathcal{F}^\bullet, f_*\mathcal{I}^\bullet)
=
\Hom_{K(\mathcal{O}_\mathcal{C})}(f^*\mathcal{F}^\bullet, \mathcal{I}^\bullet)
$$
이고(Modules on Sites, Lemma
\ref{sites-modules-lemma-adjoint-pullback-pushforward-modules}), $f$가
평탄하다고 가정했으므로 $f^*$가 완전하기 때문이다.
\end{proof}

\begin{lemma}
\label{sites-cohomology-lemma-hom-K-injective}
$\mathcal{C}$를 사이트라 하고, $\mathcal{O} \to \mathcal{O}'$를
환층의 사상이라 하자. $\mathcal{I}^\bullet$가 K-단사
$\mathcal{O}$-가군 복합체이면
$\SheafHom_\mathcal{O}(\mathcal{O}', \mathcal{I}^\bullet)$
은 K-단사 $\mathcal{O}'$-가군 복합체이다.
\end{lemma}

\begin{proof}
이는 Modules on Sites, Lemma
\ref{sites-modules-lemma-adjoint-hom-restrict}에 의해
$\Hom_{K(\mathcal{O}')}(\mathcal{G}^\bullet,
\Hom_\mathcal{O}(\mathcal{O}', \mathcal{I}^\bullet)) =
\Hom_{K(\mathcal{O})}(\mathcal{G}^\bullet, \mathcal{I}^\bullet)$
이기 때문이다.
\end{proof}






\section{국소화와 코호몰로지}
\label{sites-cohomology-section-localization}

\noindent
$\mathcal{C}$를 사이트라 하고, $f : X \to Y$를 $\mathcal{C}$의 사상이라
하자. 그러면 토포스의 사상
$$
j_{X/Y} : \Sh(\mathcal{C}/X) \longrightarrow \Sh(\mathcal{C}/Y)
$$
을 얻는다. Sites, Sections \ref{sites-section-localize}와
\ref{sites-section-more-localize}를 보라. 이러한 유형의 토포스 사상에서는
코호몰로지에 관한 몇몇 문제를 더 쉽게 다룰 수 있다. 자명한 유형의
밑변환 정리를 얻는 예를 하나 들겠다.

\begin{lemma}
\label{sites-cohomology-lemma-localize-cartesian-square}
$\mathcal{C}$를 사이트라 하자. $\mathcal{C}$의 카르테시안 도식
$$
\xymatrix{
X' \ar[d] \ar[r] & X \ar[d] \\
Y' \ar[r] & Y
}
$$
을 생각하자. 그러면 함자
$j_{Y'/Y}^{-1} \circ Rj_{X/Y, *} = Rj_{X'/Y', *} \circ j_{X'/X}^{-1}$
$D(\mathcal{C}/X) \to D(\mathcal{C}/Y')$로서 위 등식이 성립한다.
\end{lemma}

\begin{proof}
$E \in D(\mathcal{C}/X)$라 하고, $E$를 나타내는
$\mathcal{C}/X$ 위의 아벨층 K-단사 복합체 $\mathcal{I}^\bullet$를
고르자. Lemma \ref{sites-cohomology-lemma-restrict-K-injective-to-open}에 의해
$j_{X'/X}^{-1}\mathcal{I}^\bullet$도 K-단사이다. 따라서
$Rj_{X'/Y'}(j_{X'/X}^{-1}E)$를
$j_{X'/Y', *}j_{X'/X}^{-1}\mathcal{I}^\bullet$로 계산할 수 있다.
그러므로 등식은 Sites, Lemma
\ref{sites-lemma-localize-cartesian-square}에서 따른다.
\end{proof}

\noindent
환 달린 사이트 $(\mathcal{C}, \mathcal{O})$와 $\mathcal{C}$의 사상
$f : X \to Y$가 주어지면 $j_{X/Y}$는 환 달린 토포스의 사상
$$
j_{X/Y} :
(\Sh(\mathcal{C}/X), \mathcal{O}_X)
\longrightarrow
(\Sh(\mathcal{C}/Y), \mathcal{O}_Y)
$$
이 된다. Modules on Sites, Lemma
\ref{sites-modules-lemma-relocalize}를 보라.

\begin{lemma}
\label{sites-cohomology-lemma-localize-cartesian-square-modules}
$(\mathcal{C}, \mathcal{O})$를 환 달린 사이트라 하자.
$\mathcal{C}$의 카르테시안 도식
$$
\xymatrix{
X' \ar[d] \ar[r] & X \ar[d] \\
Y' \ar[r] & Y
}
$$
을 생각하자. 그러면 함자
$j_{Y'/Y}^* \circ Rj_{X/Y, *} = Rj_{X'/Y', *} \circ j_{X'/X}^*$
$D(\mathcal{O}_X) \to D(\mathcal{O}_{Y'})$로서 위 등식이 성립한다.
\end{lemma}

\begin{proof}
$j_{Y'/Y}^{-1}\mathcal{O}_Y = \mathcal{O}_{Y'}$이므로
$j_{Y'/Y}^* = Lj_{Y'/Y}^* = j_{Y'/Y}^{-1}$이다. 마찬가지로
$j_{X'/X}^* = Lj_{X'/X}^* = j_{X'/X}^{-1}$이다. 따라서 Lemma
\ref{sites-cohomology-lemma-modules-abelian-unbounded}에 의해 아벨층의 유도범주에서 결과를
증명하면 충분하며, 이는 Lemma
\ref{sites-cohomology-lemma-localize-cartesian-square}에서 증명했다.
\end{proof}












\section{역계와 코호몰로지}
\label{sites-cohomology-section-inverse-systems}

\noindent
가군층의 역계에 관한 몇 가지 결과를 증명한다.

\begin{lemma}
\label{sites-cohomology-lemma-ML-general}
$I$를 환 $A$의 아이디얼이라 하고, $\mathcal{C}$를 사이트라 하자.
$\mathcal{C}$ 위의 $A$-가군 층들의 역계
$$
\ldots \to \mathcal{F}_3 \to \mathcal{F}_2 \to \mathcal{F}_1
$$
를 생각하고
$\mathcal{F}_n = \mathcal{F}_{n + 1}/I^n\mathcal{F}_{n + 1}$이라 하자.
$p \geq 0$라 하자. 등급
$\bigoplus_{n \geq 0} I^n/I^{n + 1}$-가군으로서
$$
\bigoplus\nolimits_{n \geq 0} H^{p + 1}(\mathcal{C}, I^n\mathcal{F}_{n + 1})
$$
이 오름 사슬 조건을 만족한다고 가정하자. 그러면 역계
$M_n = H^p(\mathcal{C}, \mathcal{F}_n)$은 Mittag--Leffler 조건을
만족한다\footnote{실제로 모든 $n \geq c$에 대해
$\Im(M_n \to M_{n - c})$가 안정된 상이 되게 하는 $c \geq 0$가 존재한다.}.
\end{lemma}

\begin{proof}
$N_n = H^{p + 1}(\mathcal{C}, I^n\mathcal{F}_{n + 1})$로 놓고, 짧은 완전열
$0 \to I^n\mathcal{F}_{n + 1} \to \mathcal{F}_{n + 1} \to \mathcal{F}_n \to 0$
에서 오는 코호몰로지의 경계 사상을 $\delta_n : M_n \to N_n$이라 하자.
그러면 $\bigoplus \Im(\delta_n) \subset \bigoplus N_n$은 등급
부분가군이다. 실제로 $s \in M_n$이고 $f \in I^m$이면 가환도식
$$
\xymatrix{
0 \ar[r] &
I^n\mathcal{F}_{n + 1} \ar[d]_f \ar[r] &
\mathcal{F}_{n + 1} \ar[d]_f \ar[r] &
\mathcal{F}_n \ar[d]_f \ar[r] & 0 \\
0 \ar[r] &
I^{n + m}\mathcal{F}_{n + m + 1} \ar[r] &
\mathcal{F}_{n + m + 1} \ar[r] &
\mathcal{F}_{n + m} \ar[r] & 0
}
$$
을 얻는다. 가운데 수직 사상은 $\mathcal{F}_{n + 1}$의 국소 단면을
$\mathcal{F}_{n + m + 1}$의 단면으로 올린 뒤 $f$를 곱하여 주어지며,
다른 수직 화살표들도 마찬가지이다. 따라서
$\delta_{n + m}(fs) = f \delta_n(s)$이다. 가정에 의해
$\delta_{n_j}(s_j)$들이 $\bigoplus \Im(\delta_n)$을 등급 가군으로
생성하도록 하는 $s_j \in M_{n_j}$, $j = 1, \ldots, N$을 찾을 수 있다.
$n > c = \max(n_j)$라 하고, $s \in M_n$이라 하자. 그러면
$\delta_n(s) = \sum f_j \delta_{n_j}(s_j)$가 되게 하는
$f_j \in I^{n - n_j}$를 찾을 수 있다. 따라서
$\delta(s - \sum f_j s_j) = 0$이다. 즉 $M_n$ 안에서
$s - \sum f_js_j$로 가는 $s' \in M_{n + 1}$을 찾을 수 있다. 그러므로
$$
\Im(M_{n + 1} \to M_{n - c}) = \Im(M_n \to M_{n - c})
$$
이다. 실제로 원소 $f_js_j$들은 $M_{n - c}$에서 영으로 간다.
이것이 보조정리를 증명한다.
\end{proof}

\begin{lemma}
\label{sites-cohomology-lemma-ML-general-better}
$I$를 환 $A$의 아이디얼이라 하고, $\mathcal{C}$를 사이트라 하자.
$\mathcal{C}$ 위의 $A$-가군의 역계
$$
\ldots \to \mathcal{F}_3 \to \mathcal{F}_2 \to \mathcal{F}_1
$$
를 생각하고
$\mathcal{F}_n = \mathcal{F}_{n + 1}/I^n\mathcal{F}_{n + 1}$이라 하자.
$p \geq 0$라 하자. $n$이 주어졌을 때
$$
N_n =
\bigcap\nolimits_{m \geq n}
\Im\left(
H^{p + 1}(\mathcal{C}, I^n\mathcal{F}_{m + 1}) \to
H^{p + 1}(\mathcal{C}, I^n\mathcal{F}_{n + 1})
\right)
$$
으로 정의하자. 등급 $\bigoplus_{n \geq 0} I^n/I^{n + 1}$-가군으로서
$\bigoplus N_n$이 오름 사슬 조건을 만족하면 역계
$M_n = H^p(\mathcal{C}, \mathcal{F}_n)$은 Mittag--Leffler 조건을
만족한다\footnote{실제로 모든 $n \geq c$에 대해
$\Im(M_n \to M_{n - c})$가 안정된 상이 되게 하는 $c \geq 0$가 존재한다.}.
\end{lemma}

\begin{proof}
증명은 Lemma \ref{sites-cohomology-lemma-ML-general}의 증명과 완전히 같다. 실제로
$\bigoplus N_n$이
$\bigoplus H^{p + 1}(\mathcal{C}, I^n\mathcal{F}_{n + 1})$의 등급
$\bigoplus_{n \geq 0} I^n/I^{n + 1}$-부분가군이고, 경계 사상
$\delta_n : M_n \to H^{p + 1}(\mathcal{C}, I^n\mathcal{F}_{n + 1})$의
상이 $N_n$에 포함됨을 보이기만 하면 거기서 제시한 논증으로 결과가
따른다.

\medskip\noindent
$\xi \in N_n$이고 $f \in I^k$라고 하자. $m \gg n + k$를 고르고,
$\xi$로 가는
$\xi' \in H^{p + 1}(\mathcal{C}, I^n\mathcal{F}_{m + 1})$를 고르자.
Lemma \ref{sites-cohomology-lemma-ML-general}의 증명에서와 같이 구성한 도식
$$
\xymatrix{
0 \ar[r] &
I^n\mathcal{F}_{m + 1} \ar[d]_f \ar[r] &
\mathcal{F}_{m + 1} \ar[d]_f \ar[r] &
\mathcal{F}_n \ar[d]_f \ar[r] & 0 \\
0 \ar[r] &
I^{n + k}\mathcal{F}_{m + 1} \ar[r] &
\mathcal{F}_{m + 1} \ar[r] &
\mathcal{F}_{n + k} \ar[r] & 0
}
$$
을 생각하자. 코호몰로지에서 유도된 사상을 얻고,
$f \xi' \in H^{p + 1}(\mathcal{C}, I^{n + k}\mathcal{F}_{m + 1})$
가 $f \xi$로 감을 알 수 있다. 이것은 모든 $m \gg n + k$에 대해
성립하므로 원하는 대로 $f\xi$는 $N_{n + k}$에 속한다.

\medskip\noindent
경계 사상 $\delta_n$의 상이 $N_n$에 포함됨을 보기 위해
$m \geq n$에 대해 도식
$$
\xymatrix{
0 \ar[r] &
I^n\mathcal{F}_{m + 1} \ar[d] \ar[r] &
\mathcal{F}_{m + 1} \ar[d] \ar[r] &
\mathcal{F}_n \ar[d] \ar[r] & 0 \\
0 \ar[r] &
I^n\mathcal{F}_{n + 1} \ar[r] &
\mathcal{F}_{n + 1} \ar[r] &
\mathcal{F}_n \ar[r] & 0
}
$$
들을 생각한다. 코호몰로지에서 유도된 사상들을 보면 결론이 따른다.
\end{proof}

\begin{lemma}
\label{sites-cohomology-lemma-topology-I-adic-general}
$I$를 환 $A$의 아이디얼이라 하고, $\mathcal{C}$를 사이트라 하자.
$\mathcal{C}$ 위의 $A$-가군 층들의 역계
$$
\ldots \to \mathcal{F}_3 \to \mathcal{F}_2 \to \mathcal{F}_1
$$
를 생각하고
$\mathcal{F}_n = \mathcal{F}_{n + 1}/I^n\mathcal{F}_{n + 1}$이라 하자.
$p \geq 0$라 하자. 등급
$\bigoplus_{n \geq 0} I^n/I^{n + 1}$-가군으로서
$$
\bigoplus\nolimits_{n \geq 0} H^p(\mathcal{C}, I^n\mathcal{F}_{n + 1})
$$
이 오름 사슬 조건을 만족한다고 가정하자. 그러면
$M = \lim H^p(\mathcal{C}, \mathcal{F}_n)$ 위의 극한 위상은
$I$-진 위상이다.
\end{lemma}

\begin{proof}
$n \geq 1$에 대해
$F^n = \Ker(M \to H^p(\mathcal{C}, \mathcal{F}_n))$으로 놓고
$F^0 = M$으로 놓자. $I F^n \subset F^{n + 1}$임을 관찰하자. 특히
$I^n M \subset F^n$이다. 따라서 $I$-진 위상은 극한 위상보다 더
미세하다. 역을 위해, 주어진 $n$에 대해 $F^m \subset I^nM$이 되게 하는
$m \geq n$이 존재함을 보일 것이다\footnote{실제로 모든 $n$에 대해
$F^{n + c} \subset I^nM$이 되게 하는 $c \geq 0$가 존재한다.}.
단사 사상
$$
F^n/F^{n + 1} \longrightarrow H^p(\mathcal{C}, \mathcal{F}_{n + 1})
$$
이 있고, 그 상은
$H^p(\mathcal{C}, I^n\mathcal{F}_{n + 1}) \to
H^p(\mathcal{C}, \mathcal{F}_{n + 1})$.
$F^n/F^{n + 1}$의 역상을
$$
E_n \subset H^p(\mathcal{C}, I^n\mathcal{F}_{n + 1})
$$
이라 나타내자. 그러면 $\bigoplus E_n$은
$\bigoplus H^p(\mathcal{C}, I^n\mathcal{F}_{n + 1})$의 등급
$\bigoplus I^n/I^{n + 1}$-부분가군이고,
$\bigoplus E_n \to \bigoplus F^n/F^{n + 1}$은 등급 가군의 준동형이다.
자세한 내용은 생략한다. 가정에 의해 $\bigoplus E_n$은
$\bigoplus I^n/I^{n + 1}$ 위에서 유한 개의 동차 원소들로 생성된다.
$E_n \to F^n/F^{n + 1}$가 전사이므로
$\bigoplus F^n/F^{n + 1}$에 대해서도 같은 것이 성립한다. 따라서
$r$과 $c_1, \ldots, c_r \geq 0$, 그리고 그 상들이
$\bigoplus F^n/F^{n + 1}$을 생성하는 $a_i \in F^{c_i}$를 찾을 수 있다.
$c = \max(c_i)$로 놓자.

\medskip\noindent
$n \geq c$에 대해 $I F^n = F^{n + 1}$이라고 주장한다. 이 주장은
원하는 대로 $F^{n + c} = I^nF^c \subset I^nM$임을 보인다. 주장을
증명하기 위해 $a \in F^{n + 1}$이라고 하자.
$F^{n + 1}/F^{n + 2}$에서의 $a$의 상은 $a_i$들의 선형결합이다. 따라서
어떤 $f_i \in I^{n + 1 - c_i}$에 대해
$a  - \sum f_i a_i \in F^{n + 2}$이다. $n \geq c_i$이므로
$I^{n + 1 - c_i} = I \cdot I^{n - c_i}$이고, 따라서
$g_{i, j} \in I$이고 $h_{i, j}a_i \in F^n$이도록
$f_i = \sum g_{i, j} h_{i, j}$로 쓸 수 있다. 그러므로
$F^{n + 1} = F^{n + 2} + IF^n$이다. 간단한 귀납법으로 모든
$e > 0$에 대해 $F^{n + 1} = F^{n + e} + IF^n$을 얻는다. 따라서
$IF^n$은 $F^{n + 1}$에서 조밀하다. $I$의 생성자
$k_1, \ldots, k_r$를 고르고 연속 사상
$$
u : (F^n)^{\oplus r} \longrightarrow F^{n + 1},\quad
(x_1, \ldots, x_r) \mapsto \sum k_i x_i
$$
을 극한 위상에서 생각하자. 위의 논의에 의해 모든 $m \geq n$에 대해
$u$ 아래 $(F^m)^{\oplus r}$의 상은 $F^{m + 1}$에서 조밀하다. 열린 사상
보조정리(More on Algebra, Lemma
\ref{more-algebra-lemma-open-mapping})에 의해 $u$는 열린 사상이다.
따라서 $u$는 전사이다. 그러므로 $n \geq c$에 대해
$IF^n = F^{n + 1}$이다. 이것으로 증명을 마친다.
\end{proof}

\begin{lemma}
\label{sites-cohomology-lemma-topology-I-adic-general-better}
$I$를 환 $A$의 아이디얼이라 하고, $\mathcal{C}$를 사이트라 하자.
$$
\ldots \to \mathcal{F}_3 \to \mathcal{F}_2 \to \mathcal{F}_1
$$
을 $\mathcal{C}$ 위의 $A$-가군 층들의 역계라 하고,
$\mathcal{F}_n = \mathcal{F}_{n + 1}/I^n\mathcal{F}_{n + 1}$이라 하자.
$p \geq 0$라 하자. 주어진 $n$에 대해
$$
N_n =
\bigcap\nolimits_{m \geq n}
\Im\left(
H^p(\mathcal{C}, I^n\mathcal{F}_{m + 1}) \to
H^p(\mathcal{C}, I^n\mathcal{F}_{n + 1})
\right)
$$
으로 정의하자. 등급 $\bigoplus_{n \geq 0} I^n/I^{n + 1}$-가군으로서
$\bigoplus N_n$이 오름 사슬 조건을 만족하면
$M = \lim H^p(\mathcal{C}, \mathcal{F}_n)$ 위의 극한 위상은
$I$-진 위상이다.
\end{lemma}

\begin{proof}
증명은 Lemma \ref{sites-cohomology-lemma-topology-I-adic-general}의 증명과 완전히 같다.
실제로 $\bigoplus N_n$이
$\bigoplus H^{p + 1}(\mathcal{C}, I^n\mathcal{F}_{n + 1})$의 등급
$\bigoplus_{n \geq 0} I^n/I^{n + 1}$-부분가군이고,
$F^n/F^{n + 1} \subset H^p(\mathcal{C}, \mathcal{F}_{n + 1})$가
$N_n \to H^p(\mathcal{C}, \mathcal{F}_{n + 1})$의 상에 포함됨을 보이기만 하면
거기서 제시한 논증으로 결과가 따른다.
Lemma \ref{sites-cohomology-lemma-ML-general-better}의 증명에서 가군 구조에 관한 명제를 보았다.

\medskip\noindent
$t \in F^n$이라 하자. $H^p(\mathcal{C}, \mathcal{F}_{n + 1})$에서 $t$의 상으로
가는 원소 $s \in H^p(\mathcal{C}, I^n\mathcal{F}_{n + 1})$를 고르자.
$s$가 $N_n$에 속함을 보여야 한다. 이제 $F^n$은 사상
$M \to H^p(\mathcal{C}, \mathcal{F}_n)$의 핵이므로, 모든 $m \geq n$에 대해
$t$를 원소 $t_m \in H^p(\mathcal{C}, \mathcal{F}_{m + 1})$으로 보낼 수 있으며
그 원소는 $H^p(\mathcal{C}, \mathcal{F}_n)$에서 영으로 간다. 짧은 완전열
$0 \to I^n\mathcal{F}_{m + 1} \to \mathcal{F}_{m + 1} \to \mathcal{F}_n \to 0$
에서 오는 코호몰로지 열
$$
H^{p - 1}(\mathcal{C}, \mathcal{F}_n) \to
H^p(\mathcal{C}, I^n\mathcal{F}_{m + 1}) \to
H^p(\mathcal{C}, \mathcal{F}_{m + 1}) \to
H^p(\mathcal{C}, \mathcal{F}_n)
$$
을 생각하자. $t_m$으로 가는
$s_m \in H^p(\mathcal{C}, I^n\mathcal{F}_{m + 1})$을 고를 수 있다.
위의 열과 $m = n$일 때의 열을 비교하면, $s_m$은
$H^{p - 1}(\mathcal{C}, \mathcal{F}_n) \to
H^p(\mathcal{C}, I^n\mathcal{F}_{n + 1})$의 상에 속하는 원소를 더한 것을
제외하면 $s$로 감을 알 수 있다. 그러나 이 사상은
$H^p(\mathcal{C}, I^n\mathcal{F}_{m + 1}) \to
H^p(\mathcal{C}, I^n\mathcal{F}_{n + 1})$
를 거쳐 인수분해되므로 원하는 대로 $s$가 그 상에 속함을 알 수 있다.
\end{proof}












\section{유도 극한과 호모토피 극한}
\label{sites-cohomology-section-derived-limits}

\noindent
$\mathcal{C}$를 사이트라 하자. $n > m$이면
$\Mor((U, n), (V, m)) = \emptyset$이고, 그렇지 않으면
$\Mor((U, n), (V, m)) = \Mor(U, V)$인 범주
$\mathcal{C} \times \mathbf{N}$을 생각하자.
$\{U_i \to U\}$가 $\mathcal{C}$의 덮개일 때의 족
$\{(U_i, n) \to (U, n)\}$을 덮개로 삼아 이 범주에 사이트 구조를 준다.
그러면 $\mathcal{C} \times \mathbf{N}$ 위의 층은
$\mathcal{C}$ 위의 층들의 역계와 같은 것임을 곧바로 확인할 수 있다.
특히 $\textit{Ab}(\mathcal{C} \times \mathbf{N})$은
$\mathcal{C}$ 위의 아벨층들의 역계로 이루어진 범주이다.
이제 역계를 그 극한으로 보내는 함자
$$
\lim : \textit{Ab}(\mathcal{C} \times \mathbf{N}) \to \textit{Ab}(\mathcal{C})
$$
를 생각하자. 이는 연속이고 쌍대연속인 함자
$\mathcal{C} \times \mathbf{N} \to \mathcal{C}$에 결부된 토포스의 사상
$g : \Sh(\mathcal{C} \times \mathbf{N}) \to \Sh(\mathcal{C})$의
$g_*$에 불과하다. ($g^{-1}$은 $\mathcal{C}$ 위의 층을 그에 대응하는
상수 역계로 보냄에 유의하라.)

\medskip\noindent
위에서 설명한 일반 이론으로부터 유도함자
$$
R\lim = Rg_* : D(\mathcal{C} \times \mathbf{N}) \to D(\mathcal{C}).
$$
를 얻는다. 표시한 대로 이 함자는 흔히 $R\lim$으로 쓴다.

\medskip\noindent
한편 연속이고 쌍대연속인 함자
$\mathcal{C} \to \mathcal{C} \times \mathbf{N}$,
$U \mapsto (U, n)$은 토포스의 사상
$i_n : \Sh(\mathcal{C}) \to \Sh(\mathcal{C} \times \mathbf{N})$들을 정의한다.
물론 $i_n^{-1}$은 역계의 $n$번째 항을 고르는 함자이다. 따라서 함자의 변환
$i_{n + 1}^{-1} \to i_n^{-1}$들이 있다. 그러므로
$K \in D(\mathcal{C} \times \mathbf{N})$가 주어지면
$K_n = i_n^{-1}K \in D(\mathcal{C})$와 사상 $K_{n + 1} \to K_n$을 얻는다.
Derived Categories, Definition \ref{derived-definition-derived-limit}에서
호모토피 극한
$$
R\lim K_n \in D(\mathcal{C})
$$
의 개념을 정의했다. 두 개념은 의미가 통하는 범위에서 일치한다고 주장한다.

\begin{lemma}
\label{sites-cohomology-lemma-derived-limit-is-ok}
$\mathcal{C}$를 사이트라 하고, $K$를
$D(\mathcal{C} \times \mathbf{N})$의 대상이라 하자.
위와 같이 $K_n = i_n^{-1}K$로 놓자. 그러면 $D(\mathcal{C})$ 안에서
$$
R\lim K \cong R\lim K_n
$$
이다.
\end{lemma}

\begin{proof}
$D(\mathcal{C} \times \mathbf{N})$의 대상 $K$에 대해 $R\lim$을 계산하려고
각 항이 $\textit{Ab}(\mathcal{C} \times \mathbf{N})$의 단사대상인
K-단사 대표 $\mathcal{I}^\bullet$을 고른다. Injectives, Theorem
\ref{injectives-theorem-K-injective-embedding-grothendieck}을 보라.
$\mathcal{I}^\bullet$을 복합체들의 역계
$(\mathcal{I}_n^\bullet)$로 생각할 수 있고 그렇게 하기로 하면
$$
R\lim K = \lim \mathcal{I}_n^\bullet
$$
임을 알 수 있다. 여기서 오른쪽은 항별 역극한이다.

\medskip\noindent
$\mathcal{J} = (\mathcal{J}_n)$를
$\textit{Ab}(\mathcal{C} \times \mathbf{N})$의 단사대상이라 하자.
사상 $(U, n) \to (U, n + 1)$은
$\mathcal{C} \times \mathbf{N}$의 단사사상이므로
$\mathcal{J}(U, n + 1) \to \mathcal{J}(U, n)$은 전사이다
(Lemma \ref{sites-cohomology-lemma-restriction-along-monomorphism-surjective}).
따라서 $\mathcal{J}_{n + 1} \to \mathcal{J}_n$은 준층의 사상으로서 전사이다.

\medskip\noindent
함자 $i_n^{-1}$이 완전한 왼쪽 수반 $i_{n, !}$을 가짐에 유의하자.
실제로 $i_{n, !}\mathcal{F}$는 역계
$\ldots 0 \to 0 \to \mathcal{F} \to \ldots \to \mathcal{F}$이다.
따라서 Derived Categories, Lemma
\ref{derived-lemma-adjoint-preserve-K-injectives}에 의해 복합체
$i_n^{-1}\mathcal{I}^\bullet = \mathcal{I}_n^\bullet$은 K-단사이다.

\medskip\noindent
K-단사 복합체의 각 항을 단사대상으로 골랐으므로
$$
0 \to  \lim \mathcal{I}_n^\bullet \to \prod \mathcal{I}_n^\bullet
\to \prod \mathcal{I}_n^\bullet \to 0
$$
이 아벨 준층 복합체의 짧은 완전열이고, 따라서 아벨층 복합체의 짧은
완전열임을 얻는다. 더욱이 가운데와 오른쪽의 곱들은
$D(\mathcal{C})$ 안의 곱들을 나타낸다. Injectives, Lemma
\ref{injectives-lemma-derived-products}와 그 증명을 보라
(여기서 $\mathcal{I}_n^\bullet$이 K-단사임을 사용한다).
따라서 삼각범주에서 호모토피 극한의 정의에 의해
$R\lim K$는 역계 $(K_n)$의 호모토피 극한이다.
\end{proof}

\begin{lemma}
\label{sites-cohomology-lemma-RGamma-commutes-with-Rlim}
$(\mathcal{C}, \mathcal{O})$를 환 달린 사이트라 하자.
$U \in \Ob(\mathcal{C})$에 대한 함자 $R\Gamma(\mathcal{C}, -)$와
$R\Gamma(U, -)$는 $R\lim$과 가환한다. 더욱이 $D(\mathcal{O})$ 안의
임의의 역계 $(K_n)$과 $m \in \mathbf{Z}$에 대하여 짧은 완전열
$$
0 \to
R^1\lim H^{m - 1}(U, K_n) \to H^m(U, R\lim K_n) \to
\lim H^m(U, K_n) \to 0
$$
이 있다. $H^m(\mathcal{C}, R\lim K_n)$에도 같은 결론이 성립한다.
\end{lemma}

\begin{proof}
첫째 명제는 Injectives, Lemma
\ref{injectives-lemma-RF-commutes-with-Rlim}에서 따른다. 그런 다음
$R\lim R\Gamma(U, K_n) = R\Gamma(U, R\lim K_n)$에
More on Algebra, Remark \ref{more-algebra-remark-compare-derived-limit}을
적용하면 짧은 완전열들을 얻는다.
\end{proof}

\begin{lemma}
\label{sites-cohomology-lemma-Rf-commutes-with-Rlim}
$f : (\Sh(\mathcal{C}), \mathcal{O}) \to (\Sh(\mathcal{C}'), \mathcal{O}')$를
환 달린 토포스의 사상이라 하자. 그러면 $Rf_*$는 $R\lim$과 가환한다.
즉 $Rf_*$는 유도 극한과 가환한다.
\end{lemma}

\begin{proof}
$(K_n)$을 $D(\mathcal{O})$의 대상들의 역계라 하자. $n$에 대한 귀납법으로
$\mathcal{O}$-가군의 실제 복합체 $\mathcal{K}_n^\bullet$과
$D(\mathcal{O})$ 안의 사상 $K_{n + 1} \to K_n$을 나타내는 복합체 사상
$\mathcal{K}_{n + 1}^\bullet \to \mathcal{K}_n^\bullet$을 고를 수 있다.
달리 말해, 결부된 역계가 주어진 역계인 대상 $K$가
$D(\mathcal{C} \times \mathbf{N})$ 안에 존재한다. 이제 토포스의
사상들의 가환도식
$$
\xymatrix{
\Sh(\mathcal{C} \times \mathbf{N}) \ar[r]_g \ar[d]_{f \times 1} &
\Sh(\mathcal{C}) \ar[d]_f \\
\Sh(\mathcal{C}' \times \mathbf{N}) \ar[r]^{g'} &
\Sh(\mathcal{C}')
}
$$
을 생각하자. 그러면
$R\lim R(f \times 1)_*K = Rf_* R\lim K$이다. 정의를 풀어 쓰고
Lemma \ref{sites-cohomology-lemma-derived-limit-is-ok}을 사용하면
$R\lim (Rf_*K_n) = Rf_*(R\lim K_n)$을 얻는다.

\medskip\noindent
$\mathcal{C}$가 충분한 점을 가지는 경우의 다른 증명을 제시한다. 정의
특수 삼각형
$$
R\lim K_n \to \prod K_n \to \prod K_n
$$
을 $D(\mathcal{O})$ 안에서 생각하자. 완전 함자 $Rf_*$를 적용하면
$D(\mathcal{O}')$ 안의 특수 삼각형
$$
Rf_*(R\lim K_n) \to Rf_*\left(\prod K_n\right) \to Rf_*\left(\prod K_n\right)
$$
을 얻는다. 따라서 $Rf_*$가 유도범주에서의 곱과 가환함을 증명하면
충분하다(그 곱은 복합체의 곱으로 단순히 주어지지 않는다. Injectives,
Lemma \ref{injectives-lemma-derived-products}를 보라).
그러나 Lemma \ref{sites-cohomology-lemma-adjoint}에 의해 $Rf_*$는 오른쪽 수반이므로
이는 형식적으로 따른다(Categories, Lemma
\ref{categories-lemma-adjoint-exact} 참조).
주의: $R\lim K_n$은 $D(\mathcal{O})$ 안의 극한이 아니므로 Categories,
Lemma \ref{categories-lemma-adjoint-exact}을 직접 적용할 수는 없다.
\end{proof}

\begin{remark}
\label{sites-cohomology-remark-discuss-derived-limit}
$(\mathcal{C}, \mathcal{O})$를 환 달린 사이트라 하고,
$(K_n)$을 $D(\mathcal{O})$ 안의 역계라 하자. $K = R\lim K_n$으로 놓자.
각 $n$과 $m$에 대하여 $\mathcal{H}^m_n = H^m(K_n)$을 $K_n$의
$m$차 코호몰로지 층이라 하고, 마찬가지로
$\mathcal{H}^m = H^m(K)$으로 놓자. 준층
$$
U \longmapsto \underline{\mathcal{H}}^m_n(U) = H^m(U, K_n)
$$
을 $\underline{\mathcal{H}}^m_n$이라 쓰자. 마찬가지로
$\underline{\mathcal{H}}^m(U) = H^m(U, K)$로 놓는다. Lemma
\ref{sites-cohomology-lemma-sheafification-cohomology}에 의해
$\mathcal{H}^m_n$은 $\underline{\mathcal{H}}^m_n$의 층화이고,
$\mathcal{H}^m$은 $\underline{\mathcal{H}}^m$의 층화이다.
다음 도식이 있다.
$$
\xymatrix{
K \ar@{=}[d] &
\underline{\mathcal{H}}^m \ar[d] \ar[r] & 
\mathcal{H}^m \ar[d] \\
R\lim K_n &
\lim \underline{\mathcal{H}}^m_n \ar[r] & 
\lim \mathcal{H}^m_n
}
$$
일반적으로 $\lim \mathcal{H}^m_n$이
$\lim \underline{\mathcal{H}}^m_n$의 층화인 것은 아닐 수 있다.
$U \in \mathcal{C}$이면 Lemma
\ref{sites-cohomology-lemma-RGamma-commutes-with-Rlim}에 의해 짧은 완전열
\begin{equation}
\label{sites-cohomology-equation-ses-Rlim-over-U}
0 \to
R^1\lim \underline{\mathcal{H}}^{m - 1}_n(U) \to
\underline{\mathcal{H}}^m(U) \to
\lim \underline{\mathcal{H}}^m_n(U) \to 0
\end{equation}
을 얻는다.
\end{remark}

\noindent
다음 보조정리는 대수공간이나 대수스택 위에서 전이 사상이 전사인
준연접 가군의 역계에 적용된다.

\begin{lemma}
\label{sites-cohomology-lemma-inverse-limit-is-derived-limit}
$(\mathcal{C}, \mathcal{O})$를 환 달린 사이트라 하고,
$(\mathcal{F}_n)$을 $\mathcal{O}$-가군의 역계라 하자.
$\mathcal{B} \subset \Ob(\mathcal{C})$를 부분집합이라 하자. 다음을 가정하자.
\begin{enumerate}
\item $\mathcal{C}$의 모든 대상은 그 원소들이 $\mathcal{B}$에 속하는
덮개를 갖는다.
\item $p > 0$ 및 $U \in \mathcal{B}$에 대하여
$H^p(U, \mathcal{F}_n) = 0$이다.
\item $U \in \mathcal{B}$에 대하여 역계 $\mathcal{F}_n(U)$의
$R^1\lim$은 영이다.
\end{enumerate}
그러면 $R\lim \mathcal{F}_n = \lim \mathcal{F}_n$이고,
$p > 0$ 및 $U \in \mathcal{B}$에 대하여
$H^p(U, \lim \mathcal{F}_n) = 0$이다.
\end{lemma}

\begin{proof}
$K_n = \mathcal{F}_n$ 및 $K = R\lim \mathcal{F}_n$으로 놓자.
Remark \ref{sites-cohomology-remark-discuss-derived-limit}의 표기와 가정 (2)에 의해
$U \in \mathcal{B}$에 대하여 $m \not = 0$이면
$\underline{\mathcal{H}}_n^m(U) = 0$이고,
$\underline{\mathcal{H}}_n^0(U) = \mathcal{F}_n(U)$이다.
Equation (\ref{sites-cohomology-equation-ses-Rlim-over-U})과 가정 (3)으로부터
$m \not = 0$이면 $\underline{\mathcal{H}}^m(U) = 0$이고,
$m = 0$이면 $\lim \mathcal{F}_n(U)$와 같음을 알 수 있다.
(1)을 사용하여 층화하면 $m \not = 0$이면
$\mathcal{H}^m = 0$이고, $m = 0$이면 $\lim \mathcal{F}_n$과 같음을
얻는다. 따라서 $K = \lim \mathcal{F}_n$이다.
$H^m(U, K) = \underline{\mathcal{H}}^m(U) = 0$이 $m > 0$에 대하여
성립하므로(위 참조) 둘째 명제도 성립한다.
\end{proof}

\begin{lemma}
\label{sites-cohomology-lemma-cohomology-derived-limit-injective}
$(\mathcal{C}, \mathcal{O})$를 환 달린 사이트라 하고, $(K_n)$을
$D(\mathcal{O})$ 안의 역계라 하자. $V \in \Ob(\mathcal{C})$ 및
$m \in \mathbf{Z}$라 하자. 정수 $n(V)$와 $V$의 덮개들의 공종계
$\text{Cov}_V$가 존재하여
$\{V_i \to V\} \in \text{Cov}_V$에 대하여 다음이 성립한다고 가정하자.
\begin{enumerate}
\item $R^1\lim H^{m - 1}(V_i, K_n) = 0$이다.
\item $n \geq n(V)$에 대하여
$H^m(V_i, K_n) \to H^m(V_i, K_{n(V)})$가 단사이다.
\end{enumerate}
그러면 단면 위의 사상
$H^m(R\lim K_n)(V) \to H^m(K_{n(V)})(V)$는 단사이다.
\end{lemma}

\begin{proof}
$\gamma \in H^m(R\lim K_n)(V)$가 $H^m(K_{n(V)})(V)$에서 영으로 간다고
하자. $H^m(R\lim K_n)$은 준층 $U \mapsto H^m(U, R\lim K_n)$의
층화이므로(Lemma \ref{sites-cohomology-lemma-sheafification-cohomology}에 의한다),
$\{V_i \to V\} \in \text{Cov}_V$와 $\gamma|_{V_i}$로 가는 원소
$\tilde\gamma_i \in H^m(V_i, R\lim K_n)$를 고를 수 있다.
그러면 $\tilde\gamma_i$는
$\tilde\gamma_{i, n(V)} \in H^m(V_i, K_{n(V)})$로 간다.
$H^m(K_{n(V)})$도 준층 $U \mapsto H^m(U, K_{n(V)})$의 층화이므로
(다시 Lemma \ref{sites-cohomology-lemma-sheafification-cohomology}에 의한다),
$\{V_i \to V\}$를 세분으로 바꾼 뒤 모든 $i$에 대해
$\tilde\gamma_{i, n(V)} = 0$이라고 가정해도 된다.
이 덮개에 대하여 Lemma \ref{sites-cohomology-lemma-RGamma-commutes-with-Rlim}의
짧은 완전열
$$
0 \to
R^1\lim H^{m - 1}(V_i, K_n) \to H^m(V_i, R\lim K_n) \to
\lim H^m(V_i, K_n) \to 0
$$
을 생각하자. 가정 (1)에 의해 왼쪽 군은 영이고, 가정 (2)에 의해
오른쪽 군은 $H^m(V_i, K_{n(V)})$ 안으로 단사로 간다.
따라서 $\tilde\gamma_i = 0$이고, 원하는 대로 $\gamma = 0$이다.
\end{proof}

\begin{lemma}
\label{sites-cohomology-lemma-is-limit-per-object}
$(\mathcal{C}, \mathcal{O})$를 환 달린 사이트라 하고,
$E \in D(\mathcal{O})$라 하자.
$\mathcal{B} \subset \Ob(\mathcal{C})$를 부분집합이라 하자. 다음을 가정하자.
\begin{enumerate}
\item $\mathcal{C}$의 모든 대상은 그 원소들이 $\mathcal{B}$에 속하는
덮개를 갖는다.
\item 모든 $V \in \mathcal{B}$에 대하여 함수
$p(V, -) : \mathbf{Z} \to \mathbf{Z}$와 $V$의 덮개들의 공종계
$\text{Cov}_V$가 존재하여
$$
H^p(V_i, H^{m - p}(E)) = 0
$$
$\{V_i \to V\} \in \text{Cov}_V$ 및 $p > p(V, m)$을 만족하는
모든 정수 $p, m$에 대해 성립한다.
\end{enumerate}
그러면 Derived Categories, Remark
\ref{derived-remark-map-into-derived-limit-truncations}의 사상
$E \to R\lim \tau_{\geq -n} E$는 $D(\mathcal{O})$ 안의 동형이다.
\end{lemma}

\begin{proof}
$K_n = \tau_{\geq -n}E$ 및 $K = R\lim K_n$으로 놓자.
표준 사상 $E \to K$는 표준 사상
$E \to K_n = \tau_{\geq -n}E$들에서 온다. $E \to K$가 코호몰로지
층들의 동형 $H^m(E) \to H^m(K)$을 유도함을 보여야 한다.
나머지 증명에서는 $m$을 고정한다. $n \geq -m$이면 사상
$E \to \tau_{\geq -n}E = K_n$은 동형
$H^m(E) \to H^m(K_n)$을 유도한다.
증명을 마치려면 모든 $V \in \mathcal{B}$에 대하여 사상
$H^m(K)(V) \to H^m(K_{n(V)})(V)$가 단사가 되게 하는 정수
$n(V) \geq -m$이 존재함을 보이면 충분하다. 실제로 그러면 합성
$$
H^m(E)(V) \to H^m(K)(V) \to H^m(K_{n(V)})(V)
$$
은 전단사이고 둘째 화살표는 단사이므로 첫째 화살표는 전단사이다.
조건 (1)에 의해 이는 $H^m(E) \to H^m(K)$가 동형임을 함의한다.
다음과 같이 놓자.
$$
n(V) = 1 + \max\{-m, p(V, m - 1) - m, -1 + p(V, m) - m, -2 + p(V, m + 1) - m\}.
$$
그러면 언제나 $n(V) \geq -m$이다. 주장:
$n \geq n(V)$ 및 $\{V_i \to V\} \in \text{Cov}_V$에 대하여 사상
$$
H^{m - 1}(V_i, K_{n + 1}) \to H^{m - 1}(V_i, K_n)
\quad\text{and}\quad
H^m(V_i, K_{n + 1}) \to H^m(V_i, K_n)
$$
은 동형이다. 이 주장은 Lemma
\ref{sites-cohomology-lemma-cohomology-derived-limit-injective}의 조건 (1)과 (2)가
성립함을 뜻하며, 따라서 원하는 단사성을 함의한다.
다음 특수 삼각형들이 있음을 상기하자(Derived Categories, Remark
\ref{derived-remark-truncation-distinguished-triangle}).
$$
H^{-n - 1}(E)[n + 1] \to
K_{n + 1} \to K_n \to H^{-n - 1}(E)[n + 2]
$$
이에 연관된 긴 완전 코호몰로지 열을 보면,
$n \geq n(V)$ 및 $\{V_i \to V\} \in \text{Cov}_V$에 대하여
$$
H^{m + n}(V_i, H^{-n - 1}(E)),\quad
H^{m + n + 1}(V_i, H^{-n - 1}(E)),\quad
H^{m + n + 2}(V_i, H^{-n - 1}(E))
$$
가 영이면 주장이 따른다. 이는 $n(V)$의 선택과 보조정리의 가정에서 따른다.
\end{proof}

\begin{lemma}
\label{sites-cohomology-lemma-is-limit-spaltenstein}
$(\mathcal{C}, \mathcal{O})$를 환 달린 사이트라 하고,
$E \in D(\mathcal{O})$라 하자.
$\mathcal{B} \subset \Ob(\mathcal{C})$를 부분집합이라 하자. 다음을 가정하자.
\begin{enumerate}
\item $\mathcal{C}$의 모든 대상은 그 원소들이 $\mathcal{B}$에 속하는
덮개를 갖는다.
\item 모든 $V \in \mathcal{B}$에 대하여 정수 $d_V \geq 0$와
$V$의 덮개들의 공종계 $\text{Cov}_V$가 존재하여
$$
H^p(V_i, H^q(E)) = 0 \text{ for }
\{V_i \to V\} \in \text{Cov}_V,\ p > d_V, \text{ and }q < 0
$$
이 성립한다.
\end{enumerate}
그러면 Derived Categories, Remark
\ref{derived-remark-map-into-derived-limit-truncations}의 사상
$E \to R\lim \tau_{\geq -n} E$는 $D(\mathcal{O})$ 안의 동형이다.
\end{lemma}

\begin{proof}
이는 $p(V, m) = d_V + \max(0, m)$로 둔
Lemma \ref{sites-cohomology-lemma-is-limit-per-object}에서 따른다.
\end{proof}

\begin{lemma}
\label{sites-cohomology-lemma-is-limit}
$(\mathcal{C}, \mathcal{O})$를 환 달린 사이트라 하고,
$E \in D(\mathcal{O})$라 하자. 함수
$p(-) : \mathbf{Z} \to \mathbf{Z}$와 부분집합
$\mathcal{B} \subset \Ob(\mathcal{C})$가 존재하여 다음을 만족한다고
가정하자.
\begin{enumerate}
\item $\mathcal{C}$의 모든 대상은 그 원소들이 $\mathcal{B}$에 속하는
덮개를 갖는다.
\item $p > p(m)$ 및 $V \in \mathcal{B}$에 대하여
$H^p(V, H^{m - p}(E)) = 0$이다.
\end{enumerate}
그러면 Derived Categories, Remark
\ref{derived-remark-map-into-derived-limit-truncations}의 사상
$E \to R\lim \tau_{\geq -n} E$는 $D(\mathcal{O})$ 안의 동형이다.
\end{lemma}

\begin{proof}
$p(V, m) = p(m)$으로 놓고, $\text{Cov}_V$를 모든 $i$에 대해
$V_i \in \mathcal{B}$인 덮개 $\{V_i \to V\}$들의 집합으로 두어
Lemma \ref{sites-cohomology-lemma-is-limit-per-object}을 적용한다.
\end{proof}

\begin{lemma}
\label{sites-cohomology-lemma-is-limit-dimension}
$(\mathcal{C}, \mathcal{O})$를 환 달린 사이트라 하고,
$E \in D(\mathcal{O})$라 하자. 정수 $d \geq 0$와 부분집합
$\mathcal{B} \subset \Ob(\mathcal{C})$가 존재하여 다음을 만족한다고
가정하자.
\begin{enumerate}
\item $\mathcal{C}$의 모든 대상은 그 원소들이 $\mathcal{B}$에 속하는
덮개를 갖는다.
\item $p > d$, $q < 0$, $V \in \mathcal{B}$에 대하여
$H^p(V, H^q(E)) = 0$이다.
\end{enumerate}
그러면 Derived Categories, Remark
\ref{derived-remark-map-into-derived-limit-truncations}의 사상
$E \to R\lim \tau_{\geq -n} E$는 $D(\mathcal{O})$ 안의 동형이다.
\end{lemma}

\begin{proof}
$d_V = d$로 놓고, $\text{Cov}_V$를 모든 $i$에 대해
$V_i \in \mathcal{B}$인 덮개 $\{V_i \to V\}$들의 집합으로 두어
Lemma \ref{sites-cohomology-lemma-is-limit-spaltenstein}을 적용한다.
\end{proof}

\noindent
위의 보조정리들을 사용하면 어떤 상황에서 코호몰로지를 계산할 수 있다.

\begin{lemma}
\label{sites-cohomology-lemma-cohomology-over-U-trivial}
$(\mathcal{C}, \mathcal{O})$를 환 달린 사이트라 하고,
$K$를 $D(\mathcal{O})$의 대상이라 하자.
$\mathcal{B} \subset \Ob(\mathcal{C})$를 부분집합이라 하자. 다음을 가정하자.
\begin{enumerate}
\item $\mathcal{C}$의 모든 대상은 그 원소들이 $\mathcal{B}$에 속하는
덮개를 갖는다.
\item 모든 $p > 0$, $q \in \mathbf{Z}$, $U \in \mathcal{B}$에 대하여
$H^p(U, H^q(K)) = 0$이다.
\end{enumerate}
그러면 $q \in \mathbf{Z}$ 및 $U \in \mathcal{B}$에 대하여
$H^q(U, K) = H^0(U, H^q(K))$이다.
\end{lemma}

\begin{proof}
$d = 0$으로 둔 Lemma \ref{sites-cohomology-lemma-is-limit-dimension}에 의해
$K = R\lim \tau_{\geq -n} K$임을 관찰하자.
$U \in \mathcal{B}$라 하자. Equation
(\ref{sites-cohomology-equation-ses-Rlim-over-U})에 의해 짧은 완전열
$$
0 \to R^1\lim H^{q - 1}(U, \tau_{\geq -n}K) \to
H^q(U, K) \to \lim H^q(U, \tau_{\geq -n}K) \to 0
$$
을 얻는다. Derived Categories, Lemma
\ref{derived-lemma-two-ss-complex-functor}의 스펙트럼열을 사용하면
조건 (2)는 모든 $q$에 대하여
$H^q(U, \tau_{\geq -n} K) = H^0(U, H^q(\tau_{\geq -n} K))$
임을 함의한다. $\tau_{\geq -n}K$는 아래로 유계이므로 스펙트럼열은
수렴한다. $n > -q$이면 $H^q(\tau_{\geq -n}K) = H^q(K)$이다.
따라서 표시한 짧은 완전열의 왼쪽 계와 오른쪽 계는 각각
$H^0(U, H^{q - 1}(K))$와 $H^0(U, H^q(K))$를 값으로 하여 결국
상수이다. 보조정리가 따른다.
\end{proof}

\noindent
유도 극한을 기술할 수 있는 또 다른 경우가 있다.

\begin{lemma}
\label{sites-cohomology-lemma-derived-limit-suitable-system}
$(\mathcal{C}, \mathcal{O})$를 환 달린 사이트라 하고, $(K_n)$을
$D(\mathcal{O})$의 대상들의 역계라 하자.
$\mathcal{B} \subset \Ob(\mathcal{C})$를 부분집합이라 하자. 다음을 가정하자.
\begin{enumerate}
\item $\mathcal{C}$의 모든 대상은 그 원소들이 $\mathcal{B}$에 속하는
덮개를 갖는다.
\item 모든 $U \in \mathcal{B}$와 모든 $q \in \mathbf{Z}$에 대하여
\begin{enumerate}
\item $p > 0$에 대하여 $H^p(U, H^q(K_n)) = 0$이다.
\item 역계 $H^0(U, H^q(K_n))$의 $R^1\lim$은 영이다.
\end{enumerate}
\end{enumerate}
그러면 $q \in \mathbf{Z}$에 대하여
$H^q(R\lim K_n) = \lim H^q(K_n)$이다.
\end{lemma}

\begin{proof}
$K = R\lim K_n$으로 놓고 Remark
\ref{sites-cohomology-remark-discuss-derived-limit}의 표기를 사용하자.
$U \in \mathcal{B}$라 하자.
Lemma \ref{sites-cohomology-lemma-cohomology-over-U-trivial}과 (2)(a)에 의해
$H^q(U, K_n) = H^0(U, H^q(K_n))$이다.
함자 $R\Gamma(U, -)$가 유도 극한과 가환한다는 사실을 사용하면
$$
H^q(U, K) = H^q(R\lim R\Gamma(U, K_n)) = \lim H^0(U, H^q(K_n))
$$
이다. 여기서 마지막 등식은 More on Algebra, Remark
\ref{more-algebra-remark-compare-derived-limit}과 가정 (2)(b)에서 따른다.
따라서 $H^q(U, K)$는 $U$ 위에서의 층 $H^q(K_n)$의 단면들의 역극한이다.
$\lim H^q(K_n)$은 층이므로, 가정 (1)을 사용하면 준층
$U \mapsto H^q(U, K)$의 층화인 $H^q(K)$가
$\lim H^q(K_n)$과 같음을 얻는다. 이것이 보조정리를 증명한다.
\end{proof}







\section{K-단사 분해의 구성}
\label{sites-cohomology-section-K-injective}

\noindent
$(\mathcal{C}, \mathcal{O})$를 환 달린 사이트라 하자.
$\mathcal{F}^\bullet$을 $\mathcal{O}$-가군의 복합체라 하자.
범주 $\textit{Mod}(\mathcal{O})$는 충분한 단사대상을 가지므로
Derived Categories, Lemma
\ref{derived-lemma-special-inverse-system}을 사용하여 다음 도식을
구성할 수 있다.
$$
\xymatrix{
\ldots \ar[r] &
\tau_{\geq -2}\mathcal{F}^\bullet \ar[r] \ar[d] &
\tau_{\geq -1}\mathcal{F}^\bullet \ar[d] \\
\ldots \ar[r] & \mathcal{I}_2^\bullet \ar[r] & \mathcal{I}_1^\bullet
}
$$
이 도식은 $\mathcal{O}$-가군 복합체의 범주 안에 있으며 다음을 만족한다.
\begin{enumerate}
\item 세로 화살표들은 유사동형이다.
\item $\mathcal{I}_n^\bullet$은 단사대상들로 이루어진 아래로 유계인
복합체이다.
\item 화살표 $\mathcal{I}_{n + 1}^\bullet \to \mathcal{I}_n^\bullet$은
항별 분할 전사이다.
\end{enumerate}
$\mathcal{O}$-가군의 범주는 극한을 가지며(준층 수준에서 계산된다),
따라서 항별 극한
$\mathcal{I}^\bullet = \lim_n \mathcal{I}_n^\bullet$을 만들 수 있다.
Derived Categories, Lemmas
\ref{derived-lemma-bounded-below-injectives-K-injective} 및
\ref{derived-lemma-limit-K-injectives}
에 의해 이는 K-단사 복합체이다. 일반적으로 표준 사상
\begin{equation}
\label{sites-cohomology-equation-into-candidate-K-injective}
\mathcal{F}^\bullet \to \mathcal{I}^\bullet
\end{equation}
은 유사동형이 아닐 수 있다. 다음 보조정리에서는 이것이 유사동형이 되는
몇 가지 조건을 기술한다.

\begin{lemma}
\label{sites-cohomology-lemma-K-injective}
위에서 기술한 상황에 있다고 하자.
$\mathcal{H}^m = H^m(\mathcal{F}^\bullet)$을 $m$차 코호몰로지 층이라 쓰자.
$\mathcal{B} \subset \Ob(\mathcal{C})$를 부분집합이라 하고,
$d \in \mathbf{N}$이라 하자. 다음을 가정하자.
\begin{enumerate}
\item $\mathcal{C}$의 모든 대상은 그 원소들이 $\mathcal{B}$에 속하는
덮개를 갖는다.
\item 모든 $U \in \mathcal{B}$에 대하여 $p > d$ 및 $q < 0$이면
$H^p(U, \mathcal{H}^q) = 0$이다\footnote{
$\forall m$, $\exists p(m)$에 대하여 $p > p(m)$이면
$H^p(U. \mathcal{H}^{m - p}) = 0$인 것으로 충분하다. Lemma
\ref{sites-cohomology-lemma-is-limit}을 보라.}.
\end{enumerate}
그러면 (\ref{sites-cohomology-equation-into-candidate-K-injective})은 유사동형이다.
\end{lemma}

\begin{proof}
Derived Categories, Lemma
\ref{derived-lemma-difficulty-K-injectives}에 의해 사상
$\mathcal{F}^\bullet \to R\lim \tau_{\geq -n} \mathcal{F}^\bullet$
이 동형임을 보이면 충분하다. 이는 Lemma
\ref{sites-cohomology-lemma-is-limit-dimension}에서 따른다.
\end{proof}

\noindent
다음은 가군 복합체들의 역계의 항별 극한이 갖는 코호몰로지 층에 관한
기술적 보조정리이다. 이 보조정리의 사용은 되도록 피하고, 그 대신 유도
역극한을 이용한 논증을 사용해야 한다.

\begin{lemma}
\label{sites-cohomology-lemma-inverse-limit-complexes}
$(\mathcal{C}, \mathcal{O})$를 환 달린 사이트라 하고,
$(\mathcal{F}_n^\bullet)$을 $\mathcal{O}$-가군 복합체의 역계라 하자.
$m \in \mathbf{Z}$라 하자. 부분집합
$\mathcal{B} \subset \Ob(\mathcal{C})$와 정수 $n_0$가 주어져 다음을
만족한다고 하자.
\begin{enumerate}
\item $\mathcal{C}$의 모든 대상은 그 원소들이 $\mathcal{B}$에 속하는
덮개를 갖는다.
\item 모든 $U \in \mathcal{B}$에 대하여
\begin{enumerate}
\item 아벨군의 계 $\mathcal{F}_n^{m - 2}(U)$ 및
$\mathcal{F}_n^{m - 1}(U)$의 $R^1\lim$은 영이다
(예를 들어 이 계들은 Mittag--Leffler 조건을 만족한다).
\item 아벨군의 계 $H^{m - 1}(\mathcal{F}_n^\bullet(U))$의
$R^1\lim$은 영이다(예를 들어 이 계는 Mittag--Leffler 조건을 만족한다).
\item 모든 $n \geq n_0$에 대하여
$H^m(\mathcal{F}_n^\bullet(U)) = H^m(\mathcal{F}_{n_0}^\bullet(U))$
이다.
\end{enumerate}
\end{enumerate}
그러면 사상 $H^m(\mathcal{F}^\bullet) \to \lim H^m(\mathcal{F}_n^\bullet)
\to H^m(\mathcal{F}_{n_0}^\bullet)$들은 층의 동형이다. 여기서
$\mathcal{F}^\bullet = \lim \mathcal{F}_n^\bullet$은 항별 역극한이다.
\end{lemma}

\begin{proof}
$U \in \mathcal{B}$라 하자. $H^m(\mathcal{F}^\bullet(U))$은
$$
\lim_n \mathcal{F}_n^{m - 2}(U) \to
\lim_n \mathcal{F}_n^{m - 1}(U) \to
\lim_n \mathcal{F}_n^m(U) \to
\lim_n \mathcal{F}_n^{m + 1}(U)
$$
의 왼쪽에서 셋째 자리에서의 코호몰로지임을 주목하자.
가정 (2)(a)와 (2)(b)에 의해 More on Algebra, Lemma
\ref{more-algebra-lemma-apply-Mittag-Leffler-again}을 적용하여
$$
H^m(\mathcal{F}^\bullet(U)) = \lim H^m(\mathcal{F}_n^\bullet(U))
$$
을 얻는다. 가정 (2)(c)에 의해 모든 $n \geq n_0$에 대하여
$$
H^m(\mathcal{F}^\bullet(U)) = H^m(\mathcal{F}_n^\bullet(U))
$$
이다. 가정 (1)에 의해 모든 $n \geq n_0$에 대하여 준층
$U \mapsto H^m(\mathcal{F}^\bullet(U))$의 층화와
$U \mapsto H^m(\mathcal{F}_n^\bullet(U))$의 층화가 같음을 얻는다.
따라서 층의 역계 $H^m(\mathcal{F}_n^\bullet)$은 $n \geq n_0$에서
$H^m(\mathcal{F}^\bullet)$을 값으로 하여 상수이다. 이것이 보조정리를
증명한다.
\end{proof}







\section{유계 코호몰로지 차원}
\label{sites-cohomology-section-bounded}

\noindent
이 절에서는 함자 $Rf_*$가 언제 유계 코호몰로지 차원을 가지는지 묻는다.
비유계 복합체를 생각할 때 이는 상당히 미묘한 문제이다.

\begin{situation}
\label{sites-cohomology-situation-olsson-laszlo}
$\mathcal{C}$를 사이트라 하고, $\mathcal{O}$를 $\mathcal{C}$ 위의
환의 층이라 하자. $\mathcal{A} \subset \textit{Mod}(\mathcal{O})$를
약한 Serre 부분범주라 하자. 다음이 성립한다고 가정한다. 즉 부분집합
$\mathcal{B} \subset \Ob(\mathcal{C})$가 존재하여 다음을 만족한다.
\begin{enumerate}
\item $\mathcal{C}$의 모든 대상은 그 원소들이 $\mathcal{B}$에 속하는
덮개를 갖는다.
\item 모든 $V \in \mathcal{B}$에 대하여 정수 $d_V$와
$V$의 덮개들의 공종계 $\text{Cov}_V$가 존재하여
$$
H^p(V_i, \mathcal{F}) = 0 \text{ for }
\{V_i \to V\} \in \text{Cov}_V,\ p > d_V, \text{ and }
\mathcal{F} \in \Ob(\mathcal{A})
$$
이 성립한다.
\end{enumerate}
\end{situation}

\begin{lemma}
\label{sites-cohomology-lemma-olsson-laszlo}
\begin{reference}
이는 가정을 조금 바꾼 \cite[Proposition 2.1.4]{six-I}이며, 사이트에 대한
\cite[Proposition 3.13]{Spaltenstein}의 유사물이다.
\end{reference}
Situation \ref{sites-cohomology-situation-olsson-laszlo}에서 임의의
$E \in D_\mathcal{A}(\mathcal{O})$에 대하여 Derived Categories, Remark
\ref{derived-remark-map-into-derived-limit-truncations}의 사상
$E \to R\lim \tau_{\geq -n} E$는 $D(\mathcal{O})$ 안의 동형이다.
\end{lemma}

\begin{proof}
Lemma \ref{sites-cohomology-lemma-is-limit-spaltenstein}에서 곧바로 따른다.
\end{proof}

\begin{lemma}
\label{sites-cohomology-lemma-olsson-laszlo-modified}
Situation \ref{sites-cohomology-situation-olsson-laszlo}에서 $(K_n)$을
$D_\mathcal{A}^+(\mathcal{O})$ 안의 역계라 하자. 모든 $j$에 대하여
$\mathcal{A}$ 안의 역계 $(H^j(K_n))$이 결국 $\mathcal{H}^j$을 값으로
하여 상수라고 가정하자. 그러면 모든 $j$에 대하여
$H^j(R\lim K_n) = \mathcal{H}^j$이다.
\end{lemma}

\begin{proof}
$V \in \mathcal{B}$라 하자. $\{V_i \to V\}$를 Situation
\ref{sites-cohomology-situation-olsson-laszlo}의 집합 $\text{Cov}_V$에 속하는 덮개라 하자.
$K_n$은 아래로 유계이므로
$$
E_2^{p, q} = H^p(V_i, H^q(K_n))
$$
에서 $H^{p + q}(V_i, K_n)$으로 수렴하는 스펙트럼열이 있다.
Derived Categories, Lemma \ref{derived-lemma-two-ss-complex-functor}을 보라.
가정에 의해 $p > d_V$이면 $E_2^{p, q} = 0$임에 유의하자. 모든
$n \geq n_0$에 대하여 다음이 성립하도록 $n_0$를 고르자.
$$
\begin{matrix}
\mathcal{H}^{j + 1} & = & H^{j + 1}(K_n), \\
\mathcal{H}^j & = & H^j(K_n), \\
\ldots, \\
\mathcal{H}^{j - d_V - 2} & = & H^{j - d_V - 2}(K_n)
\end{matrix}
$$
$K_n$과 $K_{n_0}$에 대한 위의 스펙트럼열들을 비교하면,
$n \geq n_0$일 때 코호몰로지 군 $H^{j - 1}(V_i, K_n)$과
$H^j(V_i, K_n)$은 $n$과 무관함을 알 수 있다. 따라서 $n$이 충분히
크면($V$에 따라 달라진다) 단면 위의 사상
$H^j(R\lim K_n)(V) \to H^j(K_n)(V)$은 단사이다.
Lemma \ref{sites-cohomology-lemma-cohomology-derived-limit-injective}을 보라.
$\mathcal{C}$의 모든 대상은 $\mathcal{B}$의 원소들로 덮을 수 있으므로
사상 $H^j(R\lim K_n) \to \mathcal{H}^j$는 단사이다.

\medskip\noindent
전사성도 비슷하게 보인다. 즉 $U \in \Ob(\mathcal{C})$와
$\gamma \in \mathcal{H}^j(U)$를 고르자. $U$를 덮개의 원소들로 바꾼 뒤
$\gamma$를 $H^j(R\lim K_n)$의 단면으로 올리고자 한다. 따라서
Situation \ref{sites-cohomology-situation-olsson-laszlo}의 성질 (1)에 의해
$U = V \in \mathcal{B}$라고 가정해도 된다. 모든 $n \geq n_0$에 대하여
다음이 성립하도록 $n_0$를 고르자.
$$
\begin{matrix}
\mathcal{H}^{j + 1} & = & H^{j + 1}(K_n), \\
\mathcal{H}^j & = & H^j(K_n), \\
\ldots, \\
\mathcal{H}^{j - d_V - 2} & = & H^{j - d_V - 2}(K_n)
\end{matrix}
$$
다음을 만족하는 $\text{Cov}_V$의 원소 $\{V_i \to V\}$를 고르자.
$\gamma|_{V_i} \in \mathcal{H}^j(V_i) = H^j(K_{n_0})(V_i)$
가 원소 $\gamma_{n_0, i} \in H^j(V_i, K_{n_0})$으로 올라가도록 한다.
Lemma \ref{sites-cohomology-lemma-sheafification-cohomology}에 의해 $H^j(K_{n_0})$은
$U \mapsto H^j(U, K_{n_0})$의 층화이므로 이는 가능하다.
증명의 첫 문단의 논의에 의해 $H^{j - 1}(V_i, K_n)$과
$H^j(V_i, K_n)$은 $n \geq n_0$과 무관하다. 따라서 Lemma
\ref{sites-cohomology-lemma-RGamma-commutes-with-Rlim}에 의해 $\gamma_{n_0, i}$는
원소 $\gamma_i \in H^j(V_i, R\lim K_n)$으로 올라간다.
이것으로 증명을 마친다.
\end{proof}

\begin{lemma}
\label{sites-cohomology-lemma-olsson-laszlo-map-version-one}
\begin{reference}
이는 가정을 조금 바꾼 \cite[Lemma 2.1.10]{six-I}의 한 형태이다.
\end{reference}
$f : (\Sh(\mathcal{C}), \mathcal{O}) \to (\Sh(\mathcal{C}'), \mathcal{O}')$를
환 달린 토포스의 사상이라 하자.
$\mathcal{A} \subset \textit{Mod}(\mathcal{O})$와
$\mathcal{A}' \subset \textit{Mod}(\mathcal{O}')$를 약한 Serre
부분범주라 하자. 정수 $N$이 존재하여 다음을 만족한다고 가정하자.
\begin{enumerate}
\item $\mathcal{C}, \mathcal{O}, \mathcal{A}$는 Situation
\ref{sites-cohomology-situation-olsson-laszlo}의 가정을 만족한다.
\item $\mathcal{C}', \mathcal{O}', \mathcal{A}'$는 Situation
\ref{sites-cohomology-situation-olsson-laszlo}의 가정을 만족한다.
\item $p \geq 0$ 및 $\mathcal{F} \in \Ob(\mathcal{A})$에 대하여
$R^pf_*\mathcal{F} \in \Ob(\mathcal{A}')$이다.
\item $p > N$ 및 $\mathcal{F} \in \Ob(\mathcal{A})$에 대하여
$R^pf_*\mathcal{F} = 0$이다.
\end{enumerate}
그러면 $D_\mathcal{A}(\mathcal{O})$ 안의 $K$에 대하여 다음이 성립한다.
\begin{enumerate}
\item[(a)] $Rf_*K$는 $D_{\mathcal{A}'}(\mathcal{O}')$에 속한다.
\item[(b)] $j \geq N - n$에 대하여 사상
$H^j(Rf_*K) \to H^j(Rf_*(\tau_{\geq -n}K))$은 동형이다.
\end{enumerate}
\end{lemma}

\begin{proof}
Lemma \ref{sites-cohomology-lemma-olsson-laszlo}에 의해
$K = R\lim \tau_{\geq -n}K$이고, Lemma
\ref{sites-cohomology-lemma-Rf-commutes-with-Rlim}에 의해
$Rf_*K = R\lim Rf_*\tau_{\geq -n}K$이다.
복합체 $Rf_*\tau_{\geq -n}K$는 아래로 유계이다. 스펙트럼열
$$
E_2^{p, q} = R^pf_*H^q(\tau_{\geq -n}K)
$$
은 $H^{p + q}(Rf_*\tau_{\geq -n}K)$으로 수렴한다
(Derived Categories, Lemma \ref{derived-lemma-two-ss-complex-functor}).
이 스펙트럼열과 가정 (3)에 의해
$Rf_*\tau_{\geq -n}K$는 $D^+_{\mathcal{A}'}(\mathcal{O}')$에 속한다.
Homology, Lemma \ref{homology-lemma-biregular-ss-converges}를 보라.
$m \geq n$일 때 사상
$$
Rf_*(\tau_{\geq -m}K) \longrightarrow Rf_*(\tau_{\geq -n}K)
$$
은 위의 스펙트럼열들에 의해 차수 $j \geq -n + N$에서 코호몰로지 층의
동형을 유도함에 유의하자. 따라서 Lemma
\ref{sites-cohomology-lemma-olsson-laszlo-modified}을 적용하면 결론을 얻는다.
\end{proof}

\noindent
가정을 조금 달리한 위 보조정리의 변형이 때때로 필요하다.

\begin{situation}
\label{sites-cohomology-situation-olsson-laszlo-prime}
$f : (\mathcal{C}, \mathcal{O}) \to (\mathcal{C}', \mathcal{O}')$를
환 달린 사이트의 사상이라 하고, $u : \mathcal{C}' \to \mathcal{C}$를
그에 대응하는 사이트의 연속 함자라 하자.
$\mathcal{A} \subset \textit{Mod}(\mathcal{O})$를 약한 Serre
부분범주라 하자. 다음이 성립한다고 가정한다. 즉 부분집합
$\mathcal{B}' \subset \Ob(\mathcal{C}')$가 존재하여 다음을 만족한다.
\begin{enumerate}
\item $\mathcal{C}'$의 모든 대상은 그 원소들이 $\mathcal{B}'$에 속하는
덮개를 갖는다.
\item 모든 $V' \in \mathcal{B}'$에 대하여 정수 $d_{V'}$와
$V'$의 덮개들의 공종계 $\text{Cov}_{V'}$가 존재하여
$$
H^p(u(V'_i), \mathcal{F}) = 0 \text{ for }
\{V'_i \to V'\} \in \text{Cov}_{V'},\ p > d_{V'}, \text{ and }
\mathcal{F} \in \Ob(\mathcal{A})
$$
이 성립한다.
\end{enumerate}
\end{situation}

\begin{lemma}
\label{sites-cohomology-lemma-olsson-laszlo-map-version-two}
\begin{reference}
이는 가정을 조금 바꾼 \cite[Lemma 2.1.10]{six-I}의 한 형태이다.
\end{reference}
$f : (\mathcal{C}, \mathcal{O}) \to (\mathcal{C}', \mathcal{O}')$를
환 달린 사이트의 사상이라 하자. 또한 정수 $N$이 존재하여 다음을
만족한다고 가정하자.
\begin{enumerate}
\item $\mathcal{C}, \mathcal{O}, \mathcal{A}$는 Situation
\ref{sites-cohomology-situation-olsson-laszlo}의 가정을 만족한다.
\item $f : (\mathcal{C}, \mathcal{O}) \to (\mathcal{C}', \mathcal{O}')$와
$\mathcal{A}$는 Situation \ref{sites-cohomology-situation-olsson-laszlo-prime}의
가정을 만족한다.
\item $p > N$ 및 $\mathcal{F} \in \Ob(\mathcal{A})$에 대하여
$R^pf_*\mathcal{F} = 0$이다.
\end{enumerate}
그러면 $D_\mathcal{A}(\mathcal{O})$ 안의 $K$에 대하여
$j \geq N - n$이면 사상
$H^j(Rf_*K) \to H^j(Rf_*(\tau_{\geq -n}K))$은 동형이다.
\end{lemma}

\begin{proof}
$K$를 $D_\mathcal{A}(\mathcal{O})$의 대상이라 하자. Lemma
\ref{sites-cohomology-lemma-olsson-laszlo}에 의해
$K = R\lim \tau_{\geq -n}K$이고, Lemma
\ref{sites-cohomology-lemma-Rf-commutes-with-Rlim}에 의해
$Rf_*K = R\lim Rf_*(\tau_{\geq -n}K)$이다.
$V' \in \mathcal{B}'$라 하고, $\{V'_i \to V'\}$를
$\text{Cov}_{V'}$의 원소라 하자. 다음을 생각한다.
$$
H^j(V'_i, Rf_*K) = H^j(u(V'_i), K)
\quad\text{and}\quad
H^j(V'_i, Rf_*(\tau_{\geq -n}K)) = H^j(u(V'_i), \tau_{\geq -n}K)
$$
Situation \ref{sites-cohomology-situation-olsson-laszlo-prime}의 가정은 $j$와 $d_{V'}$에
따라 $n$이 충분히 클 때 마지막 군이 $n$과 무관함을 함의한다.
몇몇 세부사항은 생략한다. 이를 $j$와 $j - 1$에 적용하고 Lemma
\ref{sites-cohomology-lemma-RGamma-commutes-with-Rlim}을 사용하면
$$
H^j(V'_i, Rf_*K) = \lim H^j(V'_i, Rf_*(\tau_{\geq -n} K))
$$
이며, 오른쪽 계는 $d_{V'}$와 $j$에만 의존하는 어떤 상수보다 큰
$n$에 대해 상수임을 얻는다. 따라서 Lemma
\ref{sites-cohomology-lemma-cohomology-derived-limit-injective}은
$$
H^j(Rf_*K)(V') \longrightarrow
\left(\lim H^j(Rf_*(\tau_{\geq -n}K))\right)(V')
$$
이 단사임을 함의한다. 원소 $V' \in \mathcal{B}'$가
$\mathcal{C}'$의 모든 대상을 덮으므로 사상
$H^j(Rf_*K) \to \lim H^j(Rf_*(\tau_{\geq -n}K))$은 단사이다.
스펙트럼열
$$
E_2^{p, q} = R^pf_*H^q(\tau_{\geq -n}K)
$$
은 $H^{p + q}(Rf_*(\tau_{\geq -n}K))$으로 수렴한다
(Derived Categories, Lemma \ref{derived-lemma-two-ss-complex-functor}).
이 스펙트럼열과 가정 (3)에 의해
$H^j(Rf_*(\tau_{\geq -n}K))$는 $n \geq N - j$에 대해 상수이다.
따라서 $j \geq N - n$이면
$H^j(Rf_*K) \to H^j(Rf_*(\tau_{\geq -n}K))$은 단사이다.

\medskip\noindent
따라서 보조정리의 마지막 줄에서 ``동형''을 ``단사''로 바꾼 명제를
증명했다. 이제 $j$와 $n$을 고르고
$j \geq N - n$라고 하자. 다음 특수 삼각형을 생각하자.
$$
\tau_{\leq -n - 1}K \to K \to \tau_{\geq -n}K \to (\tau_{\leq -n - 1}K)[1]
$$
Derived Categories, Remark
\ref{derived-remark-truncation-distinguished-triangle}을 보라.
$\tau_{\geq -n}\tau_{\leq -n -1}K = 0$이므로,
$\tau_{-n - 1}K$에 대해 이미 증명한 단사성은
$$
0 = H^j(Rf_*(\tau_{\leq -n - 1}K)) =
H^{j + 1}(Rf_*(\tau_{\leq -n - 1}K)) =
H^{j + 2}(Rf_*(\tau_{\leq -n - 1}K)) = \ldots
$$
특수 삼각형
$$
Rf_*(\tau_{\leq -n - 1}K) \to Rf_*K \to Rf_*(\tau_{\geq -n}K) \to
Rf_*(\tau_{\leq -n - 1}K)[1]
$$
에 연관된 긴 완전 코호몰로지 열에 의해 이는
$H^j(Rf_*K) \to H^j(Rf_*(\tau_{\geq -n}K))$이 동형임을 함의한다.
\end{proof}








\section{Mayer-Vietoris}
\label{sites-cohomology-section-mayer-vietoris}

\noindent
통상적인 Mayer--Vietoris의 명제와 증명은 Cohomology, Section
\ref{cohomology-section-mayer-vietoris}을 보라.

\medskip\noindent
$(\mathcal{C}, \mathcal{O})$를 환 달린 사이트라 하고, 범주
$\mathcal{C}$ 안의 가환도식
$$
\xymatrix{
E \ar[d] \ar[r] & Y \ar[d] \\
Z \ar[r] & X
}
$$
을 생각하자. 이 상황에서 $D(\mathcal{O})$의 대상 $K$가 주어지면
특수 삼각형의 시작처럼 보이는 열
$$
R\Gamma(X, K) \to
R\Gamma(Z, K) \oplus
R\Gamma(Y, K) \to
R\Gamma(E, K)
$$
을 얻는다. 다음 보조정리에서 이를 더 정확히 표현한다.

\begin{lemma}
\label{sites-cohomology-lemma-c-square}
위의 상황에서 $K$를 나타내는 $\mathcal{O}$-가군의 K-단사 복합체
$\mathcal{I}^\bullet$을 고르자. 네 화살표 중 하나에는 표준 사상의
$-1$배를 사용하여 복합체의 사상
$$
\mathcal{I}^\bullet(X) \xrightarrow{\alpha}
\mathcal{I}^\bullet(Z) \oplus
\mathcal{I}^\bullet(Y) \xrightarrow{\beta}
\mathcal{I}^\bullet(E)
$$
을 얻으며 $\beta \circ \alpha = 0$이다. 따라서 표준 사상
$$
c^K_{X, Z, Y, E} :
\mathcal{I}^\bullet(X)
\longrightarrow
C(\beta)^\bullet[-1]
$$
을 얻는다. 이 사상이 표준이라는 것은 $K$를 나타내는 K-단사 복합체를
다르게 고르면 아벨군의 유도범주에서 동형인 화살표가 정해진다는 뜻이다.
$c^K_{X, Z, Y, E}$가 동형이면 그 역을 사용하여 표준 특수 삼각형
$$
R\Gamma(X, K) \to
R\Gamma(Z, K) \oplus
R\Gamma(Y, K) \to
R\Gamma(E, K) \to
R\Gamma(X, K)[1]
$$
을 얻는다. 이 모든 구성은 $K$에 대해 함자적이다.
\end{lemma}

\begin{proof}
이 보조정리는 사실상 자명하다. 예를 들어 $\mathcal{J}^\bullet$이
$K$를 나타내는 두 번째 K-단사 복합체이면 유사동형
$\mathcal{I}^\bullet \to \mathcal{J}^\bullet$을 고를 수 있고, 이는
등장하는 모든 복합체 사이의 유사동형을 정한다. 세부사항은 생략한다.
원뿔의 구성과 특수 삼각형과의 관계는 Derived Categories, Sections
\ref{derived-section-cones} 및
\ref{derived-section-homotopy-triangulated}를 보라.
\end{proof}

\begin{lemma}
\label{sites-cohomology-lemma-two-out-of-three-blow-up-square}
위의 상황에서 $K_1 \to K_2 \to K_3 \to K_1[1]$을
$D(\mathcal{O})$ 안의 특수 삼각형이라 하자.
$\{1, 2, 3\}$ 가운데 두 $i$에 대하여
$c^{K_i}_{X, Z, Y, E}$가 유사동형이면, 셋째 $i$에 대해서도 유사동형이다.
\end{lemma}

\begin{proof}
삼각형을 회전하여 $c^{K_1}_{X, Z, Y, E}$와
$c^{K_2}_{X, Z, Y, E}$가 유사동형이라고 가정해도 된다.
$K_1 \to K_2$를 나타내는 $\mathcal{O}$-가군 K-단사 복합체의 사상
$f : \mathcal{I}^\bullet_1 \to \mathcal{I}^\bullet_2$를 고르자.
그러면 K-단사 복합체 $C(f)^\bullet$이 $K_3$을 나타낸다.
Derived Categories, Lemma \ref{derived-lemma-triangle-K-injective}을 보라.
이제 사상 $c^{K_3}_{X, Z, Y, E}$는 $K(\textit{Ab})$ 안의 특수 삼각형
사이 사상의 셋째 다리이고, 다른 두 다리는 유사동형이므로 이 사상도
동형이다. 몇몇 세부사항은 생략한다. Derived Categories, Lemma
\ref{derived-lemma-third-isomorphism-triangle}을 사용하라.
\end{proof}

\noindent
이 구성이 실제로 특수 삼각형을 주는 한 판정법을 제시하자.

\begin{lemma}
\label{sites-cohomology-lemma-square-triangle}
위의 상황에서 다음을 가정하자.
\begin{enumerate}
\item $h_X^\# = h_Y^\# \amalg_{h_E^\#} h_Z^\#$이다.
\item $h_E^\# \to h_Y^\#$는 단사이다.
\end{enumerate}
그러면 Lemma \ref{sites-cohomology-lemma-c-square}의 구성은 특수 삼각형
$$
R\Gamma(X, K) \to
R\Gamma(Z, K) \oplus
R\Gamma(Y, K) \to
R\Gamma(E, K) \to R\Gamma(X, K)[1]
$$
을 주며, 이는 $D(\mathcal{C})$ 안의 $K$에 대해 함자적이다.
\end{lemma}

\begin{proof}
$K$를 각 항이 단사 아벨층인 K-단사 복합체로 나타낼 수 있다.
Section \ref{sites-cohomology-section-unbounded}를 보라. 따라서 $\mathcal{I}$가 단사
아벨층일 때
$$
0 \to \mathcal{I}(X) \to
\mathcal{I}(Z) \oplus \mathcal{I}(Y) \to
\mathcal{I}(E) \to 0
$$
이 짧은 완전열임을 보이면 충분하다. 첫 화살표는 단사이다. 실제로 조건
(1)에 의해 사상 $h_Y \amalg h_Z \to h_X$는 층화한 뒤 전사가 되는데,
이는 $\{Y \to X, Z \to X\}$가 $X$의 한 덮개에 의해 세분될 수 있다는
뜻이다. 마지막 화살표도 전사인데,
$\mathcal{I}(Y) \to \mathcal{I}(E)$가 전사이기 때문이다. 구체적으로
$\mathcal{I}(E) = \Hom(\mathbf{Z}_E^\#, \mathcal{I})$,
$\mathcal{I}(Y) = \Hom(\mathbf{Z}_Y^\#, \mathcal{I})$,
이고, (2)에 의해 사상 $\mathbf{Z}_E^\# \to \mathbf{Z}_Y^\#$는
단사이며, $\mathcal{I}$는 단사 아벨층이다. Modules on Sites, Section
\ref{sites-modules-section-free-abelian-sheaf}과 비교하라.
마지막으로 $s \in \mathcal{I}(Y)$와 $t \in \mathcal{F}(Z)$가
$\mathcal{I}(E)$의 같은 원소로 간다고 하자. 그러면 $s$와 $t$는 사상
$$
s \amalg t : h_Y^\# \amalg h_Z^\# \longrightarrow \mathcal{I}
$$
을 정의하고, 이 사상은 가정에 의해
$h_Y^\# \amalg_{h_E^\#} h_Z^\#$를 거쳐 인수분해된다. 따라서 가정
(1)에 의해 유일한 사상 $h_X^\# \to \mathcal{I}$를 얻는다. 이는
$Y$ 위에서 $s$로, $Z$ 위에서 $t$로 제한되는 $\mathcal{I}(X)$의 원소에
대응한다.
\end{proof}

\begin{lemma}
\label{sites-cohomology-lemma-square-triangle-general}
$\mathcal{C}$를 사이트라 하자. $\mathcal{C}$ 위의 집합의 준층들로 이루어진
가환도식
$$
\xymatrix{
\mathcal{D} \ar[r] \ar[d] & \mathcal{F} \ar[d] \\
\mathcal{E} \ar[r] & \mathcal{G}
}
$$
을 생각하고 다음을 가정하자.
\begin{enumerate}
\item $\mathcal{G}^\# =
\mathcal{E}^\# \amalg_{\mathcal{D}^\#} \mathcal{F}^\#$이다.
\item $\mathcal{D}^\# \to \mathcal{F}^\#$는 단사이다.
\end{enumerate}
그러면 표준 특수 삼각형
$$
R\Gamma(\mathcal{G}, K) \to
R\Gamma(\mathcal{E}, K) \oplus
R\Gamma(\mathcal{F}, K) \to
R\Gamma(\mathcal{D}, K) \to
R\Gamma(\mathcal{G}, K)[1]
$$
이 있으며, 이는 $K \in D(\mathcal{C})$에 대해 함자적이다. 여기서
$R\Gamma(\mathcal{G}, -)$는 Section \ref{sites-cohomology-section-limp}에서 논의한
코호몰로지이다.
\end{lemma}

\begin{proof}
층화는 완전하고
$R\Gamma(\mathcal{G}, -) = R\Gamma(\mathcal{G}^\#, -)$
이므로 $\mathcal{D}, \mathcal{E}, \mathcal{F}, \mathcal{G}$가
집합의 층이라고 가정해도 된다. 더욱이 코호몰로지
$R\Gamma(\mathcal{G}, -)$는 기저 사이트가 아니라 토포스에만 의존한다.
따라서 Sites, Lemma \ref{sites-lemma-topos-good-site}에 의해
$\mathcal{C}$를 준표준 위상을 가진 ``더 큰'' 사이트로 바꾸어
어떤 $\mathcal{C}$의 대상 $X, Y, Z, E$에 대하여
$\mathcal{G} = h_X$, $\mathcal{F} = h_Y$, $\mathcal{E} = h_Z$,
그리고 $\mathcal{D} = h_E$라고 가정해도 된다. 이 경우 결과는
Lemma \ref{sites-cohomology-lemma-square-triangle}에서 따른다.
\end{proof}







\section{두 위상의 비교}
\label{sites-cohomology-section-compare}

\noindent
$\mathcal{C}$를 범주라 하자.
$\text{Cov}(\mathcal{C}) \supset \text{Cov}'(\mathcal{C})$
를 $\mathcal{C}$에 사이트 구조를 주는 두 방식이라 하자.
$\text{Cov}(\mathcal{C})$에 대응하는 위상을 $\tau$라 하고,
$\text{Cov}'(\mathcal{C})$에 대응하는 위상을 $\tau'$라 하자.
그러면 $\mathcal{C}$의 항등함자는 사이트의 사상
$$
\epsilon : \mathcal{C}_\tau \longrightarrow \mathcal{C}_{\tau'}
$$
을 정의한다. 여기서 $\epsilon_*$는 기저 준층들에 대한 항등함자이고,
$\epsilon^{-1}$은 $\tau'$-층의 $\tau$-층화이다.
Sites, Examples \ref{sites-example-finer-topology} 및
\ref{sites-example-finer-topology-bis}를 보라.
위의 상황에서 다음이 성립한다.
\begin{enumerate}
\item $\epsilon_* : \Sh(\mathcal{C}_\tau) \to \Sh(\mathcal{C}_{\tau'})$
는 충실충만이고 $\epsilon^{-1} \circ \epsilon_* = \text{id}$이다.
\item $\epsilon_* : \textit{Ab}(\mathcal{C}_\tau) \to
\textit{Ab}(\mathcal{C}_{\tau'})$
는 충실충만이고 $\epsilon^{-1} \circ \epsilon_* = \text{id}$이다.
\item $R\epsilon_* : D(\mathcal{C}_\tau) \to D(\mathcal{C}_{\tau'})$
는 충실충만이고 $\epsilon^{-1} \circ R\epsilon_* = \text{id}$이다.
\item $\mathcal{O}$가 $\tau$-위상에 대한 환의 층이면,
$\mathcal{O}$는 $\tau'$-위상에 대해서도 층이고 $\epsilon$은 환 달린
사이트의 평탄 사상
$$
\epsilon :
(\mathcal{C}_\tau, \mathcal{O}_\tau)
\longrightarrow
(\mathcal{C}_{\tau'}, \mathcal{O}_{\tau'})
$$
이 된다.
\item $\epsilon_* : \textit{Mod}(\mathcal{O}_\tau) \to
\textit{Mod}(\mathcal{O}_{\tau'})$는 충실충만이고
$\epsilon^* \circ \epsilon_* = \text{id}$이다.
\item $R\epsilon_* : D(\mathcal{O}_\tau) \to D(\mathcal{O}_{\tau'})$
는 충실충만이고 $\epsilon^* \circ R\epsilon_* = \text{id}$이다.
\end{enumerate}
몇 가지 설명을 덧붙인다.

\medskip\noindent
(1)에 관하여. $\mathcal{F}$를 $\tau$-위상에서의 집합의 층이라 하자.
그러면 $\epsilon_*\mathcal{F}$는 $\mathcal{F}$를 $\tau'$-위상의 층으로
본 것일 뿐이다. $\epsilon^{-1}$을 적용한다는 것은 $\mathcal{F}$의
$\tau$-층화를 취한다는 뜻인데, $\mathcal{F}$는 이미 $\tau$-층이므로
아무것도 바꾸지 않는다. 따라서
$\epsilon^{-1}(\epsilon_*\mathcal{F})) = \mathcal{F}$.
충실충만성은 Categories, Lemma
\ref{categories-lemma-adjoint-fully-faithful}에서 따른다.

\medskip\noindent
(2)에 관하여. 아벨층의 당김과 직접상은 기저 집합의 층에 그 연산들을
하는 것과 같으므로 이는 (1)의 결과이다.

\medskip\noindent
(3)에 관하여. $K$를 $D(\mathcal{C}_\tau)$의 대상이라 하자.
$R\epsilon_*K$를 계산하기 위해 $K$를 나타내는 K-단사 복합체
$\mathcal{I}^\bullet$을 고르고
$R\epsilon_*K = \epsilon_*\mathcal{I}^\bullet$로 놓는다.
대상 $L$에 대한
$\epsilon^{-1} : D(\mathcal{C}_{\tau'}) \to D(\mathcal{C}_\tau)$의 값은
$L$을 나타내는 임의의 아벨층 복합체에 완전 함자 $\epsilon^{-1}$을
적용하여 계산하므로, $\epsilon^{-1}R\epsilon_*K$는
$\epsilon^{-1}\epsilon_*\mathcal{I}^\bullet$로 나타내어진다.
(1)에 의해
$\mathcal{I}^\bullet = \epsilon^{-1}\epsilon_*\mathcal{I}^\bullet$.
달리 말해 $\epsilon^{-1} \circ R\epsilon_* = \text{id}$이고, 앞과 같이
결론을 얻는다.

\medskip\noindent
(4)에 관하여. $\epsilon^{-1}\mathcal{O}_{\tau'} = \mathcal{O}_\tau$임에
유의하자. (1)의 논의를 보라. 따라서 $\epsilon$은 환 달린 사이트의
평탄 사상이다. Modules on Sites, Definition
\ref{sites-modules-definition-flat-morphism}을 보라. 그뿐 아니라
$\mathcal{O}_{\tau'}$-가군에 대해 $\epsilon^* = \epsilon^{-1}$임도
분명하다(가군으로서의 당김은 기저 아벨층의 당김과 같은 기저 아벨층을
갖는다).

\medskip\noindent
(5)에 관하여. 이는 (2)와 (4)에서 말한 것으로부터 분명하다.

\medskip\noindent
(6)에 관하여. 이는 (3)과 유사하다. 세부사항은 생략한다.












\section{코호몰로지 강하의 형식적 성질}
\label{sites-cohomology-section-formal-cohomological-descent}

\noindent
이 절에서는 환 달린 토포스의 사상이 아래쪽 유도범주에서 위쪽
유도범주로의 매입을 어느 정도까지 정하는지만 논의한다.
다음은 전형적인 결과이다.

\begin{lemma}
\label{sites-cohomology-lemma-downstairs}
$f : (\Sh(\mathcal{C}), \mathcal{O}_\mathcal{C}) \to
(\Sh(\mathcal{D}), \mathcal{O}_\mathcal{D})$를 환 달린 토포스의
사상이라 하자. 다음 사상이 동형인 대상 $K$들로 이루어진 충만 부분범주
$D' \subset D(\mathcal{O}_\mathcal{D})$를 생각하자.
$$
K \longrightarrow Rf_*Lf^*K
$$
그러면 $D'$은 $D(\mathcal{O}_\mathcal{D})$의 포화된 엄밀 충만
삼각부분범주이고, 함자 $Lf^* : D' \to D(\mathcal{O}_\mathcal{C})$는
충실충만하다.
\end{lemma}

\begin{proof}
이 상황에서 포화의 정의는 Derived Categories, Definition
\ref{derived-definition-saturated}을 보라. 삼각부분범주에 관한 논의는
Derived Categories, Lemma \ref{derived-lemma-triangulated-subcategory}를
보라. 보조정리의 표준 사상은 수반함자쌍 $(Lf^*, Rf_*)$의 단위이다.
Lemma \ref{sites-cohomology-lemma-adjoint}를 보라. $D'$이 포화 삼각부분범주라는 증명은
생략한다. 이는 $Lf^*$와 $Rf_*$가 삼각범주의 완전 함자라는 사실에서
형식적으로 따른다. 마지막 부분은 $Lf^*$와 $Rf_*$가 수반이라는
사실에서 형식적으로 따른다. Categories, Lemma
\ref{categories-lemma-adjoint-fully-faithful}과 비교하라.
\end{proof}

\begin{lemma}
\label{sites-cohomology-lemma-upstairs}
$f : (\Sh(\mathcal{C}), \mathcal{O}_\mathcal{C}) \to
(\Sh(\mathcal{D}), \mathcal{O}_\mathcal{D})$를 환 달린 토포스의
사상이라 하자. 다음 사상이 동형인 대상 $K$들로 이루어진 충만 부분범주
$D' \subset D(\mathcal{O}_\mathcal{C})$를 생각하자.
$$
Lf^*Rf_*K \longrightarrow K
$$
그러면 $D'$은 $D(\mathcal{O}_\mathcal{C})$의 포화된 엄밀 충만
삼각부분범주이고, 함자 $Rf_* : D' \to D(\mathcal{O}_\mathcal{D})$는
충실충만하다.
\end{lemma}

\begin{proof}
이 상황에서 포화의 정의는 Derived Categories, Definition
\ref{derived-definition-saturated}을 보라. 삼각부분범주에 관한 논의는
Derived Categories, Lemma \ref{derived-lemma-triangulated-subcategory}를
보라. 보조정리의 표준 사상은 수반함자쌍 $(Lf^*, Rf_*)$의 쌍대단위이다.
Lemma \ref{sites-cohomology-lemma-adjoint}를 보라. $D'$이 포화 삼각부분범주라는 증명은
생략한다. 이는 $Lf^*$와 $Rf_*$가 삼각범주의 완전 함자라는 사실에서
형식적으로 따른다. 마지막 부분은 $Lf^*$와 $Rf_*$가 수반이라는
사실에서 형식적으로 따른다. Categories, Lemma
\ref{categories-lemma-adjoint-fully-faithful}과 비교하라.
\end{proof}

\begin{lemma}
\label{sites-cohomology-lemma-bounded-in-image-upstairs}
$f : (\Sh(\mathcal{C}), \mathcal{O}_\mathcal{C}) \to
(\Sh(\mathcal{D}), \mathcal{O}_\mathcal{D})$를 환 달린 토포스의
사상이라 하자. $K$를 $D(\mathcal{O}_\mathcal{C})$의 대상이라 하고
다음을 가정하자.
\begin{enumerate}
\item $f$는 평탄하다.
\item $K$는 아래로 유계이다.
\item $f^*Rf_*H^q(K) \to H^q(K)$는 동형이다.
\end{enumerate}
그러면 $f^*Rf_*K \to K$는 동형이다.
\end{lemma}

\begin{proof}
$f^*Rf_*K \to K$가 동형일 필요충분조건은 코호몰로지 층 $H^j$에서
동형인 것임에 유의하자. 또한
$H^j(f^*Rf_*K) = f^*H^j(Rf_*K) = f^*H^j(Rf_*\tau_{\leq j}K) =
H^j(f^*Rf_*\tau_{\leq j}K)$.
따라서 $K$가 유계라고 가정해도 된다. 성질 (3)에 의해 코호몰로지 층들은
Lemma \ref{sites-cohomology-lemma-upstairs}의 삼각부분범주
$D' \subset D(\mathcal{O}_\mathcal{C})$에 속한다. 따라서 $K$도
그 부분범주에 속한다.
\end{proof}

\begin{lemma}
\label{sites-cohomology-lemma-bounded-in-image-downstairs}
$f : (\Sh(\mathcal{C}), \mathcal{O}_\mathcal{C}) \to
(\Sh(\mathcal{D}), \mathcal{O}_\mathcal{D})$를 환 달린 토포스의
사상이라 하자. $K$를 $D(\mathcal{O}_\mathcal{D})$의 대상이라 하고
다음을 가정하자.
\begin{enumerate}
\item $f$는 평탄하다.
\item $K$는 아래로 유계이다.
\item $H^q(K) \to Rf_*f^*H^q(K)$는 동형이다.
\end{enumerate}
그러면 $K \to Rf_*f^*K$는 동형이다.
\end{lemma}

\begin{proof}
$K \to Rf_*f^*K$가 동형일 필요충분조건은 코호몰로지 층 $H^j$에서
동형인 것임에 유의하자. 또한
$H^j(Rf_*f^*K) = H^j(Rf_*\tau_{\leq j}f^*K) = H^j(Rf_*f^*\tau_{\leq j}K)$.
따라서 $K$가 유계라고 가정해도 된다. 성질 (3)에 의해 코호몰로지 층들은
Lemma \ref{sites-cohomology-lemma-downstairs}의 삼각부분범주
$D' \subset D(\mathcal{O}_\mathcal{D})$에 속한다. 따라서 $K$도
그 부분범주에 속한다.
\end{proof}

\begin{lemma}
\label{sites-cohomology-lemma-equivalence-bounded}
$f : (\Sh(\mathcal{C}), \mathcal{O}) \to (\Sh(\mathcal{C}'), \mathcal{O}')$를
환 달린 토포스의 사상이라 하자.
$\mathcal{A} \subset \textit{Mod}(\mathcal{O})$와
$\mathcal{A}' \subset \textit{Mod}(\mathcal{O}')$를 약한 Serre
부분범주라 하자. 다음을 가정하자.
\begin{enumerate}
\item $f$는 평탄하다.
\item $f^*$는 범주의 동치 $\mathcal{A}' \to \mathcal{A}$를 유도한다.
\item $\mathcal{F}' \in \Ob(\mathcal{A}')$에 대하여
$\mathcal{F}' \to Rf_*f^*\mathcal{F}'$은 동형이다.
\end{enumerate}
그러면
$f^* : D_{\mathcal{A}'}^+(\mathcal{O}') \to D_\mathcal{A}^+(\mathcal{O})$
는 범주의 동치이고, 그 준역은
$Rf_* : D_\mathcal{A}^+(\mathcal{O}) \to D_{\mathcal{A}'}^+(\mathcal{O}')$이다.
\end{lemma}

\begin{proof}
가정 (2), (3)과 Lemmas \ref{sites-cohomology-lemma-bounded-in-image-downstairs} 및
\ref{sites-cohomology-lemma-downstairs}에 의해
$f^* : D_{\mathcal{A}'}^+(\mathcal{O}') \to D_\mathcal{A}^+(\mathcal{O})$
가 충실충만함을 알 수 있다.
$\mathcal{F} \in \Ob(\mathcal{A})$라 하자. 그러면
$\mathcal{F} = f^*\mathcal{F}'$로 쓸 수 있고
$Rf_*\mathcal{F} = Rf_* f^*\mathcal{F}' = \mathcal{F}'$이다.
특히 $p > 0$에 대하여 $R^pf_*\mathcal{F} = 0$이고
$f_*\mathcal{F} \in \Ob(\mathcal{A}')$이다.
따라서 임의의 $K \in D^+_\mathcal{A}(\mathcal{O})$에 대하여
$R^{p + q}f_*K$로 수렴하는 스펙트럼열
$E_2^{p, q} = R^pf_*H^q(K)$을 사용하면
$Rf_*K$가 $D^+_{\mathcal{A}'}(\mathcal{O}')$에 속함을 알 수 있다.
물론 Lemmas \ref{sites-cohomology-lemma-bounded-in-image-upstairs} 및
\ref{sites-cohomology-lemma-upstairs}에 의해
$Rf_* : D_\mathcal{A}^+(\mathcal{O}) \to D_{\mathcal{A}'}^+(\mathcal{O}')$
도 충실충만하다. $f^*$와 $Rf_*$가 수반이므로, 예를 들어 Categories,
Lemma \ref{categories-lemma-adjoint-fully-faithful}에 의해 보조정리의
결론을 얻는다.
\end{proof}

\begin{lemma}
\label{sites-cohomology-lemma-equivalence-unbounded-one}
\begin{reference}
이는 \cite[Theorem 2.2.3]{six-I}와 유사하다.
\end{reference}
$f : (\Sh(\mathcal{C}), \mathcal{O}) \to (\Sh(\mathcal{C}'), \mathcal{O}')$를
환 달린 토포스의 사상이라 하자.
$\mathcal{A} \subset \textit{Mod}(\mathcal{O})$와
$\mathcal{A}' \subset \textit{Mod}(\mathcal{O}')$를 약한 Serre
부분범주라 하자. 다음을 가정하자.
\begin{enumerate}
\item $f$는 평탄하다.
\item $f^*$는 범주의 동치 $\mathcal{A}' \to \mathcal{A}$를 유도한다.
\item $\mathcal{F}' \in \Ob(\mathcal{A}')$에 대하여
$\mathcal{F}' \to Rf_*f^*\mathcal{F}'$은 동형이다.
\item $\mathcal{C}, \mathcal{O}, \mathcal{A}$는 Situation
\ref{sites-cohomology-situation-olsson-laszlo}의 가정을 만족한다.
\item $\mathcal{C}', \mathcal{O}', \mathcal{A}'$는 Situation
\ref{sites-cohomology-situation-olsson-laszlo}의 가정을 만족한다.
\end{enumerate}
그러면 $f^* : D_{\mathcal{A}'}(\mathcal{O}') \to D_\mathcal{A}(\mathcal{O})$
는 범주의 동치이고, 그 준역은
$Rf_* : D_\mathcal{A}(\mathcal{O}) \to D_{\mathcal{A}'}(\mathcal{O}')$이다.
\end{lemma}

\begin{proof}
$f^*$가 완전하므로 $f^*$가 명제에 적힌 함자
$f^* : D_{\mathcal{A}'}(\mathcal{O}') \to D_\mathcal{A}(\mathcal{O})$
를 정의하고 이 함자가 절단함자 $\tau_{\geq -n}$과 가환함은 분명하다.
대응하는 아래로 유계인 범주들에서 $f^*$와 $Rf_*$가 서로 준역인
동치임은 이미 알고 있다. Lemma \ref{sites-cohomology-lemma-equivalence-bounded}를 보라.
$N = 0$으로 둔 Lemma \ref{sites-cohomology-lemma-olsson-laszlo-map-version-one}에 의해
$Rf_*$가 실제로 함자
$Rf_* : D_\mathcal{A}(\mathcal{O}) \to D_{\mathcal{A}'}(\mathcal{O}')$
를 정의하고 이 함자도 절단함자 $\tau_{\geq -n}$과 가환함을 알 수 있다.
따라서 $D_\mathcal{A}(\mathcal{O})$ 안의 $K$에 대하여 사상
$f^*Rf_*K \to K$는 동형이다. 절단들에서 이것이 성립하기 때문이다.
마찬가지로 $D_{\mathcal{A}'}(\mathcal{O}')$ 안의 $K'$에 대하여 사상
$K' \to Rf_*f^*K'$은 동형이다. 절단들에서 이것이 성립하기 때문이다.
이것으로 증명을 마친다.
\end{proof}

\begin{lemma}
\label{sites-cohomology-lemma-equivalence-unbounded-two}
\begin{reference}
이는 \cite[Theorem 2.2.3]{six-I}와 유사하다.
\end{reference}
$f : (\mathcal{C}, \mathcal{O}) \to (\mathcal{C}', \mathcal{O}')$를
환 달린 사이트의 사상이라 하자.
$\mathcal{A} \subset \textit{Mod}(\mathcal{O})$와
$\mathcal{A}' \subset \textit{Mod}(\mathcal{O}')$를 약한 Serre
부분범주라 하자. 다음을 가정하자.
\begin{enumerate}
\item $f$는 평탄하다.
\item $f^*$는 범주의 동치 $\mathcal{A}' \to \mathcal{A}$를 유도한다.
\item $\mathcal{F}' \in \Ob(\mathcal{A}')$에 대하여
$\mathcal{F}' \to Rf_*f^*\mathcal{F}'$은 동형이다.
\item $\mathcal{C}, \mathcal{O}, \mathcal{A}$는 Situation
\ref{sites-cohomology-situation-olsson-laszlo}의 가정을 만족한다.
\item $f : (\mathcal{C}, \mathcal{O}) \to (\mathcal{C}', \mathcal{O}')$와
$\mathcal{A}$는 Situation \ref{sites-cohomology-situation-olsson-laszlo-prime}의
가정을 만족한다.
\end{enumerate}
그러면 $f^* : D_{\mathcal{A}'}(\mathcal{O}') \to D_\mathcal{A}(\mathcal{O})$
는 범주의 동치이고, 그 준역은
$Rf_* : D_\mathcal{A}(\mathcal{O}) \to D_{\mathcal{A}'}(\mathcal{O}')$이다.
\end{lemma}

\begin{proof}
이 보조정리의 증명은 Lemma \ref{sites-cohomology-lemma-equivalence-unbounded-one}의
증명과 완전히 같고, 다만 Lemma
\ref{sites-cohomology-lemma-olsson-laszlo-map-version-one}에 대한 참조를 Lemma
\ref{sites-cohomology-lemma-olsson-laszlo-map-version-two}에 대한 참조로 바꾼다.
\end{proof}








\section{두 위상의 비교, II}
\label{sites-cohomology-section-compare-II}

\noindent
$\mathcal{C}$를 범주라 하자.
$\text{Cov}(\mathcal{C}) \supset \text{Cov}'(\mathcal{C})$
를 $\mathcal{C}$에 사이트 구조를 주는 두 방식이라 하자.
$\text{Cov}(\mathcal{C})$에 대응하는 위상을 $\tau$라 하고,
$\text{Cov}'(\mathcal{C})$에 대응하는 위상을 $\tau'$라 하자.
그러면 $\mathcal{C}$의 항등함자는 사이트의 사상
$$
\epsilon : \mathcal{C}_\tau \longrightarrow \mathcal{C}_{\tau'}
$$
을 정의한다. 여기서 $\epsilon_*$는 기저 준층들에 대한 항등함자이고,
$\epsilon^{-1}$은 $\tau'$-층의 $\tau$-층화이므로 분명 완전하다.
$\mathcal{O}$를 $\tau$-위상에 대한 환의 층이라 하자. 그러면
$\mathcal{O}$는 $\tau'$-위상에 대해서도 층이고, $\epsilon$은
환 달린 사이트의 사상
$$
\epsilon :
(\mathcal{C}_\tau, \mathcal{O}_\tau)
\longrightarrow
(\mathcal{C}_{\tau'}, \mathcal{O}_{\tau'})
$$
이 된다. 더 자세한 논의는 Section \ref{sites-cohomology-section-compare}을 보라.

\begin{lemma}
\label{sites-cohomology-lemma-compare-topologies-derived-adequate-modules}
위와 같은 $\epsilon : (\mathcal{C}_\tau, \mathcal{O}_\tau) \to
(\mathcal{C}_{\tau'}, \mathcal{O}_{\tau'})$를 생각하자.
$\mathcal{B} \subset \Ob(\mathcal{C})$를 부분집합이라 하고,
$\mathcal{A} \subset \textit{PMod}(\mathcal{O})$를 충만 부분범주라 하자.
다음을 가정하자.
\begin{enumerate}
\item $\mathcal{A}$의 모든 대상은 $\tau$-위상에 대한 층이다.
\item $\mathcal{A}$는 $\textit{Mod}(\mathcal{O}_\tau)$의
약한 Serre 부분범주이다.
\item $\mathcal{C}$의 모든 대상은 그 원소들이 $\mathcal{B}$에 속하는
$\tau'$-덮개를 갖는다.
\item 모든 $U \in \mathcal{B}$와 모든 $\mathcal{F} \in \mathcal{A}$에
대해 $p > 0$이면 $H^p_\tau(U, \mathcal{F}) = 0$이다.
\end{enumerate}
그러면 $\mathcal{A}$는 $\textit{Mod}(\mathcal{O}_{\tau'})$의
약한 Serre 부분범주이고, $\epsilon^*$와 $R\epsilon_*$가 주는
삼각범주의 동치
$D_\mathcal{A}(\mathcal{O}_\tau) = D_\mathcal{A}(\mathcal{O}_{\tau'})$
가 있다.
\end{lemma}

\begin{proof}
$\epsilon^{-1}\mathcal{O}_{\tau'} = \mathcal{O}_\tau$이므로
$\epsilon$은 환 달린 사이트의 평탄 사상이고, 가군의 층에 대해서는
실제로 $\epsilon^{-1} = \epsilon^*$이다. 성질 (1)에 의해
$\mathcal{A}$의 모든 대상을 $\mathcal{O}_\tau$-가군의 층인 동시에
$\mathcal{O}_{\tau'}$-가군의 층으로 생각할 수 있다. 달리 말해,
충실충만한 포함함자
$$
\mathcal{A} \to \textit{Mod}(\mathcal{O}_\tau) \to
\textit{Mod}(\mathcal{O}_{\tau'})
$$
들이 있다. 혼동을 피하기 위해 $\mathcal{A}$의 상을
$\mathcal{A}' \subset \textit{Mod}(\mathcal{O}_{\tau'})$로 쓰겠다.
그러면 $\epsilon_* : \mathcal{A} \to \mathcal{A}'$와
$\epsilon^* : \mathcal{A}' \to \mathcal{A}$가 서로 준역인 동치임은
분명하다(보조정리 앞의 논의를 보고, $\mathcal{A}'$의 대상들이
$\tau$-위상에서 층이라는 사실을 사용하라).

\medskip\noindent
조건 (3)과 (4)에 의해 $p > 0$이고
$\mathcal{F} \in \Ob(\mathcal{A})$이면
$R^p\epsilon_*\mathcal{F} = 0$이다. 실제로 $R^p\epsilon_*$는 준층
$U \mapsto H^p_\tau(U, \mathcal{F})$에 연관된 층이다.
Lemma \ref{sites-cohomology-lemma-higher-direct-images}를 보라. 따라서
$\mathcal{A}$의 임의의 완전 복합체에 함자 $\epsilon_*$를 적용해도
완전하다. 여기서 그러한 복합체란 그 항들이
$\mathcal{A}$에 속하는 $\textit{Mod}(\mathcal{O}_\tau)$의 완전
복합체와 같은 것이다. Homology, Lemma
\ref{homology-lemma-characterize-weak-serre-subcategory}를 보라.

\medskip\noindent
다음 $\textit{Mod}(\mathcal{O}_{\tau'})$의 완전열을 생각하자.
$$
\mathcal{F}'_0 \to \mathcal{F}'_1 \to
\mathcal{F}'_2 \to \mathcal{F}'_3 \to
\mathcal{F}'_4
$$
여기서 $\mathcal{F}'_0, \mathcal{F}'_1, \mathcal{F}'_3,
\mathcal{F}'_4$는 $\mathcal{A}'$에 속한다. 완전 함자
$\epsilon^*$를 적용하면 다음 $\textit{Mod}(\mathcal{O}_\tau)$의
완전열을 얻는다.
$$
\epsilon^*\mathcal{F}'_0 \to \epsilon^*\mathcal{F}'_1 \to
\epsilon^*\mathcal{F}'_2 \to \epsilon^*\mathcal{F}'_3 \to
\epsilon^*\mathcal{F}'_4
$$
실제로 $\mathcal{A}$는 약한 Serre 부분범주이고
$\epsilon^*\mathcal{F}'_0, \epsilon^*\mathcal{F}'_1,
\epsilon^*\mathcal{F}'_3, \epsilon^*\mathcal{F}'_4$는
$\mathcal{A}$에 속하므로, Homology, Definition
\ref{homology-definition-serre-subcategory}에 의해
$\epsilon^*\mathcal{F}_2$도 $\mathcal{A}$에 속한다.
다음 열들의 사상을 생각하자.
$$
\xymatrix{
\mathcal{F}'_0 \ar[r] \ar[d] &
\mathcal{F}'_1 \ar[r] \ar[d] &
\mathcal{F}'_2 \ar[r] \ar[d] &
\mathcal{F}'_3 \ar[r] \ar[d] &
\mathcal{F}'_4 \ar[d] \\
\epsilon_*\epsilon^*\mathcal{F}'_0 \ar[r] &
\epsilon_*\epsilon^*\mathcal{F}'_1 \ar[r] &
\epsilon_*\epsilon^*\mathcal{F}'_2 \ar[r] &
\epsilon_*\epsilon^*\mathcal{F}'_3 \ar[r] &
\epsilon_*\epsilon^*\mathcal{F}'_4
}
$$
앞 문단의 논의에 의해 아래 행은 완전하다. 첫 문단의 논의에 의해
지표가 $0$, $1$, $3$, $4$인 수직 화살표들은 동형이다.
$5$-보조정리(Homology, Lemma \ref{homology-lemma-five-lemma})에 의해
$\mathcal{F}'_2 \cong \epsilon_*\epsilon^*\mathcal{F}'_2$이고, 따라서
$\mathcal{F}'_2$는 $\mathcal{A}'$에 속한다. 이로써
$\mathcal{A}'$가 $\textit{Mod}(\mathcal{O}_{\tau'})$의 약한 Serre
부분범주임을 알 수 있다. Homology, Definition
\ref{homology-definition-serre-subcategory}를 보라.

\medskip\noindent
이제 유도범주 $D_\mathcal{A}(\mathcal{O}_\tau)$와
$D_{\mathcal{A}'}(\mathcal{O}_{\tau'})$를 말할 수 있다.
Derived Categories, Section \ref{derived-section-triangulated-sub}을 보라.
증명을 끝내기 위해 Lemma \ref{sites-cohomology-lemma-equivalence-unbounded-two}의
조건 (1) -- (5)를 적용할 수 있음을 보이자. 위에서 이미 (1), (2),
(3)을 보았다. 모든 대상이 $\mathcal{B}$의 대상들로 이루어진
$\tau'$-덮개를 가지므로, 더욱이 모든 대상은 $\mathcal{B}$의
대상들로 이루어진 $\tau$-덮개도 갖는다. 따라서 Lemma
\ref{sites-cohomology-lemma-equivalence-unbounded-two}의 조건 (4)가 성립한다.
마찬가지로 조건 (5)도 성립한다.
\end{proof}

\begin{lemma}
\label{sites-cohomology-lemma-descent-squares}
위와 같은 $\epsilon : (\mathcal{C}_\tau, \mathcal{O}_\tau) \to
(\mathcal{C}_{\tau'}, \mathcal{O}_{\tau'})$를 생각하자.
$A$를 집합이라 하고, 각 $\alpha \in A$에 대해
$$
\xymatrix{
E_\alpha \ar[d] \ar[r] & Y_\alpha \ar[d] \\
Z_\alpha \ar[r] & X_\alpha
}
$$
가 범주 $\mathcal{C}$의 가환 그림이라 하자. 다음을 가정하자.
\begin{enumerate}
\item $\tau'$-층 $\mathcal{F}'$가 모든 $\alpha$에 대해
$\mathcal{F}'(X_\alpha) =
\mathcal{F}'(Z_\alpha) \times_{\mathcal{F}'(E_\alpha)} 
\mathcal{F}'(Y_\alpha)$를 만족하면 $\tau$-층이다.
\item $R\epsilon_*$의 본질적 상에 속하는
$D(\mathcal{O}_{\tau'})$의 $K'$에 대해, Lemma
\ref{sites-cohomology-lemma-c-square}의 사상
$c^{K'}_{X_\alpha, Z_\alpha, Y_\alpha, E_\alpha}$는 모든
$\alpha$에 대해 동형이다.
\end{enumerate}
그러면 $K' \in D^+(\mathcal{O}_{\tau'})$가 $R\epsilon_*$의 본질적
상에 속할 필요충분조건은 사상
$c^{K'}_{X_\alpha, Z_\alpha, Y_\alpha, E_\alpha}$가 모든
$\alpha$에 대해 동형인 것이다.
\end{lemma}

\begin{proof}
“필요” 방향은 가정 (2)에서 따른다. 한편 $K'$가 영이 아닌
코호몰로지 층을 하나만 가지면 “충분” 방향은 가정 (1)에서 따른다.
일반적으로는 귀납법으로 “충분” 방향을 증명하겠다.
$D^+(\mathcal{O}_{\tau'})$의 대상 $K'$가 모든 $\alpha \in A$에 대해
사상 $c^{K'}_{X_\alpha, Z_\alpha, Y_\alpha, E_\alpha}$가 동형일 때
(P)를 만족한다고 하자.

\medskip\noindent
구체적으로, (P)를 만족하는 $D^+(\mathcal{O}_{\tau'})$의 대상
$K'$를 택하자. 첫 화살표가 수반사상인 특수 삼각형
$$
K' \to R\epsilon_*\epsilon^{-1}K' \to M' \to K'[1]
$$
을 $D^+(\mathcal{O}_{\tau'})$에서 택하자. (2)와 Lemma
\ref{sites-cohomology-lemma-two-out-of-three-blow-up-square}에 의해 $M'$는 (P)를
만족한다. 한편 $\epsilon^{-1}$을 적용하고 Section
\ref{sites-cohomology-section-compare}의 $\epsilon^{-1}R\epsilon_* = \text{id}$를
사용하면 $\epsilon^{-1}M' = 0$을 얻는다. 다음 문단에서
$M' = 0$임을 보이면 증명이 끝난다.

\medskip\noindent
(P)를 만족하고 $\epsilon^{-1}K' = 0$인
$D^+(\mathcal{O}_{\tau'})$의 대상 $K'$를 택하자.
$K' = 0$임을 보이겠다. 즉, $i < n$에 대해 $H^i(K') = 0$인
$n \in \mathbf{Z}$가 주어졌을 때 $H^n(K') = 0$임을 보이겠다.
각 $\alpha \in A$에 대해 Lemma \ref{sites-cohomology-lemma-c-square}에 의해
특수 삼각형
$$
R\Gamma_{\tau'}(X_\alpha, K') \to
R\Gamma_{\tau'}(Z_\alpha, K') \oplus
R\Gamma_{\tau'}(Y_\alpha, K') \to
R\Gamma_{\tau'}(E_\alpha, K') \to
R\Gamma_{\tau'}(X_\alpha, K')[1]
$$
이 있다. 차수 $n$의 코호몰로지를 취하고 $K'$의 코호몰로지 층에 대한
가정된 소멸을 사용하면 완전열
$$
0 \to
H^n_{\tau'}(X_\alpha, K') \to
H^n_{\tau'}(Z_\alpha, K') \oplus
H^n_{\tau'}(Y_\alpha, K') \to
H^n_{\tau'}(E_\alpha, K')
$$
을 얻는데, 이는 완전열
$$
0 \to
\Gamma(X_\alpha, H^n(K')) \to
\Gamma(Z_\alpha, H^n(K')) \oplus
\Gamma(Y_\alpha, H^n(K')) \to
\Gamma(E_\alpha, H^n(K'))
$$
과 같다. 가정 (1)에 의해 $H^n(K')$는 $\tau$-층이다.
그러나 $\epsilon^{-1}K' = 0$이므로 $\tau$-층화
$\epsilon^{-1}H^n(K') = H^n(\epsilon^{-1}K')$는 $0$이다.
따라서 원하는 대로 $H^n(K') = 0$이다.
\end{proof}

\begin{lemma}
\label{sites-cohomology-lemma-descent-squares-helper}
위와 같은 $\epsilon : (\mathcal{C}_\tau, \mathcal{O}_\tau) \to
(\mathcal{C}_{\tau'}, \mathcal{O}_{\tau'})$를 생각하자.
$$
\xymatrix{
E \ar[d] \ar[r] & Y \ar[d] \\
Z \ar[r] & X
}
$$
가 범주 $\mathcal{C}$의 가환 그림이고 다음을 만족한다고 하자.
\begin{enumerate}
\item $h_X^\# = h_Y^\# \amalg_{h_E^\#} h_Z^\#$이다.
\item $h_E^\# \to h_Y^\#$는 단사이다.
\end{enumerate}
여기서 ${}^\#$는 $\tau$-층화를 뜻한다. 그러면 $R\epsilon_*$의
본질적 상에 속하는 $K' \in D(\mathcal{O}_{\tau'})$에 대해
Lemma \ref{sites-cohomology-lemma-c-square}의 사상 $c^{K'}_{X, Z, Y, E}$는
동형이다(이때 $\tau'$-위상을 사용한다).
\end{lemma}

\begin{proof}
이 보조정리는 Lemma \ref{sites-cohomology-lemma-square-triangle}의 거의 즉각적인
결과이므로 독자에게 증명을 건너뛸 것을 강력히 권한다.
$K' = R\epsilon_*K$라 하자. $K$를 나타내는
$\mathcal{O}_\tau$-가군의 K-단사 복합체
$\mathcal{J}^\bullet$을 택하자. 그러면
$\epsilon_*\mathcal{J}^\bullet$은 $K'$를 나타내는
$\mathcal{O}_{\tau'}$-가군의 K-단사 복합체이다.
Lemma \ref{sites-cohomology-lemma-K-injective-flat}을 보라. 다음으로
$$
0 \to
\mathcal{J}^\bullet(X) \xrightarrow{\alpha}
\mathcal{J}^\bullet(Z) \oplus
\mathcal{J}^\bullet(Y) \xrightarrow{\beta}
\mathcal{J}^\bullet(E)
\to 0
$$
은 아벨군 복합체의 짧은 완전열이다. Lemma
\ref{sites-cohomology-lemma-square-triangle}와 그 증명을 보라. 이는
$c^{K'}_{X, Z, Y, E}$를 정의하는 데 쓰이는 아벨군 복합체의 열과
같으므로 결론을 얻는다.
\end{proof}







\section{코호몰로지의 비교}
\label{sites-cohomology-section-compare-general}

\noindent
서로 다른 위상에서의 코호몰로지를 비교하는 데 도움이 될 일반론을
전개한다. Section \ref{sites-cohomology-section-compare}과 같은
$\mathcal{C}$, $\tau$, $\tau'$ 및 $\mathcal{C}$의 사상
$f : X \to Y$가 주어지면 토포스 사상들의 가환 그림
\begin{equation}
\label{sites-cohomology-equation-commutative-epsilon}
\vcenter{
\xymatrix{
\Sh(\mathcal{C}_\tau/X) \ar[r]_{f_\tau} \ar[d]_{\epsilon_X} &
\Sh(\mathcal{C}_\tau/Y) \ar[d]^{\epsilon_Y} \\
\Sh(\mathcal{C}_{\tau'}/X) \ar[r]^{f_{\tau'}} &
\Sh(\mathcal{C}_{\tau'}/X)
}
}
\end{equation}
을 얻는다. 여기서 사상 $\epsilon_X$, 각각 $\epsilon_Y$는 두 위상
$\tau$와 $\tau'$를 갖춘 범주 $\mathcal{C}/X$에 대한 Section
\ref{sites-cohomology-section-compare}의 비교사상이다. 사상 $f_\tau$와 $f_{\tau'}$는
“재국소화” 사상이다(Sites, Lemma \ref{sites-lemma-relocalize}).
그림의 가환성은 Sites, Lemma \ref{sites-lemma-localize-morphism}의
특수한 경우이다. 여기서는 $\mathcal{C} = \mathcal{C}_\tau/Y$,
$\mathcal{D} = \mathcal{C}_{\tau'}/Y$, $u = \text{id}$, $U = X$,
$V = X$를 적용한다. 또한 그 보조정리에서, 또는 자명하게,
$\epsilon_{X, *} \circ f_\tau^{-1} = f_{\tau'}^{-1} \circ \epsilon_{Y, *}$
를 얻는다.

\begin{situation}
\label{sites-cohomology-situation-compare}
Section \ref{sites-cohomology-section-compare}과 같은 $\mathcal{C}$, $\tau$, $\tau'$를
생각하자. 부분집합 $\mathcal{P} \subset \text{Arrows}(\mathcal{C})$와
$\mathcal{C}$의 각 대상 $X$마다 약한 Serre 부분범주
$\mathcal{A}'_X \subset \textit{Ab}(\mathcal{C}_{\tau'}/X)$가
주어졌다고 하자. 다음을 가정한다.
\begin{enumerate}
\item
\label{sites-cohomology-item-base-change-P}
$f : X \to Y$가 $\mathcal{P}$에 속하고 $Y' \to Y$가 임의이면
$X \times_Y Y'$가 존재하고 $X \times_Y Y' \to Y'$는
$\mathcal{P}$에 속한다.
\item
\label{sites-cohomology-item-restriction-A}
$\mathcal{C}$의 임의의 사상 $f : X \to Y$에 대해
$f_{\tau'}^{-1}$은 $\mathcal{A}'_Y$를 $\mathcal{A}'_X$ 안으로 보낸다.
\item
\label{sites-cohomology-item-A-sheaf}
주어진 $\mathcal{C}$의 $X$와 $\mathcal{A}'_X$의
$\mathcal{F}'$에 대해, $\mathcal{F}'$는 $\tau$-덮개에 대한
층 조건을 만족한다. 즉,
$\mathcal{F}' = \epsilon_{X, *}\epsilon_X^{-1}\mathcal{F}'$,
\item
\label{sites-cohomology-item-A-and-P}
만일 $f : X \to Y$가 $\mathcal{P}$에 속하고
$\mathcal{F}' \in \Ob(\mathcal{A}'_X)$이면,
$R^if_{\tau', *}\mathcal{F}' \in \Ob(\mathcal{A}'_Y)$
가 모든 $i \geq 0$에 대해 성립한다.
\item
\label{sites-cohomology-item-refine-tau-by-P}
만일 $\{U_i \to U\}_{i \in I}$가 $\tau$-덮개이면 다음이 존재한다.
\begin{enumerate}
\item $\tau'$-덮개 $\{V_j \to U\}_{j \in J}$.
\item 하나의 $f_j \in \mathcal{P}$로 이루어진 $\tau$-덮개
$\{f_j : W_j \to V_j\}$.
\item $\tau'$-덮개 $\{W_{jk} \to W_j\}_{k \in K_j}$.
\end{enumerate}
이들은 $\{W_{jk} \to U\}_{j \in J, k \in K_j}$가
$\{U_i \to U\}_{i \in I}$의 세분이 되도록 택할 수 있다.
\end{enumerate}
\end{situation}

\begin{lemma}
\label{sites-cohomology-lemma-A}
Situation \ref{sites-cohomology-situation-compare}에서 $\mathcal{C}$의 $X$에 대해
$\mathcal{A}_X$를 $\mathcal{A}'_X$의 $\mathcal{F}'$에 대한
$\epsilon_X^{-1}\mathcal{F}'$ 꼴의
$\textit{Ab}(\mathcal{C}_\tau/X)$의 대상들로 이루어진 범주라 하자.
그러면 다음이 성립한다.
\begin{enumerate}
\item $\textit{Ab}(\mathcal{C}_\tau/X)$의 $\mathcal{F}$에 대해
$\mathcal{F} \in \mathcal{A}_X \Leftrightarrow
\epsilon_{X, *}\mathcal{F} \in \mathcal{A}'_X$이다.
\item $\mathcal{C}$의 임의의 사상 $f : X \to Y$에 대해
$f_\tau^{-1}$은 $\mathcal{A}_Y$를 $\mathcal{A}_X$ 안으로 보낸다.
\end{enumerate}
\end{lemma}

\begin{proof}
부분 (1)은 (\ref{sites-cohomology-item-A-sheaf})에서 따른다. 부분 (2)는
(\ref{sites-cohomology-item-restriction-A})와, 다음을 주는
(\ref{sites-cohomology-equation-commutative-epsilon})의 가환성에서 따른다.
$\epsilon_X^{-1} \circ f_{\tau'}^{-1} = f_\tau^{-1} \circ \epsilon_Y^{-1}$.
\end{proof}

\noindent
다음 목표는 Lemmas \ref{sites-cohomology-lemma-compare-cohomology-general}와
\ref{sites-cohomology-lemma-cohomological-descent-general}를 증명하는 것이다.
다음 귀납가설을 사용하는 귀납논법으로 이를 수행하겠다.

\medskip\noindent
$(V_n)$ $\mathcal{C}$의 $X$와 $\mathcal{A}_X$의 $\mathcal{F}$에 대해,
$1 \leq i \leq n$이면 $R^i\epsilon_{X, *}\mathcal{F} = 0$이다.

\begin{lemma}
\label{sites-cohomology-lemma-V-implies-C-general}
Situation \ref{sites-cohomology-situation-compare}에서 $(V_n)$이 성립한다고 하자.
$\mathcal{P}$의 $f : X \to Y$와 $\mathcal{A}_X$의 $\mathcal{F}$에
대해 $i \leq n$이면 $R^if_{\tau', *}\epsilon_{X, *}\mathcal{F} =
\epsilon_{Y, *}R^if_{\tau, *}\mathcal{F}$이다.
\end{lemma}

\begin{proof}
가환 그림 (\ref{sites-cohomology-equation-commutative-epsilon})를 더 언급하지 않고
사용하겠다. 특히
$$
Rf_{\tau', *}R\epsilon_{X, *}\mathcal{F} =
R\epsilon_{Y, *}Rf_{\tau, *}\mathcal{F}
$$
가 성립한다. 가정 $(V_n)$에 의해
$\epsilon_{X, *}\mathcal{F} \to R\epsilon_{X, *}\mathcal{F}$
는 차수 $\leq n$에서 동형이다. 따라서
$Rf_{\tau', *}\epsilon_{X, *}\mathcal{F} \to
Rf_{\tau', *}R\epsilon_{X, *}\mathcal{F}$는 차수 $\leq n$에서
동형이다. 그러므로
$$
R^if_{\tau', *}\epsilon_{X, *}\mathcal{F} \to
H^i(R\epsilon_{Y, *}Rf_{\tau, *}\mathcal{F})
$$
는 $i \leq n$에서 동형이다. $R\epsilon_{Y, *}Rf_{\tau, *}\mathcal{F}$에
대한 Lemma \ref{sites-cohomology-lemma-relative-Leray}의 스펙트럼열의 둘째 쪽을
살펴보아 보조정리를 증명하겠다. 그 모양은 다음과 같다.
$$
\begin{matrix}
\ldots &
\ldots &
\ldots &
\ldots \\
\epsilon_{Y, *}R^2f_{\tau, *}\mathcal{F} &
R^1\epsilon_{Y, *}R^2f_{\tau, *}\mathcal{F} &
R^2\epsilon_{Y, *}R^2f_{\tau, *}\mathcal{F} &
\ldots \\
\epsilon_{Y, *}R^1f_{\tau, *}\mathcal{F} &
R^1\epsilon_{Y, *}R^1f_{\tau, *}\mathcal{F} &
R^2\epsilon_{Y, *}R^1f_{\tau, *}\mathcal{F} &
\ldots \\
\epsilon_{Y, *}f_{\tau, *}\mathcal{F} &
R^1\epsilon_{Y, *}f_{\tau, *}\mathcal{F} &
R^2\epsilon_{Y, *}f_{\tau, *}\mathcal{F} &
\ldots
\end{matrix}
$$
$(C_m)$을 다음 가설이라 하자: $R^if_{\tau', *}\epsilon_{X, *}\mathcal{F} =
\epsilon_{Y, *}R^if_{\tau, *}\mathcal{F}$가 $i \leq m$에 대해 성립한다.
$(C_0)$은 성립한다. $m < n$일 때
$(C_{m - 1}) \Rightarrow (C_m)$임을 보이자. 실제로
$(C_{m - 1})$이 성립하면 $n \geq p > 0$ 및 $q \leq m - 1$에 대해
\begin{align*}
R^p\epsilon_{Y, *}R^qf_{\tau, *}\mathcal{F}
& =
R^p\epsilon_{Y, *}
\epsilon_Y^{-1} \epsilon_{Y, *} R^qf_{\tau, *}\mathcal{F} \\
& =
R^p\epsilon_{Y, *}
\epsilon_Y^{-1}R^qf_{\tau', *}\epsilon_{X, *}\mathcal{F} = 0
\end{align*}
이다. 첫째 등식은 $\epsilon_Y^{-1}\epsilon_{Y, *} = \text{id}$에서,
둘째는 $(C_{m - 1})$에서, 마지막은 $(V_n)$에서 따른다. 실제로
(\ref{sites-cohomology-item-A-and-P})에 의해
$\epsilon_Y^{-1}R^qf_{\tau', *}\epsilon_{X, *}\mathcal{F}$는
$\mathcal{A}_Y$에 속한다. 스펙트럼열을 살펴보면
$E_2^{0, m} = \epsilon_{Y, *}R^mf_{\tau, *}\mathcal{F}$
가 $p + q = m$인 영이 아닌 유일한 항 $E_2^{p, q}$이다.
$\text{d}_r^{p, q} : E_r^{p, q} \to E_r^{p + r, q - r + 1}$임을
상기하자. 따라서 이 자리에서 나가거나 이 자리로 들어오는 영이 아닌
미분 $\text{d}_r^{p, q}$, $r \geq 2$는 없다. 그러므로
$H^m(R\epsilon_{Y, *}Rf_{\tau, *}\mathcal{F}) =
\epsilon_{Y, *}R^mf_{\tau, *}\mathcal{F}$이고, 위 논의에 의해
$(C_m)$이 성립한다.

\medskip\noindent
마지막으로 $(C_{n - 1})$을 가정하자. 같은 분석에 의해
$E_2^{0, n} = \epsilon_{Y, *}R^nf_{\tau, *}\mathcal{F}$
가 $p + q = n$인 영이 아닌 유일한 항 $E_2^{p, q}$이다.
여전히 이 자리로 들어오는 영이 아닌 미분은 없지만, 이 자리에서
나가는 영이 아닌 미분은 있을 수 있다. 즉, 사상
$d_{n + 1}^{0, n} : \epsilon_{Y, *}R^nf_{\tau, *}\mathcal{F} \to
R^{n + 1}\epsilon_{Y, *}f_{\tau, *}\mathcal{F}$.
따라서 완전열
$$
0 \to R^nf_{\tau', *}\epsilon_{X, *}\mathcal{F} \to 
\epsilon_{Y, *}R^nf_{\tau, *}\mathcal{F} \to
R^{n + 1}\epsilon_{Y, *}f_{\tau, *}\mathcal{F}
$$
이 있다. (\ref{sites-cohomology-item-A-and-P})와 (\ref{sites-cohomology-item-A-sheaf})에 의해 층
$R^nf_{\tau', *}\epsilon_{X, *}\mathcal{F}$
은 $\tau$-덮개에 대한 층 조건을 만족하며,
$\epsilon_{Y, *}R^nf_{\tau, *}\mathcal{F}$도 그러하다(Section
\ref{sites-cohomology-section-compare}의 $\epsilon_*$의 설명을 사용하라).
그러나 $\tau'$-층
$R^{n + 1}\epsilon_{Y, *}f_{\tau, *}\mathcal{F}$
의 $\tau$-층화는 영이다. 이는 코호몰로지의 국소성에서 따르며,
Lemmas \ref{sites-cohomology-lemma-kill-cohomology-class-on-covering}와
\ref{sites-cohomology-lemma-higher-direct-images}를 사용한다.
따라서 $R^nf_{\tau', *}\epsilon_{X, *}\mathcal{F} \to 
\epsilon_{Y, *}R^nf_{\tau, *}\mathcal{F}$는 동형이어야 하며,
증명이 끝난다.
\end{proof}

\noindent
$E'$, 각각 $E$가 $D(\mathcal{C}_{\tau'}/X)$, 각각
$D(\mathcal{C}_\tau/X)$의 대상이면, $\mathcal{C}/X$의 대상 $U$ 위에서
$E'$, 각각 $E$의 코호몰로지를 $H^n_{\tau'}(U, E')$, 각각
$H^n_\tau(U, E)$로 쓰겠다.

\begin{lemma}
\label{sites-cohomology-lemma-V-implies-cohomology-general}
Situation \ref{sites-cohomology-situation-compare}에서 $(V_n)$이 성립한다고 하자.
$\mathcal{C}$의 $X$와 $L \in D(\mathcal{C}_{\tau'}/X)$가
$i < 0$에 대해 $H^i(L) = 0$이고, $0 \leq i \leq n$에 대해
$H^i(L)$이 $\mathcal{A}'_X$에 속하면
$H^n_{\tau'}(X, L) = H^n_\tau(X, \epsilon_X^{-1}L)$.
\end{lemma}

\begin{proof}
Lemma \ref{sites-cohomology-lemma-Leray-unbounded}에 의해
$H^n_\tau(X, \epsilon_X^{-1}L) =
H^n_{\tau'}(X, R\epsilon_{X, *}\epsilon_X^{-1}L)$.
다음으로 수렴하는 스펙트럼열
$$
E_2^{p, q} = R^p\epsilon_{X, *}\epsilon_X^{-1}H^q(L)
$$
이 있다: $H^{p + q}(R\epsilon_{X, *}\epsilon_X^{-1}L)$.
$(V_n)$에 의해 $0 < p \leq n$ 및 $0 \leq q \leq n$에 대해
$E_2^{p, q}$가 소멸한다. 따라서
$E_2^{0, q} = \epsilon_{X, *}\epsilon_X^{-1}H^q(L) = H^q(L)$
들이 $p + q \leq n$인 영이 아닌 유일한 항 $E_2^{p, q}$들이다.
따라서 사상
$$
L \longrightarrow R\epsilon_{X, *}\epsilon_X^{-1}L
$$
은 차수 $\leq n$에서 동형이다(작은 세부사항은 생략한다).
그러므로
$H^i_{\tau'}(X, L) = H^i_{\tau'}(X, R\epsilon_{X, *}\epsilon_X^{-1}L)$
가 $i \leq n$에 대해 성립하며, 보조정리가 증명되었다.
\end{proof}

\begin{lemma}
\label{sites-cohomology-lemma-V-implies-cohomology-extra-general}
Situation \ref{sites-cohomology-situation-compare}에서 $(V_n)$이 성립한다고 하자.
$\mathcal{C}$의 $X$와 $\mathcal{A}_X$의 $\mathcal{F}$에 대해 사상
$H^{n + 1}_{\tau'}(X, \epsilon_{X, *}\mathcal{F}) \to
H^{n + 1}_\tau(X, \mathcal{F})$
은 단사이고, 그 상은 $X$의 어떤 $\tau'$-덮개 위에서 자명해지는
유들로 이루어진다.
\end{lemma}

\begin{proof}
$\epsilon_X^{-1}\epsilon_{X, *}\mathcal{F} = \mathcal{F}$임을 상기하자.
따라서 이 사상은 코호몰로지 유들을 $\epsilon_X$로 당겨서 주어진다.
Leray 스펙트럼열(Lemma \ref{sites-cohomology-lemma-Leray})
$$
E_2^{p, q} = H^p_{\tau'}(X, R^q\epsilon_{X, *}\mathcal{F})
\Rightarrow
H^{p + q}_\tau(X, \mathcal{F})
$$
과 가정된 소멸을 결합하면 완전열
$$
0 \to
H^{n + 1}_{\tau'}(X, \epsilon_{X, *}\mathcal{F}) \to
H^{n + 1}_\tau(X, \mathcal{F}) \to
H^0_{\tau'}(X, R^{n + 1}\epsilon_{X, *}\mathcal{F})
$$
을 얻는다. 이는 보조정리를 다시 서술한 것이다.
\end{proof}

\begin{lemma}
\label{sites-cohomology-lemma-make-class-zero-general}
Situation \ref{sites-cohomology-situation-compare}에서 $f : X \to Y$가
$\mathcal{P}$에 속하고 $\{X \to Y\}$가 $\tau$-덮개라고 하자.
$\mathcal{F}'$가 $\mathcal{A}'_Y$에 속한다고 하자.
$n \geq 0$이고
$$
\theta \in
\text{Equalizer}\left(
\xymatrix{
H^{n + 1}_{\tau'}(X, \mathcal{F}')
\ar@<1ex>[r] \ar@<-1ex>[r] &
H^{n + 1}_{\tau'}(X \times_Y X, \mathcal{F}')
}
\right)
$$
이면, $\theta$의
$H^{n + 1}_{\tau'}(Y_i \times_Y X, \mathcal{F}')$로의 제한이
영이 되게 하는 $\tau'$-덮개 $\{Y_i \to Y\}$가 존재한다.
\end{lemma}

\begin{proof}
(\ref{sites-cohomology-item-base-change-P})에 의해 $X \times_Y X$가 존재한다.
$\mathcal{C}/Y$의 $Z$에 대해, $\mathcal{F}'$의
$\mathcal{C}_{\tau'}/Z$로의 제한을 $\mathcal{F}'|_Z$로 쓰자.
$H^{n + 1}_{\tau'}(X, \mathcal{F}') =
H^{n + 1}(\mathcal{C}_{\tau'}/X, \mathcal{F}'|_X)$임을 상기하자.
Lemma \ref{sites-cohomology-lemma-cohomology-of-open}을 보라. 보조정리의 주장은
$\theta$의 상
$\overline{\theta} \in H^0(Y, R^{n + 1}f_{\tau', *}\mathcal{F}'|_X)$가
영이라는 것이다. 다음 데카르트 그림을 생각하자.
$$
\xymatrix{
X \times_Y X \ar[d]_{\text{pr}_1} \ar[r]_{\text{pr}_2} &
X \ar[d]^f \\
X \ar[r]^f & Y
}
$$
자명한 기저변환(Lemma \ref{sites-cohomology-lemma-localize-cartesian-square})에 의해
$$
f_{\tau'}^{-1}R^{n + 1}f_{\tau', *}(\mathcal{F}'|_X) =
R^{n + 1}\text{pr}_{1, \tau', *}(\mathcal{F}'|_{X \times_Y X})
$$
가 성립한다. 만일
$\text{pr}_1^{-1}\theta = \text{pr}_2^{-1}\theta$이면,
$f_{\tau'}^{-1}R^{n + 1}f_{\tau', *}(\mathcal{F}'|_X)$의 단면
$f_{\tau'}^{-1}\overline{\theta}$는 영이다. 실제로
$\text{pr}_1^{-1}\theta$가
$H^0(X, R^{n + 1}\text{pr}_{1, \tau', *}(\mathcal{F}'|_{X \times_Y X}))$의
영원소로 가는 것은 분명하다. (\ref{sites-cohomology-item-restriction-A})에 의해
$\mathcal{F}'|_X$는 $\mathcal{A}'_X$에 속한다. 따라서
(\ref{sites-cohomology-item-A-and-P})에 의해
$\mathcal{G}' = R^{n + 1}f_{\tau', *}(\mathcal{F}'|_X)$는
$\mathcal{A}'_Y$의 대상이다. 그러므로 (\ref{sites-cohomology-item-A-sheaf})에 의해
$\mathcal{G}'$는 $\tau$-덮개에 대한 층 조건을 만족한다.
$\{X \to Y\}$가 $\tau$-덮개이므로 제한사상
$\mathcal{G}'(Y) \to \mathcal{G}'(X)$는 단사이다.
따라서 $\overline{\theta}$는 영이다.
\end{proof}

\begin{lemma}
\label{sites-cohomology-lemma-induction-step-V-C-general}
Situation \ref{sites-cohomology-situation-compare}에서
$(V_n) \Rightarrow (V_{n + 1})$이 성립한다.
\end{lemma}

\begin{proof}
$\mathcal{C}$의 $X$와 $\mathcal{A}_X$의 $\mathcal{F}$를 택하자.
어떤 $U/X$에 대해 $\xi \in H^{n + 1}_\tau(U, \mathcal{F})$라 하자.
$\xi$가 $U$의 어떤 $\tau'$-덮개의 각 원소 위에서 영으로 제한됨을
보여야 한다. Lemma \ref{sites-cohomology-lemma-higher-direct-images}를 보라.
따라서 $U$를 $U$의 어떤 $\tau'$-덮개의 원소들로 바꾸어도 된다.

\medskip\noindent
코호몰로지의 국소성(Lemma
\ref{sites-cohomology-lemma-kill-cohomology-class-on-covering})에 의해
$\xi$가 각 $U_i$ 위에서 영으로 제한되게 하는 $\tau$-덮개
$\{U_i \to U\}$를 택할 수 있다. (\ref{sites-cohomology-item-refine-tau-by-P})와 같은
$\{V_j \to V\}$, $\{f_j : W_j \to V_j\}$,
$\{W_{jk} \to W_j\}$를 택하자. (\ref{sites-cohomology-item-base-change-P})에 의해
허용되는 대로 $U$를 $V_j$로, $\mathcal{F}$를 그
$\mathcal{C}_\tau/V_j$로의 제한으로 바꾸면 다음 문단의 경우로
귀착된다.

\medskip\noindent
여기서 $f : X \to Y$는 $\mathcal{P}$의 원소이고
$\{X \to Y\}$는 $\tau$-덮개이며, $\mathcal{F}$는
$\mathcal{A}_Y$의 대상이고,
$\xi \in H^{n + 1}_\tau(Y, \mathcal{F})$는 모든 $i \in I$에 대해
$\xi$가 $X_i$ 위에서 영으로 제한되게 하는 $\tau'$-덮개
$\{X_i \to X\}_{i \in I}$가 존재하는 유이다.
문제: $\xi$가 $Y$의 어떤 $\tau'$-덮개 위에서 영으로 제한됨을 보이라.

\medskip\noindent
Lemma \ref{sites-cohomology-lemma-V-implies-cohomology-extra-general}에 의해 상이
$\xi|_X$인 유일한 $\tau'$-코호몰로지 유
$\theta \in H^{n + 1}_{\tau'}(X, \epsilon_{X, *}\mathcal{F})$
가 존재한다. $\xi|_X$는 두 사영을 통해 $X \times_Y X$ 위의 같은
유로 당겨지므로, 유일성에 의해 $\theta$에 대해서도 마찬가지이다.
Lemma \ref{sites-cohomology-lemma-make-class-zero-general}에 의해 $Y$를 어떤
$\tau'$-덮개의 원소들로 바꾼 뒤 $\theta = 0$이라 가정해도 된다.
따라서 $\xi|_X$가 영이라고 가정해도 된다.

\medskip\noindent
이제 $f : X \to Y$가 $\mathcal{P}$의 원소이고
$\{X \to Y\}$가 $\tau$-덮개이며, $\mathcal{F}$가
$\mathcal{A}_Y$의 대상이고,
$\xi \in H^{n + 1}_\tau(Y, \mathcal{F})$
가 $H^{n + 1}_\tau(X, \mathcal{F})$에서 영으로 간다고 하자.
문제: $\xi$가 $Y$의 어떤 $\tau'$-덮개 위에서 영으로 제한됨을 보이라.

\medskip\noindent
가정에 의해 $\xi$는 사상
$$
\mathcal{F} \longrightarrow Rf_{\tau, *}f_\tau^{-1}\mathcal{F}
$$
에서 영으로 간다. Lemma \ref{sites-cohomology-lemma-Leray-unbounded}를 사용하라.
절단의 특수 삼각형(Derived Categories, Remark
\ref{derived-remark-truncation-distinguished-triangle})을 사용하는
간단한 논증으로 $\xi$가 사상
$$
\mathcal{F} \longrightarrow \tau_{\leq n}Rf_{\tau, *}f_\tau^{-1}\mathcal{F}
$$
에서도 영으로 감을 알 수 있다. 이를 다음에서 정의되는
$\epsilon_{Y, *}\mathcal{F} \to K$와 비교하겠다.
$$
K = \tau_{\leq n}Rf_{\tau', *}f_{\tau'}^{-1}\epsilon_{Y, *}\mathcal{F} =
\tau_{\leq n}Rf_{\tau', *}\epsilon_{X, *}f_{\tau}^{-1}\mathcal{F}
$$
등식
$\epsilon_{X, *} f_\tau^{-1} = f_{\tau'}^{-1} \epsilon_{Y, *}$
은 (\ref{sites-cohomology-equation-commutative-epsilon})의 성질이다. 사상
$$
Rf_{\tau', *}\epsilon_{X, *}f_{\tau}^{-1}\mathcal{F} \longrightarrow
Rf_{\tau', *}R\epsilon_{X, *}f_{\tau}^{-1}\mathcal{F} =
R\epsilon_{Y, *}Rf_{\tau, *}f_\tau^{-1}\mathcal{F}
$$
을 생각하자. 이는 Lemma \ref{sites-cohomology-lemma-V-implies-C-general}의 증명에서
사용되었고, 수반에 의해 사상
$$
\epsilon_Y^{-1} Rf_{\tau', *}\epsilon_{X, *}f_{\tau}^{-1}\mathcal{F} \to
Rf_{\tau, *}f_\tau^{-1}\mathcal{F}
$$
을 유도한다. 절단을 취하면 사상
$$
\epsilon_Y^{-1}K
\longrightarrow
\tau_{\leq n}Rf_{\tau, *}f_\tau^{-1}\mathcal{F}
$$
을 얻는다. 이는 Lemma \ref{sites-cohomology-lemma-V-implies-C-general}에 의해
동형이다. 실제로 Lemma \ref{sites-cohomology-lemma-A}에 의해
$f_\tau^{-1}\mathcal{F}$는 $\mathcal{A}_X$에 속하므로 그 보조정리를
적용할 수 있다. 특수 삼각형
$$
\epsilon_{Y, *}\mathcal{F} \to K \to L \to \epsilon_{Y, *}\mathcal{F}[1]
$$
을 택하자. $\{X \to Y\}$가 $\tau$-덮개이므로 사상
$\mathcal{F} \to f_{\tau, *}f_\tau^{-1}\mathcal{F}$는 단사이다.
따라서
$\epsilon_{Y, *}\mathcal{F} \to
\epsilon_{Y, *}f_{\tau, *}f_\tau^{-1}\mathcal{F} =
f_{\tau', *}f_{\tau'}^{-1}\epsilon_{Y, *}\mathcal{F}$
도 단사이다. 그러므로 $L$은 차수 $0, \ldots, n$에서만 영이 아닌
코호몰로지 층을 갖는다.
(\ref{sites-cohomology-item-restriction-A})와 (\ref{sites-cohomology-item-A-and-P})에 의해
$f_{\tau', *}f_{\tau'}^{-1}\epsilon_{Y, *}\mathcal{F}$는
$\mathcal{A}'_Y$에 속하므로
$$
H^0(L) = \Coker(\epsilon_{Y, *}\mathcal{F} \to
f_{\tau', *}f_{\tau'}^{-1}\epsilon_{Y, *}\mathcal{F})
$$
은 약한 Serre 부분범주 $\mathcal{A}'_Y$에 속한다.
$1 \leq i \leq n$에 대해서도 (\ref{sites-cohomology-item-restriction-A})와
(\ref{sites-cohomology-item-A-and-P})에 의해
$H^i(L) = R^if_{\tau', *}f_{\tau'}^{-1}\epsilon_{Y, *}\mathcal{F}$는
$\mathcal{A}'_Y$에 속한다. 위 특수 삼각형을 $\epsilon_Y$로 당기면
특수 삼각형
$$
\mathcal{F} \to \tau_{\leq n}Rf_{\tau, *}f_\tau^{-1}\mathcal{F}
\to \epsilon_Y^{-1}L \to \mathcal{F}[1]
$$
을 얻는다. $\xi$가 가운데 항에서 영으로 가므로, $\xi$는 어떤 원소
$\xi' \in H^n_\tau(Y, \epsilon_Y^{-1}L)$의 상이다.
Lemma \ref{sites-cohomology-lemma-V-implies-cohomology-general}에 의해
$$
H^n_{\tau'}(Y, L) = H^n_\tau(Y, \epsilon_Y^{-1}L),
$$
이다. 따라서 $\xi'$를 $H^n_{\tau'}(Y, L)$의 원소로 올린 뒤
$H^{n + 1}_{\tau'}(Y, \epsilon_{Y, *}\mathcal{F})$로 가는 경계를
취하면, $\xi$가 표준사상
$H^{n + 1}_{\tau'}(Y, \epsilon_{Y, *}\mathcal{F}) \to
H^{n + 1}_\tau(Y, \mathcal{F})$의 상에 속함을 알 수 있다.
$H^{n + 1}_{\tau'}(Y,\epsilon_{Y, *}\mathcal{F})$에 대한 코호몰로지의
국소성으로 결론을 얻는다. Lemma
\ref{sites-cohomology-lemma-kill-cohomology-class-on-covering}를 보라.
\end{proof}

\begin{lemma}
\label{sites-cohomology-lemma-V-C-all-n-general}
Situation \ref{sites-cohomology-situation-compare}에서 $(V_n)$은 모든 $n$에 대해
성립한다. 더욱이 다음이 성립한다.
\begin{enumerate}
\item $\mathcal{C}$의 $X$와
$K' \in D^+_{\mathcal{A}'_X}(\mathcal{C}_{\tau'}/X)$에 대해 사상
$K' \to R\epsilon_{X, *}(\epsilon_X^{-1}K')$는 동형이다.
\item $\mathcal{P}$의 $f : X \to Y$와
$K' \in D^+_{\mathcal{A}'_X}(\mathcal{C}_{\tau'}/X)$에 대해
$Rf_{\tau', *}K' \in D^+_{\mathcal{A}'_X}(\mathcal{C}_{\tau'}/Y)$이고
$\epsilon_Y^{-1}(Rf_{\tau', *}K') = Rf_{\tau, *}(\epsilon_X^{-1}K')$.
\end{enumerate}
\end{lemma}

\begin{proof}
$(V_0)$은 빈 조건이므로 성립한다. 따라서 Lemma
\ref{sites-cohomology-lemma-induction-step-V-C-general}에 의해 $(V_n)$은 모든
$n$에 대해 성립한다.

\medskip\noindent
(1)의 증명. 대상 $K = \epsilon_X^{-1}K'$의 코호몰로지 층
$H^i(K) = \epsilon_X^{-1}H^i(K')$는 $\mathcal{A}_X$에 속한다.
따라서 스펙트럼열
$$
E_2^{p, q} = R^p\epsilon_{X, *} H^q(K) \Rightarrow
H^{p + q}(R\epsilon_{X, *}K)
$$
은 모든 $n$에 대한 $(V_n)$에 의해 퇴화하고
$$
H^n(R\epsilon_{X, *}K) = \epsilon_{X, *}H^n(K) =
\epsilon_{X, *}\epsilon_X^{-1}H^i(K') = H^i(K').
$$
를 얻는다. 여기서도 $H^i(K')$가 $\mathcal{A}'_X$에 속한다는
사실을 사용했다. 따라서 표준사상
$K' \to R\epsilon_{X, *}(\epsilon_X^{-1}K')$는 동형이다.

\medskip\noindent
(2)의 증명. 스펙트럼열
$$
E_2^{p, q} = R^pf_{\tau', *}H^q(K') \Rightarrow R^{p + q}f_{\tau', *}K'
$$
과, (\ref{sites-cohomology-item-A-and-P})에 의해
$R^pf_{\tau', *}H^q(K')$가 $\mathcal{A}'_Y$에 속한다는 사실,
$\mathcal{A}'_Y$가 $\textit{Ab}(\mathcal{C}_{\tau'}/Y)$의 약한 Serre
부분범주라는 사실, 그리고 Homology, Lemma
\ref{homology-lemma-biregular-ss-converges}를 사용하면
$Rf_{\tau', *}K' \in D^+_{\mathcal{A}'_X}(\mathcal{C}_{\tau'}/X)$.
증명을 끝내려면 기저변환 사상
$$
\epsilon_Y^{-1}(Rf_{\tau', *}K')
\longrightarrow
Rf_{\tau, *}(\epsilon_X^{-1}K')
$$
이 동형임을 보여야 한다. 위 스펙트럼열을 스펙트럼열
$$
E_2^{p, q} = R^pf_{\tau, *}H^q(\epsilon_X^{-1}K')
\Rightarrow R^{p + q}f_{\tau, *}\epsilon_X^{-1}K'
$$
과 비교하면, $K'$가 $\mathcal{A}'_X$에 속하는 영이 아닌 코호몰로지
층 $\mathcal{F}'$ 하나만 갖는 경우로 귀착된다. 세부사항은 생략한다.
그러면 Lemma \ref{sites-cohomology-lemma-V-implies-C-general}에 의해
$\epsilon_Y^{-1}R^if_{\tau', *}\mathcal{F}' =
R^if_{\tau, *}\epsilon_X^{-1}\mathcal{F}'$가 모든 $i$에 대해
성립하며 증명이 끝난다.
\end{proof}

\begin{lemma}
\label{sites-cohomology-lemma-cohomological-descent-general}
Situation \ref{sites-cohomology-situation-compare}에서, 임의의 $\mathcal{C}$의 $X$에
대해 범주
$\mathcal{A}_X \subset \textit{Ab}(\mathcal{C}_\tau/X)$
는 약한 Serre 부분범주이고, 함자
$$
R\epsilon_{X, *} :
D^+_{\mathcal{A}_X}(\mathcal{C}_\tau/X)
\longrightarrow
D^+_{\mathcal{A}'_X}(\mathcal{C}_{\tau'}/X)
$$
는 $\epsilon_X^{-1}$을 준역으로 갖는 동치이다.
\end{lemma}

\begin{proof}
$\mathcal{A}_X$에 대해 Homology, Lemma
\ref{homology-lemma-characterize-weak-serre-subcategory}에 나열된
조건들을 확인해야 한다. $\varphi : \mathcal{F} \to \mathcal{G}$가
$\mathcal{A}_X$의 사상이면
$\epsilon_{X, *}\varphi :
\epsilon_{X, *}\mathcal{F} \to \epsilon_{X, *}\mathcal{G}$는
$\mathcal{A}'_X$의 사상이다. 따라서
$\Ker(\epsilon_{X, *}\varphi)$와
$\Coker(\epsilon_{X, *}\varphi)$는 $\mathcal{A}'_X$의 대상이다.
실제로 이는 $\textit{Ab}(\mathcal{C}_{\tau'}/X)$의 약한 Serre
부분범주이다. $\epsilon_X^{-1}$을 적용하면 완전열
$$
0 \to
\epsilon_X^{-1}\Ker(\epsilon_{X, *}\varphi) \to
\mathcal{F} \to \mathcal{G} \to
\epsilon_X^{-1}\Coker(\epsilon_{X, *}\varphi) \to 0
$$
을 얻고 $\Ker(\varphi)$와 $\Coker(\varphi)$가
$\mathcal{A}_X$에 속함을 알 수 있다. 마지막으로
$$
0 \to \mathcal{F}_1 \to \mathcal{F}_2 \to \mathcal{F}_3 \to 0
$$
이 $\textit{Ab}(\mathcal{C}_\tau/X)$의 짧은 완전열이고
$\mathcal{F}_1$과 $\mathcal{F}_3$가 $\mathcal{A}_X$에 속한다고 하자.
$\epsilon_{X, *}$를 적용하면 완전열
$$
0 \to 
\epsilon_{X, *}\mathcal{F}_1 \to
\epsilon_{X, *}\mathcal{F}_2 \to
\epsilon_{X, *}\mathcal{F}_3 \to
R^1\epsilon_{X, *}\mathcal{F}_1 = 0
$$
을 얻으며, 마지막 소멸은 Lemma
\ref{sites-cohomology-lemma-V-C-all-n-general}에서 따른다. 그 범주는
$\textit{Ab}(\mathcal{C}_{\tau'}/X)$의 약한 Serre 부분범주이므로
$\epsilon_{X, *}\mathcal{F}_2$는 $\mathcal{A}'_X$에 속한다.
$\epsilon_X$로 당기면 $\mathcal{F}_2$가 $\mathcal{A}_X$에
속한다는 결론을 얻는다.

\medskip\noindent
따라서 $\mathcal{A}_X$는 $\textit{Ab}(\mathcal{C}_\tau/X)$의 약한
Serre 부분범주이고, 범주
$D^+_{\mathcal{A}_X}(\mathcal{C}_\tau/X)$를 생각할 수 있다.
$\epsilon_X^{-1} : \mathcal{A}'_X \to \mathcal{A}_X$가 동치이고
$\mathcal{F}' \to R\epsilon_{X, *}\epsilon_X^{-1}\mathcal{F}'$
가 $\mathcal{A}'_X$의 $\mathcal{F}'$에 대해 동형임에 유의하자.
실제로 Lemma \ref{sites-cohomology-lemma-V-C-all-n-general}에 의해 모든 $n$에 대해
$(V_n)$이 성립한다. 따라서 Lemma \ref{sites-cohomology-lemma-equivalence-bounded}로
결론을 얻는다.
\end{proof}

\begin{lemma}
\label{sites-cohomology-lemma-compare-cohomology-general}
Situation \ref{sites-cohomology-situation-compare}에서 $X$가 $\mathcal{C}$에 속한다고 하자.
\begin{enumerate}
\item $\mathcal{A}'_X$의 $\mathcal{F}'$에 대해
$H^n_{\tau'}(X, \mathcal{F}') = H^n_\tau(X, \epsilon_X^{-1}\mathcal{F}')$,
\item $K' \in D^+_{\mathcal{A}'_X}(\mathcal{C}_{\tau'}/X)$에 대해
$H^n_{\tau'}(X, K') = H^n_\tau(X, \epsilon_X^{-1}K')$.
\end{enumerate}
\end{lemma}

\begin{proof}
이는 Lemma \ref{sites-cohomology-lemma-V-C-all-n-general}에서
Remark \ref{sites-cohomology-remark-before-Leray}에 의해 따른다.
\end{proof}
























\section{Hausdorff이고 국소 준콤팩트인 공간 위의 코호몰로지}
\label{sites-cohomology-section-cohomology-LC}

\noindent
문헌에서 흔히 “국소 콤팩트”라고 하는 대신 “Hausdorff이고 국소
준콤팩트”라고 말하는 관례를 계속 사용한다. 대상이 Hausdorff이고
국소 준콤팩트인 위상공간이며 사상이 연속사상인 범주를
$\textit{LC}$로 쓰자.

\begin{lemma}
\label{sites-cohomology-lemma-LC-basic}
범주 $\textit{LC}$는 올곱과 끝대상을 가지므로 임의의 유한 극한을
갖는다. $\textit{LC}$의 사상 $X \to Z$와 $Y \to Z$가 주어지고
$X$와 $Y$가 준콤팩트이면 $X \times_Z Y$도 준콤팩트이다.
\end{lemma}

\begin{proof}
끝대상은 한원소 공간이다. $\textit{LC}$의 사상 $X \to Z$와
$Y \to Z$가 주어지면 올곱 $X \times_Z Y$는 $X \times Y$의
부분공간이다. Topology, Section \ref{topology-section-Hausdorff}에
의해 $X \times Y$가 Hausdorff이므로 $X \times_Z Y$도 Hausdorff이다.

\medskip\noindent
$X$와 $Y$가 준콤팩트이면 Topology, Theorem
\ref{topology-theorem-tychonov}에 의해 $X \times Y$는 준콤팩트이다.
$X \times_Z Y$는 $X \times Y$의 닫힌 부분집합이므로(Topology,
Lemma \ref{topology-lemma-fibre-product-closed}), Topology, Lemma
\ref{topology-lemma-closed-in-quasi-compact}에 의해
$X \times_Z Y$는 준콤팩트이다.

\medskip\noindent
마지막으로 일반적인 경우로 돌아가자. $x \in X$와 $y \in Y$에 대해
준콤팩트 근방 $x \in E \subset X$와 $y \in F \subset Y$를 택하면,
위 결과에 의해 $E \times_Z F$는 $(x, y)$의 준콤팩트 근방이다.
따라서 Topology, Lemma
\ref{topology-lemma-locally-quasi-compact-Hausdorff}에 의해
$X \times_Z Y$는 $\textit{LC}$의 대상이다.
\end{proof}

\noindent
$\textit{LC}$에는 보통 위상보다 더 강한 위상을 줄 수 있다.

\begin{definition}
\label{sites-cohomology-definition-covering-LC}
$\{f_i : X_i \to X\}$를 범주 $\textit{LC}$에서 목표가 고정된
사상들의 족이라 하자. 모든 $x \in X$에 대해
$i_1, \ldots, i_n \in I$ 및 준콤팩트 부분집합
$E_j \subset X_{i_j}$가 존재하여 $\bigcup f_{i_j}(E_j)$가
$x$의 근방이 되면, 이 족을 {\it qc 덮개}\footnote{이는 표준적이지
않은 표기이다. 스킴의 fpqc 덮개를 떠올리게 하려고 택했다.}라 한다.
\end{definition}

\noindent
$\textit{LC}$의 대상의 열린 덮개 $X = \bigcup U_i$는
$X$가 국소 준콤팩트이므로 qc 덮개 $\{U_i \to X\}$를 준다.
필수적인 보조정리부터 시작한다.

\begin{lemma}
\label{sites-cohomology-lemma-qc}
$X$를 Hausdorff이고 국소 준콤팩트인 공간, 즉
$\textit{LC}$의 대상이라 하자.
\begin{enumerate}
\item $X' \to X$가 $\textit{LC}$에서 동형이면
$\{X' \to X\}$는 qc 덮개이다.
\item $\{f_i : X_i \to X\}_{i\in I}$가 qc 덮개이고 각 $i$에 대해
$\{g_{ij} : X_{ij} \to X_i\}_{j\in J_i}$가 qc 덮개이면
$\{X_{ij} \to X\}_{i \in I, j\in J_i}$도 qc 덮개이다.
\item $\{X_i \to X\}_{i\in I}$가 qc 덮개이고
$X' \to X$가 $\textit{LC}$의 사상이면
$\{X' \times_X X_i \to X'\}_{i\in I}$도 qc 덮개이다.
\end{enumerate}
\end{lemma}

\begin{proof}
열린 덮개가 qc 덮개라는 위의 언급으로 부분 (1)이 성립한다.

\medskip\noindent
(2)의 증명. $x \in X$라 하자. $\bigcup f_{i_a}(E_a)$가 $x$의
근방이 되게 하는 $i_1, \ldots, i_n \in I$와 준콤팩트 부분집합
$E_a \subset X_{i_a}$를 택하자. 각 $e \in E_a$에 대해 유한 부분집합
$J_e \subset J_{i_a}$와 준콤팩트 부분집합
$F_{e, j} \subset X_{ij}$, $j \in J_e$를 택하여
$\bigcup g_{ij}(F_{e, j})$가 $e$의 근방이 되게 할 수 있다.
$E_a$가 준콤팩트이므로 다음을 만족하는 유한한 원소들
$e_1, \ldots, e_{m_a}$를 택할 수 있다.
$$
E_a \subset
\bigcup\nolimits_{k = 1, \ldots, m_a}
\bigcup\nolimits_{j \in J_{e_k}} g_{ij}(F_{e_k, j})
$$
그러면
$$
\bigcup\nolimits_{a = 1, \ldots, n}
\bigcup\nolimits_{k = 1, \ldots, m_a}
\bigcup\nolimits_{j \in J_{e_k}} f_i(g_{ij}(F_{e_k, j}))
$$
는 $x$의 근방이다.

\medskip\noindent
(3)의 증명. $x' \in X'$를 점이라 하고 그 상을 $x \in X$라 하자.
$\bigcup f_{i_j}(E_j)$가 $x$의 근방이 되게 하는
$i_1, \ldots, i_n \in I$와 준콤팩트 부분집합
$E_j \subset X_{i_j}$를 택하자. $x'$의 준콤팩트 근방
$F \subset X'$를 택하여 $x$의 준콤팩트 근방
$\bigcup f_{i_j}(E_j)$ 안으로 가게 하자. 그러면
$F \times_X E_j \subset X' \times_X X_{i_j}$는 준콤팩트 부분집합이고,
$F$는 사상 $\coprod F \times_X E_j \to F$의 상이다.
따라서 기저변환은 qc 덮개이고 증명이 끝난다.
\end{proof}

\noindent
$\textit{LC}$의 모든 대상이 Hausdorff이므로 $\textit{LC}$의 임의의
사상 $f : X \to Y$는 위상공간의 분리 연속사상이다. 따라서
$f$가 위상공간의 고유사상일 필요충분조건은 $f$가 보편적으로
닫힌 것이다. Topology, Section \ref{topology-section-proper}의
논의를 보라.

\begin{lemma}
\label{sites-cohomology-lemma-proper-surjective-is-qc-covering}
$f : X \to Y$를 $\textit{LC}$의 사상이라 하자.
$f$가 고유이고 전사이면 $\{f : X \to Y\}$는 qc 덮개이다.
\end{lemma}

\begin{proof}
$y \in Y$를 점이라 하자. 각 $x \in X_y$에 대해 준콤팩트 근방
$E_x \subset X$를 택하고 열린집합 $x \in U_x \subset E_x$를 택하자.
$f$가 고유이므로 올 $X_y$는 준콤팩트이고, 다음을 만족하는
$x_1, \ldots, x_n \in X_y$를 택할 수 있다.
$X_y \subset U_{x_1} \cup \ldots \cup U_{x_n}$.
$f(E_{x_1}) \cup \ldots \cup f(E_{x_n})$가 $y$의 근방임을 주장한다.
실제로 $f$가 닫힌 사상이므로(Topology, Theorem
\ref{topology-theorem-characterize-proper}),
$Z = f(X \setminus U_{x_1} \cup \ldots \cup U_{x_n})$는 $y$를
포함하지 않는 $Y$의 닫힌 부분집합이다. $f$가 전사이므로 원하는 대로
$Y \setminus Z$는 $f(E_{x_1}) \cup \ldots \cup f(E_{x_n})$에 포함된다.
\end{proof}

\noindent
몇 가지 집합론적 문제를 제외하면 Lemma \ref{sites-cohomology-lemma-qc}는
$\textit{LC}$가 qc 덮개들의 모음과 함께 사이트를 이룸을 보여준다.
집합론적 문제를 피하도록 적절히 수정한 이 사이트를
$\textit{LC}_{qc}$로 쓰겠다.

\begin{remark}[집합론적 문제]
\label{sites-cohomology-remark-set-theoretic-LC}
$\textit{LC}$의 대상들은 진정한 모임을 이루므로 이 범주는 “큰”
범주이다. 마찬가지로 덮개들도 진정한 모임을 이룬다. 위상공간
$X$의 {\it 크기}를 $X$의 점들의 집합의 기수로 정의하자.
기수들 위의 함수 $Bound$를, 예를 들어 Sets, Equation
(\ref{sets-equation-bound})에서처럼 택하자. 마지막으로
$S_0$를 $\textit{LC}$의 대상들의 초기 집합이라 하자. 예를 들면
$S_0 = \{(\mathbf{R}, \text{euclidean topology})\}$로 둘 수 있다.
Sets, Lemma \ref{sets-lemma-construct-category}에서와 꼭 마찬가지로
다음을 만족하는 극한 순서수 $\alpha$를 택할 수 있다.
$\textit{LC}_\alpha = \textit{LC} \cap V_\alpha$
는 $S_0$를 포함하고, $\textit{LC}$에 존재하는 모든 가산 극한과
쌍대극한 아래 보존된다. 더욱이 $X \in \textit{LC}_\alpha$이고
$Y \in \textit{LC}$이며
$\text{size}(Y) \leq Bound(\text{size}(X))$이면, $Y$는
$\textit{LC}_\alpha$의 어떤 대상과 동형이다. 다음으로 Sets, Lemma
\ref{sets-lemma-coverings-site}를 적용하여 $\textit{LC}_\alpha$ 위의
qc 덮개들의 집합 $\text{Cov}$를 택하되, $\textit{LC}_\alpha$의
모든 qc 덮개가 이 집합의 어떤 덮개와 조합적으로 동치가 되게 한다.
이렇게 사이트 $(\textit{LC}_\alpha, \text{Cov})$를 얻고 이를
$\textit{LC}_{qc}$로 쓰겠다.
\end{remark}

\noindent
Remark \ref{sites-cohomology-remark-set-theoretic-LC}의 사이트 $\textit{LC}_{qc}$에는
둘째 위상이 있다. 즉, 대상 $X$가 주어지면 $X_i \to X$가 열린
몰입인 $\textit{LC}_{qc}$의 모든 덮개 $\{X_i \to X\}$를 생각할 수
있다. 이 사이트를 $\textit{LC}_{Zar}$로 쓰자. 항등함자
$\textit{LC}_{Zar} \to \textit{LC}_{qc}$는 연속이고 사이트의 사상
$$
\epsilon : \textit{LC}_{qc} \longrightarrow \textit{LC}_{Zar}
$$
을 정의한다. Section \ref{sites-cohomology-section-compare}을 보라.
Hausdorff이고 국소 준콤팩트인 위상공간 $X$, 더 정확히는
$X \in \Ob(\textit{LC}_{qc})$에 대해 유도된 사상을
$$
\epsilon_X : \textit{LC}_{qc}/X \longrightarrow \textit{LC}_{Zar}/X
$$
로 쓰겠다(Sites, Lemma \ref{sites-lemma-localize-morphism}를 보라).
$X_{Zar}$를 그 대상들이 $X$의 열린집합인 사이트라 하자.
Sites, Example \ref{sites-example-site-topological}을 보라.
다음 사이트의 사상이 있다.
$$
\pi_X : \textit{LC}_{Zar}/X \longrightarrow X_{Zar}
$$
이는 연속함자
$X_{Zar} \to \textit{LC}_{Zar}/X$, $U \mapsto U$.
로 주어진다. 실제로 $X_{Zar}$는 올곱과 끝대상을 가지고 위 함자는
이들과 가환하므로 Sites, Proposition
\ref{sites-proposition-get-morphism}을 적용할 수 있다. 흔히
$\pi$를 토포스의 사상
$$
\pi_X : \Sh(\textit{LC}_{Zar}/X) \longrightarrow \Sh(X)
$$
으로 생각하며, 이때 등식 $\Sh(X_{Zar}) = \Sh(X)$를 사용한다.

\begin{lemma}
\label{sites-cohomology-lemma-describe-pullback-pi}
$X$를 $\textit{LC}_{qc}$의 대상이라 하고 $\mathcal{F}$를
$X$ 위의 층이라 하자. 규칙
$$
\textit{LC}_{qc}/X \longrightarrow \textit{Sets},\quad
(f : Y \to X) \longmapsto \Gamma(Y, f^{-1}\mathcal{F})
$$
은 층이고, 따라서 더욱이 $\textit{LC}_{Zar}/X$ 위의 층이기도 하다.
이 층은 $\textit{LC}_{Zar}/X$ 위에서는
$\pi_X^{-1}\mathcal{F}$와 같고, $\textit{LC}_{qc}/X$ 위에서는
$\epsilon_X^{-1}\pi_X^{-1}\mathcal{F}$와 같다.
\end{lemma}

\begin{proof}
보조정리의 공식으로 주어지는 준층을 $\mathcal{G}$로 쓰자.
물론 공식의 당김 $f^{-1}$은 위상공간 위 층들의 보통 당김을 뜻한다.
정의에서 곧바로 $\mathcal{G}$가 Zar 위상에 대한 층임을 알 수 있다.

\medskip\noindent
$Y \to X$를 $\textit{LC}_{qc}$의 사상이라 하고,
$\mathcal{V} = \{g_i : Y_i \to Y\}_{i \in I}$를 qc 덮개라 하자.
$\mathcal{G}$가 qc 위상에 대한 층임을 보이려면 사상
$\mathcal{G}(Y) \to H^0(\mathcal{V}, \mathcal{G})$
이 동형임을 보이면 충분하다. Sites, Section
\ref{sites-section-sheafification}을 보라. qc 덮개는 전사이고 줄기에서
단면들의 같음을 판별할 수 있으므로 이 사상은 단사이다(Sheaves,
Lemmas \ref{sheaves-lemma-sheaf-subset-stalks}와
\ref{sheaves-lemma-stalk-pullback-presheaf}를 사용하라). 따라서
$\mathcal{G}$는 $\textit{LC}_{qc}$ 위의 분리 준층이다. 그러므로
임의의 원소
$(s_i) \in H^0(\mathcal{V}, \mathcal{G})$
가 $\mathcal{V}$를 어떤 세분으로 바꾼 뒤 $\mathcal{G}(Y)$의 상에
속하는 원소로 감을 보이면 충분하다(Sites, Theorem
\ref{sites-theorem-plus}).

\medskip\noindent
$Y_{i, Zar}$ 위의 층과 $Y_i$ 위의 층을 동일시하면,
$\mathcal{G}|_{Y_{i, Zar}}$는 연속사상 $g_i : Y_i \to Y$에 의한
$f^{-1}\mathcal{F}$의 당김이다. 따라서 각 $j$에 대해 열린집합
$W_{ij} \subset Y$와 단면 $t_{ij} \in \mathcal{G}(W_{ij})$가 존재하고,
$V_{ij}$가 $W_{ij}$ 안으로 가며 $s|_{V_{ij}}$가 $t_{ij}$의 당김이
되도록 열린 덮개 $Y_i = \bigcup V_{ij}$를 택할 수 있다. 달리 말해
덮개 $\{Y_i \to Y\}$를 세분한 뒤, 열린집합
$W_i \subset Y$가 존재하여 $Y_i \to Y$가 $W_i$를 통해 분해되고
$W_i$ 위의 $\mathcal{G}$의 단면
$t_i$가 주어진 단면 $s_i$로 제한된다고 가정해도 된다.
더욱이 $y \in Y$가 $Y_i \to Y$와 $Y_j \to Y$의 상에 모두 속하면,
줄기 $f^{-1}\mathcal{F}_y$ 안의 상 $t_{i, y}$와 $t_{j, y}$는 같다.
실제로 $s_i$와 $s_j$는 $Y_i \times_Y Y_j$ 위에서 일치한다.
따라서 각 $y \in Y$에 대해 잘 정의된 원소 $t_y$가
$f^{-1}\mathcal{F}_y$에 있으며, $y$가 $Y_i \to Y$의 상에
속할 때 이는 $t_{i, y}$와 일치한다. 원소 $(t_y)$가 $Y$ 위의
$f^{-1}\mathcal{F}$의 대역단면에서 옴을 보이겠다. 그러면
보조정리의 증명이 끝난다.

\medskip\noindent
이를 $Y$ 위에서 국소적으로 보이면 충분하다. Sheaves, Section
\ref{sheaves-section-sheafification}을 보라. $y_0 \in Y$라 하자.
$i_1, \ldots, i_n \in I$와 준콤팩트 부분집합
$E_j \subset Y_{i_j}$를 택하여 $\bigcup g_{i_j}(E_j)$가
$y_0$의 근방이 되게 하자. $V \subset Y$를 $y_0$의 열린 근방으로
택하되, $\bigcup g_{i_j}(E_j)$에 포함되고
$W_{i_1} \cap \ldots \cap W_{i_n}$에도 포함되게 하자.
$t_{i_1, y_0} = \ldots = t_{i_n, y_0}$이므로 $V$를 줄인 뒤에는
$f^{-1}\mathcal{F}$의 단면 $t_{i_j}|_V$, $j = 1, \ldots, n$이
서로 일치한다고 가정할 수 있다. $V \subset \bigcup g_{i_j}(E_j)$이므로
$(t_y)_{y \in V}$가 이 단면에서 옴을 알 수 있다.

\medskip\noindent
이제 $\mathcal{G}$가 $\textit{LC}_{qc}$ 위에서는
$\epsilon_X^{-1}\pi_X^{-1}\mathcal{F}$와 같고, 각각
$\textit{LC}_{Zar}$ 위에서는 $\pi_X^{-1}\mathcal{F}$와 같음을
보여야 한다. 두 경우 모두 먼저 준층
$$
(f : Y \to X)
\longmapsto
\colim_{f(Y) \subset U \subset X} \mathcal{F}(U)
$$
을 취한 뒤 층화함으로써 당김을 정의한다. Zar 위상에서 층화하면
바로 우리의 층 $\mathcal{G}$가 나오며, $\mathcal{G}$가 qc 층이라는
사실은 같은 구성이 qc 위상에서도 성립함을 보여준다.
\end{proof}

\noindent
$X \in \Ob(\textit{LC}_{Zar})$라 하고 $\mathcal{H}$를
$\textit{LC}_{Zar}/X$ 위의 아벨 층이라 하자. 이제
$\textit{LC}_{Zar}/X$의 대상 $U$ 위에서 $\mathcal{H}$의
코호몰로지를 $H^n_{Zar}(U, \mathcal{H})$로 쓰겠다.

\begin{lemma}
\label{sites-cohomology-lemma-collect-true-things-Zar}
$X$를 $\textit{LC}_{Zar}$의 대상이라 하자. 그러면
\begin{enumerate}
\item $\mathcal{F} \in \textit{Ab}(X)$에 대해
$H^n_{Zar}(X, \pi_X^{-1}\mathcal{F}) = H^n(X, \mathcal{F})$,
\item $\pi_{X, *} : \textit{Ab}(\textit{LC}_{Zar}/X) \to \textit{Ab}(X)$
는 완전함자이다.
\item 수반의 단위원 $\text{id} \to \pi_{X, *} \circ \pi_X^{-1}$은
동형이다.
\item $K \in D(X)$에 대해 표준사상
$K \to R\pi_{X, *} \pi_X^{-1}K$는 동형이다.
\end{enumerate}
$f : X \to Y$를 $\textit{LC}_{Zar}$의 사상이라 하자. 그러면
\begin{enumerate}
\item[(5)] 토포스들의 가환도식
$$
\xymatrix{
\Sh(\textit{LC}_{Zar}/X) \ar[r]_{f_{Zar}} \ar[d]_{\pi_X} &
\Sh(\textit{LC}_{Zar}/Y) \ar[d]^{\pi_Y} \\
\Sh(X_{Zar}) \ar[r]^f &
\Sh(Y_{Zar})
}
$$
이 있다.
\item[(6)] $L \in D^+(Y)$에 대해
$H^n_{Zar}(X, \pi_Y^{-1}L) = H^n(X, f^{-1}L)$,
\item[(7)] $f$가 고유이면 다음이 성립한다.
\begin{enumerate}
\item $\pi_Y^{-1} \circ f_* = f_{Zar, *} \circ \pi_X^{-1}$는
$\Sh(X) \to \Sh(\textit{LC}_{Zar}/Y)$인 함자들의 등식이다.
\item $\pi_Y^{-1} \circ Rf_* = Rf_{Zar, *} \circ \pi_X^{-1}$는
$D^+(X) \to D^+(\textit{LC}_{Zar}/Y)$인 함자들의 등식이다.
\end{enumerate}
\end{enumerate}
\end{lemma}

\begin{proof}
(1)의 증명. 등식 $H^n_{Zar}(X, \pi_X^{-1}\mathcal{F}) = H^n(X, \mathcal{F})$
은 $\textit{LC}_{Zar}$에서 $X$의 덮개가 $X$의 열린 덮개와 같다는
자명한 관찰에서 나오는 일반적인 사실이다. 자세한 증명을 보고 싶은
독자는 함자 $X_{Zar} \to \textit{LC}_{Zar}$에 Lemma
\ref{sites-cohomology-lemma-cohomology-bigger-site}를 적용하면 된다.

\medskip\noindent
(2)의 증명. $\pi_X$를 정의할 때 사용한 함자
$X_{Zar} \to \textit{LC}_{Zar}/X$에 Sites, Lemma
\ref{sites-lemma-bigger-site}를 적용하면 어떤 토포스의 사상
$\tau_X : \Sh(X_{Zar}) \to \Sh(\textit{LC}_{Zar})$에 대해
$\pi_{X, *} = \tau_X^{-1}$임을 알 수 있으므로 성립한다.

\medskip\noindent
(3)의 증명. Sites, Lemma \ref{sites-lemma-bigger-site}에 의해
$\tau_X^{-1} \circ \pi_X^{-1}$이 항등함자이므로 성립한다.
또는 Lemma \ref{sites-cohomology-lemma-describe-pullback-pi}의
$\pi_X^{-1}$에 대한 구체적인 기술로부터 이를 이끌어낼 수도 있다.

\medskip\noindent
(4)의 증명. $K$를 나타내는 아벨 층들의 복합체에 (3)을 적용하라.

\medskip\noindent
(5)의 증명. 토포스의 사상 $f_{Zar}$는 Sites, Lemma
\ref{sites-lemma-relocalize}를 적용하여 얻는다. 우리의 경우에는
Sites, Lemma \ref{sites-lemma-relocalize-given-fibre-products}에 의해
연속함자 $Z/Y \mapsto Z \times_Y X/X$에서 얻는다. 대응하는
연속함자들이 올바르게 합성되므로 도식은 가환한다(Sites, Lemma
\ref{sites-lemma-composition-morphisms-sites}를 보라).

\medskip\noindent
(6)의 증명. $\mathcal{G}$가 $\textit{Ab}(Y)$에 속할 때
$H^n_{Zar}(X, \pi_Y^{-1}\mathcal{G}) =
H^n_{Zar}(X, f_{Zar}^{-1}\pi_Y^{-1}\mathcal{G})$
이다. Lemma \ref{sites-cohomology-lemma-cohomology-of-open}을 보라.
(5)의 도식이 가환하므로 이는
$H^n_{Zar}(X, \pi_X^{-1}f^{-1}\mathcal{G})$와 같다.
따라서 $L$이 차수 $0$의 층 하나로 이루어진 경우에는 (1)에서
결론이 나온다. 일반적인 경우에는 $L$을 아래로 유계인 아벨 층들의
복합체로 나타내면 된다.

\medskip\noindent
(7a)의 증명. $\mathcal{F}$를 $X$ 위의 층이라 하자.
$g : Z \to Y$를 $\textit{LC}_{Zar}/Y$의 대상이라 하자. 올곱
$$
\xymatrix{
Z' \ar[r]_{f'} \ar[d]_{g'} & Z \ar[d]^g \\
X \ar[r]^f & Y
}
$$
을 생각하자. 그러면
$$
(f_{Zar, *}\pi_X^{-1}\mathcal{F})(Z/Y) =
(\pi_X^{-1}\mathcal{F})(Z'/X)  =
\Gamma(Z', (g')^{-1}\mathcal{F})  =
\Gamma(Z, f'_*(g')^{-1}\mathcal{F})
$$
이며, 둘째 등식에는 Lemma \ref{sites-cohomology-lemma-describe-pullback-pi}를
사용했다. 한편
$$
(\pi_Y^{-1}f_*\mathcal{F})(Z/Y) = \Gamma(Z, g^{-1}f_*\mathcal{F})
$$
도 Lemma \ref{sites-cohomology-lemma-describe-pullback-pi}에 의해 성립한다.
따라서 집합의 층에 대한 고유 기저변환(Cohomology, Lemma
\ref{cohomology-lemma-proper-base-change-sheaves-of-sets})에 의해
두 집합은 표준적으로 동형이다. 이 동형은 제한사상과 양립하고
동형
$\pi_Y^{-1}f_*\mathcal{F} = f_{Zar, *}\pi_X^{-1}\mathcal{F}$.
를 정의한다. 따라서 함자들의 동형
$\pi_Y^{-1} \circ f_* = f_{Zar, *} \circ \pi_X^{-1}$.
을 얻는다.

\medskip\noindent
(7b)의 증명. $K \in D^+(X)$라 하자. Lemma
\ref{sites-cohomology-lemma-unbounded-describe-higher-direct-images}에 의해
$Rf_{Zar, *}\pi_X^{-1}K$의 $n$차 코호몰로지 층은 준층
$$
(g : Z \to Y) \longmapsto H^n_{Zar}(Z', \pi_X^{-1}K)
$$
에 결합된 층이다. 표기는 위와 같다. 다음을 관찰하자.
\begin{align*}
H^n_{Zar}(Z', \pi_X^{-1}K)
& =
H^n(Z', (g')^{-1}K) \\
& =
H^n(Z, Rf'_*(g')^{-1}K) \\
& =
H^n(Z, g^{-1}Rf_*K) \\
& =
H^n_{Zar}(Z, \pi_Y^{-1}Rf_*K)
\end{align*}
첫째 등식은 $K$와 $g' : Z' \to X$에 (6)을 적용한 것이다.
둘째 등식은 $f' : Z' \to Z$에 대한 Leray이다(Cohomology, Lemma
\ref{cohomology-lemma-before-Leray}). 셋째 등식은 고유 기저변환 정리
(Cohomology, Theorem \ref{cohomology-theorem-proper-base-change})이다.
넷째 등식은 $g : Z \to Y$와 $Rf_*K$에 (6)을 적용한 것이다.
따라서 $Rf_{Zar, *}\pi_X^{-1}K$와 $\pi_Y^{-1}Rf_*K$는 같은
코호몰로지 층들을 갖는다. 표준 기저변환 사상
$\pi_Y^{-1}Rf_*K \to Rf_{Zar, *}\pi_X^{-1}K$가 이 동형을 유도함을
확인하는 일은 생략한다.
\end{proof}

\noindent
Lemma \ref{sites-cohomology-lemma-describe-pullback-pi}의 상황에서 $\epsilon$과
$\pi$의 합성 및 등식 $\Sh(X) = \Sh(X_{Zar})$는 토포스의 사상
$$
a_X : \Sh(\textit{LC}_{qc}/X) \longrightarrow \Sh(X)
$$
을 결정한다.

\begin{lemma}
\label{sites-cohomology-lemma-push-pull-LC}
$f : X \to Y$를 $\textit{LC}_{qc}$의 사상이라 하자.
그러면 토포스들의 가환도식
$$
\vcenter{
\xymatrix{
\Sh(\textit{LC}_{qc}/X) \ar[r]_{f_{qc}} \ar[d]_{\epsilon_X} &
\Sh(\textit{LC}_{qc}/Y) \ar[d]^{\epsilon_Y} \\
\Sh(\textit{LC}_{Zar}/X) \ar[r]^{f_{Zar}} &
\Sh(\textit{LC}_{Zar}/Y)
}
}
\quad\text{and}\quad
\vcenter{
\xymatrix{
\Sh(\textit{LC}_{qc}/X) \ar[r]_{f_{qc}} \ar[d]_{a_X} &
\Sh(\textit{LC}_{qc}/Y) \ar[d]^{a_Y} \\
\Sh(X) \ar[r]^f &
\Sh(Y)
}
}
$$
이 있고 $a_X = \pi_X \circ \epsilon_X$,
$a_Y = \pi_X \circ \epsilon_X$이다. $f$가 고유이면
$a_Y^{-1} \circ f_* = f_{qc, *} \circ a_X^{-1}$이다.
\end{lemma}

\begin{proof}
토포스의 사상 $f_{qc}$는 Sites, Lemma
\ref{sites-lemma-relocalize}에서 얻는 사상이다. 우리의 경우에는
Sites, Lemma \ref{sites-lemma-relocalize-given-fibre-products}에 의해
연속함자 $Z/Y \mapsto Z \times_Y X/X$에서 얻는다. 대응하는
연속함자들이 올바르게 합성되므로 왼쪽 도식은 가환한다(Sites, Lemma
\ref{sites-lemma-composition-morphisms-sites}를 보라).
왼쪽 도식이 가환하고 Lemma \ref{sites-cohomology-lemma-collect-true-things-Zar}의
부분 (5)가 성립하므로 오른쪽 도식도 가환한다.

\medskip\noindent
마지막 주장의 증명. Lemma \ref{sites-cohomology-lemma-collect-true-things-Zar}의
부분 (7a)의 증명을 되풀이할 수도 있지만, 여기서는 그 결과로부터
이끌어내겠다. $\epsilon_{Y, *}$는 바탕 준층들 위에서 항등함자이므로
동형을 반영한다. Lemma \ref{sites-cohomology-lemma-describe-pullback-pi}의 기술은
$\epsilon_{Y, *} \circ a_Y^{-1} = \pi_Y^{-1}$임을 보여주며,
$X$에 대해서도 마찬가지이다. 표준사상
$a_Y^{-1}f_*\mathcal{F} \to f_{qc, *}a_X^{-1}\mathcal{F}$
이 동형임을 보이려면
$$
\pi_Y^{-1}f_*\mathcal{F} =
\epsilon_{Y, *}a_Y^{-1}f_*\mathcal{F} \to
\epsilon_{Y, *}f_{qc, *}a_X^{-1}\mathcal{F} =
f_{Zar, *}\epsilon_{X, *}a_X^{-1}\mathcal{F} =
f_{Zar, *}\pi_X^{-1}\mathcal{F}
$$
가 동형임을 보이면 충분하다. 이는 Lemma
\ref{sites-cohomology-lemma-collect-true-things-Zar}의 부분 (7a)이다.
\end{proof}

\begin{lemma}
\label{sites-cohomology-lemma-compare-qc-zar}
비교사상 $\epsilon : \textit{LC}_{qc} \to \textit{LC}_{Zar}$를
생각하자. 위상공간의 고유사상들의 모임을 $\mathcal{P}$로 쓰자.
$\textit{LC}_{Zar}$의 $X$에 대해
$\mathcal{A}'_X \subset \textit{Ab}(\textit{LC}_{Zar}/X)$를,
$\textit{Ab}(X)$의 $\mathcal{F}$에 대해
$\pi_X^{-1}\mathcal{F}$ 꼴인 층들로 이루어진 충만한 부분범주라 하자.
그러면
(\ref{sites-cohomology-item-base-change-P}),
(\ref{sites-cohomology-item-restriction-A}),
(\ref{sites-cohomology-item-A-sheaf}),
(\ref{sites-cohomology-item-A-and-P}) 및
(\ref{sites-cohomology-item-refine-tau-by-P})
는 Situation \ref{sites-cohomology-situation-compare}에서 성립한다.
\end{lemma}

\begin{proof}
먼저 Homology, Lemma
\ref{homology-lemma-characterize-weak-serre-subcategory}의 조건
(1), (2), (3), (4)를 확인하여
$\mathcal{A}'_X \subset \textit{Ab}(\textit{LC}_{Zar}/X)$가 약한
Serre 부분범주임을 보이자. Lemma \ref{sites-cohomology-lemma-collect-true-things-Zar}의
부분 (3)에 의해 $\pi_X^{-1}$은 완전하고 충실충만하므로 부분
(1), (2), (3)은 곧바로 성립한다.
$0 \to \pi_X^{-1}\mathcal{F} \to \mathcal{G} \to \pi_X^{-1}\mathcal{F}' \to 0$
가 $\textit{Ab}(\textit{LC}_{Zar}/X)$의 짧은 완전열이면,
Lemma \ref{sites-cohomology-lemma-collect-true-things-Zar}의 부분 (2)에 의해
$0 \to \mathcal{F} \to \pi_{X, *}\mathcal{G} \to \mathcal{F}' \to 0$
는 완전하다. 따라서
$\mathcal{G} = \pi_X^{-1}\pi_{X, *}\mathcal{G}$는
$\mathcal{A}'_X$에 속하며, 이로써 마지막 조건도 확인된다.

\medskip\noindent
성질 (\ref{sites-cohomology-item-base-change-P})는 Lemma \ref{sites-cohomology-lemma-LC-basic} 및
고유사상의 기저변환이 고유사상이라는 사실에 의해 성립한다
(Topology, Theorem \ref{topology-theorem-characterize-proper}과
Lemma \ref{topology-lemma-base-change-separated}를 보라).

\medskip\noindent
성질 (\ref{sites-cohomology-item-restriction-A})는 Lemma
\ref{sites-cohomology-lemma-collect-true-things-Zar}의 가환도식 (5)에서 나온다.

\medskip\noindent
성질 (\ref{sites-cohomology-item-A-sheaf})는 Lemma
\ref{sites-cohomology-lemma-describe-pullback-pi}이다.

\medskip\noindent
성질 (\ref{sites-cohomology-item-A-and-P})는 Lemma
\ref{sites-cohomology-lemma-collect-true-things-Zar}의 부분 (7)(b)이다.

\medskip\noindent
(\ref{sites-cohomology-item-refine-tau-by-P})의 증명. qc 덮개
$\{U_i \to U\}$가 주어졌다고 하자. $u \in U$에 대해
$i_1, \ldots, i_m \in I$와 준콤팩트 부분집합
$E_j \subset U_{i_j}$를 택하여 $\bigcup f_{i_j}(E_j)$가
$u$의 근방이 되게 하자. Hausdorff 준콤팩트 공간들 사이의
연속사상 $Y = \coprod_{j = 1, \ldots, m} E_j \to U$는 고유임을
관찰하라(Topology, Lemma \ref{topology-lemma-closed-map}).
$\bigcup f_{i_j}(E_j)$에 포함되는 열린 근방 $u \in V$를 택하자.
그러면 $Y \times_U V \to V$는 전사인 고유사상이므로 Lemma
\ref{sites-cohomology-lemma-proper-surjective-is-qc-covering}에 의해 $qc$ 덮개이다.
모든 $u \in U$에 대해 이렇게 할 수 있으므로
(\ref{sites-cohomology-item-refine-tau-by-P})가 성립한다.
\end{proof}

\begin{lemma}
\label{sites-cohomology-lemma-V-C-all-n}
위와 같은 표기를 사용하자.
\begin{enumerate}
\item $X \in \Ob(\textit{LC}_{qc})$이고 $\mathcal{F}$가 $X$ 위의
아벨 층이면 $\epsilon_{X, *}a_X^{-1}\mathcal{F} = \pi_X^{-1}\mathcal{F}$
이고, $i > 0$에 대해 $R^i\epsilon_{X, *}(a_X^{-1}\mathcal{F}) = 0$이다.
\item $f : X \to Y$가 $\textit{LC}_{qc}$의 고유사상이고
$\mathcal{F}$가 $X$ 위의 아벨 층이면 모든 $i$에 대해
$a_Y^{-1}(R^if_*\mathcal{F}) = R^if_{qc, *}(a_X^{-1}\mathcal{F})$
이다.
\item $X \in \Ob(\textit{LC}_{qc})$이고 $K$가 $D^+(X)$에 속하면
사상 $\pi_X^{-1}K \to R\epsilon_{X, *}(a_X^{-1}K)$는 동형이다.
\item $f : X \to Y$가 $\textit{LC}_{qc}$의 고유사상이고
$K$가 $D^+(X)$에 속하면
$a_Y^{-1}(Rf_*K) = Rf_{qc, *}(a_X^{-1}K)$이다.
\end{enumerate}
\end{lemma}

\begin{proof}
Lemma \ref{sites-cohomology-lemma-compare-qc-zar}에 의해 Section
\ref{sites-cohomology-section-compare-general}의 모든 보조정리를 현재 상황에
적용할 수 있다. 그 결과들을 옮겨 쓰기 위해, Lemma \ref{sites-cohomology-lemma-A}의
범주 $\mathcal{A}_X$가
$a_X^{-1} : \textit{Ab}(X) \to \textit{Ab}(\textit{LC}_{qc}/X)$의
본질적 상임을 관찰하자.

\medskip\noindent
부분 (1)은 모든 $n$에 대한 $(V_n)$과 동치이며, 이는 Lemma
\ref{sites-cohomology-lemma-V-C-all-n-general}에 의해 성립한다.

\medskip\noindent
부분 (2)는 Lemma \ref{sites-cohomology-lemma-V-implies-C-general}의 결론에
$\epsilon_Y^{-1}$을 적용하면 나온다.

\medskip\noindent
부분 (3)은 Lemma \ref{sites-cohomology-lemma-V-C-all-n-general}의 부분 (1)에서
나온다. 실제로 $\pi_X^{-1}K$는
$D^+_{\mathcal{A}'_X}(\textit{LC}_{Zar}/X)$에 속하고
$a_X^{-1} = \epsilon_X^{-1} \circ a_X^{-1}$이다.

\medskip\noindent
같은 이유로 부분 (4)는 Lemma
\ref{sites-cohomology-lemma-V-C-all-n-general}의 부분 (2)에서 나온다.
\end{proof}

\begin{lemma}
\label{sites-cohomology-lemma-cohomological-descent-LC}
$X$를 $\textit{LC}_{qc}$의 대상이라 하자. $K \in D^+(X)$에 대해 사상
$$
K \longrightarrow Ra_{X, *}a_X^{-1}K
$$
은 동형이다. 여기서 $a_X : \Sh(\textit{LC}_{qc}/X) \to \Sh(X)$는
위에서 정의한 사상이다.
\end{lemma}

\begin{proof}
먼저 명제를 $K$가 아벨 층 하나로 주어지는 경우로 환원하자.
즉, $K$를 아래로 유계인 복합체 $\mathcal{F}^\bullet$로 나타내자.
층의 경우에 의해
$\mathcal{F}^n = a_{X, *} a_X^{-1} \mathcal{F}^n$
이고 $q > 0$에 대해 층 $R^qa_{X, *}a_X^{-1}\mathcal{F}^n$이
영임을 안다. $a_X^{-1}\mathcal{F}^\bullet$와 함자 $a_{X, *}$에
Leray의 비순환성 보조정리(Derived Categories, Lemma
\ref{derived-lemma-leray-acyclicity})를 적용하면 결론이 나온다.
이제부터 $K = \mathcal{F}$라 가정하자.

\medskip\noindent
Lemma \ref{sites-cohomology-lemma-describe-pullback-pi}에 의해
$a_{X, *}a_X^{-1}\mathcal{F} = \mathcal{F}$이다. 따라서
$q > 0$에 대해 $R^qa_{X, *}a_X^{-1}\mathcal{F} = 0$임을 보이면 충분하다.
이를 위해 $a_X = \epsilon_X \circ \pi_X$와 Leray 스펙트럼열
Lemma \ref{sites-cohomology-lemma-relative-Leray}를 사용할 수 있다. Lemma
\ref{sites-cohomology-lemma-V-C-all-n}에 의해 $i > 0$에 대해
$R^i\epsilon_{X, *}(a_X^{-1}\mathcal{F}) = 0$이고
$\epsilon_{X, *}a_X^{-1}\mathcal{F} = \pi_X^{-1}\mathcal{F}$이다.
Lemma \ref{sites-cohomology-lemma-collect-true-things-Zar}에 의해 $j > 0$에 대해
$R^j\pi_{X, *}(\pi_X^{-1}\mathcal{F}) = 0$이다. 이로써 증명이 끝난다.
\end{proof}

\begin{lemma}
\label{sites-cohomology-lemma-compare-cohomology-LC}
$X \in \Ob(\textit{LC}_{qc})$와 위에서 정의한
$a_X : \Sh(\textit{LC}_{qc}/X) \to \Sh(X)$에 대해 다음이 성립한다.
\begin{enumerate}
\item $X$ 위의 아벨 층 $\mathcal{F}$에 대해
$H^n(X, \mathcal{F}) = H^n_{qc}(X, a_X^{-1}\mathcal{F})$,
\item $K \in D^+(X)$에 대해
$H^n(X, K) = H^n_{qc}(X, a_X^{-1}K)$이다.
\end{enumerate}
예를 들어 $A$가 아벨 군이면
$H^n(X, \underline{A}) = H^n_{qc}(X, \underline{A})$이다.
\end{lemma}

\begin{proof}
Lemma \ref{sites-cohomology-lemma-cohomological-descent-LC}에서 나오며, 여기에는
Remark \ref{sites-cohomology-remark-before-Leray}를 사용한다.
\end{proof}











\section{Ext에 대한 스펙트럼열}
\label{sites-cohomology-section-spectral-sequence-ext}

\noindent
이 절에서는 Ext 함자를 생각할 때 나타나는 여러 스펙트럼열을 모은다.
환 달린 사이트 $(\mathcal{C}, \mathcal{O})$ 위의 가군 복합체의
임의의 쌍 $\mathcal{G}^\bullet, \mathcal{F}^\bullet$에 대해
$$
\Ext^n_\mathcal{O}(\mathcal{G}^\bullet, \mathcal{F}^\bullet)
=
\Hom_{D(\mathcal{O})}(\mathcal{G}^\bullet, \mathcal{F}^\bullet[n])
$$
로 쓰겠다. 이는 Derived Categories, Section
\ref{derived-section-ext}의 일반적인 규약을 따른다.

\begin{example}
\label{sites-cohomology-example-hom-complex-into-sheaf}
$(\mathcal{C}, \mathcal{O})$를 환 달린 사이트라 하자.
$\mathcal{K}^\bullet$를 위로 유계인 $\mathcal{O}$-가군의 복합체라 하고,
$\mathcal{F}$를 $\mathcal{O}$-가군이라 하자. 그러면 $E_2$-쪽이
$$
E_2^{i, j} =
\Ext_\mathcal{O}^i(H^{-j}(\mathcal{K}^\bullet), \mathcal{F})
\Rightarrow
\Ext_\mathcal{O}^{i + j}(\mathcal{K}^\bullet, \mathcal{F})
$$
인 스펙트럼열과, $E_1$-쪽이
$$
E_1^{i, j} =
\Ext_\mathcal{O}^j(\mathcal{K}^{-i}, \mathcal{F})
\Rightarrow
\Ext_\mathcal{O}^{i + j}(\mathcal{K}^\bullet, \mathcal{F}).
$$
인 또 하나의 스펙트럼열이 있다. 이 스펙트럼열들을 구성하려면 단사분해
$\mathcal{F} \to \mathcal{I}^\bullet$를 택하고 이중복합체
$\Hom_\mathcal{O}(\mathcal{K}^\bullet, \mathcal{I}^\bullet)$에서
나오는 두 스펙트럼열을 생각한다. Homology, Section
\ref{homology-section-double-complex}을 보라.
\end{example}








\section{컵곱}
\label{sites-cohomology-section-cup-product}

\noindent
$(\mathcal{C}, \mathcal{O})$를 환 달린 사이트라 하자.
$K, M$을 $D(\mathcal{O})$의 대상이라 하고
$A = \Gamma(\mathcal{C}, \mathcal{O})$로 두자. 이 상황에서
(대역) 컵곱은 $D(A)$ 안의 사상
$$
R\Gamma(\mathcal{C}, K) \otimes_A^\mathbf{L}
R\Gamma(\mathcal{C}, M)
\longrightarrow
R\Gamma(\mathcal{C}, K \otimes_\mathcal{O}^\mathbf{L} M)
$$
이다. 이를 Remark \ref{sites-cohomology-remark-cup-product}에서와 같이 환 달린
토포스의 사상
$(\Sh(\mathcal{C}), \mathcal{O}) \to (\Sh(pt), A)$
에 대한 상대 컵곱으로 정의한다.

\medskip\noindent
상대 컵곱의 자연스러운 양립성을 정식화하고 증명하자. 즉, 환 달린
토포스의 사상
$f : (\Sh(\mathcal{C}), \mathcal{O}_\mathcal{C}) \to
(\Sh(\mathcal{D}), \mathcal{O}_\mathcal{D})$
가 있다고 하자. $\mathcal{K}^\bullet$와 $\mathcal{M}^\bullet$를
$\mathcal{O}_\mathcal{C}$-가군의 복합체라 하자. 소박한 컵곱
$$
\text{Tot}(
f_*\mathcal{K}^\bullet
\otimes_{\mathcal{O}_\mathcal{D}}
f_*\mathcal{M}^\bullet)
\longrightarrow
f_*\text{Tot}(\mathcal{K}^\bullet
\otimes_{\mathcal{O}_\mathcal{C}}
\mathcal{M}^\bullet)
$$
이 있다. 이것이 상대 컵곱과 관련됨을 주장한다.

\begin{lemma}
\label{sites-cohomology-lemma-cup-compatible-with-naive}
위 상황에서 다음 도식은 가환한다.
$$
\xymatrix{
f_*\mathcal{K}^\bullet
\otimes_{\mathcal{O}_\mathcal{D}}^\mathbf{L}
f_*\mathcal{M}^\bullet \ar[r] \ar[d]
&
Rf_*\mathcal{K}^\bullet
\otimes_{\mathcal{O}_\mathcal{D}}^\mathbf{L}
Rf_*\mathcal{M}^\bullet \ar[d]^{\text{비고 \ref{sites-cohomology-remark-cup-product}}} \\
\text{Tot}(
f_*\mathcal{K}^\bullet
\otimes_{\mathcal{O}_\mathcal{D}}
f_*\mathcal{M}^\bullet) \ar[d]_{\text{소박한 컵곱}} &
Rf_*(\mathcal{K}^\bullet
\otimes_{\mathcal{O}_\mathcal{C}}^\mathbf{L}
\mathcal{M}^\bullet) \ar[d] \\
f_*\text{Tot}(\mathcal{K}^\bullet
\otimes_{\mathcal{O}_\mathcal{C}}
\mathcal{M}^\bullet) \ar[r] &
Rf_*\text{Tot}(\mathcal{K}^\bullet
\otimes_{\mathcal{O}_\mathcal{C}}
\mathcal{M}^\bullet)
}
$$
\end{lemma}

\begin{proof}
Remark \ref{sites-cohomology-remark-cup-product}의 구성에 의해 도식을 시계방향으로
돌아 얻는 사상
$$
f_*\mathcal{K}^\bullet
\otimes_{\mathcal{O}_\mathcal{D}}^\mathbf{L}
f_*\mathcal{M}^\bullet 
\longrightarrow
Rf_*\text{Tot}(\mathcal{K}^\bullet
\otimes_{\mathcal{O}_\mathcal{C}}
\mathcal{M}^\bullet)
$$
은 사상
\begin{align*}
Lf^*(f_*\mathcal{K}^\bullet
\otimes_{\mathcal{O}_\mathcal{D}}^\mathbf{L}
f_*\mathcal{M}^\bullet)
& =
Lf^*f_*\mathcal{K}^\bullet
\otimes_{\mathcal{O}_\mathcal{D}}^\mathbf{L}
Lf^*f_*\mathcal{M}^\bullet \\
& \to
Lf^*Rf_*\mathcal{K}^\bullet
\otimes_{\mathcal{O}_\mathcal{D}}^\mathbf{L}
Lf^*Rf_*\mathcal{M}^\bullet \\
& \to
\mathcal{K}^\bullet
\otimes_{\mathcal{O}_\mathcal{D}}^\mathbf{L}
\mathcal{M}^\bullet \\
& \to
\text{Tot}(\mathcal{K}^\bullet
\otimes_{\mathcal{O}_\mathcal{C}}
\mathcal{M}^\bullet)
\end{align*}
에 수반된다. Lemma \ref{sites-cohomology-lemma-adjoints-push-pull-compatibility}에
의해 이는 또한
\begin{align*}
Lf^*(f_*\mathcal{K}^\bullet
\otimes_{\mathcal{O}_\mathcal{D}}^\mathbf{L}
f_*\mathcal{M}^\bullet)
& =
Lf^*f_*\mathcal{K}^\bullet
\otimes_{\mathcal{O}_\mathcal{D}}^\mathbf{L}
Lf^*f_*\mathcal{M}^\bullet \\
& \to
f^*f_*\mathcal{K}^\bullet
\otimes_{\mathcal{O}_\mathcal{D}}^\mathbf{L}
f^*f_*\mathcal{M}^\bullet \\
& \to
\mathcal{K}^\bullet
\otimes_{\mathcal{O}_\mathcal{D}}^\mathbf{L}
\mathcal{M}^\bullet \\
& \to
\text{Tot}(\mathcal{K}^\bullet
\otimes_{\mathcal{O}_\mathcal{C}}
\mathcal{M}^\bullet)
\end{align*}
와 같다. 반시계방향으로 돌면 사상
\begin{align*}
Lf^*(f_*\mathcal{K}^\bullet
\otimes_{\mathcal{O}_\mathcal{D}}^\mathbf{L}
f_*\mathcal{M}^\bullet)
& \to
Lf^*\text{Tot}(
f_*\mathcal{K}^\bullet
\otimes_{\mathcal{O}_\mathcal{D}}
f_*\mathcal{M}^\bullet) \\
& \to
Lf^*f_*\text{Tot}(\mathcal{K}^\bullet
\otimes_{\mathcal{O}_\mathcal{C}}
\mathcal{M}^\bullet) \\
& \to
Lf^*Rf_*\text{Tot}(\mathcal{K}^\bullet
\otimes_{\mathcal{O}_\mathcal{C}}
\mathcal{M}^\bullet) \\
& \to
\text{Tot}(\mathcal{K}^\bullet
\otimes_{\mathcal{O}_\mathcal{C}}
\mathcal{M}^\bullet)
\end{align*}
에 수반되는 사상을 얻는다. Lemma
\ref{sites-cohomology-lemma-adjoints-push-pull-compatibility}에 의해 이는 또한
\begin{align*}
Lf^*(f_*\mathcal{K}^\bullet
\otimes_{\mathcal{O}_\mathcal{D}}^\mathbf{L}
f_*\mathcal{M}^\bullet)
& \to
Lf^*\text{Tot}(
f_*\mathcal{K}^\bullet
\otimes_{\mathcal{O}_\mathcal{D}}
f_*\mathcal{M}^\bullet) \\
& \to
Lf^*f_*\text{Tot}(\mathcal{K}^\bullet
\otimes_{\mathcal{O}_\mathcal{C}}
\mathcal{M}^\bullet) \\
& \to
f^*f_*\text{Tot}(\mathcal{K}^\bullet
\otimes_{\mathcal{O}_\mathcal{C}}
\mathcal{M}^\bullet) \\
& \to
\text{Tot}(\mathcal{K}^\bullet
\otimes_{\mathcal{O}_\mathcal{C}}
\mathcal{M}^\bullet)
\end{align*}
와 같다. 이제 다음 도식을 살펴보면 증명이 끝난다.
$$
\xymatrix{
Lf^*(f_*\mathcal{K}^\bullet
\otimes_{\mathcal{O}_\mathcal{D}}^\mathbf{L}
f_*\mathcal{M}^\bullet) \ar[d] \ar[rr] & &
Lf^*f_*\mathcal{K}^\bullet \otimes_{\mathcal{O}_\mathcal{C}}^\mathbf{L}
Lf^*f_*\mathcal{M}^\bullet \ar[d] \\
Lf^*\text{Tot}(
f_*\mathcal{K}^\bullet
\otimes_{\mathcal{O}_\mathcal{D}}
f_*\mathcal{M}^\bullet) \ar[d]_{naive} \ar[r] &
f^*\text{Tot}(
f_*\mathcal{K}^\bullet
\otimes_{\mathcal{O}_\mathcal{D}}
f_*\mathcal{M}^\bullet) \ar[ldd]^{naive} \ar[dd] &
f^*f_*\mathcal{K}^\bullet \otimes_{\mathcal{O}_\mathcal{C}}^\mathbf{L}
f^*f_*\mathcal{M}^\bullet \ar[dd] \ar[ldd] \\
Lf^*f_*\text{Tot}(\mathcal{K}^\bullet
\otimes_{\mathcal{O}_\mathcal{C}}
\mathcal{M}^\bullet) \ar[d] \\
f^*f_*\text{Tot}(\mathcal{K}^\bullet \otimes_{\mathcal{O}_\mathcal{C}}
\mathcal{M}^\bullet) \ar[rd] &
\text{Tot}(f^*f_*\mathcal{K}^\bullet \otimes_{\mathcal{O}_\mathcal{C}}
f^*f_*\mathcal{M}^\bullet) \ar[d] &
\mathcal{K}^\bullet \otimes_{\mathcal{O}_\mathcal{C}}^\mathbf{L}
\mathcal{M}^\bullet \ar[ld] \\
& \text{Tot}(\mathcal{K}^\bullet
\otimes_{\mathcal{O}_\mathcal{C}}
\mathcal{M}^\bullet)
}
$$
이 도식의 모든 다각형은 가환한다. 위쪽 다각형은 Lemma
\ref{sites-cohomology-lemma-tensor-pull-compatibility}에 의해 가환한다.
두 소박한 컵곱을 포함하는 정사각형은 $Lf^* \to f^*$가
가군의 복합체에 대해 함자적이므로 가환한다. 두 사상
$\mathcal{A}^\bullet \otimes^\mathbf{L} \mathcal{B}^\bullet \to
\text{Tot}(\mathcal{A}^\bullet \otimes \mathcal{B}^\bullet)$
을 포함하는 정사각형도 마찬가지이다. 마지막으로 남은 정사각형의
가환성은 복합체의 수준에서 성립하며, $f^*$와 $f_*$의 수반성에 의한
소박한 컵곱의 정의로 볼 수 있다. 도식의 바깥을 따라 도는 두 사상이
위에서 주어진 두 사상이므로 증명이 끝난다.
\end{proof}

\begin{lemma}
\label{sites-cohomology-lemma-cup-product-associative}
$f : (\Sh(\mathcal{C}), \mathcal{O}) \to (\Sh(\mathcal{C}'), \mathcal{O}')$
를 환 달린 토포스의 사상이라 하자. Remark
\ref{sites-cohomology-remark-cup-product}의 상대 컵곱은 다음 도식이
$$
\xymatrix{
Rf_*K \otimes_{\mathcal{O}'}^\mathbf{L}
Rf_*L \otimes_{\mathcal{O}'}^\mathbf{L}
Rf_*M \ar[r] \ar[d] &
Rf_*(K \otimes_\mathcal{O}^\mathbf{L} L)
\otimes_{\mathcal{O}'}^\mathbf{L} Rf_*M \ar[d] \\
Rf_*K \otimes_{\mathcal{O}'}^\mathbf{L}
Rf_*(L \otimes_\mathcal{O}^\mathbf{L} M) \ar[r] &
Rf_*(K \otimes_\mathcal{O}^\mathbf{L} 
L \otimes_\mathcal{O}^\mathbf{L} M)
}
$$
$D(\mathcal{O})$의 모든 $K, L, M$에 대해 $D(\mathcal{O}')$에서
가환한다는 의미에서 결합적이다.
\end{lemma}

\begin{proof}
어느 쪽으로 돌아도 $D(\mathcal{O})$ 안의 자명한 사상
\begin{align*}
Lf^*(Rf_*K \otimes_{\mathcal{O}'}^\mathbf{L}
Rf_*L \otimes_{\mathcal{O}'}^\mathbf{L}
Rf_*M) & =
Lf^*(Rf_*K) \otimes_\mathcal{O}^\mathbf{L}
Lf^*(Rf_*L) \otimes_\mathcal{O}^\mathbf{L}
Lf^*(Rf_*M) \\
& \to
K \otimes_\mathcal{O}^\mathbf{L} 
L \otimes_\mathcal{O}^\mathbf{L} M
\end{align*}
에 수반되는 사상을 얻는다.
\end{proof}

\begin{lemma}
\label{sites-cohomology-lemma-cup-product-commutative}
$f : (\Sh(\mathcal{C}), \mathcal{O}) \to (\Sh(\mathcal{C}'), \mathcal{O}')$
를 환 달린 토포스의 사상이라 하자. Remark
\ref{sites-cohomology-remark-cup-product}의 상대 컵곱은 다음 도식이
$$
\xymatrix{
Rf_*K \otimes_{\mathcal{O}'}^\mathbf{L} Rf_*L \ar[r] \ar[d]_\psi &
Rf_*(K \otimes_\mathcal{O}^\mathbf{L} L) \ar[d]^{Rf_*\psi} \\
Rf_*L \otimes_{\mathcal{O}'}^\mathbf{L} Rf_*K \ar[r] &
Rf_*(L \otimes_\mathcal{O}^\mathbf{L} K)
}
$$
$D(\mathcal{O})$의 모든 $K, L$에 대해 $D(\mathcal{O}')$에서
가환한다는 의미에서 가환적이다. 여기서 $\psi$는 유도범주 위의
가환성 제약이다(Lemma \ref{sites-cohomology-lemma-symmetric-monoidal-derived}).
\end{lemma}

\begin{proof}
생략한다.
\end{proof}

\begin{lemma}
\label{sites-cohomology-lemma-compose-cup-product}
$f : (\Sh(\mathcal{C}), \mathcal{O}) \to (\Sh(\mathcal{C}'), \mathcal{O}')$와
$f' : (\Sh(\mathcal{C}'), \mathcal{O}') \to
(\Sh(\mathcal{C}''), \mathcal{O}'')$
을 환 달린 토포스의 사상들이라 하자. Remark
\ref{sites-cohomology-remark-cup-product}의 상대 컵곱은 다음 도식이
$$
\xymatrix{
R(f' \circ f)_*K \otimes_{\mathcal{O}''}^\mathbf{L} R(f' \circ f)_*L
\ar@{=}[rr] \ar[d] & &
Rf'_*Rf_*K \otimes_{\mathcal{O}''}^\mathbf{L} Rf'_*Rf_*L \ar[d] \\
R(f' \circ f)_*(K \otimes_\mathcal{O}^\mathbf{L} L) \ar@{=}[r] &
Rf'_*Rf_*(K \otimes_\mathcal{O}^\mathbf{L} L) &
Rf'_*(Rf_*K \otimes_{\mathcal{O}'}^\mathbf{L}  Rf_*L) \ar[l]
}
$$
$D(\mathcal{O})$의 모든 $K, L$에 대해 $D(\mathcal{O}'')$에서
가환한다는 의미에서 합성과 양립한다.
\end{lemma}

\begin{proof}
도식을 어느 쪽으로 돌아도 사상
\begin{align*}
& L(f' \circ f)^*\left(R(f' \circ f)_*K
\otimes_{\mathcal{O}''}^\mathbf{L}
R(f' \circ f)_*L\right) \\
& =
L(f' \circ f)^*R(f' \circ f)_*K
\otimes_\mathcal{O}^\mathbf{L}
L(f' \circ f)^*R(f' \circ f)_*L) \\
& \to
K \otimes_\mathcal{O}^\mathbf{L} L
\end{align*}
에 수반되는 사상을 얻으므로 성립한다. 이 사상은 $D(\mathcal{O})$에
속한다. 이를 보려면 다음 쌍대단위들의 합성
$$
L(f' \circ f)^*R(f' \circ f)_* =
Lf^* L(f')^* Rf'_* Rf_*  \to
Lf^* Rf_* \to \text{id}
$$
이 $L(f' \circ f)^*$와 $R(f' \circ f)_*$에 대한 쌍대단위임을 사용한다.
Categories, Lemma \ref{categories-lemma-compose-counits}를 보라.
\end{proof}

\begin{lemma}
\label{sites-cohomology-lemma-base-change-cup-product}
환 달린 토포스들의 가환 정사각형
$$
\xymatrix{
(\Sh(\mathcal{C}'), \mathcal{O}_{\mathcal{C}'}) \ar[r]_{g'} \ar[d]_{f'} &
(\Sh(\mathcal{C}, \mathcal{O}_{\mathcal{C}}) \ar[d]^f \\
(\Sh(\mathcal{D}'), \mathcal{O}_{\mathcal{D}'}) \ar[r]^g &
(\Sh(\mathcal{D}), \mathcal{O}_{\mathcal{D}}) 
}
$$
을 생각하자. $K, L$이 $D(\mathcal{O}_{\mathcal{C}})$에 속한다고 하자.
상대 컵곱은 다음 도식이
$$
\xymatrix{
Lg^*(Rf_*K \otimes_{\mathcal{O}_{\mathcal{D}}}^\mathbf{L} Rf_*L)
\ar[r] \ar@{=}[d] &
Lg^*(Rf_*(K \otimes_{\mathcal{O}_{\mathcal{C}}}^\mathbf{L} L)) \ar[d] \\
Lg^*Rf_*K \otimes_{\mathcal{O}_{\mathcal{D}'}}^\mathbf{L} Lg^*Rf_*L \ar[d] &
R(f')_*L(g')^*(K \otimes_{\mathcal{O}_{\mathcal{C}}}^\mathbf{L} L) \ar@{=}[d] \\
R(f')_*(L(g')^*K \otimes_{\mathcal{O}_{\mathcal{D}'}} R(f')_*(L(g')^*L \ar[r] &
R(f')_*(L(g')^*K \otimes_{\mathcal{O}_{\mathcal{C}'}}^\mathbf{L} L(g')^*L)
}
$$
$D(\mathcal{O}_{\mathcal{D}'})$에서 가환한다는 의미에서 이 정사각형과
양립한다. 가로 화살표들은 상대 컵곱(Remark
\ref{sites-cohomology-remark-cup-product})으로 주어지고, 세로 화살표들은 기저변환 사상
(Remark \ref{sites-cohomology-remark-base-change})과 Lemma
\ref{sites-cohomology-lemma-pullback-tensor-product}로 주어진다.
\end{lemma}

\begin{proof}
생략한다.
\end{proof}






\section{Hom 복합체}
\label{sites-cohomology-section-hom-complexes}

\noindent
$(\mathcal{C}, \mathcal{O})$를 환 달린 사이트라 하자.
$\mathcal{L}^\bullet$와 $\mathcal{M}^\bullet$를
$\mathcal{O}$-가군의 두 복합체라 하자. $\mathcal{O}$-가군의 복합체
$\SheafHom^\bullet(\mathcal{L}^\bullet, \mathcal{M}^\bullet)$를
구성하겠다. 구체적으로 각 $n$에 대해
$$
\SheafHom^n(\mathcal{L}^\bullet, \mathcal{M}^\bullet) =
\prod\nolimits_{n = p + q}
\SheafHom_\mathcal{O}(\mathcal{L}^{-q}, \mathcal{M}^p)
$$
로 둔다. $\SheafHom^n$은 $\mathcal{L}^\bullet$에서
$\mathcal{M}^\bullet$로 가는 모든 차수 $n$의 동차
$\mathcal{O}$-선형사상들로 이루어진 $\mathcal{O}$-가군의 층이라고
생각하면 좋다. 여기서 두 복합체는 등급화된 $\mathcal{O}$-가군으로
본다. 이 용어 아래 미분을 규칙
$$
\text{d}(f) =
\text{d}_\mathcal{M} \circ f - (-1)^n f \circ \text{d}_\mathcal{L}
$$
으로 정의한다. 여기서
$f \in \SheafHom^n_\mathcal{O}(\mathcal{L}^\bullet, \mathcal{M}^\bullet)$이다.
$\text{d}^2 = 0$의 확인은 생략한다. 이 구성은 Differential Graded
Algebra, Example \ref{dga-example-category-complexes}의 특수한 경우이다.
구성으로부터 모든 $n \in \mathbf{Z}$와 각
$U \in \Ob(\mathcal{C})$에 대해 곧바로
\begin{equation}
\label{sites-cohomology-equation-cohomology-hom-complex}
H^n(\Gamma(U, \SheafHom^\bullet(\mathcal{L}^\bullet, \mathcal{M}^\bullet))) =
\Hom_{K(\mathcal{O}_U)}(\mathcal{L}^\bullet|_U, \mathcal{M}^\bullet[n]|_U)
\end{equation}
를 얻는다. 마찬가지로 대역단면의 복합체에 대해
\begin{equation}
\label{sites-cohomology-equation-global-cohomology-hom-complex}
H^n(\Gamma(\mathcal{C},
\SheafHom^\bullet(\mathcal{L}^\bullet, \mathcal{M}^\bullet))) =
\Hom_{K(\mathcal{O})}(\mathcal{L}^\bullet, \mathcal{M}^\bullet[n])
\end{equation}
를 얻는다.

\begin{lemma}
\label{sites-cohomology-lemma-compose}
$(\mathcal{C}, \mathcal{O})$를 환 달린 사이트라 하자.
$\mathcal{O}$-가군의 복합체
$\mathcal{K}^\bullet, \mathcal{L}^\bullet, \mathcal{M}^\bullet$가
주어지면 동형
$$
\SheafHom^\bullet(\mathcal{K}^\bullet,
\SheafHom^\bullet(\mathcal{L}^\bullet, \mathcal{M}^\bullet))
=
\SheafHom^\bullet(\text{Tot}(\mathcal{K}^\bullet \otimes_\mathcal{O}
\mathcal{L}^\bullet), \mathcal{M}^\bullet)
$$
이 있다. 이는 $\mathcal{K}^\bullet, \mathcal{L}^\bullet,
\mathcal{M}^\bullet$에 대해 함자적인 $\mathcal{O}$-가군 복합체들의
동형이다.
\end{lemma}

\begin{proof}
생략한다. 힌트: More on Algebra, Lemma
\ref{more-algebra-lemma-compose}와 꼭 같은 방식으로 증명한다.
\end{proof}

\begin{lemma}
\label{sites-cohomology-lemma-composition}
$(\mathcal{C}, \mathcal{O})$를 환 달린 사이트라 하자.
$\mathcal{O}$-가군의 복합체
$\mathcal{K}^\bullet, \mathcal{L}^\bullet, \mathcal{M}^\bullet$가
주어지면 표준사상
$$
\text{Tot}\left(
\SheafHom^\bullet(\mathcal{L}^\bullet, \mathcal{M}^\bullet)
\otimes_\mathcal{O}
\SheafHom^\bullet(\mathcal{K}^\bullet, \mathcal{L}^\bullet)
\right)
\longrightarrow
\SheafHom^\bullet(\mathcal{K}^\bullet, \mathcal{M}^\bullet)
$$
이 있다. 이는 $\mathcal{O}$-가군 복합체들의 사상이다.
\end{lemma}

\begin{proof}
생략한다. 힌트: More on Algebra, Lemma
\ref{more-algebra-lemma-composition}과 꼭 같은 방식으로 증명한다.
\end{proof}

\begin{lemma}
\label{sites-cohomology-lemma-diagonal-better}
$(\mathcal{C}, \mathcal{O})$를 환 달린 사이트라 하자.
$\mathcal{O}$-가군의 복합체
$\mathcal{K}^\bullet, \mathcal{L}^\bullet, \mathcal{M}^\bullet$가
주어지면 표준사상
$$
\text{Tot}\left(
\mathcal{K}^\bullet \otimes_\mathcal{O}
\SheafHom^\bullet(\mathcal{M}^\bullet, \mathcal{L}^\bullet)
\right)
\longrightarrow
\SheafHom^\bullet(\mathcal{M}^\bullet,
\text{Tot}(\mathcal{K}^\bullet \otimes_\mathcal{O} \mathcal{L}^\bullet))
$$
이 있다. 이는 세 복합체 모두에 대해 함자적인
$\mathcal{O}$-가군 복합체들의 사상이다.
\end{lemma}

\begin{proof}
생략한다. 힌트: More on Algebra, Lemma
\ref{more-algebra-lemma-diagonal-better}와 꼭 같은 방식으로 증명한다.
\end{proof}

\begin{lemma}
\label{sites-cohomology-lemma-diagonal}
$(\mathcal{C}, \mathcal{O})$를 환 달린 사이트라 하자.
$\mathcal{O}$-가군의 복합체
$\mathcal{K}^\bullet, \mathcal{L}^\bullet, \mathcal{M}^\bullet$가
주어지면 표준사상
$$
\mathcal{K}^\bullet
\longrightarrow
\SheafHom^\bullet(\mathcal{L}^\bullet,
\text{Tot}(\mathcal{K}^\bullet \otimes_\mathcal{O} \mathcal{L}^\bullet))
$$
이 있다. 이는 두 복합체에 대해 함자적인
$\mathcal{O}$-가군 복합체들의 사상이다.
\end{lemma}

\begin{proof}
생략한다. 힌트: More on Algebra, Lemma
\ref{more-algebra-lemma-diagonal}과 꼭 같은 방식으로 증명한다.
\end{proof}

\begin{lemma}
\label{sites-cohomology-lemma-evaluate-and-more}
$(\mathcal{C}, \mathcal{O})$를 환 달린 사이트라 하자.
$\mathcal{O}$-가군의 복합체
$\mathcal{K}^\bullet, \mathcal{L}^\bullet, \mathcal{M}^\bullet$가
주어지면 표준사상
$$
\text{Tot}(\SheafHom^\bullet(\mathcal{L}^\bullet,
\mathcal{M}^\bullet) \otimes_\mathcal{O} \mathcal{K}^\bullet)
\longrightarrow
\SheafHom^\bullet(\SheafHom^\bullet(\mathcal{K}^\bullet,
\mathcal{L}^\bullet), \mathcal{M}^\bullet)
$$
이 있다. 이는 세 복합체 모두에 대해 함자적인
$\mathcal{O}$-가군 복합체들의 사상이다.
\end{lemma}

\begin{proof}
생략한다. 힌트: More on Algebra, Lemma
\ref{more-algebra-lemma-evaluate-and-more}와 꼭 같은 방식으로 증명한다.
\end{proof}

\begin{lemma}
\label{sites-cohomology-lemma-RHom-into-K-injective}
$(\mathcal{C}, \mathcal{O})$를 환 달린 사이트라 하자.
$L$과 $M$을 $D(\mathcal{O})$의 대상이라 하자.
$\mathcal{I}^\bullet$을 $M$을 나타내는 $\mathcal{O}$-가군의
K-단사 복합체라 하고, $\mathcal{L}^\bullet$를 $L$을 나타내는
$\mathcal{O}$-가군의 복합체라 하자. 그러면 모든
$U \in \Ob(\mathcal{C})$에 대해
$$
H^0(\Gamma(U, \SheafHom^\bullet(\mathcal{L}^\bullet, \mathcal{I}^\bullet))) =
\Hom_{D(\mathcal{O}_U)}(L|_U, M|_U)
$$
이다. 마찬가지로
$H^0(\Gamma(\mathcal{C},
\SheafHom^\bullet(\mathcal{L}^\bullet, \mathcal{I}^\bullet))) =
\Hom_{D(\mathcal{O})}(L, M)$.
\end{lemma}

\begin{proof}
다음이 성립한다.
\begin{align*}
H^0(\Gamma(U, \SheafHom^\bullet(\mathcal{L}^\bullet, \mathcal{I}^\bullet)))
& =
\Hom_{K(\mathcal{O}_U)}(\mathcal{L}^\bullet|_U, \mathcal{I}^\bullet|_U) \\
& =
\Hom_{D(\mathcal{O}_U)}(L|_U, M|_U)
\end{align*}
첫째 등식은 (\ref{sites-cohomology-equation-cohomology-hom-complex})이다.
둘째 등식은 Lemma \ref{sites-cohomology-lemma-restrict-K-injective-to-open}에 의해
$\mathcal{I}^\bullet|_U$가 K-단사이므로 성립한다.
마지막 등식의 증명도 비슷하며, 대신
(\ref{sites-cohomology-equation-global-cohomology-hom-complex})를 사용한다.
\end{proof}

\begin{lemma}
\label{sites-cohomology-lemma-RHom-well-defined}
$(\mathcal{C}, \mathcal{O})$를 환 달린 사이트라 하자.
$(\mathcal{I}')^\bullet \to \mathcal{I}^\bullet$
를 $\mathcal{O}$-가군의 K-단사 복합체들 사이의 준동형이라 하고,
$(\mathcal{L}')^\bullet \to \mathcal{L}^\bullet$를
$\mathcal{O}$-가군의 복합체들 사이의 준동형이라 하자. 그러면
$$
\SheafHom^\bullet(\mathcal{L}^\bullet, (\mathcal{I}')^\bullet)
\longrightarrow
\SheafHom^\bullet((\mathcal{L}')^\bullet, \mathcal{I}^\bullet)
$$
는 준동형이다.
\end{lemma}

\begin{proof}
$M$을 $\mathcal{I}^\bullet$과 $(\mathcal{I}')^\bullet$이 나타내는
$D(\mathcal{O})$의 대상이라 하고, $L$을 $\mathcal{L}^\bullet$와
$(\mathcal{L}')^\bullet$가 나타내는 $D(\mathcal{O})$의 대상이라 하자.
Lemma \ref{sites-cohomology-lemma-RHom-into-K-injective}에 의해 두 층
$$
H^0(\SheafHom^\bullet(\mathcal{L}^\bullet, (\mathcal{I}')^\bullet))
\quad\text{and}\quad
H^0(\SheafHom^\bullet((\mathcal{L}')^\bullet, \mathcal{I}^\bullet))
$$
은 모두 준층
$$
U \longmapsto \Hom_{D(\mathcal{O}_U)}(L|_U, M|_U)
$$
에 결합된 층과 같다. 따라서 이 사상은 준동형이다.
\end{proof}

\begin{lemma}
\label{sites-cohomology-lemma-RHom-from-K-flat-into-K-injective}
$(\mathcal{C}, \mathcal{O})$를 환 달린 사이트라 하자.
$\mathcal{I}^\bullet$을 $\mathcal{O}$-가군의 K-단사 복합체라 하고,
$\mathcal{L}^\bullet$을 $\mathcal{O}$-가군의 K-평탄 복합체라 하자.
그러면 $\SheafHom^\bullet(\mathcal{L}^\bullet, \mathcal{I}^\bullet)$은
$\mathcal{O}$-가군의 K-단사 복합체이다.
\end{lemma}

\begin{proof}
실제로 $\mathcal{K}^\bullet$이 $\mathcal{O}$-가군의 비순환 복합체이면
다음이 성립한다.
\begin{align*}
\Hom_{K(\mathcal{O})}(\mathcal{K}^\bullet,
\SheafHom^\bullet(\mathcal{L}^\bullet, \mathcal{I}^\bullet))
& =
H^0(\Gamma(\mathcal{C},
\SheafHom^\bullet(\mathcal{K}^\bullet,
\SheafHom^\bullet(\mathcal{L}^\bullet, \mathcal{I}^\bullet)))) \\
& =
H^0(\Gamma(\mathcal{C}, \SheafHom^\bullet(\text{Tot}(
\mathcal{K}^\bullet \otimes_\mathcal{O} \mathcal{L}^\bullet),
\mathcal{I}^\bullet))) \\
& =
\Hom_{K(\mathcal{O})}(
\text{Tot}(\mathcal{K}^\bullet \otimes_\mathcal{O} \mathcal{L}^\bullet),
\mathcal{I}^\bullet) \\
& =
0
\end{align*}
첫 번째 등식은 (\ref{sites-cohomology-equation-global-cohomology-hom-complex})에서 따르고,
두 번째 등식은 Lemma \ref{sites-cohomology-lemma-compose}에서 따르며,
세 번째 등식은 (\ref{sites-cohomology-equation-global-cohomology-hom-complex})에서 따른다.
마지막 등식은
$\text{Tot}(\mathcal{K}^\bullet \otimes_\mathcal{O} \mathcal{L}^\bullet)$
이 $\mathcal{L}^\bullet$의 K-평탄성
(Definition \ref{sites-cohomology-definition-K-flat})에 의해 비순환이고,
$\mathcal{I}^\bullet$이 K-단사이기 때문에 성립한다.
\end{proof}








\section{유도 범주에서의 내부 Hom}
\label{sites-cohomology-section-internal-hom}

\noindent
$(\mathcal{C}, \mathcal{O})$를 환 달린 사이트라 하고, $L, M$을
$D(\mathcal{O})$의 대상이라 하자. $D(\mathcal{O})$의 임의의 세 번째
대상 $K$에 대해 표준 전단사
\begin{equation}
\label{sites-cohomology-equation-internal-hom}
\Hom_{D(\mathcal{O})}(K, R\SheafHom(L, M))
=
\Hom_{D(\mathcal{O})}(K \otimes_\mathcal{O}^\mathbf{L} L, M)
\end{equation}
가 존재하도록 $D(\mathcal{O})$의 대상 $R\SheafHom(L, M)$을 구성하고자
한다. 요네다 보조정리(Categories, Lemma
\ref{categories-lemma-yoneda})에 의해 이 식은 $R\SheafHom(L, M)$을
유일한 동형사상까지 결정한다.

\medskip\noindent
그러한 대상을 구성하기 위해 $M$을 나타내는 $\mathcal{O}$-가군의
K-단사 복합체 $\mathcal{I}^\bullet$와 $L$을 나타내는 임의의
$\mathcal{O}$-가군 복합체 $\mathcal{L}^\bullet$를 택한다. 다음과 같이
둔다.
$$
R\SheafHom(L, M) = \SheafHom^\bullet(\mathcal{L}^\bullet, \mathcal{I}^\bullet)
$$
여기서 오른쪽은 Section \ref{sites-cohomology-section-hom-complexes}에서 구성한
$\mathcal{O}$-가군 복합체이다. 이는 Lemma
\ref{sites-cohomology-lemma-RHom-well-defined}에 의해 잘 정의된다. 따라서 함자
$$
D(\mathcal{O})^{opp} \times D(\mathcal{O}) \longrightarrow D(\mathcal{O}),
\quad
(K, L) \longmapsto R\SheafHom(K, L)
$$
를 얻는다. (\ref{sites-cohomology-equation-internal-hom})을 증명하기에 앞서
$R\SheafHom(K, L)$의 코호몰로지 군을 계산하자.

\begin{lemma}
\label{sites-cohomology-lemma-section-RHom-over-U}
$(\mathcal{C}, \mathcal{O})$를 환 달린 사이트라 하고, $L, M$을
$D(\mathcal{O})$의 대상이라 하자. $\mathcal{C}$의 모든 대상 $U$에 대해
$$
H^0(U, R\SheafHom(L, M)) =
\Hom_{D(\mathcal{O}_U)}(L|_U, M|_U)
$$
이고, 또한 $H^0(\mathcal{C}, R\SheafHom(L, M)) =
\Hom_{D(\mathcal{O})}(L, M)$이다.
\end{lemma}

\begin{proof}
$M$을 나타내는 $\mathcal{O}$-가군의 K-단사 복합체
$\mathcal{I}^\bullet$와 $L$을 나타내는 K-평탄 복합체
$\mathcal{L}^\bullet$를 택하자. 그러면 Lemma
\ref{sites-cohomology-lemma-RHom-from-K-flat-into-K-injective}에 의해
$\SheafHom^\bullet(\mathcal{L}^\bullet, \mathcal{I}^\bullet)$
는 K-단사이다. 따라서 $U$ 위의 단면을 취하는 것만으로 $U$ 위의
코호몰로지를 계산할 수 있고, Lemma
\ref{sites-cohomology-lemma-RHom-into-K-injective}에서 결론이 따른다.
\end{proof}

\begin{lemma}
\label{sites-cohomology-lemma-internal-hom}
$(\mathcal{C}, \mathcal{O})$를 환 달린 사이트라 하고, $K, L, M$을
$D(\mathcal{O})$의 대상이라 하자. 위에서 설명한 구성에 대해
$D(\mathcal{O})$ 안의 표준 동형사상
$$
R\SheafHom(K, R\SheafHom(L, M)) =
R\SheafHom(K \otimes_\mathcal{O}^\mathbf{L} L, M)
$$
이 있다. 이는 $K, L, M$에 대해 함자적이며, $H^0(\mathcal{C}, -)$를
취하면 (\ref{sites-cohomology-equation-internal-hom})을 준다.
\end{lemma}

\begin{proof}
$M$을 나타내는 K-단사 복합체 $\mathcal{I}^\bullet$와 $L$을 나타내는
$\mathcal{O}$-가군의 K-평탄 복합체 $\mathcal{L}^\bullet$를 택하자.
임의의 $\mathcal{O}$-가군 복합체 $\mathcal{K}^\bullet$에 대해
다음이 성립한다.
$$
\SheafHom^\bullet(\mathcal{K}^\bullet,
\SheafHom^\bullet(\mathcal{L}^\bullet, \mathcal{I}^\bullet))
=
\SheafHom^\bullet(
\text{Tot}(\mathcal{K}^\bullet \otimes_\mathcal{O} \mathcal{L}^\bullet),
\mathcal{I}^\bullet)
$$
이다(Lemma \ref{sites-cohomology-lemma-compose}). 왼쪽은
$R\SheafHom(K, R\SheafHom(L, M))$을 나타내고(Lemma
\ref{sites-cohomology-lemma-RHom-from-K-flat-into-K-injective} 사용), 오른쪽은
$R\SheafHom(K \otimes_\mathcal{O}^\mathbf{L} L, M)$을 나타낸다.
이로써 보조정리의 표시된 식이 증명된다. 대역 단면을 취하고 Lemma
\ref{sites-cohomology-lemma-section-RHom-over-U}를 사용하면
(\ref{sites-cohomology-equation-internal-hom})을 얻는다.
\end{proof}

\begin{lemma}
\label{sites-cohomology-lemma-restriction-RHom-to-U}
$(\mathcal{C}, \mathcal{O})$를 환 달린 사이트라 하고, $K, L$을
$D(\mathcal{O})$의 대상이라 하자. $R\SheafHom(K, L)$의 구성은 제한과
가환한다. 즉 $\mathcal{C}$의 모든 대상 $U$에 대해
$R\SheafHom(K|_U, L|_U) = R\SheafHom(K, L)|_U$이다.
\end{lemma}

\begin{proof}
이는 구성과 Lemma \ref{sites-cohomology-lemma-restrict-K-injective-to-open}에서
명백하다.
\end{proof}

\begin{lemma}
\label{sites-cohomology-lemma-RHom-triangulated}
$(\mathcal{C}, \mathcal{O})$를 환 달린 사이트라 하자. 이중함자
$R\SheafHom(- , -)$는 두 변수 각각에서 특수 삼각형을 특수
삼각형으로 보낸다.
\end{lemma}

\begin{proof}
대응
$$
(\mathcal{L}^\bullet, \mathcal{M}^\bullet) \longmapsto
\SheafHom^\bullet(\mathcal{L}^\bullet, \mathcal{M}^\bullet)
$$
은 어느 한 변수에서든 복합체들의 항별 분할 짧은 완전열을 항별
분할 짧은 완전열로 보낸다는 관찰에서 결론이 따른다. 세부 사항은
생략한다.
\end{proof}

\begin{lemma}
\label{sites-cohomology-lemma-internal-hom-evaluate}
$(\mathcal{C}, \mathcal{O})$를 환 달린 사이트라 하고, $K, L, M$을
$D(\mathcal{O})$의 대상이라 하자. 표준 사상
$$
R\SheafHom(L, M) \otimes_\mathcal{O}^\mathbf{L} K
\longrightarrow
R\SheafHom(R\SheafHom(K, L), M)
$$
이 있으며, 이는 $D(\mathcal{O})$ 안에서 $K, L, M$에 대해 함자적이다.
\end{lemma}

\begin{proof}
$M$을 나타내는 K-단사 복합체 $\mathcal{I}^\bullet$, $L$을 나타내는
K-단사 복합체 $\mathcal{J}^\bullet$, 그리고 $K$를 나타내는 K-평탄
복합체 $\mathcal{K}^\bullet$를 택한다. 이 사상은 다음 사상을
사용하여 정의한다.
$$
\text{Tot}(\SheafHom^\bullet(\mathcal{J}^\bullet,
\mathcal{I}^\bullet) \otimes_\mathcal{O} \mathcal{K}^\bullet)
\longrightarrow
\SheafHom^\bullet(\SheafHom^\bullet(\mathcal{K}^\bullet,
\mathcal{J}^\bullet), \mathcal{I}^\bullet)
$$
이는 Lemma \ref{sites-cohomology-lemma-evaluate-and-more}의 사상이다. 복합체들을 이렇게
특별히 택했으므로 왼쪽은
$R\SheafHom(L, M) \otimes_\mathcal{O}^\mathbf{L} K$를 나타내고,
오른쪽은 $R\SheafHom(R\SheafHom(K, L), M)$을 나타낸다. 이것이
$D(\mathcal{O})$의 세 대상 모두에 대해 함자적이라는 증명은 생략한다.
\end{proof}

\begin{lemma}
\label{sites-cohomology-lemma-internal-hom-composition}
\begin{slogan}
RSheafHom에서의 합성.
\end{slogan}
$(\mathcal{C}, \mathcal{O})$를 환 달린 사이트라 하자. $D(\mathcal{O})$의
$K, L, M$이 주어지면 표준 사상
$$
R\SheafHom(L, M) \otimes_\mathcal{O}^\mathbf{L} R\SheafHom(K, L)
\longrightarrow R\SheafHom(K, M)
$$
이 $D(\mathcal{O})$ 안에 존재한다.
\end{lemma}

\begin{proof}
$M$을 나타내는 K-단사 복합체 $\mathcal{I}^\bullet$와 $L$을 나타내는
K-단사 복합체 $\mathcal{J}^\bullet$, 그리고 $K$를 나타내는 임의의
$\mathcal{O}$-가군 복합체 $\mathcal{K}^\bullet$를 택하자. Lemma
\ref{sites-cohomology-lemma-composition}에 의해 복합체들의 사상
$$
\text{Tot}\left(
\SheafHom^\bullet(\mathcal{J}^\bullet, \mathcal{I}^\bullet)
\otimes_\mathcal{O}
\SheafHom^\bullet(\mathcal{K}^\bullet, \mathcal{J}^\bullet)
\right)
\longrightarrow
\SheafHom^\bullet(\mathcal{K}^\bullet, \mathcal{I}^\bullet)
$$
이 있다. $\mathcal{O}$-가군 복합체
$\SheafHom^\bullet(\mathcal{J}^\bullet, \mathcal{I}^\bullet)$,
$\SheafHom^\bullet(\mathcal{K}^\bullet, \mathcal{J}^\bullet)$, 그리고
$\SheafHom^\bullet(\mathcal{K}^\bullet, \mathcal{I}^\bullet)$
는 각각 $R\SheafHom(L, M)$, $R\SheafHom(K, L)$, 그리고
$R\SheafHom(K, M)$을 나타낸다. K-평탄 복합체
$\mathcal{H}^\bullet$와 유사동형
$\mathcal{H}^\bullet \to
\SheafHom^\bullet(\mathcal{K}^\bullet, \mathcal{J}^\bullet)$를 택하면,
다음 사상이 있다.
$$
\text{Tot}\left(
\SheafHom^\bullet(\mathcal{J}^\bullet, \mathcal{I}^\bullet)
\otimes_\mathcal{O} \mathcal{H}^\bullet
\right)
\longrightarrow
\text{Tot}\left(
\SheafHom^\bullet(\mathcal{J}^\bullet, \mathcal{I}^\bullet)
\otimes_\mathcal{O}
\SheafHom^\bullet(\mathcal{K}^\bullet, \mathcal{J}^\bullet)
\right)
$$
이 사상의 정의역은
$R\SheafHom(L, M) \otimes_\mathcal{O}^\mathbf{L} R\SheafHom(K, L)$을
나타낸다. 표시된 두 화살표를 합성하면 원하는 사상을 얻는다. 이
구성이 함자적이라는 증명은 생략한다.
\end{proof}

\begin{lemma}
\label{sites-cohomology-lemma-internal-hom-diagonal-better}
$(\mathcal{C}, \mathcal{O})$를 환 달린 사이트라 하자. $D(\mathcal{O})$의
$K, L, M$이 주어지면 표준 사상
$$
K \otimes_\mathcal{O}^\mathbf{L} R\SheafHom(M, L)
\longrightarrow
R\SheafHom(M, K \otimes_\mathcal{O}^\mathbf{L} L)
$$
이 있으며, 이는 $D(\mathcal{O})$ 안에서 $K, L, M$에 대해 함자적이다.
\end{lemma}

\begin{proof}
$K$를 나타내는 K-평탄 복합체 $\mathcal{K}^\bullet$와 $L$을 나타내는
K-단사 복합체 $\mathcal{I}^\bullet$를 택하고, $M$을 나타내는 임의의
$\mathcal{O}$-가군 복합체 $\mathcal{M}^\bullet$를 택하자.
$\mathcal{J}^\bullet$가 K-단사인 유사동형
$\text{Tot}(\mathcal{K}^\bullet \otimes_{\mathcal{O}_X} \mathcal{I}^\bullet)
\to \mathcal{J}^\bullet$
을 택한다. 그러면 사상
$$
\text{Tot}\left(
\mathcal{K}^\bullet \otimes_\mathcal{O}
\SheafHom^\bullet(\mathcal{M}^\bullet, \mathcal{I}^\bullet)
\right)
\to
\SheafHom^\bullet(\mathcal{M}^\bullet,
\text{Tot}(\mathcal{K}^\bullet \otimes_\mathcal{O} \mathcal{I}^\bullet))
\to
\SheafHom^\bullet(\mathcal{M}^\bullet, \mathcal{J}^\bullet)
$$
을 사용한다. 여기서 첫째 사상은 Lemma
\ref{sites-cohomology-lemma-diagonal-better}의 사상이다.
\end{proof}

\begin{lemma}
\label{sites-cohomology-lemma-internal-hom-diagonal}
$(\mathcal{C}, \mathcal{O})$를 환 달린 사이트라 하자. $D(\mathcal{O})$의
$K, L$이 주어지면 표준 사상
$$
K \longrightarrow R\SheafHom(L, K \otimes_\mathcal{O}^\mathbf{L} L)
$$
이 있으며, 이는 $D(\mathcal{O})$ 안에서 $K$와 $L$ 모두에 대해
함자적이다.
\end{lemma}

\begin{proof}
$K$를 나타내는 K-평탄 복합체 $\mathcal{K}^\bullet$와 $L$을 나타내는
임의의 $\mathcal{O}$-가군 복합체 $\mathcal{L}^\bullet$를 택하자.
K-단사 복합체 $\mathcal{J}^\bullet$와 유사동형
$\text{Tot}(\mathcal{K}^\bullet \otimes_\mathcal{O} \mathcal{L}^\bullet)
\to \mathcal{J}^\bullet$를 택한다. 그러면
$$
\mathcal{K}^\bullet \to
\SheafHom^\bullet(\mathcal{L}^\bullet,
\text{Tot}(\mathcal{K}^\bullet \otimes_\mathcal{O} \mathcal{L}^\bullet))
\to
\SheafHom^\bullet(\mathcal{L}^\bullet, \mathcal{J}^\bullet)
$$
을 사용한다. 여기서 첫째 사상은 Lemma \ref{sites-cohomology-lemma-diagonal}에서 온다.
\end{proof}

\begin{lemma}
\label{sites-cohomology-lemma-dual}
$(\mathcal{C}, \mathcal{O})$를 환 달린 사이트라 하고, $L$을
$D(\mathcal{O})$의 대상이라 하자.
$L^\vee = R\SheafHom(L, \mathcal{O})$라 두자. $D(\mathcal{O})$의
$M$에 대해 표준 사상
\begin{equation}
\label{sites-cohomology-equation-eval}
M \otimes^\mathbf{L}_\mathcal{O} L^\vee \longrightarrow R\SheafHom(L, M)
\end{equation}
이 있으며, 이는 표준 사상
$$
H^0(\mathcal{C}, M \otimes_\mathcal{O}^\mathbf{L} L^\vee)
\longrightarrow
\Hom_{D(\mathcal{O})}(L, M)
$$
을 유도한다. 이 사상은 $D(\mathcal{O})$의 $M$에 대해 함자적이다.
\end{lemma}

\begin{proof}
(\ref{sites-cohomology-equation-eval})의 사상은 식별
$M = R\SheafHom(\mathcal{O}, M)$을 사용하는 Lemma
\ref{sites-cohomology-lemma-internal-hom-composition}의 특수한 경우이다.
\end{proof}

\begin{remark}
\label{sites-cohomology-remark-projection-formula-for-internal-hom}
$f : (\Sh(\mathcal{C}), \mathcal{O}_\mathcal{C}) \to
(\Sh(\mathcal{D}), \mathcal{O}_\mathcal{D})$를 환 달린 토포스의 사상이라
하고, $K, L$을 $D(\mathcal{O}_\mathcal{C})$의 대상이라 하자. 다음과
같은 표준 사상이 있다고 주장한다.
$$
Rf_*R\SheafHom(L, K) \longrightarrow R\SheafHom(Rf_*L, Rf_*K)
$$
즉 (\ref{sites-cohomology-equation-internal-hom})에 의해 이는 사상
$Rf_*R\SheafHom(L, K) \otimes_{\mathcal{O}_\mathcal{D}}^\mathbf{L} Rf_*L
\to Rf_*K$을 주는 것과 같다. 이를 위해 합성
$$
Rf_*R\SheafHom(L, K) \otimes_{\mathcal{O}_\mathcal{D}}^\mathbf{L} Rf_*L \to
Rf_*(R\SheafHom(L, K) \otimes_{\mathcal{O}_\mathcal{C}}^\mathbf{L} L) \to
Rf_*K
$$
을 사용할 수 있다. 여기서 첫째 화살표는 상대 컵곱
(Remark \ref{sites-cohomology-remark-cup-product})이고, 둘째 화살표는 표준 사상
$R\SheafHom(L, K) \otimes_{\mathcal{O}_\mathcal{C}}^\mathbf{L} L \to K$
에 $Rf_*$를 적용한 것이다. 이 표준 사상은 Lemma
\ref{sites-cohomology-lemma-internal-hom-composition}에서 한 자리에
$\mathcal{O}_\mathcal{C}$를 넣어 얻는다.
\end{remark}

\begin{remark}
\label{sites-cohomology-remark-prepare-fancy-base-change}
$h : (\Sh(\mathcal{C}), \mathcal{O}) \to
(\Sh(\mathcal{C}'), \mathcal{O}')$를 환 달린 토포스의 사상이라 하고,
$K, L$을 $D(\mathcal{O}')$의 대상이라 하자. 표준 사상
$$
Lh^*R\SheafHom(K, L) \longrightarrow R\SheafHom(Lh^*K, Lh^*L)
$$
이 $D(\mathcal{O})$ 안에 존재한다고 주장한다.
(\ref{sites-cohomology-equation-internal-hom})은 Lemma
\ref{sites-cohomology-lemma-internal-hom}에서 증명했으며, 이에 의해 그러한 사상을 주는
것은 사상
$$
Lh^*R\SheafHom(K, L) \otimes^\mathbf{L} Lh^*K \longrightarrow Lh^*L
$$
을 주는 것과 같다. Lemma \ref{sites-cohomology-lemma-pullback-tensor-product}에 의해 이
화살표의 정의역은 $Lh^*(\SheafHom(K, L) \otimes^\mathbf{L} K)$이므로,
표준 사상
$$
R\SheafHom(K, L) \otimes^\mathbf{L} K \longrightarrow L.
$$
을 구성하면 충분하다. 이를 위해
$$
\text{id} :
R\SheafHom(K, L)
\longrightarrow
R\SheafHom(K, L)
$$
에 (\ref{sites-cohomology-equation-internal-hom})으로 대응하는 화살표를 취한다.
\end{remark}

\begin{remark}
\label{sites-cohomology-remark-fancy-base-change}
환 달린 토포스들의 가환 도식
$$
\xymatrix{
(\Sh(\mathcal{C}'), \mathcal{O}_{\mathcal{C}'})
\ar[r]_h \ar[d]_{f'} &
(\Sh(\mathcal{C}), \mathcal{O}_\mathcal{C}) \ar[d]^f \\
(\Sh(\mathcal{D}'), \mathcal{O}_{\mathcal{D}'})
\ar[r]^g &
(\Sh(\mathcal{D}), \mathcal{O}_\mathcal{D})
}
$$
이 주어졌다고 하자. $K, L$을 $D(\mathcal{O}_\mathcal{C})$의 대상이라
하자. 표준 기저변환 사상
$$
Lg^*Rf_*R\SheafHom(K, L)
\longrightarrow
R(f')_*R\SheafHom(Lh^*K, Lh^*L)
$$
이 $D(\mathcal{O}_{\mathcal{D}'})$ 안에 존재한다고 주장한다. 구체적으로
합성
\begin{align*}
L(f')^*Lg^*Rf_*R\SheafHom(K, L)
& =
Lh^*Lf^*Rf_*R\SheafHom(K, L) \\
& \to
Lh^*R\SheafHom(K, L) \\
& \to
R\SheafHom(Lh^*K, Lh^*L)
\end{align*}
에 수반되는 사상을 취한다. 여기서 첫째 화살표는 수반 사상
$Lf^*Rf_* \to \text{id}$를 사용하고, 둘째 화살표는 Remark
\ref{sites-cohomology-remark-prepare-fancy-base-change}에서 구성한 표준 사상이다.
\end{remark}




\section{대역 유도 Hom}
\label{sites-cohomology-section-global-RHom}

\noindent
$(\Sh(\mathcal{C}), \mathcal{O})$를 환 달린 토포스라 하고,
$K, L \in D(\mathcal{O})$라 하자. 유도 범주에서의 내부 Hom 구성을
사용하면 $D(\Gamma(\mathcal{C}, \mathcal{O}))$의 잘 정의된 대상
$$
R\Hom_\mathcal{O}(K, L) = R\Gamma(\mathcal{C}, R\SheafHom(K, L))
$$
을 얻는다. Lemma \ref{sites-cohomology-lemma-section-RHom-over-U}에 의해
$$
H^0(R\Hom_\mathcal{O}(K, L)) = \Hom_{D(\mathcal{O})}(K, L)
$$
이고,
$$
H^p(R\Hom_\mathcal{O}(K, L)) = \Ext_{D(\mathcal{O})}^p(K, L)
$$
이다. $f : (\mathcal{C}', \mathcal{O}') \to
(\mathcal{C}, \mathcal{O})$가 환 달린 토포스의 사상이면 표준 사상
$$
R\Hom_\mathcal{O}(K, L) \longrightarrow R\Hom_{\mathcal{O}'}(Lf^*K, Lf^*L)
$$
이 $D(\Gamma(\mathcal{O}))$ 안에 존재한다. 이는 Remark
\ref{sites-cohomology-remark-prepare-fancy-base-change}에서 정의한 사상의 대역 단면을
취하여 얻는다.









\section{아래 첨자 느낌표의 유도 함자}
\label{sites-cohomology-section-derived-lower-shriek}

\noindent
이 절에서는 $Lg^*$와 $Rg_*$뿐 아니라 유도 함자 $Lg_!$도 존재하는
환 달린 토포스의 사상 $g$를 연구한다.

\begin{lemma}
\label{sites-cohomology-lemma-pullback-injective-pre-limp}
$u : \mathcal{C} \to \mathcal{D}$를 사이트들의 연속이고 쌍대연속인
함자라 하고, $g : \Sh(\mathcal{C}) \to \Sh(\mathcal{D})$를 이에
대응하는 토포스의 사상이라 하자. $\mathcal{O}_\mathcal{D}$를 환의
층이라 하고, $\mathcal{I}$를 단사 $\mathcal{O}_\mathcal{D}$-가군이라
하자. 그러면 모든 $p > 0$ 및 $U \in \Ob(\mathcal{C})$에 대해
$H^p(U, g^{-1}\mathcal{I}) = 0$이다.
\end{lemma}

\begin{proof}
임의의 $\mathcal{C}$의 덮개 $\mathcal{U} = \{U_i \to U\}$에 대해 모든
고차 체흐 코호몰로지 군
$\check H^p(\mathcal{U}, g^{-1}\mathcal{I})$
이 소멸함을 증명하면, 보조정리의 소멸은 Lemma
\ref{sites-cohomology-lemma-cech-vanish-collection}에서 따른다. $u$가 연속이므로
$u(\mathcal{U}) = \{u(U_i) \to u(U)\}$는 $\mathcal{D}$의 덮개이고,
$u(U_{i_0} \times_U \ldots \times_U U_{i_n}) =
u(U_{i_0}) \times_{u(U)} \ldots \times_{u(U)} u(U_{i_n})$.
따라서
$$
\check H^p(\mathcal{U}, g^{-1}\mathcal{I}) =
\check H^p(u(\mathcal{U}), \mathcal{I})
$$
이다. 실제로 Sites, Lemma \ref{sites-lemma-when-shriek}에 의해
$g^{-1} = u^p$이다. $\mathcal{I}$가 단사
$\mathcal{O}_\mathcal{D}$-가군이므로 이 체흐 코호몰로지 군들은
소멸한다. Lemma \ref{sites-cohomology-lemma-injective-module-trivial-cech}를 보라.
\end{proof}

\begin{lemma}
\label{sites-cohomology-lemma-existence-derived-lower-shriek}
$u : \mathcal{C} \to \mathcal{D}$를 사이트들의 연속이고 쌍대연속인
함자라 하고, $g : \Sh(\mathcal{C}) \to \Sh(\mathcal{D})$를 이에
대응하는 토포스의 사상이라 하자. $\mathcal{O}_\mathcal{D}$를 환의
층이라 하고,
$\mathcal{O}_\mathcal{C} = g^{-1}\mathcal{O}_\mathcal{D}$.
라 두자. 함자 $g_! : \textit{Mod}(\mathcal{O}_\mathcal{C}) \to
\textit{Mod}(\mathcal{O}_\mathcal{D})$
(Modules on Sites, Lemma
\ref{sites-modules-lemma-lower-shriek-modules} 참조)는 왼쪽 유도 함자
$$
Lg_! : D(\mathcal{O}_\mathcal{C}) \longrightarrow D(\mathcal{O}_\mathcal{D})
$$
를 가지며, 이는 $g^*$의 왼쪽 수반이다. 또한
$U \in \Ob(\mathcal{C})$에 대해
$$
Lg_!(j_{U!}\mathcal{O}_U) =
g_!j_{U!}\mathcal{O}_U =
j_{u(U)!} \mathcal{O}_{u(U)}.
$$
여기서 $j_{U!}$와 $j_{u(U)!}$는 국소화 사상
$j_U : \mathcal{C}/U \to \mathcal{C}$ 및
$j_{u(U)} : \mathcal{D}/u(U) \to \mathcal{D}$에 결부된 영에 의한
연장이다.
\end{lemma}

\begin{proof}
Derived Categories, Proposition
\ref{derived-proposition-left-derived-exists}를 사용하여 $Lg_!$를
구성한다. 이를 위해 그 명제의 가정 (1), (2), (3), (4), (5)를 확인해야
한다. 먼저 $g_!$는 왼쪽 수반이므로 오른쪽 완전이고 모든 여극한과
가환한다. 따라서 (5)가 성립한다. 환 달린 사이트 위의 가군 범주는
Grothendieck 아벨 범주이므로 조건 (3)과 (4)가 성립한다.
$\mathcal{P} \subset \Ob(\textit{Mod}(\mathcal{O}_\mathcal{C}))$를
$j_{U!}\mathcal{O}_U$ 꼴 가군들의 직합인
$\mathcal{O}_\mathcal{C}$-가군들의 모임이라 하자. Modules on Sites,
Lemma \ref{sites-modules-lemma-lower-shriek-modules}의 증명에서 보듯이
$g_!j_{U!}\mathcal{O}_U = j_{u(U)!} \mathcal{O}_{u(U)}$이다. 모든
$\mathcal{O}_\mathcal{C}$-가군은 $\mathcal{P}$의 한 대상의 몫이다.
Modules on Sites, Lemma
\ref{sites-modules-lemma-module-quotient-flat}을 보라. 따라서 (1)이
성립한다. 마지막으로 (2)를 증명해야 한다. 모든 $n$에 대해
$\mathcal{K}^n \in \mathcal{P}$인 $\mathcal{O}_\mathcal{C}$-가군의
위로 유계인 비순환 복합체 $\mathcal{K}^\bullet$를 택하자.
$g_!\mathcal{K}^\bullet$가 완전임을 보여야 한다. 이를 위해 모든 단사
$\mathcal{O}_\mathcal{D}$-가군 $\mathcal{I}$에 대해
$$
\Hom_{D(\mathcal{O}_\mathcal{D})}(
g_!\mathcal{K}^\bullet, \mathcal{I}[n]) = 0
$$
가 모든 $n \in \mathbf{Z}$에 대해 성립함을 보이면 충분하다.
$\mathcal{I}$가 단사이므로
\begin{align*}
\Hom_{D(\mathcal{O}_\mathcal{D})}(
g_!\mathcal{K}^\bullet, \mathcal{I}[n])
& =
\Hom_{K(\mathcal{O}_\mathcal{D})}(
g_!\mathcal{K}^\bullet, \mathcal{I}[n]) \\
& =
H^n(\Hom_{\mathcal{O}_\mathcal{D}}(
g_!\mathcal{K}^\bullet, \mathcal{I})) \\
& =
H^n(\Hom_{\mathcal{O}_\mathcal{C}}(
\mathcal{K}^\bullet, g^{-1}\mathcal{I}))
\end{align*}
이고, 마지막 등식은 $g_!$와 $g^{-1}$의 수반성에서 따른다.

\medskip\noindent
$g^{-1}\mathcal{I}$가 단사 $\mathcal{O}_\mathcal{C}$-가군이라면 이
군의 소멸은 명백할 것이다. 그러나 일반적으로 $g_!$가 완전이 아니므로
$g^{-1}\mathcal{I}$가 반드시 단사 $\mathcal{O}_\mathcal{C}$-가군인
것은 아니다. 한편 다음은 알고 있다.
$$
\Ext^p_{\mathcal{O}_\mathcal{C}}(
j_{U!}\mathcal{O}_U, g^{-1}\mathcal{I}) =
H^p(U, g^{-1}\mathcal{I}) = 0 \text{ for }p \geq 1
$$
여기서 첫 번째 등식은
$\Hom_{\mathcal{O}_\mathcal{C}}(j_{U!}\mathcal{O}_U, \mathcal{H}) =
\mathcal{H}(U)$에 유도 함자를 취하여 얻고, $p > 0$ 및
$U \in \Ob(\mathcal{C})$에 대한
$H^p(U, g^{-1}\mathcal{I})$의 소멸은 Lemma
\ref{sites-cohomology-lemma-pullback-injective-pre-limp}에서 따른다. 각
$\mathcal{K}^{-q}$가 $j_{U!}\mathcal{O}_U$ 꼴 가군들의 직합이므로
$$
\Ext^p_{\mathcal{O}_\mathcal{C}}(\mathcal{K}^{-q}, g^{-1}\mathcal{I}) = 0
\text{ for }p \geq 1\text{ and all }q
$$
이다. 다음 스펙트럼 열을 사용하자(Example
\ref{sites-cohomology-example-hom-complex-into-sheaf} 참조).
$$
E_1^{p, q} = \Ext^q_{\mathcal{O}_\mathcal{C}}(
\mathcal{K}^{-p}, g^{-1}\mathcal{I})
\Rightarrow
\Ext^{p + q}_{\mathcal{O}_\mathcal{C}}(
\mathcal{K}^\bullet, g^{-1}\mathcal{I}) = 0.
$$
복합체 $\mathcal{K}^\bullet$가 비순환이므로 유도 범주에서 소멸하고,
따라서 Ext 군들이 소멸한다. 그러므로 이 스펙트럼 열은 0으로
수렴한다. 위에서 증명한 고차 Ext들의 소멸에 의해 $E_1$ 쪽에서
영이 아닐 수 있는 항은
$E_1^{p, 0} = \Hom_{\mathcal{O}_\mathcal{C}}(
\mathcal{K}^{-p}, g^{-1}\mathcal{I})$뿐이다. 따라서 복합체
$\Hom_{\mathcal{O}_\mathcal{C}}(
\mathcal{K}^\bullet, g^{-1}\mathcal{I})$
가 원하는 대로 비순환임을 얻는다.

\medskip\noindent
따라서 왼쪽 유도 함자 $Lg_!$가 존재한다. 이는
$g^{-1} = g^* = Rg^* = Lg^*$의 왼쪽 수반이다. 즉
\begin{equation}
\label{sites-cohomology-equation-to-prove}
\Hom_{D(\mathcal{O}_\mathcal{C})}(K, g^*L) =
\Hom_{D(\mathcal{O}_\mathcal{D})}(Lg_!K, L)
\end{equation}
가 Derived Categories, Lemma
\ref{derived-lemma-derived-adjoint-functors}에 의해 성립한다. 이로써
증명이 끝난다.
\end{proof}

\begin{remark}
\label{sites-cohomology-remark-when-derived-shriek-equal}
경고! $u : \mathcal{C} \to \mathcal{D}$, $g$,
$\mathcal{O}_\mathcal{D}$, $\mathcal{O}_\mathcal{C}$가 Lemma
\ref{sites-cohomology-lemma-existence-derived-lower-shriek}에서와 같다고 하자. 일반적으로
다음 도식은 {\bf 가환하지 않는다}.
$$
\xymatrix{
D(\mathcal{O}_\mathcal{C}) \ar[r]_{Lg_!} \ar[d]_{forget} &
D(\mathcal{O}_\mathcal{D}) \ar[d]^{forget} \\
D(\mathcal{C}) \ar[r]^{Lg^{Ab}_!} &
D(\mathcal{D})
}
$$
여기서 함자 $Lg_!^{Ab}$는 $\mathcal{C}$와 $\mathcal{D}$ 양쪽에서 상수층
$\mathbf{Z}$를 구조층으로 사용하여 Lemma
\ref{sites-cohomology-lemma-existence-derived-lower-shriek}에서 구성한 함자이다.
일반적으로는 $g_! = g_!^{Ab}$조차 성립하지 않는다(Modules on Sites,
Remark \ref{sites-modules-remark-when-shriek-equal} 참조). 그러나 이
현상은 {\bf $g_! = g_!^{Ab}$인 경우에도 일어날 수 있다}! 실제로 Lemma
\ref{sites-cohomology-lemma-existence-derived-lower-shriek}의 증명에서 $Lg_!$의 구성을
살펴보면, $Lg_!$가 $Lg_!^{\textit{Ab}}$와 일치할 필요충분조건은 모든
$\mathcal{C}$의 대상 $U$에 대해 표준 사상
$$
Lg^{Ab}_!j_{U!}\mathcal{O}_U \longrightarrow j_{u(U)!}\mathcal{O}_{u(U)}
$$
이 $D(\mathcal{D})$에서 동형사상인 것이다. 일반적으로 말할 수 있는
것은 자연변환
$$
Lg_!^{Ab} \circ forget \longrightarrow forget \circ Lg_!
$$
이 존재한다는 것뿐이다.
\end{remark}

\begin{lemma}
\label{sites-cohomology-lemma-pullback-injective-limp}
$u : \mathcal{C} \to \mathcal{D}$를 사이트들의 연속이고 쌍대연속인
함자라 하고, $g : \Sh(\mathcal{C}) \to \Sh(\mathcal{D})$를 이에
대응하는 토포스의 사상이라 하자. $\mathcal{O}_\mathcal{D}$를 환의
층이라 하고, $\mathcal{I}$를 단사 $\mathcal{O}_\mathcal{D}$-가군이라
하자. 만약
$g_!^{Sh} : \Sh(\mathcal{C}) \to \Sh(\mathcal{D})$
가 올곱과 가환하면\footnote{$\mathcal{C}$가 유한 연결 극한을 가지고
$u$가 그것들과 가환하면 성립한다. Sites, Lemma
\ref{sites-lemma-preserve-equalizers}를 보라.},
$g^{-1}\mathcal{I}$는 완전 비순환이다.
\end{lemma}

\begin{proof}
Lemma \ref{sites-cohomology-lemma-characterize-limp}의 판정법을 사용한다. 조건 (1)은
Lemma \ref{sites-cohomology-lemma-pullback-injective-pre-limp}에 의해 성립한다.
$K' \to K$를 $\mathcal{C}$ 위 집합층들의 전사 사상이라 하자.
$g_!^{Sh}$가 왼쪽 수반이므로 $g_!^{Sh}K' \to g_!^{Sh}K$는 전사이다.
또한
\begin{align*}
H^0(K' \times_K \ldots \times_K K', g^{-1}\mathcal{I}) 
& =
H^0(g_!^{Sh}(K' \times_K \ldots \times_K K'), \mathcal{I}) \\
& =
H^0(g_!^{Sh}K' \times_{g_!^{Sh}K} \ldots \times_{g_!^{Sh}K} g_!^{Sh}K',
\mathcal{I})
\end{align*}
이다. 이는 $g_!^{Sh}$에 대한 가정에서 따른다. $\mathcal{I}$는 단사
가군이므로 항등사상에 Lemma
\ref{sites-cohomology-lemma-direct-image-injective-sheaf}를 적용하면 완전 비순환이다.
따라서 Lemma \ref{sites-cohomology-lemma-characterize-limp}의 역을 사용하면 복합체
$$
0 \to H^0(K, g^{-1}\mathcal{I}) \to H^0(K', g^{-1}\mathcal{I}) \to
H^0(K' \times_K K', g^{-1}\mathcal{I}) \to \ldots
$$
가 원하는 대로 완전임을 알 수 있다.
\end{proof}

\begin{lemma}
\label{sites-cohomology-lemma-pullback-same-cohomology}
$u : \mathcal{C} \to \mathcal{D}$를 사이트들의 연속이고 쌍대연속인
함자라 하고, $g : \Sh(\mathcal{C}) \to \Sh(\mathcal{D})$를 이에
대응하는 토포스의 사상이라 하자. $U \in \Ob(\mathcal{C})$라 하자.
\begin{enumerate}
\item $D(\mathcal{D})$의 $M$에 대해
$R\Gamma(U, g^{-1}M) = R\Gamma(u(U), M)$.
\item $\mathcal{O}_\mathcal{D}$가 환의 층이고
$\mathcal{O}_\mathcal{C} = g^{-1}\mathcal{O}_\mathcal{D}$이면
$D(\mathcal{O}_\mathcal{D})$의 $M$에 대해
$R\Gamma(U, g^*M) = R\Gamma(u(U), M)$.
\end{enumerate}
\end{lemma}

\begin{proof}
아래로 유계인 경우에는 $K$를 아래로 유계인 단사 복합체로 나타내고
Lemma \ref{sites-cohomology-lemma-pullback-injective-pre-limp}와 Leray의 비순환성
보조정리를 사용하면 (1)과 (2)를 알 수 있다. 비유계인 경우에는 먼저
(1)이 (2)의 특수한 경우임을 주목하자. (2)에 대해서는
$$
R\Gamma(U, g^*M) =
R\Hom_{\mathcal{O}_\mathcal{C}}(j_{U!}\mathcal{O}_U, g^*M) =
R\Hom_{\mathcal{O}_\mathcal{D}}(j_{u(U)!}\mathcal{O}_{u(U)}, M) =
R\Gamma(u(U), M)
$$
를 사용할 수 있다. 가운데 등식은 Lemma
\ref{sites-cohomology-lemma-existence-derived-lower-shriek}의 귀결이다.
\end{proof}

\begin{lemma}
\label{sites-cohomology-lemma-special-square-cocontinuous}
환 달린 토포스들의 가환 도식
$$
\xymatrix{
(\Sh(\mathcal{C}'), \mathcal{O}_{\mathcal{C}'})
\ar[r]_{(g', (g')^\sharp)} \ar[d]_{(f', (f')^\sharp)} &
(\Sh(\mathcal{C}), \mathcal{O}_\mathcal{C}) \ar[d]^{(f, f^\sharp)} \\
(\Sh(\mathcal{D}'), \mathcal{O}_{\mathcal{D}'}) \ar[r]^{(g, g^\sharp)} &
(\Sh(\mathcal{D}), \mathcal{O}_\mathcal{D})
}
$$
이 주어졌다고 하자. 다음을 가정하자.
\begin{enumerate}
\item $f$, $f'$, $g$, $g'$는 Sites, Lemma
\ref{sites-lemma-cocontinuous-morphism-topoi}에서와 같이 각각
쌍대연속 함자 $u$, $u'$, $v$, $v'$에 대응한다.
\item $v \circ u' = u \circ v'$,
\item $v$와 $v'$는 쌍대연속일 뿐 아니라 연속이다.
\item $\mathcal{D}'$의 모든 대상 $V'$에 대해 $v$가 주는 함자
${}^{u'}_{V'}\mathcal{I} \to {}^{\ \ \ u}_{v(V')}\mathcal{I}$
는 공종이다.
\item $g^{-1}\mathcal{O}_{\mathcal{D}} = \mathcal{O}_{\mathcal{D}'}$
이고 $(g')^{-1}\mathcal{O}_{\mathcal{C}} =
\mathcal{O}_{\mathcal{C}'}$이다.
\item $g'_! : \textit{Ab}(\mathcal{C}') \to \textit{Ab}(\mathcal{C})$
는 완전이다\footnote{$\mathcal{C}'$에 올곱과 등화자가 존재하고 $v'$가
그것들과 가환하면 성립한다. Modules on Sites, Lemma
\ref{sites-modules-lemma-exactness-lower-shriek}을 보라.}.
\end{enumerate}
그러면 함자 $D(\mathcal{O}_\mathcal{C}) \to
D(\mathcal{O}_{\mathcal{D}'})$로서
$Rf'_* \circ (g')^* = g^* \circ Rf_*$이다.
\end{lemma}

\begin{proof}
조건 (5)에 의해 $g^* = Lg^* = g^{-1}$이고
$(g')^* = L(g')^* = (g')^{-1}$이다. Lemma
\ref{sites-cohomology-lemma-modules-abelian-unbounded}에 의해 아벨 층들의 유도 범주
$D(\mathcal{C})$에서 결과를 증명하면 충분하다. 대상
$K \in D(\mathcal{C})$를 택하고, $K$를 나타내는 $\mathcal{C}$ 위
아벨 층들의 K-단사 복합체 $\mathcal{I}^\bullet$를 택하자. Derived
Categories, Lemma \ref{derived-lemma-adjoint-preserve-K-injectives}와 가정
(6)에 의해 $(g')^{-1}\mathcal{I}^\bullet$는 $\mathcal{C}'$ 위 아벨
층들의 K-단사 복합체이다. Modules on Sites, Lemma
\ref{sites-modules-lemma-special-square-cocontinuous}
에 의해 $f'_*(g')^{-1}\mathcal{I}^\bullet =
g^{-1}f_*\mathcal{I}^\bullet$이다. $f_*\mathcal{I}^\bullet$는 $Rf_*K$를
나타내고 $f'_*(g')^{-1}\mathcal{I}^\bullet$는
$Rf'_*(g')^{-1}K$를 나타내므로 결론이 따른다.
\end{proof}

\begin{lemma}
\label{sites-cohomology-lemma-special-square-continuous}
환 달린 토포스들의 가환 도식
$$
\xymatrix{
(\Sh(\mathcal{C}'), \mathcal{O}_{\mathcal{C}'})
\ar[r]_{(g', (g')^\sharp)} \ar[d]_{(f', (f')^\sharp)} &
(\Sh(\mathcal{C}), \mathcal{O}_\mathcal{C}) \ar[d]^{(f, f^\sharp)} \\
(\Sh(\mathcal{D}'), \mathcal{O}_{\mathcal{D}'}) \ar[r]^{(g, g^\sharp)} &
(\Sh(\mathcal{D}), \mathcal{O}_\mathcal{D})
}
$$
과 함자들의 도식
$$
\xymatrix{
\mathcal{C}' \ar[r]_{v'} &
\mathcal{C} \\
\mathcal{D}' \ar[r]^v \ar[u]^{u'} &
\mathcal{D} \ar[u]_u
}
$$
이 주어졌다고 하자. Sites, Sections
\ref{sites-section-morphism-sites}와
\ref{sites-section-cocontinuous-morphism-topoi}의 표기 아래 다음을
가정하자.
\begin{enumerate}
\item $u$와 $u'$는 연속이며 사상 $f$와 $f'$를 유도한다.
\item $v$와 $v'$는 쌍대연속이며 사상 $g$와 $g'$를 유도한다.
\item $u \circ v = v' \circ u'$,
\item $v$와 $v'$는 쌍대연속일 뿐 아니라 연속이다.
\item $g^{-1}\mathcal{O}_{\mathcal{D}} = \mathcal{O}_{\mathcal{D}'}$
이고 $(g')^{-1}\mathcal{O}_{\mathcal{C}} =
\mathcal{O}_{\mathcal{C}'}$이다.
\end{enumerate}
그러면 함자 $D^+(\mathcal{O}_\mathcal{C}) \to
D^+(\mathcal{O}_{\mathcal{D}'})$로서
$Rf'_* \circ (g')^* = g^* \circ Rf_*$이다. 더 나아가
\begin{enumerate}
\item[(6)] $g'_! : \textit{Ab}(\mathcal{C}') \to \textit{Ab}(\mathcal{C})$
가 완전이라고 가정하자\footnote{$\mathcal{C}'$에 올곱과 등화자가
존재하고 $v'$가 그것들과 가환하면 성립한다. Modules on Sites, Lemma
\ref{sites-modules-lemma-exactness-lower-shriek}을 보라.}.
\end{enumerate}
그러면 함자 $D(\mathcal{O}_\mathcal{C}) \to
D(\mathcal{O}_{\mathcal{D}'})$로서도
$Rf'_* \circ (g')^* = g^* \circ Rf_*$이다.
\end{lemma}

\begin{proof}
조건 (5)에 의해 $g^* = Lg^* = g^{-1}$이고
$(g')^* = L(g')^* = (g')^{-1}$이다. Lemma
\ref{sites-cohomology-lemma-modules-abelian-unbounded}에 의해 아벨 층들의 유도 범주
$D^+(\mathcal{C})$ 또는 $D(\mathcal{C})$에서 결과를 증명하면 충분하다.

\medskip\noindent
대상 $K \in D^+(\mathcal{C})$를 택하자. $K$를 나타내는
$\mathcal{C}$ 위 단사 아벨 층들의 아래로 유계인 복합체
$\mathcal{I}^\bullet$를 택한다. Lemma
\ref{sites-cohomology-lemma-pullback-injective-pre-limp}에 의해 모든 $p > 0$, 모든 $q$,
모든 $U' \in \Ob(\mathcal{C}')$에 대해
$H^p(U', (g')^{-1}\mathcal{I}^q) = 0$이다. Lemma
\ref{sites-cohomology-lemma-higher-direct-images}에서 보듯이
$R^pf'_*(g')^{-1}\mathcal{I}^q$는 준층
$V' \mapsto H^p(u'(V'), (g')^{-1}\mathcal{I}^q)$에 결합된 층이다.
따라서 $(g')^{-1}\mathcal{I}^q$는 함자 $f'_*$에 대해 오른쪽
비순환이다. Leray의 비순환성 보조정리(Derived Categories, Lemma
\ref{derived-lemma-leray-acyclicity})에 의해
$f'_*(g')^*\mathcal{I}^\bullet$는 $Rf'_*(g')^{-1}K$를 나타낸다.
Modules on Sites, Lemma
\ref{sites-modules-lemma-special-square-continuous}
에 의해 $f'_*(g')^{-1}\mathcal{I}^\bullet =
g^{-1}f_*\mathcal{I}^\bullet$이다. $g^{-1}f_*\mathcal{I}^\bullet$는
$g^{-1}Rf_*K$를 나타내므로 결론이 따른다.

\medskip\noindent
대상 $K \in D(\mathcal{C})$를 택하고, $K$를 나타내는 $\mathcal{C}$ 위
아벨 층들의 K-단사 복합체 $\mathcal{I}^\bullet$를 택하자. Derived
Categories, Lemma \ref{derived-lemma-adjoint-preserve-K-injectives}와 가정
(6)에 의해 $(g')^{-1}\mathcal{I}^\bullet$는 $\mathcal{C}'$ 위 아벨
층들의 K-단사 복합체이다. Modules on Sites, Lemma
\ref{sites-modules-lemma-special-square-continuous}
에 의해 $f'_*(g')^{-1}\mathcal{I}^\bullet =
g^{-1}f_*\mathcal{I}^\bullet$이다. $f_*\mathcal{I}^\bullet$는 $Rf_*K$를
나타내고 $f'_*(g')^{-1}\mathcal{I}^\bullet$는
$Rf'_*(g')^{-1}K$를 나타내므로 결론이 따른다.
\end{proof}








\section{올범주에 대한 아래 첨자 느낌표의 유도 함자}
\label{sites-cohomology-section-derived-lower-shriek-fibred}

\noindent
이 절에서는 Section \ref{sites-cohomology-section-derived-lower-shriek}에서 논의한 상황의
몇몇 특수한 경우를 자세히 다룬다. 가군에 대한 아래 첨자 느낌표와
아벨 군층에 대한 아래 첨자 느낌표가 서로 같음을 확인한다. 처음 읽을
때에는 이 절을 건너뛰기를 권한다.

\begin{situation}
\label{sites-cohomology-situation-fibred-category}
$(\mathcal{D}, \mathcal{O}_\mathcal{D})$를 환 달린 사이트라 하고,
$p : \mathcal{C} \to \mathcal{D}$를 올범주라 하자. $\mathcal{C}$에
$\mathcal{D}$에서 물려받은 위상을 준다(Stacks, Section
\ref{stacks-section-topology}). $p$에 결부된 토포스의 사상을
$\pi : \Sh(\mathcal{C}) \to \Sh(\mathcal{D})$로 나타낸다(Stacks,
Lemma \ref{stacks-lemma-topology-inherited-functorial}).
$\mathcal{O}_\mathcal{C} = \pi^{-1}\mathcal{O}_\mathcal{D}$라 두면 환
달린 토포스의 사상
$$
\pi :
(\Sh(\mathcal{C}), \mathcal{O}_\mathcal{C})
\longrightarrow
(\Sh(\mathcal{D}), \mathcal{O}_\mathcal{D})
$$
을 얻는다.
\end{situation}

\begin{lemma}
\label{sites-cohomology-lemma-fibred-category-with-object}
Situation \ref{sites-cohomology-situation-fibred-category}의 가정과 표기를 사용하자.
$U \in \Ob(\mathcal{C})$에 대해 유도된 토포스의 사상
$$
\pi_U : \Sh(\mathcal{C}/U) \longrightarrow \Sh(\mathcal{D}/p(U))
$$
을 생각하자. 그러면 토포스의 사상
$$
\sigma : \Sh(\mathcal{D}/p(U)) \to \Sh(\mathcal{C}/U)
$$
이 존재하여 $\pi_U \circ \sigma = \text{id}$이고
$\sigma^{-1} = \pi_{U, *}$이다.
\end{lemma}

\begin{proof}
$\pi_U$는 $\pi$를 국소화들에 제한한 것이다. Sites, Lemma
\ref{sites-lemma-localize-cocontinuous}를 보라. $\mathcal{D}/p(U)$의 대상
$V \to p(U)$에 대해, $p$가 올범주이므로 존재하는 $\mathcal{D}$ 위
$\mathcal{C}$의 강하게 데카르트인 사상을
$V \times_{p(U)} U \to U$로 나타내자. 함자
$$
v : \mathcal{D}/p(U) \to \mathcal{C}/U,\quad
V/p(U) \mapsto V \times_{p(U)} U/U
$$
는 $\mathcal{C}$ 위상의 정의에 의해 연속이다. 또한 강하게 데카르트인
사상의 정의에 의해 $p$의 오른쪽 수반이다. 따라서 Sites, Section
\ref{sites-section-cocontinuous-adjoint}에서 논의한 상황에 있으며, 층
$\pi_{U, *}\mathcal{F}$는 $V \mapsto
\mathcal{F}(V \times_{p(U)} U)$와 같다. 특히 Sites, Lemma
\ref{sites-lemma-have-functor-other-way-morphism}를 보라.

\medskip\noindent
여기서는 더 많은 사실이 성립한다. 즉 $\mathcal{C}$의 덮개에 나타나는
모든 사상이 강하게 데카르트이므로 함자 $v$는 쌍대연속이기도 하다.
따라서 $v$는 보조정리에 표시한 사상 $\sigma$를 정의한다. 등식
$\sigma^{-1} = \pi_{U, *}$는 정의에서 바로 따른다.
$\pi_U^{-1}\mathcal{G}$는 규칙
$U'/U \mapsto \mathcal{G}(p(U')/p(U))$로 주어지므로
$\sigma^{-1} \circ \pi_U^{-1} = \text{id}$이고, 따라서
$\pi_U \circ \sigma = \text{id}$이다.
\end{proof}

\begin{situation}
\label{sites-cohomology-situation-morphism-fibred-categories}
$(\mathcal{D}, \mathcal{O}_\mathcal{D})$를 환 달린 사이트라 하자.
$u : \mathcal{C}' \to \mathcal{C}$를 $\mathcal{D}$ 위 올범주들의
$1$-사상이라 하자(Categories, Definition
\ref{categories-definition-fibred-categories-over-C}). $\mathcal{C}$와
$\mathcal{C}'$에 각각 물려받은 위상을 주고(Stacks, Definition
\ref{stacks-definition-topology-inherited}),
$\pi : \Sh(\mathcal{C}) \to \Sh(\mathcal{D})$,
$\pi' : \Sh(\mathcal{C}') \to \Sh(\mathcal{D})$, 그리고
$g : \Sh(\mathcal{C}') \to \Sh(\mathcal{C})$
를 대응하는 토포스의 사상들이라 하자(Stacks, Lemma
\ref{stacks-lemma-topology-inherited-functorial}).
$\mathcal{O}_\mathcal{C} = \pi^{-1}\mathcal{O}_\mathcal{D}$ 및
$\mathcal{O}_{\mathcal{C}'} = (\pi')^{-1}\mathcal{O}_\mathcal{D}$라 두자.
$g^{-1}\mathcal{O}_\mathcal{C} = \mathcal{O}_{\mathcal{C}'}$이므로
$$
\xymatrix{
(\Sh(\mathcal{C}'), \mathcal{O}_{\mathcal{C}'}) \ar[rd]_{\pi'} \ar[rr]_g & &
(\Sh(\mathcal{C}), \mathcal{O}_\mathcal{C}) \ar[ld]^\pi \\
& (\Sh(\mathcal{D}), \mathcal{O}_\mathcal{D})
}
$$
은 환 달린 토포스 사상들의 가환 도식이다.
\end{situation}

\begin{lemma}
\label{sites-cohomology-lemma-morphism-fibred-categories-with-object}
Situation \ref{sites-cohomology-situation-morphism-fibred-categories}의 가정과 표기를
사용하자. $U' \in \Ob(\mathcal{C}')$에 대해 $U = u(U')$ 및
$V = p'(U')$라 두고, 유도된 환 달린 토포스의 사상들을 생각하자.
$$
\xymatrix{
(\Sh(\mathcal{C}'/U'), \mathcal{O}_{U'}) \ar[rd]_{\pi'_{U'}} \ar[rr]_{g'} & &
(\Sh(\mathcal{C}), \mathcal{O}_U) \ar[ld]^{\pi_U} \\
& (\Sh(\mathcal{D}/V), \mathcal{O}_V)
}
$$
그러면 토포스의 사상
$$
\sigma' : \Sh(\mathcal{D}/V) \to \Sh(\mathcal{C}'/U'),
$$
이 존재하여, $\sigma = g' \circ \sigma'$라 두면
$\pi'_{U'} \circ \sigma' = \text{id}$, $\pi_U \circ \sigma = \text{id}$,
$(\sigma')^{-1} = \pi'_{U', *}$, 그리고
$\sigma^{-1} = \pi_{U, *}$이다.
\end{lemma}

\begin{proof}
$v' : \mathcal{D}/V \to \mathcal{C}'/U'$를
$p' : \mathcal{C}' \to \mathcal{D}'$와 대상 $U'$에서 시작하여 Lemma
\ref{sites-cohomology-lemma-fibred-category-with-object}의 증명에서 구성한 함자라 하자.
$u$는 $\mathcal{D}$ 위 올범주들의 $1$-사상이므로 강하게 데카르트인
사상을 강하게 데카르트인 사상으로 보낸다. 따라서 함자
$v = u \circ v'$는 $p : \mathcal{C} \to \mathcal{D}$ 및 $U$에 대한
Lemma \ref{sites-cohomology-lemma-fibred-category-with-object}의 증명에 나오는 함자이다.
그러므로 이 보조정리는 그 보조정리에서 따른다.
\end{proof}

\begin{lemma}
\label{sites-cohomology-lemma-properties-lower-shriek-fibred-category}
Situation \ref{sites-cohomology-situation-morphism-fibred-categories}의 가정과 표기를
사용하자.
\begin{enumerate}
\item 가군과 아벨 층 각각에서 $g^* = g^{-1}$의 왼쪽 수반
$g_! : \textit{Mod}(\mathcal{O}_{\mathcal{C}'}) \to
\textit{Mod}(\mathcal{O}_\mathcal{C})$ 및
$g_!^{\textit{Ab}} : \textit{Ab}(\mathcal{C}') \to \textit{Ab}(\mathcal{C})$
가 존재한다.
\item 도식
$$
\xymatrix{
\textit{Mod}(\mathcal{O}_{\mathcal{C}'}) \ar[d] \ar[r]_{g_!} &
\textit{Mod}(\mathcal{O}_\mathcal{C}) \ar[d] \\
\textit{Ab}(\mathcal{C}') \ar[r]^{g_!^{\textit{Ab}}} &
\textit{Ab}(\mathcal{C})
}
$$
은 가환한다.
\item 가군과 아벨 층의 유도 범주 각각에서 $g^* = g^{-1}$의 왼쪽 수반
$Lg_! : D(\mathcal{O}_{\mathcal{C}'}) \to D(\mathcal{O}_\mathcal{C})$
및
$Lg_!^{\textit{Ab}} : D(\mathcal{C}') \to D(\mathcal{C})$
이 존재한다.
\item 도식
$$
\xymatrix{
D(\mathcal{O}_{\mathcal{C}'}) \ar[d] \ar[r]_{Lg_!} &
D(\mathcal{O}_\mathcal{C}) \ar[d] \\
D(\mathcal{C}') \ar[r]^{Lg_!^{\textit{Ab}}} &
D(\mathcal{C})
}
$$
은 가환한다.
\end{enumerate}
\end{lemma}

\begin{proof}
Stacks, Lemma \ref{stacks-lemma-topology-inherited-functorial}에 의해 함자
$u$는 연속이고 쌍대연속이다. 따라서 함자 $g_!$,
$g_!^{\textit{Ab}}$, $Lg_!$, $Lg_!^{\textit{Ab}}$의 존재는 Modules on
Sites, Sections
\ref{sites-modules-section-exactness-lower-shriek},
\ref{sites-modules-section-lower-shriek-modules}
및 Section \ref{sites-cohomology-section-derived-lower-shriek}에서 알 수 있다.

\medskip\noindent
(2)를 증명하려면 모든 $\mathcal{C}'$의 대상 $U'$에 대해 표준 사상
$$
g_!^{\textit{Ab}}j_{U'!}\mathcal{O}_{U'} \to j_{u(U')!}\mathcal{O}_{u(U')}
$$
이 동형사상임을 보이면 충분하다. Modules on Sites, Remark
\ref{sites-modules-remark-when-shriek-equal}을 보라. 마찬가지로 (4)를
증명하려면 모든 $\mathcal{C}'$의 대상 $U'$에 대해 표준 사상
$$
Lg_!^{\textit{Ab}}j_{U'!}\mathcal{O}_{U'} \to j_{u(U')!}\mathcal{O}_{u(U')}
$$
이 $D(\mathcal{C})$에서 동형사상임을 보이면 충분하다. Remark
\ref{sites-cohomology-remark-when-derived-shriek-equal}을 보라. 이는 앞의 식도 함의하므로
이 사실을 증명하겠다.

\medskip\noindent
국소화 사상 $j$에 대해 함자 $j_!$와 $j_!^{\textit{Ab}}$가 일치한다는
사실(Modules on Sites, Remark
\ref{sites-modules-remark-localize-shriek-equal} 참조)과 $j_!$가
완전이라는 사실(Modules on Sites, Lemma
\ref{sites-modules-lemma-extension-by-zero-exact})을 사용한다. Lemma
\ref{sites-cohomology-lemma-morphism-fibred-categories-with-object}의 표기를 채택하자.
Sites, Lemma \ref{sites-lemma-localize-cocontinuous}의 가환성과 수반
함자의 유일성에 의해
$Lg_!^{\textit{Ab}} \circ j_{U'!} =
j_{U!} \circ L(g')^{\textit{Ab}}_!$이므로
$L(g')^{\textit{Ab}}_!\mathcal{O}_{U'} = \mathcal{O}_U$임을 증명하면
충분하다. Lemma \ref{sites-cohomology-lemma-morphism-fibred-categories-with-object}의
결과를 사용하면 $D(\mathcal{C}/u(U'))$의 임의의 대상 $E$에 대해 다음
등식열을 얻는다.
\begin{align*}
\Hom_{D(\mathcal{C}/U)}(L(g')_!^{\textit{Ab}}\mathcal{O}_{U'}, E)
& =
\Hom_{D(\mathcal{C}'/U')}(\mathcal{O}_{U'}, (g')^{-1}E) \\
& =
\Hom_{D(\mathcal{C}'/U')}((\pi'_{U'})^{-1}\mathcal{O}_V, (g')^{-1}E) \\
& =
\Hom_{D(\mathcal{D}/V)}(\mathcal{O}_V, R\pi'_{U', *}(g')^{-1}E) \\
& =
\Hom_{D(\mathcal{D}/V)}(\mathcal{O}_V, (\sigma')^{-1}(g')^{-1}E) \\
& =
\Hom_{D(\mathcal{D}/V)}(\mathcal{O}_V, \sigma^{-1}E) \\
& =
\Hom_{D(\mathcal{D}/V)}(\mathcal{O}_V, \pi_{U, *}E) \\
& =
\Hom_{D(\mathcal{C}/U)}(\pi_U^{-1}\mathcal{O}_V, E) \\
& =
\Hom_{D(\mathcal{C}/U)}(\mathcal{O}_U, E)
\end{align*}
요네다 보조정리에 의해 결론을 얻는다.
\end{proof}

\begin{remark}
\label{sites-cohomology-remark-fibred-category}
Situation \ref{sites-cohomology-situation-fibred-category}의 가정과 표기를 사용하자.
$\mathcal{C}' = \mathcal{D}$라 두고 $u$를 $\mathcal{C}$의 구조 함자로
두면 Situation \ref{sites-cohomology-situation-morphism-fibred-categories}의 상황을 얻는다.
따라서 Lemma \ref{sites-cohomology-lemma-properties-lower-shriek-fibred-category}에 의해
함자 $\pi_!$, $\pi_!^{\textit{Ab}}$, $L\pi_!$,
$L\pi_!^{\textit{Ab}}$가 존재하며,
$forget \circ \pi_! = \pi_!^{\textit{Ab}} \circ forget$이고
$forget \circ L\pi_! = L\pi_!^{\textit{Ab}} \circ forget$이다.
\end{remark}

\begin{remark}
\label{sites-cohomology-remark-morphism-fibred-categories}
Situation \ref{sites-cohomology-situation-morphism-fibred-categories}의 가정과 표기를
사용하자. $\mathcal{F}$를 $\mathcal{C}$ 위의 아벨 층,
$\mathcal{F}'$를 $\mathcal{C}'$ 위의 아벨 층이라 하고,
$t : \mathcal{F}' \to g^{-1}\mathcal{F}$를 사상이라 하자. 그러면 표준
사상
$$
L\pi'_!(\mathcal{F}') \longrightarrow L\pi_!(\mathcal{F})
$$
을 얻는다. 여기에는 $t$에 수반되는 사상
$g_!\mathcal{F}' \to \mathcal{F}$, 사상
$Lg_!(\mathcal{F}') \to g_!\mathcal{F}'$, 그리고 등식
$L\pi'_! = L\pi_! \circ Lg_!$를 사용한다.
\end{remark}

\begin{lemma}
\label{sites-cohomology-lemma-compute-pi-shriek}
Situation \ref{sites-cohomology-situation-fibred-category}의 가정과 표기를 사용하자.
$\textit{Ab}(\mathcal{C})$의 $\mathcal{F}$에 대해 층
$\pi_!\mathcal{F}$는 준층
$$
V \longmapsto \colim_{\mathcal{C}_V^{opp}} \mathcal{F}|_{\mathcal{C}_V}
$$
에 결합된 층이다. 제한 사상은 증명에서 설명한다.
\end{lemma}

\begin{proof}
보조정리의 규칙을 $\mathcal{H}$로 나타내자. $\mathcal{D}$의 사상
$h : V' \to V$에 대해 올범주들의 당김 함자
$h^* : \mathcal{C}_V \to \mathcal{C}_{V'}$가 존재한다(Categories,
Definition \ref{categories-definition-pullback-functor-fibred-category}).
또한 $U \in \Ob(\mathcal{C}_V)$에 대해 $h$ 위에 놓이는 강하게
데카르트인 사상 $h^*U \to U$가 존재한다. 이 강하게 데카르트인
사상들을 따른 제한은 함자들의 변환
$$
\mathcal{F}|_{\mathcal{C}_V}
\longrightarrow
\mathcal{F}|_{\mathcal{C}_{V'}} \circ h^*.
$$
을 정의한다. 따라서 여극한 사이의 사상
$\mathcal{H}(V) \to \mathcal{H}(V')$를 얻는다. Categories, Lemma
\ref{categories-lemma-functorial-colimit}을 보라.

\medskip\noindent
보조정리를 증명하기 위해 모든 $\mathcal{C}$ 위의 층 $\mathcal{G}$에
대해
$$
\Mor_{\textit{PSh}(\mathcal{D})}(\mathcal{H}, \mathcal{G}) =
\Mor_{\Sh(\mathcal{C})}(\mathcal{F}, \pi^{-1}\mathcal{G})
$$
임을 보이자. 왼쪽의 한 원소는 모든 $\mathcal{C}$의 $U$에 대한 사상
$\mathcal{F}(U) \to \mathcal{G}(p(U))$들의 양립하는 계이다.
$\mathcal{C}$ 위상을 택한 방식에 의해
$\pi^{-1}\mathcal{G}(U) = \mathcal{G}(p(U))$이므로 오른쪽도 같은
것을 나타내며, 결론을 얻는다.
\end{proof}





\section{범주 위의 호몰로지}
\label{sites-cohomology-section-homology}

\noindent
점 위의 범주의 경우에는 아래 첨자 느낌표 함자의 왼쪽 유도 함자들을
호몰로지 함자라고 부르겠다.

\begin{example}[점 위의 범주]
\label{sites-cohomology-example-category-to-point}
$\mathcal{C}$를 범주라 하자. $\mathcal{C}$에 무질서 위상을 주자
(Sites, Example \ref{sites-example-indiscrete}). 그러면 $\mathcal{C}$에서
준층과 층은 일치한다. 유일한 대상 하나와 유일한 사상 하나를 갖는 범주를
$*$라 할 때 함자 $p : \mathcal{C} \to *$는 쌍대연속이고 연속이다. 대응하는
토포스의 사상을 $\pi : \Sh(\mathcal{C}) \to \Sh(*)$라 하자. $B$를 환이라
하자. $*$에는 환의 층 $B$를 주고, $\mathcal{C}$에는
$\mathcal{O}_\mathcal{C} = \pi^{-1}B$를 주며 이를 $\underline{B}$로 나타낸다.
그러면
$$
\pi : (\Sh(\mathcal{C}), \underline{B}) \to (\Sh(*), B)
$$
는 Situation \ref{sites-cohomology-situation-fibred-category}의 한 예이다.
Remark \ref{sites-cohomology-remark-fibred-category}에 의해 가군과 아벨 층에 대한
$\pi_!$를 구별할 필요가 없다. Lemma
\ref{sites-cohomology-lemma-compute-pi-shriek}에 의해
$\pi_!\mathcal{F} = \colim_{\mathcal{C}^{opp}} \mathcal{F}$.
따라서 $L_n\pi_!$는 여극한을 취하는 함자의 $n$번째 왼쪽 유도 함자이다.
이하에서는
$$
H_n(\mathcal{C}, \mathcal{F}) = L_n\pi_!(\mathcal{F})
$$
로 쓰고, 이를 $\mathcal{C}$ 위에서 {\it $\mathcal{F}$의 $n$번째 호몰로지 군}
이라 부르겠다.
\end{example}

\begin{example}[호몰로지의 계산]
\label{sites-cohomology-example-left-derived-colimits}
Example \ref{sites-cohomology-example-category-to-point}의 함자 $H_n(\mathcal{C}, -)$는
다음과 같이 계산할 수 있다. $\mathcal{F} \in
\Ob(\textit{Ab}(\mathcal{C}))$라 하자. 사슬 복합체
$$
K_\bullet(\mathcal{F}) :
\ \ldots \to
\bigoplus\nolimits_{U_2 \to U_1 \to U_0} \mathcal{F}(U_0)
\to
\bigoplus\nolimits_{U_1 \to U_0} \mathcal{F}(U_0)
\to
\bigoplus\nolimits_{U_0} \mathcal{F}(U_0)
$$
를 생각하자. 여기서 전이 사상들은
$$
(U_2 \to U_1 \to U_0, s)
\longmapsto
(U_1 \to U_0, s) - (U_2 \to U_0, s) + (U_2 \to U_1, s|_{U_1})
$$
로 주어지며, 다른 차수에서도 이와 비슷하다. 구성에 의해
$$
H_0(\mathcal{C}, \mathcal{F}) =
\colim_{\mathcal{C}^{opp}} \mathcal{F} =
H_0(K_\bullet(\mathcal{F})),
$$
이다. Categories, Lemma
\ref{categories-lemma-colimits-coproducts-coequalizers}를 보라.
$K_\bullet(\mathcal{F})$의 구성은 $\mathcal{F}$에 대해 함자적이며,
$\textit{Ab}(\mathcal{C})$의 짧은 완전열을 복합체의 짧은 완전열로 보낸다.
따라서 함자열 $\mathcal{F} \mapsto H_n(K_\bullet(\mathcal{F}))$는
$\delta$-함자를 이룬다. Homology, Definition
\ref{homology-definition-cohomological-delta-functor} 및 Lemma
\ref{homology-lemma-long-exact-sequence-cochain}를 보라.
$\mathcal{F} = j_{U!}\mathbf{Z}_U$일 때 복합체
$K_\bullet(\mathcal{F})$는 항들이
$$
X_n = \coprod\nolimits_{U_n \to \ldots \to U_1 \to U_0}
\Mor_\mathcal{C}(U_0, U)
$$
인 심플렉셜 집합 $X_\bullet$ 위의 자유 $\mathbf{Z}$-가군에 결부된 복합체이다.
이 심플렉셜 집합은 한원소 집합 $\{*\}$ 위의 상수 심플렉셜 집합과 호모토피
동치이다. 실제로 사상 $X_\bullet \to \{*\}$는 자명하고, 사상
$\{*\} \to X_n$는 $*$를 $(U \to \ldots \to U, \text{id}_U)$로 보내어
주어지며, 두 사상 $X_\bullet \to X_\bullet$ 사이의 호모토피를 정하는 사상들
$$
h_{n, i} : X_n \longrightarrow X_n
$$
(Simplicial, Lemma \ref{simplicial-lemma-relations-homotopy})은 규칙
$$
h_{n, i} :
(U_n \to \ldots \to U_0, f)
\longmapsto
(U_n \to \ldots \to U_i \to U \to \ldots \to U, \text{id})
$$
로 주어진다. 여기서 $i > 0$이고 $h_{n, 0} = \text{id}$이다. 확인은 생략한다.
이로부터 $K_\bullet(j_{U!}\mathbf{Z}_U)$의 음의 차수 코호몰로지가 자명함을
얻는다(Simplicial, Remark \ref{simplicial-remark-homotopy-better}의 함자성과
Simplicial, Lemma \ref{simplicial-lemma-homotopy-s-N}의 결과를 사용한다).
따라서 예를 들어 Homology, Lemma
\ref{homology-lemma-efface-implies-universal}와 Derived Categories, Lemma
\ref{derived-lemma-right-derived-delta-functor}의 쌍대들을 적용하면,
$K_\bullet(\mathcal{F})$가 $H_0(\mathcal{C}, -)$의 왼쪽 유도 함자
$H_n(\mathcal{C}, -)$들을 계산함을 알 수 있다.
\end{example}

\begin{example}
\label{sites-cohomology-example-morphism-categories}
$u : \mathcal{C}' \to \mathcal{C}$를 함자라 하자. $\mathcal{C}'$와
$\mathcal{C}$에 Example \ref{sites-cohomology-example-category-to-point}에서와 같이 무질서
위상을 주자. 유일한 대상 하나와 유일한 사상 하나를 갖는 범주를 $*$라 하면
함자 $u$, $\mathcal{C}' \to *$, $\mathcal{C} \to *$는 쌍대연속이고
연속이다. 대응하는 토포스의 사상들을
$g : \Sh(\mathcal{C}') \to \Sh(\mathcal{C})$,
$\pi' : \Sh(\mathcal{C}') \to \Sh(*)$,
$\pi : \Sh(\mathcal{C}) \to \Sh(*)$,
라 하자. $B$를 환이라 하자. $*$에는 환의 층 $B$를 주고,
$\mathcal{C}'$, $\mathcal{C}$에는 상수층 $\underline{B}$를 준다. 그러면
$$
\xymatrix{
(\Sh(\mathcal{C}'), \underline{B}) \ar[rd]_{\pi'} \ar[rr]_g & &
(\Sh(\mathcal{C}), \underline{B}) \ar[ld]^\pi \\
& (\Sh(*), B)
}
$$
는 Situation \ref{sites-cohomology-situation-morphism-fibred-categories}의 한 예이다. 따라서
Lemma \ref{sites-cohomology-lemma-properties-lower-shriek-fibred-category}를 $g$에 적용할 수
있으므로 가군과 아벨 층에 대한 $g_!$를 구별할 필요가 없다.
특히 Remark \ref{sites-cohomology-remark-morphism-fibred-categories}는 표준 사상
$$
H_n(\mathcal{C}', \mathcal{F}')
\longrightarrow
H_n(\mathcal{C}, \mathcal{F})
$$
을 준다. 여기서 $\textit{Ab}(\mathcal{C})$의 $\mathcal{F}$,
$\textit{Ab}(\mathcal{C}')$의 $\mathcal{F}'$, 그리고 사상
$t : \mathcal{F}' \to g^{-1}\mathcal{F}$가 주어져 있다. Example
\ref{sites-cohomology-example-left-derived-colimits}의 호몰로지 계산으로 표현하면, 이 사상들은
복합체의 사상
$$
K_\bullet(\mathcal{F}') \longrightarrow K_\bullet(\mathcal{F})
$$
에서 오며, 이 사상은 규칙
$$
(U'_n \to \ldots \to U'_0, s') \longmapsto
(u(U'_n) \to \ldots \to u(U'_0), t(s'))
$$
으로 주어진다. 표기는 자명하다.
\end{example}

\begin{remark}
\label{sites-cohomology-remark-map-evaluation-to-derived}
Example \ref{sites-cohomology-example-category-to-point}의 표기와 가정을 사용하자.
$\mathcal{F}^\bullet$을 $\mathcal{C}$ 위 아벨 층들의 유계 복합체라 하자.
$\mathcal{C}$의 임의의 대상 $U$에 대해 $D(\textit{Ab})$ 안의 표준 사상
$$
\mathcal{F}^\bullet(U) \longrightarrow L\pi_!(\mathcal{F}^\bullet)
$$
이 있다. $\mathcal{F}^\bullet$이 $\underline{B}$-가군의 복합체이면 이 사상은
$D(B)$ 안에 있다. 이를 증명하기 위해 $\mathcal{P}^\bullet$이 사영 대상들의
복합체인 준동형사상 $\mathcal{P}^\bullet \to \mathcal{F}^\bullet$을 택하여
$L\pi_!(\mathcal{F}^\bullet)$을 계산함을 주목하자. 그런데 위상이 무질서하므로
$\mathcal{P}^\bullet(U) \to \mathcal{F}^\bullet(U)$는 준동형사상이고, 따라서
$D(\textit{Ab})$, 각각 $D(B)$에서 역으로 만들 수 있다. 여기에 $\pi_!$를
여극한으로 계산함으로써 생기는 표준 사상
$\mathcal{P}^\bullet(U) \to \pi_!(\mathcal{P}^\bullet)$을 합성하면 원하는
화살표를 얻는다.
\end{remark}

\begin{lemma}
\label{sites-cohomology-lemma-initial-final}
Example \ref{sites-cohomology-example-category-to-point}의 표기와 가정을 사용하자.
$\mathcal{C}$가 시작 대상 또는 끝 대상을 가지면 $D(\textit{Ab})$, 각각
$D(B)$ 위에서 $L\pi_! \circ \pi^{-1} = \text{id}$이다.
\end{lemma}

\begin{proof}
$\mathcal{C}$가 시작 대상을 가지면 $\pi_!$는 이 대상에서 값을 취하여
계산되므로 주장은 명백하다. $\mathcal{C}$가 끝 대상을 가지면 $R\pi_*$는 이
대상에서 값을 취하여 계산되므로 $D(\textit{Ab})$, 각각 $D(B)$ 위에서
$R\pi_* \circ \pi^{-1} \cong \text{id}$이다. 따라서
$\pi^{-1} : D(\textit{Ab}) \to D(\mathcal{C})$,
각각 $\pi^{-1} : D(B) \to D(\underline{B})$는 충실충만하다. Categories,
Lemma \ref{categories-lemma-adjoint-fully-faithful}를 보라. 같은 보조정리에
의해 원하는 대로 $L\pi_! \circ \pi^{-1} = \text{id}$이다.
\end{proof}

\begin{lemma}
\label{sites-cohomology-lemma-change-of-rings}
Example \ref{sites-cohomology-example-category-to-point}의 표기와 가정을 사용하자.
$B \to B'$를 환 사상이라 하자. 환 달린 토포스의 가환 그림
$$
\xymatrix{
(\Sh(\mathcal{C}), \underline{B}) \ar[d]_\pi &
(\Sh(\mathcal{C}), \underline{B'}) \ar[d]^{\pi'} \ar[l]^h \\
(*, B) & (*, B') \ar[l]_f
}
$$
을 생각하자. 그러면 $L\pi_! \circ Lh^* = Lf^* \circ L\pi'_!$이다.
\end{lemma}

\begin{proof}
두 함자 모두 자명한 함자 $D(B') \to D(\underline{B})$의 오른쪽 수반이다.
\end{proof}

\begin{lemma}
\label{sites-cohomology-lemma-compute-by-cosimplicial-resolution}
Example \ref{sites-cohomology-example-category-to-point}의 표기와 가정을 사용하자.
$U_\bullet$을 $\mathcal{C}$의 코심플렉셜 대상이라 하고, 모든
$U \in \Ob(\mathcal{C})$에 대해 심플렉셜 집합
$\Mor_\mathcal{C}(U_\bullet, U)$가 한원소 집합 위의 상수 심플렉셜 집합과
호모토피 동치라고 하자. 그러면
$$
L\pi_!(\mathcal{F}) = \mathcal{F}(U_\bullet)
$$
가 $\textit{Ab}(\mathcal{C})$, 각각 $\textit{Mod}(\underline{B})$의
$\mathcal{F}$에 대해 함자적으로 $D(\textit{Ab})$, 각각 $D(B)$ 안에서
성립한다.
\end{lemma}

\begin{proof}
Lemma \ref{sites-cohomology-lemma-properties-lower-shriek-fibred-category}에 의해 가군과 아벨
층에 대한 $L\pi_!$가 일치하므로 $\mathcal{F}$가 아벨 층인 경우만 증명하면
충분하다. $U \in \Ob(\mathcal{C})$에 대해 아벨 층
$j_{U!}\mathbf{Z}_U$는 $\textit{Ab}(\mathcal{C})$의 사영 대상이다. 실제로
$\Hom(j_{U!}\mathbf{Z}_U, \mathcal{F}) = \mathcal{F}(U)$
이고 위상이 무질서하므로 단면을 취하는 함자는 완전하다. Modules on Sites,
Lemma \ref{sites-modules-lemma-module-quotient-flat}에 의해 모든 아벨 층은
$j_{U!}\mathbf{Z}_U$들의 직합의 몫이다. 따라서 항들이 위 꼴의 층들의
직합인 분해
$$
\ldots \to \mathcal{G}^{-1} \to \mathcal{G}^0 \to \mathcal{F} \to 0
$$
를 택하고 $L\pi_!(\mathcal{F}) = \pi_!(\mathcal{G}^\bullet)$를 취하여
$L\pi_!(\mathcal{F})$를 계산할 수 있다. 이중 복합체
$A^{\bullet, \bullet} = \mathcal{G}^\bullet(U_\bullet)$을 생각하자. 사상
$\mathcal{G}^0 \to \mathcal{F}$는 복합체의 사상
$A^{0, \bullet} \to \mathcal{F}(U_\bullet)$을 준다. $\pi_!$는
$\mathcal{C}^{opp}$ 위의 여극한을 취하여 계산되므로(Lemma
\ref{sites-cohomology-lemma-compute-pi-shriek}), 두 합성
$\mathcal{G}^m(U_1) \to \mathcal{G}^m(U_0) \to \pi_!\mathcal{G}^m$은 같다.
따라서 복합체의 표준 사상
$$
\text{Tot}(A^{\bullet, \bullet})
\longrightarrow
\pi_!(\mathcal{G}^\bullet) = L\pi_!(\mathcal{F})
$$
을 얻는다. 보조정리를 증명하려면 복합체들
$$
\ldots \to \mathcal{G}^m(U_1) \to \mathcal{G}^m(U_0) \to
\pi_!\mathcal{G}^m \to 0
$$
이 완전함을 보이면 충분하다. Homology, Lemma
\ref{homology-lemma-double-complex-gives-resolution}를 보라. 층
$\mathcal{G}^m$들은 $j_{U!}\mathbf{Z}_U$들의 직합이므로
$\mathcal{G} = j_{U!}\mathbf{Z}_U$인 경우로 귀착된다. 복합체
$j_{U!}\mathbf{Z}_U(U_\bullet)$은 한원소 집합과 호모토피 동치라고 가정한
심플렉셜 집합 $\Mor_\mathcal{C}(U_\bullet, U)$ 위의 자유
$\mathbf{Z}$-가군에 결부된 아벨 군의 복합체이다. 따라서
$$
j_{U!}\mathbf{Z}_U(U_\bullet) \to \mathbf{Z}
$$
는 아벨 군의 호모토피 동치이고, 따라서 준동형사상이다(Simplicial, Remark
\ref{simplicial-remark-homotopy-better} 및 Lemma
\ref{simplicial-lemma-homotopy-s-N}). Lemma
\ref{sites-cohomology-lemma-properties-lower-shriek-fibred-category}의 증명에서 보인 대로
$\pi_!j_{U!}\mathbf{Z}_U = \mathbf{Z}$이므로 증명이 끝난다.
\end{proof}

\begin{lemma}
\label{sites-cohomology-lemma-get-it-now}
Example \ref{sites-cohomology-example-morphism-categories}의 표기와 가정을 사용하자.
$\mathcal{C}'$의 코심플렉셜 대상 $U'_\bullet$이 존재하여 Lemma
\ref{sites-cohomology-lemma-compute-by-cosimplicial-resolution}을 $\mathcal{C}'$의
$U'_\bullet$과 $\mathcal{C}$의 $u(U'_\bullet)$ 모두에 적용할 수 있다고 하자.
그러면 함자
$D(\mathcal{C}) \to D(\textit{Ab})$,
각각 $D(\mathcal{C}, \underline{B}) \to D(B)$로서
$L\pi'_! \circ g^{-1} = L\pi_!$이다.
\end{lemma}

\begin{proof}
Lemma \ref{sites-cohomology-lemma-compute-by-cosimplicial-resolution}과 $g^{-1}$가 $u$와의
앞합성으로 주어진다는 사실에서 바로 따른다.
\end{proof}

\begin{lemma}
\label{sites-cohomology-lemma-product-categories}
$\mathcal{C}_i$, $i = 1, 2$를 범주들이라 하자.
$u_i : \mathcal{C}_1 \times \mathcal{C}_2 \to \mathcal{C}_i$를 사영 함자들이라
하자. $B$를 환이라 하자. Example \ref{sites-cohomology-example-morphism-categories}에서와 같이
대응하는 환 달린 토포스의 사상들을
$g_i : (\Sh(\mathcal{C}_1 \times \mathcal{C}_2), \underline{B}) \to
(\Sh(\mathcal{C}_i), \underline{B})$라 하자. $K_i \in
D(\mathcal{C}_i, B)$에 대해 자명한 표기 아래
$$
L(\pi_1 \times \pi_2)_!(
g_1^{-1}K_1 \otimes_{\underline{B}}^\mathbf{L} g_2^{-1}K_2)
=
L\pi_{1, !}(K_1) \otimes_B^\mathbf{L} L\pi_{2, !}(K_2)
$$
가 $D(B)$ 안에서 성립한다.
\end{lemma}

\begin{proof}
양변은 여극한과 가환하므로 $U \in \Ob(\mathcal{C}_1)$ 및
$V \in \Ob(\mathcal{C}_2)$에 대해
$K_1 = j_{U!}\underline{B}_U$이고 $K_2 = j_{V!}\underline{B}_V$
인 경우만 증명하면 충분하다. Lemma
\ref{sites-cohomology-lemma-existence-derived-lower-shriek}의 $L\pi_!$ 구성을 보라. 이 경우
$$
g_1^{-1}K_1 \otimes_{\underline{B}}^\mathbf{L} g_2^{-1}K_2 =
g_1^{-1}K_1 \otimes_{\underline{B}} g_2^{-1}K_2 =
j_{(U, V)!}\underline{B}_{(U, V)}
$$
이다. 확인은 생략한다. 보조정리의 식의 좌변과 우변은 모두 $B$로 계산되므로
결과가 따른다. Lemma \ref{sites-cohomology-lemma-existence-derived-lower-shriek}의
$L\pi_!$ 구성을 보라.
\end{proof}

\begin{lemma}
\label{sites-cohomology-lemma-eilenberg-zilber}
Example \ref{sites-cohomology-example-category-to-point}의 표기와 가정을 사용하자. Lemma
\ref{sites-cohomology-lemma-compute-by-cosimplicial-resolution}을 적용할 수 있는
$\mathcal{C}$의 코심플렉셜 대상 $U_\bullet$이 존재하면
$$
L\pi_!(K_1 \otimes^\mathbf{L}_{\underline{B}} K_2) =
L\pi_!(K_1) \otimes^\mathbf{L}_B L\pi_!(K_2)
$$
가 모든 $K_i \in D(\underline{B})$에 대해 성립한다.
\end{lemma}

\begin{proof}
범주와 함자의 그림
$$
\xymatrix{
& & \mathcal{C} \\
\mathcal{C} \ar[r]^-u &
\mathcal{C} \times \mathcal{C} \ar[rd]^{u_2} \ar[ru]_{u_1} \\
& & \mathcal{C}
}
$$
을 생각하자. 여기서 $u$는 대각 함자이고 $u_i$들은 사영 함자이다. 이로부터
환 달린 토포스의 사상 $g$, $g_1$, $g_2$를 얻는다. $\mathcal{C}$의 임의의
대상 $(U_1, U_2)$에 대해
$$
\Mor_{\mathcal{C} \times \mathcal{C}}(u(U_\bullet), (U_1, U_2)) =
\Mor_\mathcal{C}(U_\bullet, U_1) \times \Mor_\mathcal{C}(U_\bullet, U_2)
$$
이고, 이는 Simplicial, Lemma \ref{simplicial-lemma-products-homotopy}에 의해
점과 호모토피 동치이다. 따라서 Lemma \ref{sites-cohomology-lemma-get-it-now}에 의해
$D(\mathcal{C} \times \mathcal{C}, B)$의 임의의 $K$에 대해
$L\pi_!(g^{-1}K) = L(\pi \times \pi)_!(K)$이다.
$K = g_1^{-1}K_1 \otimes_B^\mathbf{L} g_2^{-1}K_2$로 택하자. 그러면
$g^{-1} = g^* = Lg^*$가 유도 텐서곱과 가환하므로
$g^{-1}K = K_1 \otimes^\mathbf{L}_{\underline{B}} K_2$이다(Lemma
\ref{sites-cohomology-lemma-pullback-tensor-product}). 이제 Lemma
\ref{sites-cohomology-lemma-product-categories}를 적용하면 증명이 끝난다.
\end{proof}

\begin{remark}[심플렉셜 가군]
\label{sites-cohomology-remark-simplicial-modules}
$\mathcal{C} = \Delta$라 하고 $B$를 임의의 환이라 하자. 이는 Example
\ref{sites-cohomology-example-category-to-point}의 특수한 경우이며, 여기서는 Lemma
\ref{sites-cohomology-lemma-compute-by-cosimplicial-resolution}의 가정이 성립한다. 구체적으로 항등 함자로
주어지는 $\Delta$의 코심플렉셜 대상을 $U_\bullet$이라 하자. 조건을 확인하려면
$[m] \in \Ob(\Delta)$에 대해 심플렉셜 집합
$\Delta[m] : n \mapsto \Mor_\Delta([n], [m])$이 점과 호모토피 동치임을
보여야 한다. 이는 Simplicial, Example
\ref{simplicial-example-simplex-contractible}에 설명되어 있다.

\medskip\noindent
이 상황에서 범주 $\textit{Mod}(\underline{B})$는 심플렉셜 $B$-가군의 범주이고,
함자 $L\pi_!$는 심플렉셜 $B$-가군 $M_\bullet$을 이에 결부된 $B$-가군의
복합체 $s(M_\bullet)$로 보낸다. 따라서 위 결과들을 심플렉셜 가군에 관한
결과로 다시 해석할 수 있다. 예를 들어 Lemma
\ref{sites-cohomology-lemma-eilenberg-zilber}의 특수한 경우로 다음을 얻는다.
$M_\bullet$, $M'_\bullet$이 평탄한 심플렉셜 $B$-가군이면 복합체
$s(M_\bullet \otimes_B M'_\bullet)$은 이중 복합체
$s(M_\bullet) \otimes_B s(M'_\bullet)$에 결부된 전체 복합체와
준동형이다. (힌트: 평탄성을 사용하여 유도 텐서곱을 보통 텐서곱으로
바꾸어라.) 이는 \cite{Eilenberg-Zilber}에서 찾을 수 있는
Eilenberg--Zilber 정리의 특수한 경우이다.
\end{remark}

\begin{lemma}
\label{sites-cohomology-lemma-O-homology-qis}
$\mathcal{C}$를 무질서 위상을 갖춘 범주라 하자. $\mathcal{O} \to
\mathcal{O}'$를 $\mathcal{C}$ 위 환의 층들의 사상이라 하자. 다음을 가정하자.
\begin{enumerate}
\item Lemma \ref{sites-cohomology-lemma-compute-by-cosimplicial-resolution}에서와 같은
$\mathcal{C}$의 코심플렉셜 대상 $U_\bullet$이 존재한다.
\item $L\pi_!\mathcal{O} \to L\pi_!\mathcal{O}'$는 동형사상이다.
\end{enumerate}
$D(\mathcal{O})$의 $K$에 대해
$$
L\pi_!(K) = L\pi_!(K \otimes_\mathcal{O}^\mathbf{L} \mathcal{O}')
$$
가 $D(\textit{Ab})$ 안에서 성립한다.
\end{lemma}

\begin{proof}
주의: 이 증명에서 $L\pi_!$는 아벨 층에 대한 $\pi_!$의 왼쪽 유도 함자를
나타낸다. $L\pi_!$는 여극한과 가환하므로 $\mathcal{O}$-가군의 위로 유계인
복합체에 대해서만 이를 증명하면 충분하다(Derived Categories, Proposition
\ref{derived-proposition-left-derived-exists}의 논증과 비교하거나, 단지 위로
유계인 복합체만 다루어도 된다). 이러한 모든 복합체는 항들이
$U \in \Ob(\mathcal{C})$에 대한 $j_{U!}\mathcal{O}_U$들의 직합인 위로 유계인
복합체와 준동형이다. Modules on Sites, Lemma
\ref{sites-modules-lemma-module-quotient-flat}을 보라. 따라서
$j_{U!}\mathcal{O}_U$에 대해 보조정리를 증명하면 충분하다. 가정에 의해
$$
S_\bullet = \Mor_\mathcal{C}(U_\bullet, U)
$$
는 한원소 집합 위의 상수 심플렉셜 집합과 호모토피 동치인 심플렉셜 집합이다.
$P_n = \mathcal{O}(U_n)$ 및 $P'_n = \mathcal{O}'(U_n)$로 두자. 심플렉셜
아벨 군
$$
X_\bullet : n \longmapsto \bigoplus\nolimits_{s \in S_n} P_n
$$
에 결부된 복합체는 Lemma \ref{sites-cohomology-lemma-compute-by-cosimplicial-resolution}에
의해 $L\pi_!(j_{U!}\mathcal{O}_U)$를 계산한다. $j_{U!}\mathcal{O}_U$는
평탄한 $\mathcal{O}$-가군이므로
$j_{U!}\mathcal{O}_U \otimes^\mathbf{L}_\mathcal{O} \mathcal{O}' =
j_{U!}\mathcal{O}'_U$이고, 이것의 $L\pi_!$는 심플렉셜 아벨 군
$$
X'_\bullet : n \longmapsto \bigoplus\nolimits_{s \in S_n} P'_n
$$
에 결부된 복합체로 계산된다. 심플렉셜 집합 $T_\bullet$에 항들이
$\bigoplus_{t \in T_n} P_n$인 심플렉셜 아벨 군을 대응시키는 규칙은
함자이므로 $X_\bullet \to P_\bullet$은 심플렉셜 아벨 군의 호모토피
동치이다. 마찬가지로 심플렉셜 집합 $T_\bullet$에 항들이
$\bigoplus_{t \in T_n} P'_n$인 심플렉셜 아벨 군을 대응시키는 규칙도
함자이다. 따라서 $X'_\bullet \to P'_\bullet$은 심플렉셜 아벨 군의
호모토피 동치이다. 가정에 의해 $P_\bullet \to P'_\bullet$은
준동형사상이다(Lemma \ref{sites-cohomology-lemma-compute-by-cosimplicial-resolution}에 의해
$P_\bullet$, 각각 $P'_\bullet$은 $L\pi_!\mathcal{O}$, 각각
$L\pi_!\mathcal{O}'$를 계산한다). 그러므로 원하는 대로 $X_\bullet$과
$X'_\bullet$은 서로 준동형이다.
\end{proof}

\begin{remark}
\label{sites-cohomology-remark-O-homology-B-homology-general}
$\mathcal{C}$와 $B$가 Example \ref{sites-cohomology-example-category-to-point}에서와 같다고
하자. Lemma \ref{sites-cohomology-lemma-compute-by-cosimplicial-resolution}에서와 같은
코심플렉셜 대상이 존재한다고 하자. $\mathcal{O} \to \underline{B}$를
동형사상 $L\pi_!\mathcal{O} \to L\pi_!\underline{B}$를 유도하는
$\mathcal{C}$ 위 환의 층들의 사상이라 하자. 이 경우 삼각 범주의 완전 함자
$$
L\pi_! : D(\mathcal{O}) \longrightarrow D(B)
$$
를 얻는다. 실제로 $D(\mathcal{O})$의 임의의 대상 $K$에 대해 Lemma
\ref{sites-cohomology-lemma-O-homology-qis}에 의해
$L\pi^{\textit{Ab}}_!(K) =
L\pi^{\textit{Ab}}_!(K \otimes_{\mathcal{O}}^\mathbf{L} \underline{B})$
이다. 따라서 표시된 함자를 $- \otimes^\mathbf{L}_\mathcal{O} \underline{B}$와
함자 $L\pi_! : D(\underline{B}) \to D(B)$의 합성으로 정의할 수 있다. 다시
말해, 보조정리의 $L\pi_!(K)$와
$L\pi_!(K \otimes_\mathcal{O}^\mathbf{L} \underline{B})$ 사이의 표준이고
함자적인 식별로부터 $L\pi_!(K)$ 위의 $B$-가군 구조를 얻는다.
\end{remark}







\section{아래 첨자 느낌표의 유도 함자 계산}
\label{sites-cohomology-section-calculate}

\noindent
이 절에서는 Section \ref{sites-cohomology-section-homology}의 결과를 적용하여 Situation
\ref{sites-cohomology-situation-fibred-category}의 $L\pi_!$와 Situation
\ref{sites-cohomology-situation-morphism-fibred-categories}의 $Lg_!$를 계산한다.

\begin{lemma}
\label{sites-cohomology-lemma-compute-left-derived-pi-shriek-pre}
Situation \ref{sites-cohomology-situation-fibred-category}의 가정과 표기를 사용하자.
$\textit{PAb}(\mathcal{C})$의 $\mathcal{F}$와 $n \geq 0$에 대해 아래의
준층에 결부된 $\mathcal{D}$ 위의 아벨 층을 $L_n(\mathcal{F})$라 하자.
$$
V \longmapsto H_n(\mathcal{C}_V, \mathcal{F}|_{\mathcal{C}_V})
$$
제한 사상들은 증명에서 설명한다. 그러면
$L_n(\mathcal{F}) = L_n(\mathcal{F}^\#)$이다.
\end{lemma}

\begin{proof}
$\mathcal{D}$의 사상 $h : V' \to V$에 대해 올범주의 당김 함자
$h^* : \mathcal{C}_V \to \mathcal{C}_{V'}$가 있다(Categories, Definition
\ref{categories-definition-pullback-functor-fibred-category}). 또한
$U \in \Ob(\mathcal{C}_V)$에 대해 $h$ 위에 놓이는 강하게 데카르트인 사상
$h^*U \to U$가 있다. 이 강하게 데카르트인 사상들을 따르는 제한은 함자들의
변환
$$
\mathcal{F}|_{\mathcal{C}_V}
\longrightarrow
\mathcal{F}|_{\mathcal{C}_{V'}} \circ h^*.
$$
을 정의한다. Example \ref{sites-cohomology-example-morphism-categories}에 의해 원하는 제한 사상
$$
H_n(\mathcal{C}_V, \mathcal{F}|_{\mathcal{C}_V})
\longrightarrow
H_n(\mathcal{C}_{V'}, \mathcal{F}|_{\mathcal{C}_{V'}})
$$
을 얻는다. 이 준층을 $L_{n, p}(\mathcal{F})$로 나타내면
$L_n(\mathcal{F}) = L_{n, p}(\mathcal{F})^\#$이다. 표준 사상
$\gamma : \mathcal{F} \to \mathcal{F}^+$(Sites, Theorem
\ref{sites-theorem-plus})는 표준 사상
$L_{n, p}(\mathcal{F}) \to L_{n, p}(\mathcal{F}^+)$를 정의한다. 층화한 뒤에
이 사상이 동형사상이 됨을 증명해야 한다.

\medskip\noindent
Example \ref{sites-cohomology-example-left-derived-colimits}의 호몰로지 계산을 사용하자.
$\mathcal{F}$를 올범주 $\mathcal{C}_V$에 제한한 것에 결부된 복합체를
$K_\bullet(\mathcal{F}|_{\mathcal{C}_V})$로 나타내자. 위의 설명에 의해 값이
$$
V \longmapsto K_\bullet(\mathcal{F}|_{\mathcal{C}_V})
$$
인 복합체의 준층 $K_\bullet(\mathcal{F})$를 얻으며, 그 코호몰로지 준층들은
준층 $L_{n, p}(\mathcal{F})$이다. 따라서
$$
K_\bullet(\mathcal{F}) \longrightarrow K_\bullet(\mathcal{F}^+)
$$
가 층화하면 동형사상이 됨을 보이면 충분하다.

\medskip\noindent
단사성. $V$를 $\mathcal{D}$의 대상이라 하고,
$\xi \in K_n(\mathcal{F})(V)$를 $K_n(\mathcal{F}^+)(V)$에서 영으로 가는
원소라 하자. $K_n(\mathcal{F})(V_j)$에서 $\xi|_{V_j}$가 영이 되도록 하는 피복
$\{V_j \to V\}$가 존재함을 보여야 한다. 다음과 같이 쓰자.
$$
\xi = \sum (U_{i, n + 1} \to \ldots \to U_{i, 0}, \sigma_i)
$$
여기서 $\sigma_i \in \mathcal{F}(U_{i, 0})$이다. $\mathcal{C}_V$의 각 사상열
$U_n \to \ldots \to U_0$이 많아야 한 번 나타나도록 정리한다. 복합체
$K_\bullet$의 정의에 나오는 합들은 직합이므로, 이것이
$K_\bullet(\mathcal{F}^+)(V)$에서 영으로 가려면 모든 $\sigma_i$가
$\mathcal{F}^+(U_{i, 0})$에서 영으로 가야 한다. $\mathcal{F}^+$의 구성에 의해
$\sigma_i|_{U_{i, 0, j}}$가 영이 되도록 하는 피복
$\{U_{i, 0, j} \to U_{i, 0}\}$들이 존재한다. $\mathcal{C}$ 위 위상의 구성에
의해 $U_{i, 0, j} \to U_{i, 0}$는 어떤 사상 $V_{i, j} \to V$의 당김으로
쓸 수 있고(Categories, Definition
\ref{categories-definition-pullback-functor-fibred-category}), 각
$\{V_{i, j} \to V\}$는 피복이다. 각 피복 $\{V_{i, j} \to V\}$를 지배하는 피복
$\{V_j \to V\}$를 택하자. 그러면 $\xi|_{V_j} = 0$임이 명백하다.

\medskip\noindent
전사성. 증명은 생략한다. 힌트: 단사성의 증명과 같이 논증하라.
\end{proof}

\begin{lemma}
\label{sites-cohomology-lemma-compute-left-derived-pi-shriek}
Situation \ref{sites-cohomology-situation-fibred-category}의 가정과 표기를 사용하자.
$\textit{Ab}(\mathcal{C})$의 $\mathcal{F}$와 $n \geq 0$에 대해 층
$L_n\pi_!(\mathcal{F})$는 Lemma
\ref{sites-cohomology-lemma-compute-left-derived-pi-shriek-pre}에서 구성한 층
$L_n(\mathcal{F})$와 같다.
\end{lemma}

\begin{proof}
$\textit{PAb}(\mathcal{C}) \to \textit{Ab}(\mathcal{C})$에서 정의된 함자열
$\mathcal{F} \mapsto L_n(\mathcal{F})$을 생각하자. 각
$V \in \Ob(\mathcal{D})$에 대해 함자열 $H_n(\mathcal{C}_V, - )$가
$\delta$-함자를 이루므로 함자 $\mathcal{F} \mapsto L_n(\mathcal{F})$들도
그러하다. 이들이 보편 $\delta$-함자를 이룸을 보이는 것이 목표이다. 이를 위해
이 함자들이 소멸하는 몇몇 아벨 준층을 구성한다.

\medskip\noindent
각 $U' \in \Ob(\mathcal{C})$에 대해 아벨 준층
$\mathcal{F}_{U'} = j_{U'!}^{\textit{PAb}}\mathbf{Z}_{U'}$
을 생각하자(Modules on Sites, Remark
\ref{sites-modules-remark-localize-presheaves}). 다음을 상기하자.
$$
\mathcal{F}_{U'}(U) =
\bigoplus\nolimits_{\Mor_\mathcal{C}(U, U')} \mathbf{Z}
$$
$U$가 $\mathcal{D})$의 $V = p(U)$ 위에 놓이고 $U'$가 $V' = p(U')$ 위에
놓이면, 임의의 사상 $a : U \to U'$는 $U \to h^*U' \to U'$로 유일하게
분해된다. 여기서 $h = p(a) : V \to V'$이다(Categories, Definition
\ref{categories-definition-pullback-functor-fibred-category} 참조). 따라서
$$
\mathcal{F}_{U'}|_{\mathcal{C}_V}
=
\bigoplus\nolimits_{h \in \Mor_\mathcal{D}(V, V')}
j_{h^*U'!}\mathbf{Z}_{h^*U'}
$$
이다. 여기서 $j_{h^*U'} : \Sh(\mathcal{C}_V/h^*U') \to
\Sh(\mathcal{C}_V)$는 국소화 사상이다. 층
$j_{h^*U'!}\mathbf{Z}_{h^*U'}$의 고차 호몰로지 군들은 소멸한다(Example
\ref{sites-cohomology-example-left-derived-colimits} 참조). 따라서 모든 $n > 0$과 모든
$U'$에 대해 $L_n(\mathcal{F}_{U'}) = 0$이다. 이로부터 임의의 아벨 준층
$\mathcal{F}$는 모든 $n > 0$에 대해 $L_n(\mathcal{G}) = 0$인 아벨 준층
$\mathcal{G}$의 몫임을 얻는다(Modules on Sites, Lemma
\ref{sites-modules-lemma-module-quotient-flat}).
$L_n(\mathcal{F}) = L_n(\mathcal{F}^\#)$이므로 아벨 층에 대해서도 같은
사실이 성립한다. 따라서 함자열 $L_n(-)$는
$\textit{Ab}(\mathcal{C})$ 위의 보편 델타 함자이다(Homology, Lemma
\ref{homology-lemma-efface-implies-universal}). Lemma
\ref{sites-cohomology-lemma-compute-pi-shriek}에 의해 $n = 0$일 때
$H^{-n}(L\pi_!(-))$와 일치하므로, 보편 $\delta$-함자의 유일성에서 결론이
따른다(Homology, Lemma
\ref{homology-lemma-uniqueness-universal-delta-functor} 및 Derived Categories,
Lemma \ref{derived-lemma-right-derived-delta-functor}).
\end{proof}

\begin{lemma}
\label{sites-cohomology-lemma-compute-left-derived-g-shriek}
Situation \ref{sites-cohomology-situation-morphism-fibred-categories}의 가정과 표기를 사용하자.
$\mathcal{C}'$ 위의 아벨 층 $\mathcal{F}'$에 대해 층
$L_ng_!(\mathcal{F}')$는 준층
$$
U \longmapsto H_n(\mathcal{I}_U, \mathcal{F}'_U)
$$
에 결부된 층이다. 표기와 제한 사상은 증명을 보라.
\end{lemma}

\begin{proof}
$p(U) = V$라 하자. 범주 $\mathcal{I}_U$의 대상은 순서쌍
$(U', \varphi)$이며, 여기서 $\varphi : U \to u(U')$는
$p(\varphi) = \text{id}_V$를 만족하는 $\mathcal{C}$의 사상, 즉 $\varphi$는 올범주
$\mathcal{C}_V$의 사상이다. 사상 $(U'_1, \varphi_1) \to
(U'_2, \varphi_2)$는 $\varphi_2 = u(a) \circ \varphi_1$를 만족하는 올범주
$\mathcal{C}'_V$의 사상 $a : U'_1 \to U'_2$로 주어진다. 준층
$\mathcal{F}'_U$는 $(U', \varphi)$를 $\mathcal{F}'(U')$로 보낸다. 아래에서
제한 사상들을 구성하겠다.

\medskip\noindent
Categories, Lemma
\ref{categories-lemma-ameliorate-morphism-fibred-categories}에서와 같은 $u$의
분해
$$
\xymatrix{
\mathcal{C}' \ar@<1ex>[r]^{u'} &
\mathcal{C}'' \ar[r]^{u''} \ar@<1ex>[l]^w & \mathcal{C}
}
$$
를 택하자. 그러면 $g_! = g''_! \circ g'_!$이고 유도 함자에 대해서도
마찬가지이다. 한편 함자 $g'_!$는 완전하다. Modules on Sites, Lemma
\ref{sites-modules-lemma-have-left-adjoint}를 보라. 따라서
$\mathcal{F}'' = g'_!\mathcal{F}'$로 두면
$Lg_!(\mathcal{F}') = Lg''_!(\mathcal{F}'')$이다. $w$에 결부된 토포스의
사상을 $h : \Sh(\mathcal{C}'') \to \Sh(\mathcal{C}')$라 할 때
$\mathcal{F}'' = h^{-1}\mathcal{F}'$임을 주목하자. Sites, Lemma
\ref{sites-lemma-have-left-adjoint}를 보라. 함자 $u''$는 $\mathcal{C}''$를
$\mathcal{C}$ 위의 올범주로 만들므로 $L_ng''_!$의 계산에 Lemma
\ref{sites-cohomology-lemma-compute-left-derived-pi-shriek}을 적용할 수 있다. Categories,
Lemma \ref{categories-lemma-ameliorate-morphism-fibred-categories}의 증명에 나오는
$\mathcal{C}''$의 구성에 의해 올범주 $\mathcal{C}''_U$가
$\mathcal{I}_U$와 같으므로 결과가 따른다. 더욱이
$h^{-1}\mathcal{F}'|_{\mathcal{C}''_U}$는 위에서 설명한 규칙으로 주어진다
(Stacks, Lemma \ref{stacks-lemma-topology-inherited-functorial}에 의해 $w$는
연속이고 쌍대연속이므로 Sites, Lemma \ref{sites-lemma-when-shriek}을 적용할
수 있다).
\end{proof}








\section{심플렉셜 가군}
\label{sites-cohomology-section-simplicial-modules}

\noindent
$A_\bullet$을 심플렉셜 환이라 하자. $A_\bullet$을 무질서 위상을 준
$\Delta$ 위의 층으로 생각할 수 있음을 상기하자. Simplicial, Section
\ref{simplicial-section-simplicial-presheaves}를 보라. 그러면 $A_\bullet$ 위의
심플렉셜 가군 $M_\bullet$은 단지 $\Delta$ 위의 $A_\bullet$-가군층이다. 다시
말해, 모든 $n \geq 0$에 대해 $A_n$-가군 $M_n$이 있고, 모든 사상
$\varphi : [n] \to [m]$에 대해 대응하는 사상
$$
M_\bullet(\varphi) : M_m \longrightarrow M_n
$$
이 있다. 이는 $A_\bullet(\varphi)$-선형이고, 이 사상들은 통상적인 방식으로
합성된다.

\medskip\noindent
$\mathcal{C}$를 사이트라 하자. $\mathcal{C}$ 위의 {\it 환의 심플렉셜 층}
$\mathcal{A}_\bullet$은 $\mathcal{C}$ 위 환의 층들의 범주에서의 심플렉셜
대상이다. 이 경우 대응 $U \mapsto \mathcal{A}_\bullet(U)$는 심플렉셜 환의
층이고, 실제로 두 개념은 동치이다. 심플렉셜 아벨 층, Lie 대수의 심플렉셜
층 등에 대해서도 비슷한 논의가 성립한다.

\medskip\noindent
그러나 위의 심플렉셜 환의 경우처럼 심플렉셜 층을 생각하는 다른 방법이 있다.
즉 사영
$$
p : \Delta \times \mathcal{C} \longrightarrow \mathcal{C}
$$
을 생각하자. 이는 강하게 데카르트인 사상들이 정확히
$([n], U) \to ([n], V)$ 꼴인 올범주를 정의한다. 범주
$\Delta \times \mathcal{C}$에 $\mathcal{C}$에서 물려받은 위상을 준다(Stacks,
Section \ref{stacks-section-topology} 참조). $\Delta \times \mathcal{C}$의
피복들에 대한 간단한 기술(Stacks, Lemma
\ref{stacks-lemma-topology-inherited})에서 $\mathcal{C}$ 위 환의 심플렉셜 층은
$\Delta \times \mathcal{C}$ 위 환의 층과 같은 것임이 바로 따른다.

\medskip\noindent
심플렉셜 환 위의 심플렉셜 가군과 유사하게, 환의 심플렉셜 층 위의 심플렉셜
가군을 다음과 같이 정의한다.

\begin{definition}
\label{sites-cohomology-definition-simplicial-module}
$\mathcal{C}$를 사이트라 하고, $\mathcal{A}_\bullet$을 $\mathcal{C}$ 위 환의
심플렉셜 층이라 하자. {\it 심플렉셜 $\mathcal{A}_\bullet$-가군}
$\mathcal{F}_\bullet$(때로는 {\it $\mathcal{A}_\bullet$-가군의 심플렉셜 층}
이라고도 한다)은 $\mathcal{A}_\bullet$에 결부된
$\Delta \times \mathcal{C}$ 위 환의 층에 대한 가군층이다.
\end{definition}

\noindent
심플렉셜 가군의 범주 $\textit{Mod}(\mathcal{A}_\bullet)$과 이에 대응하는 유도
범주 $D(\mathcal{A}_\bullet)$를 얻는다. 환의 심플렉셜 층의 사상
$\mathcal{A}_\bullet \to \mathcal{B}_\bullet$이 주어지면 함자
$$
- \otimes^\mathbf{L}_{\mathcal{A}_\bullet} \mathcal{B}_\bullet :
D(\mathcal{A}_\bullet)
\longrightarrow
D(\mathcal{B}_\bullet)
$$
를 얻는다. 더욱이 앞 절들의 내용은 함자
$$
L\pi_! : D(\mathcal{A}_\bullet) \longrightarrow D(\mathcal{C})
$$
를 정한다. 심플렉셜 가군 $\mathcal{F}_\bullet$이 주어지면 대상
$L\pi_!(\mathcal{F}_\bullet)$은 이에 결부된 사슬 복합체
$s(\mathcal{F}_\bullet)$로 표현된다(Simplicial, Section
\ref{simplicial-section-complexes}). 이는 Lemmas
\ref{sites-cohomology-lemma-compute-left-derived-pi-shriek} 및
\ref{sites-cohomology-lemma-compute-by-cosimplicial-resolution}에서 따른다.

\begin{lemma}
\label{sites-cohomology-lemma-base-change-by-qis}
$\mathcal{C}$를 사이트라 하자. $\mathcal{A}_\bullet \to
\mathcal{B}_\bullet$을 $\mathcal{C}$ 위 환의 심플렉셜 층의 준동형이라 하자.
$L\pi_!\mathcal{A}_\bullet \to L\pi_!\mathcal{B}_\bullet$이
$D(\mathcal{C})$ 안의 동형사상이면
$$
L\pi_!(K) =
L\pi_!(K \otimes^\mathbf{L}_{\mathcal{A}_\bullet} \mathcal{B}_\bullet)
$$
가 $D(\mathcal{A}_\bullet)$의 모든 $K$에 대해 성립한다.
\end{lemma}

\begin{proof}
$([n], U)$를 $\Delta \times \mathcal{C}$의 대상이라 하자. $L\pi_!$는
여극한과 가환하므로 $\mathcal{O}$-가군의 위로 유계인 복합체에 대해서만
이를 증명하면 충분하다(Derived Categories, Proposition
\ref{derived-proposition-left-derived-exists}의 논증과 비교하거나, 단지 위로
유계인 복합체만 다루어도 된다). 이러한 모든 복합체는 항들이 평탄 가군인
위로 유계인 복합체와 준동형이다. Modules on Sites, Lemma
\ref{sites-modules-lemma-module-quotient-flat}을 보라. 따라서 평탄한
$\mathcal{A}_\bullet$-가군 $\mathcal{F}$에 대해 보조정리를 증명하면 충분하다.
이 경우 유도 텐서곱은 보통 텐서곱이고, 이 역시 층이다. 따라서 Lemma
\ref{sites-cohomology-lemma-compute-left-derived-pi-shriek}에 의해 식의 양변의 코호몰로지 층을
Lemma \ref{sites-cohomology-lemma-compute-left-derived-pi-shriek-pre}의 절차로 계산할 수 있다.
그러므로 $\mathcal{F}$를 올범주들, 즉 $\Delta \times U$에 제한한 것에 대해
결과를 증명하면 충분하다. 이 경우 결과는 Lemma
\ref{sites-cohomology-lemma-O-homology-qis}에서 따른다.
\end{proof}

\begin{remark}
\label{sites-cohomology-remark-homology-augmentation}
$\mathcal{C}$를 사이트라 하자. $\epsilon : \mathcal{A}_\bullet \to
\mathcal{O}$를 환의 층들의 범주에서의 확대 사상이라 하자(Simplicial,
Definition \ref{simplicial-definition-augmentation}). $\epsilon$이 준동형사상
$s(\mathcal{A}_\bullet) \to \mathcal{O}$를 유도한다고 하자. 이 경우 삼각
범주의 완전 함자
$$
L\pi_! : D(\mathcal{A}_\bullet) \longrightarrow D(\mathcal{O})
$$
를 얻는다. 실제로 $D(\mathcal{A}_\bullet)$의 임의의 대상 $K$에 대해 Lemma
\ref{sites-cohomology-lemma-base-change-by-qis}에 의해
$L\pi_!(K) = L\pi_!(K \otimes_{\mathcal{A}_\bullet}^\mathbf{L} \mathcal{O})$
이다. 따라서 표시된 함자를 $- \otimes^\mathbf{L}_{\mathcal{A}_\bullet}
\mathcal{O}$와 Remark \ref{sites-cohomology-remark-fibred-category}의 함자
$L\pi_! : D(\Delta \times \mathcal{C}, \pi^{-1}\mathcal{O}) \to
D(\mathcal{O})$의 합성으로 정의할 수 있다. 다시 말해, 보조정리의
$L\pi_!(K)$와
$L\pi_!(K \otimes_{\mathcal{A}_\bullet}^\mathbf{L} \mathcal{O})$ 사이의
표준이고 함자적인 식별로부터 $L\pi_!(K)$ 위의 $\mathcal{O}$-가군 구조를
얻는다.
\end{remark}






\section{범주 위의 코호몰로지}
\label{sites-cohomology-section-cohomology}

\noindent
Example \ref{sites-cohomology-example-category-to-point}의 상황에서는 유도 함자 $L\pi_!$뿐만
아니라 함자 $R\pi_*$도 있다. $\mathcal{C}$ 위의 아벨 층 $\mathcal{F}$에
대해 $H_n(\mathcal{C}, \mathcal{F}) = H^{-n}(L\pi_!\mathcal{F})$이고
$H^n(\mathcal{C}, \mathcal{F}) = H^n(R\pi_*\mathcal{F})$이다.

\begin{example}[코호몰로지의 계산]
\label{sites-cohomology-example-right-derived-limits}
Example \ref{sites-cohomology-example-category-to-point}의 함자 $H^n(\mathcal{C}, -)$는 다음과
같이 계산할 수 있다. $\mathcal{F} \in \Ob(\textit{Ab}(\mathcal{C}))$라 하자.
쌍대사슬 복합체
$$
K^\bullet(\mathcal{F}) :
\prod\nolimits_{U_0} \mathcal{F}(U_0)
\to
\prod\nolimits_{U_0 \to U_1} \mathcal{F}(U_0)
\to
\prod\nolimits_{U_0 \to U_1 \to U_2} \mathcal{F}(U_0)
\to \ldots
$$
를 생각하자. 여기서 전이 사상들은
$$
(s_{U_0 \to U_1})
\longmapsto
((U_0 \to U_1 \to U_2) \mapsto s_{U_0 \to U_1} - s_{U_0 \to U_2}
+ s_{U_1 \to U_2}|_{U_0})
$$
로 주어지며, 다른 차수에서도 이와 비슷하다. 구성에 의해
$$
H^0(\mathcal{C}, \mathcal{F}) =
\lim_{\mathcal{C}^{opp}} \mathcal{F} =
H^0(K^\bullet(\mathcal{F})),
$$
이다. Categories, Lemma \ref{categories-lemma-limits-products-equalizers}를 보라.
$K^\bullet(\mathcal{F})$의 구성은 $\mathcal{F}$에 대해 함자적이며,
$\textit{Ab}(\mathcal{C})$의 짧은 완전열을 복합체의 짧은 완전열로 보낸다.
따라서 함자열 $\mathcal{F} \mapsto H^n(K^\bullet(\mathcal{F}))$는
$\delta$-함자를 이룬다. Homology, Definition
\ref{homology-definition-cohomological-delta-functor} 및 Lemma
\ref{homology-lemma-long-exact-sequence-cochain}를 보라. $\mathcal{C}$의 대상
$U$에 대응하며 $p_U^{-1}$가 $U$에서 값을 취하는 것과 같은 점을
$p_U : \Sh(*) \to \Sh(\mathcal{C})$로 나타내자. Sites, Example
\ref{sites-example-indiscrete-points}를 보라. $A$를 아벨 군이라 하고
$\mathcal{F} = p_{U, *}A$로 두자. 이 경우 복합체
$K^\bullet(\mathcal{F})$의 항은 $\text{Map}(X_n, A)$이며, 여기서
$$
X_n = \coprod\nolimits_{U_0 \to \ldots \to U_{n - 1} \to U_n}
\Mor_\mathcal{C}(U, U_0)
$$
이다. 이 심플렉셜 집합은 한원소 집합 $\{*\}$ 위의 상수 심플렉셜 집합과
호모토피 동치이다. 실제로 사상 $X_\bullet \to \{*\}$는 자명하고, 사상
$\{*\} \to X_n$는 $*$를 $(U \to \ldots \to U, \text{id}_U)$로 보내어
주어지며, 두 사상 $X_\bullet \to X_\bullet$ 사이의 호모토피를 정하는 사상들
$$
h_{n, i} : X_n \longrightarrow X_n
$$
(Simplicial, Lemma \ref{simplicial-lemma-relations-homotopy})은 규칙
$$
h_{n, i} :
(U_0 \to \ldots \to U_n, f)
\longmapsto
(U \to \ldots \to U \to U_i \to \ldots \to U_n, \text{id})
$$
로 주어진다. 여기서 $i > 0$이고 $h_{n, 0} = \text{id}$이다. 확인은 생략한다.
$\text{Map}(-, A)$가 반변 함자이므로 $K^\bullet(p_{U, *}A)$의 양의 차수
코호몰로지는 자명하다(Simplicial, Remark
\ref{simplicial-remark-homotopy-better}의 함자성과 Simplicial, Lemma
\ref{simplicial-lemma-homotopy-s-Q}의 결과를 사용한다). 따라서
$\mathcal{F}$가 $p_{U, *}A$ 꼴 층들의 곱일 때도
$K^\bullet(\mathcal{F})$는 양의 차수에서 비순환이다. $\mathcal{C}$ 위의 모든
아벨 층은 그러한 곱에 매장되므로, 예를 들어 Homology, Lemma
\ref{homology-lemma-efface-implies-universal}와 Derived Categories, Lemma
\ref{derived-lemma-right-derived-delta-functor}에 의해
$K^\bullet(\mathcal{F})$는 $H^0(\mathcal{C}, -)$의 왼쪽 유도 함자
$H^n(\mathcal{C}, -)$들을 계산한다.
\end{example}

\begin{example}[Ext의 계산]
\label{sites-cohomology-example-computing-exts}
Example \ref{sites-cohomology-example-category-to-point}에서 $\mathcal{C}$ 위 환의 층
$\mathcal{O}$가 더 주어졌다고 하자. $\mathcal{F}$, $\mathcal{G}$를
$\mathcal{O}$-가군이라 하자. 차수 $n$의 항이
$$
\prod\nolimits_{U_0 \to U_1 \to \ldots \to U_n}
\Hom_{\mathcal{O}(U_n)}(\mathcal{G}(U_n), \mathcal{F}(U_0))
$$
인 복합체 $K^\bullet(\mathcal{G}, \mathcal{F})$를 생각하자. 전이 사상은
$$
(\varphi_{U_0 \to U_1})
\longmapsto
((U_0 \to U_1 \to U_2) \mapsto
\varphi_{U_0 \to U_1} \circ \rho^{U_2}_{U_1}
- \varphi_{U_0 \to U_2}
+ \rho^{U_1}_{U_0} \circ \varphi_{U_1 \to U_2}
$$
로 주어지고, 다른 차수에서도 이와 비슷하다. 여기서 $\rho$들은 제한 사상을
나타낸다. 구성에 의해 모든 $\mathcal{O}$-가군의 순서쌍
$\mathcal{F}, \mathcal{G}$에 대해
$$
\Hom_\mathcal{O}(\mathcal{G}, \mathcal{F}) =
H^0(K^\bullet(\mathcal{G}, \mathcal{F}))
$$
이다. 대응
$(\mathcal{G}, \mathcal{F}) \mapsto K^\bullet(\mathcal{G}, \mathcal{F})$
은 첫째 변수의 직합을 곱으로 보내고 둘째 변수의 곱과 가환하는 이함자이다.
다음을 주장한다.
$$
\Ext^i_\mathcal{O}(\mathcal{G}, \mathcal{F}) =
H^i(K^\bullet(\mathcal{G}, \mathcal{F}))
$$
이는 $i \geq 0$이고 다음 중 하나가 성립하면 참이다.
\begin{enumerate}
\item 모든 $U \in \Ob(\mathcal{C})$에 대해 $\mathcal{G}(U)$는 사영
$\mathcal{O}(U)$-가군이다.
\item 모든 $U \in \Ob(\mathcal{C})$에 대해 $\mathcal{F}(U)$는 단사
$\mathcal{O}(U)$-가군이다.
\end{enumerate}
실제로 (1)의 경우 함자 $K^\bullet(\mathcal{G}, -)$는
$\mathcal{O}$-가군의 범주에서 아벨 군의 쌍대사슬 복합체의 범주로 가는 완전
함자이다. 따라서 Example \ref{sites-cohomology-example-right-derived-limits}에서와 같이
논증하면, $\mathcal{O}(U)$-가군 $A$에 대해 $\mathcal{F}$가
$p_{U, *}A$일 때 $K^\bullet(\mathcal{G}, \mathcal{F})$가 양의 차수에서
비순환임을 보이면 충분하다. 짧은 완전열
\begin{equation}
\label{sites-cohomology-equation-split}
0 \to \mathcal{G}' \to \bigoplus j_{U_i!}\mathcal{O}_{U_i} \to
\mathcal{G} \to 0
\end{equation}
을 택하자. Modules on Sites, Lemma
\ref{sites-modules-lemma-module-quotient-flat}을 보라. (1)은 가운데 층과 오른쪽
층에 대해 성립하므로 $\mathcal{G}'$에 대해서도 성립하며,
(\ref{sites-cohomology-equation-split})에서 $\mathcal{C}$의 대상에서 값을 취하면 가군의 분할
완전열을 얻는다. 따라서 복합체의 짧은 완전열
$$
0 \to
K^\bullet(\mathcal{G}, \mathcal{F}) \to
\prod K^\bullet(j_{U_i!}\mathcal{O}_{U_i}, \mathcal{F}) \to
K^\bullet(\mathcal{G}', \mathcal{F}) \to 0
$$
을 모든 $\mathcal{F}$에 대해, 특히 $\mathcal{F} = p_{U, *}A$에 대해 얻는다.
$H^0$ 위에서는
$$
0 \to \Hom(\mathcal{G}, p_{U, *}A) \to
\Hom(\prod j_{U_i!}\mathcal{O}_{U_i}, p_{U, *}A) \to
\Hom(\mathcal{G}', p_{U, *}A) \to 0
$$
를 얻는다. 이는
$\Hom(\mathcal{H}, p_{U, *}A) = \Hom_{\mathcal{O}(U)}(\mathcal{H}(U), A)$
이고 $U$ 위에서 (\ref{sites-cohomology-equation-split})의 단면열이 분할 완전하므로 완전하다.
따라서 차원 이동을 사용하면 모든 $U, U' \in \Ob(\mathcal{C})$에 대해
$K^\bullet(j_{U'!}\mathcal{O}_{U'}, p_{U, *}A)$가 양의 차수에서 비순환임을
증명하면 충분하다. 이 경우 $K^n(j_{U'!}\mathcal{O}_{U'}, p_{U, *}A)$는
$$
\prod\nolimits_{U \to U_0 \to U_1 \to \ldots \to U_n \to U'} A
$$
와 같다. 다시 말해, $K^\bullet(j_{U'!}\mathcal{O}_{U'}, p_{U, *}A)$는 항이
$\text{Map}(X_\bullet, A)$인 복합체이며, 여기서
$$
X_n = \coprod\nolimits_{U_0 \to \ldots \to U_{n - 1} \to U_n}
\Mor_\mathcal{C}(U, U_0) \times \Mor_\mathcal{C}(U_n, U')
$$
이다. 이 심플렉셜 집합은 한원소 집합 $\{*\}$ 위의 상수 심플렉셜 집합과
호모토피 동치이며, 이는 Example \ref{sites-cohomology-example-right-derived-limits}의 해당
주장과 정확히 같은 방법으로 증명할 수 있다. 이로써 주장의 증명이 끝난다.

\medskip\noindent
(2)의 논증도 이와 비슷하다(다만 쌍대이다).
\end{example}






\section{범주 위의 가군}
\label{sites-cohomology-section-modules-cohomology}

\noindent
이 절의 내용은 스킴의 범주 위의 준군 스택에서 준연접 가군의
유도 범주에 대한 한 변형을 정의하는 데 사용된다.
Sheaves on Stacks, Section \ref{stacks-sheaves-section-QC}를 보라.

\medskip\noindent
$\mathcal{C}$를 범주라 하자. $\mathcal{C}$를 무질서 위상을 갖춘
사이트로 생각한다. Example \ref{sites-cohomology-example-computing-exts}에서와 같이
$\mathcal{O}$를 $\mathcal{C}$ 위의 환의 층이라 하자.
달리 말하면 $\mathcal{O}$는 범주 $\mathcal{C}$ 위의 환의
준층이다. Categories, Definition \ref{categories-definition-presheaf}를 보라.

\begin{definition}
\label{sites-cohomology-definition-cartesian}
위의 상황에서 $\mathit{QC}(\mathcal{C}, \mathcal{O})$, 또는 간단히
{\it $\mathit{QC}(\mathcal{O})$}로
$D(\mathcal{O}) = D(\mathcal{C}, \mathcal{O})$의 다음 충만한
부분범주를 나타낸다. 그 대상은 $\mathcal{C}$의 모든 $U \to V$에 대해
표준 사상
$$
R\Gamma(V, K) \otimes_{\mathcal{O}(V)}^\mathbf{L} \mathcal{O}(U)
\longrightarrow
R\Gamma(U, K)
$$
이 $D(\mathcal{O}(U))$에서 동형이 되는 대상 $K$이다.
\end{definition}

\begin{lemma}
\label{sites-cohomology-lemma-cartesian-good-sub}
위의 상황에서 부분범주 $\mathit{QC}(\mathcal{O})$는
$D(\mathcal{O})$의 엄밀히 충만하고 포화된 삼각 부분범주이며,
임의의 직합에 의해 보존된다.
\end{lemma}

\begin{proof}
$U$를 $\mathcal{C}$의 대상이라 하자. $\mathcal{C}$ 위의 위상이
무질서 위상이므로 함자 $\mathcal{F} \mapsto \mathcal{F}(U)$는
완전하며 직합과 가환한다. 따라서 완전 함자
$K \mapsto R\Gamma(U, K)$는 $K$를 $\mathcal{O}$-가군의 임의의 복합체
$\mathcal{F}^\bullet$로 나타내고 $\mathcal{F}^\bullet(U)$를 취하여
계산된다. 그러므로 $R\Gamma(U, -)$는 직합과 가환한다.
Injectives, Lemma \ref{injectives-lemma-derived-products}를 보라.
마찬가지로 $\mathcal{C}$의 사상 $U \to V$가 주어지면
유도 텐서곱 함자
$- \otimes_{\mathcal{O}(V)}^\mathbf{L} \mathcal{O}(U) :
D(\mathcal{O}(V)) \to D(\mathcal{O}(U))$
는 완전하며 직합과 가환한다. 이 관찰들로부터 보조정리가 곧바로
따른다. 세부 사항은 생략한다.
\end{proof}

\begin{lemma}
\label{sites-cohomology-lemma-QC-bounded-above}
위의 상황에서 $M$이 $\mathit{QC}(\mathcal{O})$의 대상이고,
모든 $i > b$에 대해 $H^i(M) = 0$인 $b \in \mathbf{Z}$가 있다고
가정하자. 그러면 $H^b(M)$은 Modules on Sites, Definition
\ref{sites-modules-definition-site-local}의 의미에서
$(\mathcal{C}, \mathcal{O})$ 위의 준연접 가군이다.
\end{lemma}

\begin{proof}
Modules on Sites, Lemma
\ref{sites-modules-lemma-chaotic-quasi-coherent}에 의해,
$\mathcal{C}$의 모든 사상 $U \to V$에 대해 사상
$$
H^p(M)(V) \otimes_{\mathcal{O}(V)} \mathcal{O}(U) \to H^b(M)(U)
$$
이 동형임을 보이면 충분하다. 사상
$$
R\Gamma(V, M) \otimes_{\mathcal{O}(V)}^\mathbf{L} \mathcal{O}(U)
\to R\Gamma(U, M)
$$
이 동형이라는 것이 주어져 있다. 따라서 예를 들어 Tor 스펙트럼
열에 의해 결론이 따른다. 세부 사항은 생략한다.
\end{proof}

\begin{lemma}
\label{sites-cohomology-lemma-QC-final-object}
위의 상황에서 $\mathcal{C}$가 끝 대상 $X$를 갖는다고 가정하자.
$R = \mathcal{O}(X)$로 놓고, 자명한 사이트 사상을
$f : (\mathcal{C}, \mathcal{O}) \to (pt, R)$로 나타내자.
그러면 $Lf^*$와 $Rf_*$에 의해
$\mathit{QC}(\mathcal{O}) = D(R)$이다.
\end{lemma}

\begin{proof}
생략한다.
\end{proof}

\begin{lemma}
\label{sites-cohomology-lemma-QC-hom-out-of}
위의 상황에서 $K$가 $\mathit{QC}(\mathcal{O})$의 대상이고
$M$은 $D(\mathcal{O})$의 임의의 대상이라고 가정하자.
$\mathcal{C}$의 모든 대상 $U$에 대해
$$
\Hom_{D(\mathcal{O}_U)}(K|_U, M|_U) =
R\Hom_{\mathcal{O}(U)}(R\Gamma(U, K), R\Gamma(U, M))
$$
이다.
\end{lemma}

\begin{proof}
$\mathcal{C}$를 $\mathcal{C}/U$로 바꾸어도 된다. 따라서
$U = X$가 $\mathcal{C}$의 끝 대상이라고 가정할 수 있다.
Lemma \ref{sites-cohomology-lemma-QC-final-object}에 의해 $K = Lf^*P$이고, 여기서
$P = R\Gamma(U, K) = R\Gamma(X, K) = Rf_*K$이다.
$Lf^*$가 $Rf_*(-) = R\Gamma(U, -)$의 왼쪽 수반이므로 결론이 따른다.
\end{proof}

\noindent
$(\mathcal{C}, \mathcal{O})$가 위와 같다고 하자. $\mathcal{O}$-가군의
복합체 $\mathcal{F}^\bullet$에 대해 $\mathcal{F}^\bullet$의
{\it 크기} $|\mathcal{F}^\bullet|$를
$$
|\mathcal{F}^\bullet| =
\left|
\coprod\nolimits_{i \in \mathbf{Z},\ U \in \Ob(\mathcal{C})} \mathcal{F}^i(U)
\right|
$$
로 정의한다. $D(\mathcal{O})$의 대상 $K$에 대해 $K$의
{\it 크기} $|K|$를 기수
$$
|K| = \min
\left\{
\left|
\mathcal{F}^\bullet
\right|
\text{ where }\mathcal{F}^\bullet\text{ represents }K
\right\}
$$
로 정의한다. 기수의 성질에 의해 이 최솟값은 존재한다.

\begin{lemma}
\label{sites-cohomology-lemma-build-from}
위의 상황에서 다음 성질을 갖는 기수 $\kappa$가 존재한다.
$\mathcal{O}$-가군의 복합체 $\mathcal{F}^\bullet$와 부분집합
$\Omega^i_U \subset \mathcal{F}^i(U)$가 주어지면,
$\Omega^i_U \subset \mathcal{H}^i(U)$이고
$|\mathcal{H}^\bullet| \leq \max(\kappa, |\bigcup \Omega^i_U|)$인
부분복합체 $\mathcal{H}^\bullet \subset \mathcal{F}^\bullet$가 존재한다.
\end{lemma}

\begin{proof}
$\mathcal{H}^i(U)$를 임의의 사상 $f : U \to V$를 따른 제한으로 얻는
$\Omega^i_V$의 상들과 $\text{d}(\Omega^{i - 1}_U)$가 생성하는
$\mathcal{F}^i(U)$의 $\mathcal{O}(U)$-부분가군으로 정의한다.
$\mathcal{H}^i(U)$의 크기는 $\aleph_0$, $\mathcal{O}(U)$의 크기,
$\text{Arrows}(\mathcal{C})$의 크기, 그리고
$|\bigcup \Omega^i_U|$의 최댓값으로 제한된다. 세부 사항은 생략한다.
\end{proof}

\begin{lemma}
\label{sites-cohomology-lemma-qis-small}
위의 상황에서 다음 성질을 갖는 기수 $\kappa$가 존재한다.
$D(\mathcal{O})$의 대상 $K$를 나타내는 $\mathcal{O}$-가군의 복합체
$\mathcal{F}^\bullet$가 주어지면, $\mathcal{H}^\bullet$가 $K$를 나타내고
$|\mathcal{H}^\bullet| \leq \max(\kappa, |K|)$를 만족하는
부분복합체 $\mathcal{H}^\bullet \subset \mathcal{F}^\bullet$가 존재한다.
\end{lemma}

\begin{proof}
먼저 모든 $i$와 $U$에 대해 부분집합
$\Omega^i_U \subset \Ker(\text{d} :
\mathcal{F}^i(U) \to \mathcal{F}^{i + 1}(U))$
를 택하여
$H^i(K)(U) = H^i(\mathcal{F}^\bullet(U))$ 위로
전단사로 보내지게 한다. $K$를 크기가 $|K|$인 복합체로 나타낼
수 있으므로 $|\Omega^i_U| \leq |K|$이다.
Lemma \ref{sites-cohomology-lemma-build-from}을 적용하면 $\Omega^i_U$를 포함하고
크기가 $\max(\kappa, |K|)$ 이하인 부분복합체
$\mathcal{S}^\bullet \subset \mathcal{F}^\bullet$를 얻는다.
따라서 $H^i(\mathcal{S}^\bullet) \to H^i(\mathcal{F}^\bullet)$는
층의 전사이다.

\medskip\noindent
부분복합체
$$
\mathcal{S}^\bullet =
\mathcal{S}_0^\bullet \subset
\mathcal{S}_1^\bullet \subset
\mathcal{S}_2^\bullet \subset \ldots \subset \mathcal{F}^\bullet
$$
를 귀납적으로 구성하되, 그 크기는 $\leq \max(\kappa, |K|)$이고
$H^i(\mathcal{S}_n^\bullet) \to H^i(\mathcal{F}^\bullet)$의 핵이
$H^i(\mathcal{S}_n^\bullet) \to H^i(\mathcal{S}_{n + 1}^\bullet)$의
핵과 같게 할 것이다. 이 구성이 끝나면
$\mathcal{H}^\bullet = \bigcup \mathcal{S}_n^\bullet$를 답으로 취할 수 있다.

\medskip\noindent
$\mathcal{S}_n^\bullet$가 주어졌을 때
$\mathcal{S}_{n + 1}^\bullet$를 다음과 같이 구성한다.
모든 $U$와 $i$에 대해 부분집합
$\Omega^{i - 1}_U \subset \mathcal{F}^{i - 1}(U)$
을 택하여 $\text{d} : \mathcal{F}^{i - 1}(U) \to \mathcal{F}^i(U)$가
$\Omega^{i - 1}_U$를
$$
\mathcal{S}_n^i(U) \cap
\text{Im}(\text{d} : \mathcal{F}^{i - 1}(U)\to \mathcal{F}^i(U))
$$
위로 전단사로 보내게 한다. $\mathcal{S}_n^i(U)$의 크기가 그렇게
제한되어 있으므로 $|\Omega^i_U| \leq |K|$임을 주목하라.
그런 다음 부분집합
$$
\mathcal{S}^i(U) \cup \Omega^i_U \subset \mathcal{F}^i(U)
$$
에 Lemma \ref{sites-cohomology-lemma-build-from}을 적용하여
$\mathcal{S}_{n + 1}^\bullet$를 얻는다. 이제 모든 것이 명백하다.
\end{proof}

\begin{lemma}
\label{sites-cohomology-lemma-cartesian-colimit-kappa}
위의 상황에서 다음 성질들을 갖는 기수 $\kappa$가 존재한다.
\begin{enumerate}
\item $\mathit{QC}(\mathcal{O})$의 모든 영이 아닌 대상 $K$에 대해,
$|E| \leq \kappa$인 $\mathit{QC}(\mathcal{O})$의 영이 아닌 사상
$E \to K$가 존재한다.
\item $|E| \leq \kappa$인 $\mathit{QC}(\mathcal{O})$의 모든 사상
$\alpha : E \to \bigoplus_n K_n$에 대해,
$|E_n| \leq \kappa$인 $\mathit{QC}(\mathcal{O})$의 사상
$E_n \to K_n$들이 존재하여 $\alpha$가
$\bigoplus E_n \to \bigoplus K_n$를 통해 분해된다.
\end{enumerate}
\end{lemma}

\begin{proof}
$\kappa$를 다음 기수들의 집합의 상계로 택한다.
\begin{enumerate}
\item 모든 $U \in \Ob(\mathcal{C})$에 대한
$|\coprod_V j_{U!}\mathcal{O}_U(V)|$,
\item $\mathcal{C}$의 모든 사상 $U \to V$에 대해
More on Algebra, Lemma \ref{more-algebra-lemma-enlarge}에서 얻는 기수
$\kappa(\mathcal{O}(V) \to \mathcal{O}(U))$,
\item Lemma \ref{sites-cohomology-lemma-qis-small}에서 얻는 기수.
\end{enumerate}
다음을 주장한다. $\mathit{QC}(\mathcal{O})$의 대상을 나타내는 임의의
복합체 $\mathcal{F}^\bullet$와, $|\mathcal{S}^\bullet| \leq \kappa$인
임의의 부분복합체
$\mathcal{S}^\bullet \subset \mathcal{F}^\bullet$
에 대해, $\mathcal{S}^\bullet$를 포함하는
$\mathcal{F}^\bullet$의 부분복합체 $\mathcal{H}^\bullet$가 존재하여
$\mathcal{H}^\bullet$는 $\mathit{QC}(\mathcal{O})$의 대상을 나타내고
$|\mathcal{H}^\bullet| \leq \kappa$이다.
다음 두 문단에서 이 주장으로부터 보조정리가 따름을 보인다.

\medskip\noindent
(1)에서처럼 $K$를 $\mathit{QC}(\mathcal{O})$의 영이 아닌 대상이라
하자. $K$가 $\mathcal{O}$-가군의 복합체
$\mathcal{F}^\bullet$로 나타난다고 하자. 그러면 어떤 $i$에 대해
$H^i(\mathcal{F}^\bullet)$가 영이 아니다. 따라서 $\mathcal{C}$의
대상 $U$와 $d(s) = 0$인 단면 $s \in \mathcal{F}^i(U)$가 존재하여,
$U$ 위의 $H^i(\mathcal{F}^\bullet)$의 영이 아닌 단면을 정한다.
그러면 $s : j_{U!}\mathcal{O}_U[-i] \to \mathcal{F}^\bullet$의 상은
부분복합체
$\mathcal{S}^\bullet \subset \mathcal{F}^\bullet$이고
$|\mathcal{S}^\bullet| \leq \kappa$이다. 위의 주장을 적용하면
$|\mathcal{H}^\bullet| \leq \kappa$인
$\mathit{QC}(\mathcal{O})$의 영이 아닌 사상
$\mathcal{H}^\bullet \to \mathcal{F}^\bullet$를 얻는다. 따라서 (1)이 성립한다.

\medskip\noindent
(2)에서와 같은 $\alpha : E \to \bigoplus K_n$을 택하자.
$K_n$을 나타내는 임의의 복합체 $\mathcal{K}_n^\bullet$를 택한다.
그러면 $\bigoplus \mathcal{K}_n^\bullet$는 $\bigoplus K_n$을 나타낸다.
유도 범주의 구성에 의해 $E$를 복합체 $\mathcal{E}^\bullet$로 나타내어
$\alpha$가 복합체의 사상
$a : \mathcal{E}^\bullet \to \bigoplus \mathcal{K}_n^\bullet$
로 나타나게 할 수 있다. Lemma \ref{sites-cohomology-lemma-qis-small}과 위에서 택한
$\kappa$에 의해 $|\mathcal{E}^\bullet| \leq \kappa$라고 가정해도 된다.
위의 주장을 적용하면
$a_n : \mathcal{E}^\bullet \to \mathcal{K}_n^\bullet$의 상을 포함하고,
$|E_n| \leq \kappa$인 $\mathit{QC}(\mathcal{O})$의 대상 $E_n$을
나타내는 부분복합체
$\mathcal{E}_n^\bullet \subset \mathcal{K}_n^\bullet$를 얻는다.
이는 원하는 바이다.

\medskip\noindent
주장의 증명. $\mathcal{F}^\bullet$를
$\mathit{QC}(\mathcal{O})$의 대상을 나타내는 복합체라 하고,
$\mathcal{S}^\bullet \subset \mathcal{F}^\bullet$를 크기가
$\leq \kappa$인 부분복합체라 하자. 부분복합체
$$
\mathcal{S}^\bullet = \mathcal{S}_0^\bullet \subset
\mathcal{S}_1^\bullet \subset
\mathcal{S}_2^\bullet \subset \ldots \subset \mathcal{F}^\bullet
$$
를 귀납적으로 구성하되, 그 크기는 $\leq \kappa$이고
$\mathcal{C}$의 모든 사상 $f : U \to V$와 모든
$i \in \mathbf{Z}$에 대해 다음을 만족하게 할 것이다.
\begin{enumerate}
\item 사상
$H^i(\mathcal{S}_n^\bullet(V)
\otimes_{\mathcal{O}(V)}^\mathbf{L} \mathcal{O}(U)) \to
H^i(\mathcal{S}_n^\bullet(U))$
의 핵은
$H^i(\mathcal{S}_{n + 1}^\bullet(V)
\otimes_{\mathcal{O}(V)}^\mathbf{L} \mathcal{O}(U))$에서 영으로 보내진다.
\item 사상
$H^i(\mathcal{S}_n^\bullet(U)) \to H^i(\mathcal{S}_{n + 1}^\bullet(U))$
의 상은
$H^i(\mathcal{S}_{n + 1}^\bullet(V)
\otimes_{\mathcal{O}(V)}^\mathbf{L} \mathcal{O}(U)) \to
H^i(\mathcal{S}_{n + 1}^\bullet(U))$의 상에 포함된다.
\end{enumerate}
이 구성이 끝나면
$\mathcal{H}^\bullet = \bigcup \mathcal{S}_n^\bullet$로 놓을 수 있다.
실제로 유도 텐서곱과 환 위의 가군 복합체의 코호몰로지를 취하는
연산은 여과 여극한과 가환하므로, 조건 (1)과 (2)를 합치면
$$
\mathcal{H}^\bullet(V) \otimes_{\mathcal{O}(V)}^\mathbf{L} \mathcal{O}(U)
\longrightarrow
\mathcal{H}^\bullet(U)
$$
이 $\mathcal{C}$의 모든 $f : U \to V$에 대해 모든 차수의
코호몰로지에서 동형이고, 따라서 $D(\mathcal{O}(U))$에서 동형임이
보장된다. 그러므로 $\mathcal{H}^\bullet$는 원하는 대로
$\mathit{QC}(\mathcal{O})$의 대상을 나타낸다.

\medskip\noindent
$\mathcal{S}_n$이 주어졌을 때 $\mathcal{S}_{n + 1}$의 구성.
$\mathcal{C}$의 모든 사상 $f : U \to V$에 대해 가환 그림
$$
\xymatrix{
\mathcal{S}_n^\bullet(V)
\ar[r] \ar[d] &
\mathcal{S}_n^\bullet(U) \ar[d] \\
\mathcal{F}^\bullet(V)
\ar[r] &
\mathcal{F}^\bullet(U)
}
$$
을 생각한다. 이것은 환 사상 $\mathcal{O}(V) \to \mathcal{O}(U)$에
대해 More on Algebra, Lemma \ref{more-algebra-lemma-enlarge}에 나오는
그림이다. 즉, 아래 행은 $D(\mathcal{O}(U))$에서 동형
$$
\mathcal{F}^\bullet(V) \otimes_{\mathcal{O}(V)}^\mathbf{L} \mathcal{O}(U)
\longrightarrow
\mathcal{F}^\bullet(U)
$$
을 유도한다. 따라서 More on Algebra, Lemma
\ref{more-algebra-lemma-enlarge}에서와 같은 부분복합체
$$
\mathcal{S}_n^\bullet(V) \subset M^\bullet_f \subset \mathcal{F}^\bullet(V)
\quad\text{and}\quad
\mathcal{S}_n^\bullet(U) \subset N^\bullet_f \subset \mathcal{F}^\bullet(U)
$$
를 택할 수 있고, 특히 $|N^i_f|, |M^i_f| \leq \kappa$임을 알 수 있다.
이제 부분집합
$$
\mathcal{S}_n^i(U) \amalg \coprod\nolimits_{f : U \to V} N^i_f
\amalg \coprod\nolimits_{g : W \to U} M^i_g
\subset
\mathcal{F}^i(U)
$$
을 사용하여 Lemma \ref{sites-cohomology-lemma-build-from}을 적용하면 부분복합체
$$
\mathcal{S}_n^\bullet \subset
\mathcal{S}_{n + 1}^\bullet \subset
\mathcal{F}^\bullet
$$
를 얻는데, 이는 그 부분집합들을 포함하며
$|\mathcal{S}_{n + 1}^\bullet| \leq \kappa$를 만족한다.
More on Algebra, Lemma \ref{more-algebra-lemma-enlarge}의 구성에 의해
$\mathcal{S}_n^\bullet(V) \subset M^\bullet_f$와
$\mathcal{S}_n^\bullet(U) \subset N^\bullet_f$에 대해 대응하는 명제가
성립하므로 조건 (1)과 (2)가 성립한다. 이로써 증명이 끝난다.
\end{proof}

\begin{proposition}
\label{sites-cohomology-proposition-cartesian-brown}
$\mathcal{C}$를 무질서 위상을 갖춘 사이트로 보는 범주라 하고,
$\mathcal{O}$를 $\mathcal{C}$ 위의 환의 층이라 하자.
$\mathit{QC}(\mathcal{O})$를
Definition \ref{sites-cohomology-definition-cartesian}에서와 같이 두면 다음이 성립한다.
\begin{enumerate}
\item $\mathit{QC}(\mathcal{O})$는 $D(\mathcal{O})$의 엄밀히 충만하고
포화된 삼각 부분범주이며, 임의의 직합에 의해 보존된다.
\item 직합을 곱으로 보내는 임의의 반변 코호몰로지적 함자
$H : \mathit{QC}(\mathcal{O}) \to \textit{Ab}$
는 표현 가능하다.
\item 직합을 직합으로 보내는 삼각 범주의 임의의 완전 함자
$F : \mathit{QC}(\mathcal{O}) \to \mathcal{D}$는
완전한 오른쪽 수반을 갖는다.
\item 포함 함자 $\mathit{QC}(\mathcal{O}) \to D(\mathcal{O})$는
완전한 오른쪽 수반을 갖는다.
\end{enumerate}
\end{proposition}

\begin{proof}
(1)은 Lemma \ref{sites-cohomology-lemma-cartesian-good-sub}이다.
(2)는 Lemma \ref{sites-cohomology-lemma-cartesian-colimit-kappa}와
Derived Categories, Lemma \ref{derived-lemma-brown-bis}에서 따른다.
(3)은 Lemma \ref{sites-cohomology-lemma-cartesian-colimit-kappa}와
Derived Categories, Proposition \ref{derived-proposition-brown-bis}에서
따른다. (4)는 (3)의 특수한 경우이다.
\end{proof}

\noindent
범주 사이의 함자 $u : \mathcal{C}' \to \mathcal{C}$가 주어졌다고 하자.
$\mathcal{C}$와 $\mathcal{C}'$를 무질서 위상을 갖춘 사이트로 보면
$u$는 연속이자 쌍연속인 함자이다. 따라서 토포스의 사상
$g : \Sh(\mathcal{C}') \to \Sh(\mathcal{C})$를 얻는다.
Sites, Lemma \ref{sites-lemma-cocontinuous-morphism-topoi}를 보라.
또한 $\mathcal{C}$ 위의 환의 층 $\mathcal{O}$와
$\mathcal{C}'$ 위의 환의 층 $\mathcal{O}'$, 그리고 사상
$g^\sharp : g^{-1}\mathcal{O} \to \mathcal{O}'$이
주어졌다고 가정하자. 대응하는 환 달린 토포스의 사상을 간단히
$g : (\Sh(\mathcal{C}'), \mathcal{O}') \to (\Sh(\mathcal{C}), \mathcal{O})$로
나타낸다. Modules on Sites, Section
\ref{sites-modules-section-ringed-topoi}를 보라.

\begin{lemma}
\label{sites-cohomology-lemma-cartesian-functorial}
$g : (\Sh(\mathcal{C}'), \mathcal{O}') \to (\Sh(\mathcal{C}), \mathcal{O})$가
위와 같다고 하자. 그러면 함자
$Lg^* : D(\mathcal{O}) \to D(\mathcal{O}')$는
$\mathit{QC}(\mathcal{O})$를 $\mathit{QC}(\mathcal{O}')$ 안으로 보낸다.
\end{lemma}

\begin{proof}
$U' \in \Ob(\mathcal{C}')$라 하고, $\mathcal{C}$에서의 그 상을
$U = u(U')$라 하자. $pt$는 대상 하나와 사상 하나를 갖는 범주를
나타낸다고 하자. 표시된 환 달린 토포스를 각각
$(\Sh(pt), \mathcal{O}'(U'))$와 $(\Sh(pt), \mathcal{O}(U))$로 나타내자.
물론 이 환 달린 토포스 위의 가군의 유도 범주를
$D(\mathcal{O}'(U'))$ 및 $D(\mathcal{O}(U))$와 동일시한다.
그러면 환 달린 토포스의 가환 그림
$$
\xymatrix{
(\Sh(pt), \mathcal{O}'(U')) \ar[rr]_{U'} \ar[d] & &
(\Sh(\mathcal{C}'), \mathcal{O}') \ar[d]^g \\
(\Sh(pt), \mathcal{O}(U)) \ar[rr]^U & &
(\Sh(\mathcal{C}), \mathcal{O})
}
$$
을 얻는다. 아래 가로 사상을 따른 당김은 $D(\mathcal{O})$의 $K$를
$R\Gamma(U, K)$로 보낸다. 왼쪽 세로 화살표에 의한 당김은 $M$을
$M \otimes_{\mathcal{O}(U)}^\mathbf{L} \mathcal{O}'(U')$로 보낸다.
그림을 어느 방향으로 돌아가도 같은 결과를 얻으므로
(Lemma \ref{sites-cohomology-lemma-derived-pullback-composition}) 다음을 얻는다.
$$
R\Gamma(U', Lg^*K) =
R\Gamma(U, K) \otimes_{\mathcal{O}(U)}^\mathbf{L} \mathcal{O}'(U')
$$
마지막으로 $f' : U' \to V'$를 $\mathcal{C}'$의 사상이라 하고,
$\mathcal{C}$에서의 그 상을
$f = u(f') : U = u(U') \to V = u(V')$로 나타내자.
$K$가 $\mathit{QC}(\mathcal{O})$에 속하면
\begin{align*}
R\Gamma(V', Lg^*K) \otimes_{\mathcal{O}'(V')}^\mathbf{L} \mathcal{O}'(U')
& =
R\Gamma(V, K) \otimes_{\mathcal{O}(V)}^\mathbf{L} \mathcal{O}'(V')
\otimes_{\mathcal{O}'(V')}^\mathbf{L} \mathcal{O}'(U') \\
& =
R\Gamma(V, K) \otimes_{\mathcal{O}(V)}^\mathbf{L} \mathcal{O}'(U') \\
& =
R\Gamma(V, K) \otimes_{\mathcal{O}(V)}^\mathbf{L} \mathcal{O}(U)
\otimes_{\mathcal{O}(U)}^\mathbf{L} \mathcal{O}'(U') \\
& =
R\Gamma(U, K) \otimes_{\mathcal{O}(U)}^\mathbf{L} \mathcal{O}'(U') \\
& =
R\Gamma(U', Lg^*K)
\end{align*}
를 얻으며, 이는 원하는 바이다. 여기서 위의 관찰을 $U'$와 $V'$에
각각 사용했다.
\end{proof}

\begin{lemma}
\label{sites-cohomology-lemma-cartesian-flat-quasi-coherent}
$\mathcal{C}$를 무질서 위상을 갖춘 사이트로 보는 범주라 하고,
$\mathcal{O}$를 $\mathcal{C}$ 위의 환의 층이라 하자.
$\mathcal{C}$의 모든 $U \to V$에 대해 제한 사상
$\mathcal{O}(V) \to \mathcal{O}(U)$가 평탄한 환 사상이라고 가정하자.
그러면 $\mathit{QC}(\mathcal{O})$는 코호몰로지 층이 준연접인
복합체들의 부분범주
$D_\QCoh(\mathcal{O}) \subset D(\mathcal{O})$와 일치한다.
\end{lemma}

\begin{proof}
우리의 가정 아래
$\QCoh(\mathcal{O}) \subset \textit{Mod}(\mathcal{O})$가
약한 Serre 부분범주임을 상기하자. Modules on Sites, Lemma
\ref{sites-modules-lemma-chaotic-flat-quasi-coherent}를 보라.
따라서 $D(\mathcal{O})$의 충만한 부분범주
$$
D_\QCoh(\mathcal{O}) = D_{\QCoh(\mathcal{O})}(\textit{Mod}(\mathcal{O}))
$$
를 취하는 것은 의미가 있다. Derived Categories, Section
\ref{derived-section-triangulated-sub}를 보라. (엄밀히 말해
보조정리의 증명에는 이것이 필요하지 않다.)

\medskip\noindent
$M$을 $\mathit{QC}(\mathcal{O})$의 대상이라 하자.
$\mathcal{C}$의 모든 사상 $U \to V$에 대해 제한 사상
$\mathcal{O}(V) \to \mathcal{O}(U)$가 평탄하므로
\begin{align*}
H^i(M)(U)
& =
H^i(R\Gamma(U, M)) \\
& =
H^i(R\Gamma(V, M) \otimes_{\mathcal{O}(V)}^\mathbf{L} \mathcal{O}(U)) \\
& =
H^i(R\Gamma(V, M)) \otimes_{\mathcal{O}(V)} \mathcal{O}(U) \\
& =
H^i(M)(V) \otimes_{\mathcal{O}(V)} \mathcal{O}(U)
\end{align*}
이다. 따라서 Modules on Sites, Lemma
\ref{sites-modules-lemma-chaotic-quasi-coherent}에 의해
$H^i(M)$은 준연접이다. 위의 첫 번째와 마지막 등식은
$\mathcal{C}$ 위의 위상이 무질서 위상이므로 $\mathcal{C}$의 대상
위에서 단면을 취하는 함자가 완전하다는 사실에서 따른다.

\medskip\noindent
반대로 $M$이 $D_\QCoh(\mathcal{O})$의 대상이면,
Modules on Sites, Lemma
\ref{sites-modules-lemma-chaotic-quasi-coherent}에 의해 사상
$R\Gamma(V, M) \to R\Gamma(U, M)$은 동형
$H^i(M)(U) \to H^i(M)(V) \otimes_{\mathcal{O}(V)} \mathcal{O}(U)$을
유도한다. 따라서 $\mathcal{O}(V) \to \mathcal{O}(U)$의 평탄성에 의해
$R\Gamma(V, K) \otimes_{\mathcal{O}(V)}^\mathbf{L} \mathcal{O}(U)
\to R\Gamma(U, K)$는 $D(\mathcal{O}(U))$에서 동형이고,
$M$이 $\mathit{QC}(\mathcal{O})$에 속한다는 결론을 얻는다.
\end{proof}

\begin{lemma}
\label{sites-cohomology-lemma-cartesion-plus-topology}
$\epsilon : (\mathcal{C}_\tau, \mathcal{O}_\tau) \to
(\mathcal{C}_{\tau'}, \mathcal{O}_{\tau'})$가
Section \ref{sites-cohomology-section-compare}에서와 같다고 하자. 다음을 가정한다.
\begin{enumerate}
\item $\tau'$는 범주 $\mathcal{C}$ 위의 무질서 위상이다.
\item 모든 $U \in \Ob(\mathcal{C})$와 $\mathcal{O}(U)$-가군의 모든
K-평탄 복합체 $M^\bullet$에 대해 사상
$$
M^\bullet \longrightarrow
R\Gamma((\mathcal{C}/U)_\tau,
(M^\bullet \otimes_{\mathcal{O}(U)} \mathcal{O}_U)^\#)
$$
은 유사동형이다(설명은 증명을 보라).
\end{enumerate}
그러면 $\epsilon^*$와 $R\epsilon_*$는 $\mathit{QC}(\mathcal{O})$와
$D(\mathcal{C}_\tau, \mathcal{O}_\tau)$의 다음 충만한 부분범주 사이에
서로 유사역인 동치를 정의한다. 그 부분범주는
$R\epsilon_*K$가 $\mathit{QC}(\mathcal{O})$에 속하는 대상 $K$들로
이루어진다.\footnote{이는 $\mathcal{C}$의 모든 $U \to V$에 대해
$R\Gamma(V, K) \otimes_{\mathcal{O}(V)}^\mathbf{L} \mathcal{O}(U)
\to R\Gamma(U, K)$가 동형이라는 뜻이다.}
\end{lemma}

\begin{proof}
Section \ref{sites-cohomology-section-compare}에서 한 관찰들을 더 언급하지 않고 사용한다.
$R\epsilon_*$가 충실충만하고
$\epsilon^* \circ R\epsilon_* = \text{id}$이므로, 보조정리를 증명하려면
$\mathit{QC}(\mathcal{O})$의 $M$에 대해
$R\epsilon_*(\epsilon^*M) = M$임을 보이면 충분하다.
조건 (2)는 바로 이를 보이는 데 필요한 조건이다.
즉, $M$을 나타내며 각 항이 평탄한 $\mathcal{O}$-가군의 K-평탄 복합체
$\mathcal{M}^\bullet$를 택한다. 그러면 $\epsilon^*M$은
$\mathcal{M}^\bullet$의 $\tau$-층화
$(\mathcal{M}^\bullet)^\#$로 나타난다.
$U \in \Ob(\mathcal{C})$라 하자. Leray에 의해
$$
R\Gamma(U, R\epsilon_*(\epsilon^*M)) =
R\Gamma((\mathcal{C}/U)_\tau, (\mathcal{M}^\bullet)^\#|_{\mathcal{C}/U}) =
R\Gamma((\mathcal{C}/U)_\tau, (\mathcal{M}^\bullet|_{\mathcal{C}/U})^\#)
$$
를 얻는다. 마지막 등식은 층화가 $\mathcal{C}/U$로의 제한과
가환하기 때문이다. 평소처럼 $\mathcal{O}$의 $\mathcal{C}/U$로의
제한을 $\mathcal{O}_U$로 나타내자. 복합체의 사상
$$
\mathcal{M}^\bullet(U) \otimes_{\mathcal{O}(U)} \mathcal{O}_U
\longrightarrow
\mathcal{M}^\bullet|_{\mathcal{C}/U}
$$
을 생각하자. 이는 ($\tau'$-위상에서) $\mathcal{O}_U$-가군의
복합체의 사상이다. $\mathcal{M}^\bullet$의 선택에 의해
$\mathcal{M}^\bullet(U)$는 $\mathcal{O}(U)$-가군의 K-평탄 복합체이다.
Lemma \ref{sites-cohomology-lemma-pullback-K-flat}와, $U$의 $\mathcal{C}$로의 포함이
환 달린 토포스의 사상
$(\Sh(pt), \mathcal{O}(U)) \to (\Sh(\mathcal{C}_{\tau'}), \mathcal{O})$을
정의한다는 사실을 사용하라. $M$이 $\mathit{QC}(\mathcal{O})$에
속하므로 표시된 화살표가 유사동형이라는 결론을 얻는다.
층화는 완전하므로 층화한 뒤에도 마찬가지이다. 따라서
$$
R\Gamma(U, R\epsilon_*(\epsilon^*M)) =
R\Gamma((\mathcal{C}/U)_\tau,
(M^\bullet \otimes_{\mathcal{O}(U)} \mathcal{O}_U)^\#)
$$
이고, 가정 (2)에 의해 이는 원하는 대로
$R\Gamma(U, M) = \mathcal{M}^\bullet(U)$와 같다.
\end{proof}

\begin{lemma}
\label{sites-cohomology-lemma-QC-hom-out-of-plus-topology}
표기와 가정은 Lemma \ref{sites-cohomology-lemma-cartesion-plus-topology}에서와 같다.
$K$가 $\mathit{QC}(\mathcal{O})$의 대상이고
$M$은 $D(\mathcal{O}_\tau)$의 임의의 대상이라고 가정하자.
$\mathcal{C}$의 모든 대상 $U$에 대해
$$
\Hom_{D((\mathcal{O}_U)_\tau)}(\epsilon^*K|_U, M|_U) =
R\Hom_{\mathcal{O}(U)}(R\Gamma(U, K), R\Gamma(U, M))
$$
이다.
\end{lemma}

\begin{proof}
수반성에 의해
$$
\Hom_{D((\mathcal{O}_U)_\tau)}(\epsilon^*K|_U, M|_U) =
\Hom_{D((\mathcal{O}_U)_{\tau'})}(K|_U, R\epsilon_*M|_U)
$$
이다. 따라서 Lemma \ref{sites-cohomology-lemma-QC-hom-out-of}와 Leray에 의해 성립하는
$$
R\Gamma(U, M) = R\Gamma(U, R\epsilon_*M)
$$
에서 결론이 따른다.
\end{proof}







\section{강한 완전 복합체}
\label{sites-cohomology-section-strictly-perfect}

\noindent
이 절은 Cohomology, Section
\ref{cohomology-section-strictly-perfect}에 대응한다.

\begin{definition}
\label{sites-cohomology-definition-strictly-perfect}
$(\mathcal{C}, \mathcal{O})$를 환 달린 사이트라 하자.
$\mathcal{E}^\bullet$를 $\mathcal{O}$-가군의 복합체라 하자.
유한 개를 제외한 모든 $i$에 대해 $\mathcal{E}^i$가 영이고,
모든 $i$에 대해 $\mathcal{E}^i$가 유한 자유
$\mathcal{O}$-가군의 직합인자이면
$\mathcal{E}^\bullet$가 {\it 강한 완전}이라고 한다.
\end{definition}

\noindent
$U$를 $\mathcal{C}$의 대상이라 하자. ``$\mathcal{E}^\bullet$를
$\mathcal{O}_U$-가군의 강한 완전 복합체라 하자''라고 자주 말할
것인데, 이는 $\mathcal{E}^\bullet$가 환 달린 사이트
$(\mathcal{C}/U, \mathcal{O}_U)$ 위의 가군의 강한 완전
복합체라는 뜻이다. Modules on Sites, Definition
\ref{sites-modules-definition-localize-ringed-site}를 보라.

\begin{lemma}
\label{sites-cohomology-lemma-cone}
강한 완전 복합체의 사상 위의 원뿔은 강한 완전이다.
\end{lemma}

\begin{proof}
이는 정의에서 곧바로 따른다.
\end{proof}

\begin{lemma}
\label{sites-cohomology-lemma-tensor}
두 강한 완전 복합체의 텐서곱에 딸린 전체 복합체는
강한 완전이다.
\end{lemma}

\begin{proof}
생략한다.
\end{proof}

\begin{lemma}
\label{sites-cohomology-lemma-strictly-perfect-pullback}
$(f, f^\sharp) : (\mathcal{C}, \mathcal{O}_\mathcal{C}) \to
(\mathcal{D}, \mathcal{O}_\mathcal{D})$
를 환 달린 토포스의 사상이라 하자. $\mathcal{F}^\bullet$가
$\mathcal{O}_\mathcal{D}$-가군의 강한 완전 복합체이면,
$f^*\mathcal{F}^\bullet$는 $\mathcal{O}_\mathcal{C}$-가군의
강한 완전 복합체이다.
\end{lemma}

\begin{proof}
유한 자유 가군의 당김이 유한 자유임을 Modules on Sites, Lemma
\ref{sites-modules-lemma-global-pullback}에서 보았다.
함자 $f^*$는 가법 함자이므로 직합인자를 보존한다. 보조정리가 따른다.
\end{proof}

\begin{lemma}
\label{sites-cohomology-lemma-local-lift-map}
$(\mathcal{C}, \mathcal{O})$를 환 달린 사이트라 하고,
$U$를 $\mathcal{C}$의 대상이라 하자.
$\mathcal{O}_U$-가군의 실선 그림
$$
\xymatrix{
\mathcal{E} \ar@{..>}[dr] \ar[r] & \mathcal{F} \\
& \mathcal{G} \ar[u]_p
}
$$
이 주어졌다고 하자. 여기서 $\mathcal{E}$는 유한 자유
$\mathcal{O}_U$-가군의 직합인자이고 $p$는 전사이다.
그러면 덮개 $\{U_i \to U\}$가 존재하여, 각 $U_i$ 위에서
그림을 가환하게 하는 점선 화살표가 존재한다.
\end{lemma}

\begin{proof}
어떤 $n$에 대해 $\mathcal{E} = \mathcal{O}_U^{\oplus n}$이라고
가정해도 된다. 이 경우 점선 화살표를 찾는 것은 기저 원소들의
$\Gamma(U, \mathcal{F})$에서의 상들을 올리는 것과 동치이다.
층의 전사에 대한 특징짓기에 의해 이는 국소적으로 가능하다
(Sites, Section \ref{sites-section-sheaves-injective}).
\end{proof}

\begin{lemma}
\label{sites-cohomology-lemma-local-homotopy}
$(\mathcal{C}, \mathcal{O})$를 환 달린 사이트라 하고,
$U$를 $\mathcal{C}$의 대상이라 하자.
\begin{enumerate}
\item $\alpha : \mathcal{E}^\bullet \to \mathcal{F}^\bullet$를
$\mathcal{O}_U$-가군의 복합체의 사상이라 하자.
$\mathcal{E}^\bullet$가 강한 완전이고 $\mathcal{F}^\bullet$가
비순환적이면, 각 $\alpha|_{U_i}$가 영과 호모토픽인 덮개
$\{U_i \to U\}$가 존재한다.
\item $\alpha : \mathcal{E}^\bullet \to \mathcal{F}^\bullet$를
$\mathcal{O}_U$-가군의 복합체의 사상이라 하자.
$\mathcal{E}^\bullet$가 강한 완전이고, $i < a$에 대해
$\mathcal{E}^i = 0$이며, $i \geq a$에 대해
$H^i(\mathcal{F}^\bullet) = 0$이면, 각 $\alpha|_{U_i}$가
영과 호모토픽인 덮개 $\{U_i \to U\}$가 존재한다.
\end{enumerate}
\end{lemma}

\begin{proof}
첫 번째 명제는 두 번째 명제에서 따르므로 (2)만 증명한다.
복합체 $\mathcal{E}^\bullet$의 길이에 관한 귀납법으로 증명하겠다.
어떤 유한 자유 $\mathcal{O}$-가군의 직합인자 $\mathcal{E}$와
정수 $n \geq a$에 대해
$\mathcal{E}^\bullet \cong \mathcal{E}[-n]$이면,
Lemma \ref{sites-cohomology-lemma-local-lift-map}과 사상
$\mathcal{F}^{n - 1} \to \Ker(\mathcal{F}^n \to \mathcal{F}^{n + 1})$
이 $H^n(\mathcal{F}^\bullet)$의 가정된 소멸에 의해 전사라는
사실에서 결론이 따른다.
$i \in [a, b]$가 아닌 경우 $\mathcal{E}^i$가 영이면,
복합체의 분할 완전열
$$
0 \to \mathcal{E}^b[-b] \to \mathcal{E}^\bullet \to
\sigma_{\leq b - 1}\mathcal{E}^\bullet \to 0
$$
을 얻으며, 이는 $K(\mathcal{O}_U)$에서 구별 삼각형을 정한다.
따라서 삼각 범주의 공리로부터 완전열
$$
\Hom_{K(\mathcal{O}_U)}(
\sigma_{\leq b - 1}\mathcal{E}^\bullet, \mathcal{F}^\bullet)
\to
\Hom_{K(\mathcal{O}_U)}(\mathcal{E}^\bullet, \mathcal{F}^\bullet)
\to
\Hom_{K(\mathcal{O}_U)}(\mathcal{E}^b[-b], \mathcal{F}^\bullet)
$$
을 얻는다. 위의 결과에 의해 합성
$\mathcal{E}^b[-b] \to \mathcal{F}^\bullet$는 $U$의 어떤 덮개의
각 원소 위에서 영과 호모토픽이다. 따라서 우리의 사상이 위에 표시한
완전열의 왼쪽 항의 원소에서 온다고 가정할 수 있다.
귀납 가정에 의해 이 원소는 $U$의 어떤 덮개의 각 원소 위에서 영이다.
\end{proof}

\begin{lemma}
\label{sites-cohomology-lemma-lift-through-quasi-isomorphism}
$(\mathcal{C}, \mathcal{O})$를 환 달린 사이트라 하고,
$U$를 $\mathcal{C}$의 대상이라 하자.
$\mathcal{O}_U$-가군 복합체의 실선 그림
$$
\xymatrix{
\mathcal{E}^\bullet \ar@{..>}[dr] \ar[r]_\alpha & \mathcal{F}^\bullet \\
& \mathcal{G}^\bullet \ar[u]_f
}
$$
이 주어졌다고 하자. $\mathcal{E}^\bullet$가 강한 완전이고,
$j < a$에 대해 $\mathcal{E}^j = 0$이며, $H^j(f)$는 $j > a$에
대해 동형이고 $j = a$에 대해 전사라고 하자.
그러면 덮개 $\{U_i \to U\}$와 각 $i$에 대해 $U_i$ 위에서
그림을 호모토피까지 가환하게 하는 점선 화살표가 존재한다.
\end{lemma}

\begin{proof}
$f$에 대한 우리의 가정에 의해 원뿔 $C(f)^\bullet$의
$\geq a$인 차수의 코호몰로지 층은 소멸한다.
따라서 Lemma \ref{sites-cohomology-lemma-local-homotopy}에 의해 덮개
$\{U_i \to U\}$가 존재하여 합성
$\mathcal{E}^\bullet \to \mathcal{F}^\bullet \to C(f)^\bullet$
이 $U_i$ 위에서 영과 호모토픽이다. 또한
$$
\mathcal{G}^\bullet \to \mathcal{F}^\bullet \to C(f)^\bullet \to
\mathcal{G}^\bullet[1]
$$
가 $K(\mathcal{O}_{U_i})$에서 구별 삼각형으로 제한되므로,
$\alpha|_{U_i}$를 호모토피까지 사상
$\alpha_i : \mathcal{E}^\bullet|_{U_i} \to \mathcal{G}^\bullet|_{U_i}$
으로 올릴 수 있다. 이는 원하는 바이다.
\end{proof}

\begin{lemma}
\label{sites-cohomology-lemma-local-actual}
$(\mathcal{C}, \mathcal{O})$를 환 달린 사이트라 하고,
$U$를 $\mathcal{C}$의 대상이라 하자.
$\mathcal{E}^\bullet$, $\mathcal{F}^\bullet$를
$\mathcal{O}_U$-가군의 복합체라 하고,
$\mathcal{E}^\bullet$가 강한 완전이라고 하자.
\begin{enumerate}
\item 임의의 원소
$\alpha \in \Hom_{D(\mathcal{O}_U)}(\mathcal{E}^\bullet, \mathcal{F}^\bullet)$
에 대해 덮개 $\{U_i \to U\}$가 존재하여
$\alpha|_{U_i}$는 복합체의 사상
$\alpha_i : \mathcal{E}^\bullet|_{U_i} \to \mathcal{F}^\bullet|_{U_i}$으로
주어진다.
\item 복합체의 사상
$\alpha : \mathcal{E}^\bullet \to \mathcal{F}^\bullet$의 군
$\Hom_{D(\mathcal{O}_U)}(\mathcal{E}^\bullet, \mathcal{F}^\bullet)$
에서의 상이 영이라고 하자. 그러면 덮개 $\{U_i \to U\}$가 존재하여
$\alpha|_{U_i}$는 영과 호모토픽이다.
\end{enumerate}
\end{lemma}

\begin{proof}
(1)의 증명. 유도 범주의 구성에 의해 유사동형
$f : \mathcal{F}^\bullet \to \mathcal{G}^\bullet$와 복합체의 사상
$\beta : \mathcal{E}^\bullet \to \mathcal{G}^\bullet$를 찾아
$\alpha = f^{-1}\beta$가 되게 할 수 있다. 따라서 결론은
Lemma \ref{sites-cohomology-lemma-lift-through-quasi-isomorphism}에서 따른다.
(2)의 증명은 생략한다.
\end{proof}

\begin{lemma}
\label{sites-cohomology-lemma-Rhom-strictly-perfect}
$(\mathcal{C}, \mathcal{O})$를 환 달린 사이트라 하자.
$\mathcal{E}^\bullet$, $\mathcal{F}^\bullet$를
$\mathcal{O}$-가군의 복합체라 하고,
$\mathcal{E}^\bullet$가 강한 완전이라고 하자.
그러면 내부 Hom
$R\SheafHom(\mathcal{E}^\bullet, \mathcal{F}^\bullet)$은 항이
$$
\mathcal{H}^n =
\bigoplus\nolimits_{n = p + q}
\SheafHom_\mathcal{O}(\mathcal{E}^{-q}, \mathcal{F}^p)
$$
이고 미분이 Section \ref{sites-cohomology-section-internal-hom}에서 기술된 복합체
$\mathcal{H}^\bullet$로 나타난다.
\end{lemma}

\begin{proof}
K-단사 복합체로 가는 유사동형
$\mathcal{F}^\bullet \to \mathcal{I}^\bullet$를 택하자.
항이
$$
(\mathcal{H}')^n =
\prod\nolimits_{n = p + q}
\SheafHom_\mathcal{O}(\mathcal{L}^{-q}, \mathcal{I}^p)
$$
인 복합체를 $(\mathcal{H}')^\bullet$라 하자. Section
\ref{sites-cohomology-section-internal-hom}의 구성에 의해 이는
$R\SheafHom(\mathcal{E}^\bullet, \mathcal{F}^\bullet)$을 나타낸다.
사상
$$
\mathcal{H}^\bullet \longrightarrow (\mathcal{H}')^\bullet
$$
이 유사동형임을 보이면 충분하다. $\mathcal{C}$의 대상 $U$가
주어지면 직접 확인하여
$$
H^0(\mathcal{H}^\bullet(U)) =
\Hom_{K(\mathcal{O}_U)}(\mathcal{E}^\bullet|_U, \mathcal{I}^\bullet|_U)
\to
H^0((\mathcal{H}')^\bullet(U)) =
\Hom_{D(\mathcal{O}_U)}(\mathcal{E}^\bullet|_U, \mathcal{I}^\bullet|_U)
$$
를 얻는다. Lemma \ref{sites-cohomology-lemma-local-actual}에 의해
$U \mapsto H^0(\mathcal{H}^\bullet(U))$
의 층화는
$U \mapsto H^0((\mathcal{H}')^\bullet(U))$의 층화와 같다.
다른 코호몰로지 층에도 비슷한 논증을 적용할 수 있다.
따라서 $\mathcal{H}^\bullet$는 $(\mathcal{H}')^\bullet$와
유사동형이며, 이로써 보조정리가 증명된다.
\end{proof}

\begin{lemma}
\label{sites-cohomology-lemma-Rhom-complex-of-direct-summands-finite-free}
$(\mathcal{C}, \mathcal{O})$를 환 달린 사이트라 하자.
$\mathcal{E}^\bullet$, $\mathcal{F}^\bullet$를 다음을 만족하는
$\mathcal{O}$-가군의 복합체라 하자.
\begin{enumerate}
\item $n \ll 0$에 대해 $\mathcal{F}^n = 0$이다.
\item $n \gg 0$에 대해 $\mathcal{E}^n = 0$이다.
\item $\mathcal{E}^n$은 유한 자유 $\mathcal{O}$-가군의
직합인자와 동형이다.
\end{enumerate}
그러면 내부 Hom
$R\SheafHom(\mathcal{E}^\bullet, \mathcal{F}^\bullet)$은 항이
$$
\mathcal{H}^n =
\bigoplus\nolimits_{n = p + q}
\SheafHom_\mathcal{O}(\mathcal{E}^{-q}, \mathcal{F}^p)
$$
이고 미분이 Section \ref{sites-cohomology-section-internal-hom}에서 기술된 복합체
$\mathcal{H}^\bullet$로 나타난다.
\end{lemma}

\begin{proof}
$\mathcal{I}^\bullet$가 아래로 유계인 단사 대상들의 복합체인
유사동형 $\mathcal{F}^\bullet \to \mathcal{I}^\bullet$를 택하자.
$\mathcal{I}^\bullet$는 K-단사임을 주목하라
(Derived Categories, Lemma
\ref{derived-lemma-bounded-below-injectives-K-injective}).
따라서 Section \ref{sites-cohomology-section-internal-hom}의 구성에 의해
$R\SheafHom(\mathcal{E}^\bullet, \mathcal{F}^\bullet)$은
항이
$$
(\mathcal{H}')^n =
\prod\nolimits_{n = p + q}
\SheafHom_\mathcal{O}(\mathcal{E}^{-q}, \mathcal{I}^p) =
\bigoplus\nolimits_{n = p + q}
\SheafHom_\mathcal{O}(\mathcal{E}^{-q}, \mathcal{I}^p)
$$
인 복합체 $(\mathcal{H}')^\bullet$로 나타난다
(영이 아닌 항이 유한 개뿐이므로 등식이 성립한다).
$\mathcal{H}^\bullet$는 항이
$\SheafHom_\mathcal{O}(\mathcal{E}^{-q}, \mathcal{F}^p)$
인 이중복합체에 딸린 전체 복합체이며,
$(\mathcal{H}')^\bullet$도 마찬가지임을 주목하라.
자연스러운 사상 $(\mathcal{H}')^\bullet \to \mathcal{H}^\bullet$은
이중복합체의 사상에서 온다.
따라서 이 사상이 유사동형임을 보이기 위해 이중복합체의 스펙트럼 열
(Homology, Lemma \ref{homology-lemma-first-quadrant-ss})
$$
{}'E_1^{p, q} =
H^p(\SheafHom_\mathcal{O}(\mathcal{E}^{-q}, \mathcal{F}^\bullet))
$$
을 사용할 수 있다. 이는 $H^{p + q}(\mathcal{H}^\bullet)$로 수렴하며
$(\mathcal{H}')^\bullet$에도 마찬가지이다. 보조정리의 증명을 끝내려면
$\mathcal{E}$가 유한 자유 $\mathcal{O}$-가군의 직합인자일 때
$\mathcal{F}^\bullet \to \mathcal{I}^\bullet$가 코호몰로지 층의 동형
$$
H^p(\SheafHom_\mathcal{O}(\mathcal{E}, \mathcal{F}^\bullet))
\longrightarrow
H^p(\SheafHom_\mathcal{O}(\mathcal{E}, \mathcal{I}^\bullet))
$$
을 유도함을 보이면 충분하다. $\mathcal{E}$가 유한 자유일 때 이는
명백하므로 결론이 따른다.
\end{proof}









\section{유사코히런트 가군}
\label{sites-cohomology-section-pseudo-coherent}

\noindent
이 절에서는 유사코히런트 복합체를 논의한다.

\begin{definition}
\label{sites-cohomology-definition-pseudo-coherent}
$(\mathcal{C}, \mathcal{O})$를 환 달린 사이트라 하자.
$\mathcal{E}^\bullet$를 $\mathcal{O}$-가군의 복합체라 하고,
$m \in \mathbf{Z}$라 하자.
\begin{enumerate}
\item $\mathcal{C}$의 모든 대상 $U$에 대해 덮개
$\{U_i \to U\}$와 각 $i$에 대해 복합체의 사상
$\alpha_i : \mathcal{E}_i^\bullet \to \mathcal{E}^\bullet|_{U_i}$
이 존재하되, $\mathcal{E}_i^\bullet$가
$\mathcal{O}_{U_i}$-가군의 강한 완전 복합체이고,
$H^j(\alpha_i)$가 $j > m$에 대해 동형이며
$H^m(\alpha_i)$가 전사이면, $\mathcal{E}^\bullet$가
{\it $m$-유사코히런트}라고 한다.
\item 모든 $m$에 대해 $m$-유사코히런트이면
$\mathcal{E}^\bullet$가 {\it 유사코히런트}라고 한다.
\item $D(\mathcal{O})$의 대상 $E$가
$m$-유사코히런트(각각 유사코히런트)인
$\mathcal{O}$-가군의 복합체로 나타내어질 수 있는 경우, 그리고
그 경우에만 이를 {\it $m$-유사코히런트}
(각각 {\it 유사코히런트})라고 한다.
\end{enumerate}
\end{definition}

\noindent
$\mathcal{C}$가 준콤팩트인 끝 대상 $X$를 가지면
(예를 들어 $X$의 모든 덮개가 유한 덮개로 세분될 수 있으면),
$D(\mathcal{O})$의 $m$-유사코히런트 대상은
$D^-(\mathcal{O})$에 속한다. 그러나 일반적으로는 그럴 필요가 없다.

\begin{lemma}
\label{sites-cohomology-lemma-pseudo-coherent-independent-representative}
$(\mathcal{C}, \mathcal{O})$를 환 달린 사이트라 하고,
$E$를 $D(\mathcal{O})$의 대상이라 하자.
\begin{enumerate}
\item $\mathcal{C}$가 끝 대상 $X$를 가지고, 덮개
$\{U_i \to X\}$, $\mathcal{O}_{U_i}$-가군의 강한 완전 복합체
$\mathcal{E}_i^\bullet$, 그리고 $j > m$에 대해
$H^j(\alpha_i)$가 동형이고 $H^m(\alpha_i)$가 전사인
$D(\mathcal{O}_{U_i})$의 사상
$\alpha_i : \mathcal{E}_i^\bullet \to E|_{U_i}$가 존재하면,
$E$는 $m$-유사코히런트이다.
\item $E$가 $m$-유사코히런트이면 $E$를 나타내는
$\mathcal{O}$-가군의 임의의 복합체가 $m$-유사코히런트이다.
\item $\mathcal{C}$의 모든 대상 $U$에 대해 덮개
$\{U_i \to U\}$가 존재하여 $E|_{U_i}$가
$m$-유사코히런트이면, $E$는 $m$-유사코히런트이다.
\end{enumerate}
\end{lemma}

\begin{proof}
임의의 $E$를 나타내는 복합체를 $\mathcal{F}^\bullet$라 하고,
$X$, $\{U_i \to X\}$, $\alpha_i : \mathcal{E}_i \to E|_{U_i}$가
(1)에서와 같다고 하자. $\mathcal{F}^\bullet$가 복합체로서
$m$-유사코히런트임을 보이겠다. 그러면 $\mathcal{C}$가 끝 대상을
갖는 경우 (1)과 (2)가 증명된다. Lemma \ref{sites-cohomology-lemma-local-actual}에 의해
덮개 $\{U_i \to X\}$를 세분한 뒤 사상 $\alpha_i$들을
복합체의 사상
$\alpha_i : \mathcal{E}_i^\bullet \to \mathcal{F}^\bullet|_{U_i}$으로
나타낼 수 있다. 가정에 의해 $H^j(\alpha_i)$는 $j > m$에
대해 동형이고 $H^m(\alpha_i)$는 전사이므로
$\mathcal{F}^\bullet$는 $m$-유사코히런트이다.

\medskip\noindent
(2)의 증명. 위에서 $\mathcal{C}$의 모든 대상 $U$에 대해
$\mathcal{F}^\bullet|_U$가 $\mathcal{O}_U$-가군의 복합체로서
$m$-유사코히런트임을 알 수 있다. 정의의 형식적 귀결로
$\mathcal{F}^\bullet$는 $m$-유사코히런트이다.

\medskip\noindent
(3)의 증명. 정의와 Sites, Definition
\ref{sites-definition-site}의 (2)에서 따른다.
\end{proof}

\begin{lemma}
\label{sites-cohomology-lemma-pseudo-coherent-pullback}
$(f, f^\sharp) : (\mathcal{C}, \mathcal{O}_\mathcal{C}) \to
(\mathcal{D}, \mathcal{O}_\mathcal{D})$
를 환 달린 사이트의 사상이라 하자. $E$를
$D(\mathcal{O}_\mathcal{C})$의 대상이라 하자.
$E$가 $m$-유사코히런트이면 $Lf^*E$도 $m$-유사코히런트이다.
\end{lemma}

\begin{proof}
$f$가 함자 $u : \mathcal{D} \to \mathcal{C}$로 주어진다고 하자.
$U$를 $\mathcal{C}$의 대상이라 하자. Sites, Lemma
\ref{sites-lemma-morphism-of-sites-covering}에 의해 덮개
$\{U_i \to U\}$와 각 $i$에 대해 $\mathcal{D}$의 어떤 대상 $V_i$로부터
사상 $U_i \to u(V_i)$를 찾을 수 있다.
Lemma \ref{sites-cohomology-lemma-pseudo-coherent-independent-representative}에 의해
$Lf^*E|_{U_i}$가 $m$-유사코히런트임을 보이면 충분하다.
이를 위해서는 $Lf^*E|_{u(V_i)}$가 $m$-유사코히런트임을 보이면
충분하다. 실제로 $Lf^*E|_{U_i}$는
$Lf^*E|_{u(V_i)}$를 $\mathcal{C}/U_i$로 제한한 것이다
(Modules on Sites, Lemma \ref{sites-modules-lemma-relocalize}).
Modules on Sites, Lemma
\ref{sites-modules-lemma-localize-morphism-ringed-sites}의 가환 그림에
의해, 환 달린 사이트의 사상
$(\mathcal{C}/u(V_i), \mathcal{O}_{u(V_i)}) \to
(\mathcal{D}/V_i, \mathcal{O}_{V_i})$에 대해
보조정리를 증명하면 충분하다. 따라서 $\mathcal{D}$가
끝 대상 $Y$를 가지고 $X = u(Y)$가 $\mathcal{C}$의 끝 대상이라고
가정해도 된다.

\medskip\noindent
$\{V_i \to Y\}$를 다음과 같은 덮개라 하자. 각 $i$에 대해
$\mathcal{O}_{V_i}$-가군의 강한 완전 복합체
$\mathcal{F}_i^\bullet$와 $D(\mathcal{O}_{V_i})$의 사상
$\alpha_i : \mathcal{F}_i^\bullet \to E|_{V_i}$가 존재하고,
$H^j(\alpha_i)$는 $j > m$에 대해 동형이며
$H^m(\alpha_i)$는 전사이다. 위와 같이 논증하면 사상
$(\mathcal{C}/u(V_i), \mathcal{O}_{u(V_i)}) \to
(\mathcal{D}/V_i, \mathcal{O}_{V_i})$에 대해 결론을 증명하면 충분하다.
따라서 $\mathcal{O}_\mathcal{D}$-가군의 강한 완전 복합체
$\mathcal{F}^\bullet$와 $D(\mathcal{O}_\mathcal{D})$의 사상
$\alpha : \mathcal{F}^\bullet \to E$가 존재하고,
$H^j(\alpha)$가 $j > m$에 대해 동형이며 $H^m(\alpha)$가
전사라고 가정해도 된다. 이 경우 구별 삼각형
$$
\mathcal{F}^\bullet \to E \to C \to \mathcal{F}^\bullet[1]
$$
을 택하자. $\alpha$에 대한 가정은 모든 $j \geq m$에 대해
코호몰로지 층 $H^j(C)$가 영이라는 것과 정확히 같다.
$Lf^*$를 적용하면 구별 삼각형
$$
Lf^*\mathcal{F}^\bullet \to Lf^*E \to Lf^*C \to Lf^*\mathcal{F}^\bullet[1]
$$
을 얻는다. $Lf^*$를 왼쪽 유도 함자로 구성한 방식에 의해
$j \geq m$에 대해 $H^j(Lf^*C) = 0$이다
(Derived Categories, Lemma
\ref{derived-lemma-negative-vanishing}의 쌍대).
따라서 $H^j(Lf^*\alpha)$는 $j > m$에 대해 동형이고,
$H^m(Lf^*\alpha)$는 전사이다. 한편 $\mathcal{F}^\bullet$는
평탄한 $\mathcal{O}_\mathcal{D}$-가군의 위로 유계인 복합체이므로
$Lf^*\mathcal{F}^\bullet = f^*\mathcal{F}^\bullet$이다.
Lemma \ref{sites-cohomology-lemma-strictly-perfect-pullback}을 적용하면 결론이 따른다.
\end{proof}

\begin{lemma}
\label{sites-cohomology-lemma-cone-pseudo-coherent}
$(\mathcal{C}, \mathcal{O})$를 환 달린 사이트라 하고
$m \in \mathbf{Z}$라 하자.
$(K, L, M, f, g, h)$를 $D(\mathcal{O})$의 구별 삼각형이라 하자.
\begin{enumerate}
\item $K$가 $(m + 1)$-유사코히런트이고 $L$이
$m$-유사코히런트이면 $M$은 $m$-유사코히런트이다.
\item $K$와 $M$이 $m$-유사코히런트이면 $L$도
$m$-유사코히런트이다.
\item $L$이 $(m + 1)$-유사코히런트이고 $M$이
$m$-유사코히런트이면 $K$는 $(m + 1)$-유사코히런트이다.
\end{enumerate}
\end{lemma}

\begin{proof}
(1)의 증명. $U$를 $\mathcal{C}$의 대상이라 하자.
덮개 $\{U_i \to U\}$와 $D(\mathcal{O}_{U_i})$의 사상
$\alpha_i : \mathcal{K}_i^\bullet \to K|_{U_i}$를 택하되,
$\mathcal{K}_i^\bullet$는 강한 완전이고
$H^j(\alpha_i)$는 $j > m + 1$에 대해 동형이며
$j = m + 1$에 대해 전사이게 한다.
$\mathcal{K}_i^\bullet$를
$\sigma_{\geq m + 1}\mathcal{K}_i^\bullet$
로 바꿀 수 있으므로, $j < m + 1$에 대해
$\mathcal{K}_i^j = 0$이라고 가정해도 된다.
덮개를 세분한 뒤 $D(\mathcal{O}_{U_i})$의 사상
$\beta_i : \mathcal{L}_i^\bullet \to L|_{U_i}$를 택하되,
$\mathcal{L}_i^\bullet$는 강한 완전이고,
$H^j(\beta)$는 $j > m$에 대해 동형이며
$j = m$에 대해 전사이게 할 수 있다.
Lemma \ref{sites-cohomology-lemma-lift-through-quasi-isomorphism}에 의해,
덮개를 세분한 뒤 복합체의 사상
$\gamma_i : \mathcal{K}^\bullet \to \mathcal{L}^\bullet$
을 찾아 그림
$$
\xymatrix{
K|_{U_i} \ar[r] & L|_{U_i} \\
\mathcal{K}_i^\bullet \ar[u]^{\alpha_i} \ar[r]^{\gamma_i} &
\mathcal{L}_i^\bullet \ar[u]_{\beta_i}
}
$$
들이 $D(\mathcal{O}_{U_i})$에서 가환하게 할 수 있다
(이를 위해서는 사상 $\alpha_i$, $\beta_i$,
$K|_{U_i} \to L|_{U_i}$를 실제 복합체의 사상으로 나타내야 한다.
일부 세부 사항은 생략한다).
원뿔 $C(\gamma_i)^\bullet$는 강한 완전이다
(Lemma \ref{sites-cohomology-lemma-cone}). 그림의 가환성에 의해 구별 삼각형의 사상
$$
(\mathcal{K}_i^\bullet, \mathcal{L}_i^\bullet, C(\gamma_i)^\bullet)
\longrightarrow
(K|_{U_i}, L|_{U_i}, M|_{U_i}).
$$
이 존재한다. 유도된 긴 완전 코호몰로지 열의 사상과
Homology, Lemmas \ref{homology-lemma-four-lemma} 및
\ref{homology-lemma-five-lemma}에 의해
$C(\gamma_i)^\bullet \to M|_{U_i}$는 $> m$인 차수의
코호몰로지에서 동형을, 차수 $m$에서 전사를 유도한다.
따라서 Lemma \ref{sites-cohomology-lemma-pseudo-coherent-independent-representative}에
의해 $M$은 $m$-유사코히런트이다.

\medskip\noindent
(2)와 (3)은 구별 삼각형을 회전시켜 (1)에서 얻는다.
\end{proof}

\begin{lemma}
\label{sites-cohomology-lemma-tensor-pseudo-coherent}
$(\mathcal{C}, \mathcal{O})$를 환 달린 사이트라 하고,
$K, L$을 $D(\mathcal{O})$의 대상이라 하자.
\begin{enumerate}
\item $K$가 $n$-유사코히런트이고 $i > a$에 대해
$H^i(K) = 0$이며, $L$이 $m$-유사코히런트이고 $j > b$에 대해
$H^j(L) = 0$이면, $t = \max(m + a, n + b)$에 대해
$K \otimes_\mathcal{O}^\mathbf{L} L$은 $t$-유사코히런트이다.
\item $K$와 $L$이 유사코히런트이면
$K \otimes_\mathcal{O}^\mathbf{L} L$도 유사코히런트이다.
\end{enumerate}
\end{lemma}

\begin{proof}
(1)의 증명. $U$를 $\mathcal{C}$의 대상이라 하자.
$U$를 어떤 덮개의 원소들로 바꾸고 $\mathcal{C}$를 국소화
$\mathcal{C}/U$로 바꾸어, 강한 완전 복합체
$\mathcal{K}^\bullet$, $\mathcal{L}^\bullet$와 사상
$\alpha : \mathcal{K}^\bullet \to K$와
$\beta : \mathcal{L}^\bullet \to L$이 존재한다고 가정해도 된다.
여기서 $H^i(\alpha)$는 $i > n$에 대해 동형이고 $i = n$에 대해
전사이며, $H^i(\beta)$는 $i > m$에 대해 동형이고 $i = m$에 대해
전사이다. 그러면 사상
$$
\alpha \otimes^\mathbf{L} \beta :
\text{Tot}(\mathcal{K}^\bullet \otimes_\mathcal{O} \mathcal{L}^\bullet)
\to K \otimes_\mathcal{O}^\mathbf{L} L
$$
은 $i > t$인 차수 $i$의 코호몰로지 층에서 동형을 유도하고,
$i = t$에서 전사를 유도한다. 이는 Tor의 스펙트럼 열에서 따른다
(세부 사항은 생략한다).

\medskip\noindent
(2)의 증명. $U$를 $\mathcal{C}$의 대상이라 하자.
먼저 $U$를 어떤 덮개의 원소들로, $\mathcal{C}$를 국소화
$\mathcal{C}/U$로 바꾸어 $K$와 $L$이 위로 유계인 경우로
귀착시킬 수 있다. 그러면 명제는 (1)에서 곧바로 따른다.
\end{proof}

\begin{lemma}
\label{sites-cohomology-lemma-summands-pseudo-coherent}
$(\mathcal{C}, \mathcal{O})$를 환 달린 사이트라 하고
$m \in \mathbf{Z}$라 하자. $K \oplus L$이 $D(\mathcal{O})$에서
$m$-유사코히런트(각각 유사코히런트)이면 $K$와 $L$도 그러하다.
\end{lemma}

\begin{proof}
$K \oplus L$이 $m$-유사코히런트라고 가정하자.
$U$를 $\mathcal{C}$의 대상이라 하자. $U$를 어떤 덮개의 원소들로
바꾸어 $K \oplus L \in D^-(\mathcal{O}_U)$라고 가정해도 되며,
따라서 $L \in D^-(\mathcal{O}_U)$이다.
구별 삼각형
$$
(K \oplus L, K \oplus L, L \oplus L[1]) =
(K, K, 0) \oplus (L, L, L \oplus L[1])
$$
이 있음을 주목하라. Derived Categories, Lemma
\ref{derived-lemma-direct-sum-triangles}를 보라.
Lemma \ref{sites-cohomology-lemma-cone-pseudo-coherent}에 의해
$L \oplus L[1]$은 $m$-유사코히런트이다.
따라서 $L[1] \oplus L[2]$도 $m$-유사코히런트이고,
귀납법으로 $L[n] \oplus L[n + 1]$은 $m$-유사코히런트이다.
$L$이 위로 유계이므로 충분히 큰 $n$에 대해 $L[n]$은
$m$-유사코히런트이다. 따라서 구별 삼각형
$$
(L[n], L[n] \oplus L[n - 1], L[n - 1])
$$
들을 사용하여 거꾸로 진행하면
$L[n - 1], L[n - 2], \ldots, L$이 $m$-유사코히런트라는
원하는 결론을 얻는다.
\end{proof}

\begin{lemma}
\label{sites-cohomology-lemma-finite-cohomology}
$(\mathcal{C}, \mathcal{O})$를 환 달린 사이트라 하고,
$K$를 $D(\mathcal{O})$의 대상이라 하자.
$m \in \mathbf{Z}$라 하자.
\begin{enumerate}
\item $K$가 $m$-유사코히런트이고 $i > m$에 대해
$H^i(K) = 0$이면, $H^m(K)$는 유한형 $\mathcal{O}$-가군이다.
\item $K$가 $m$-유사코히런트이고 $i > m + 1$에 대해
$H^i(K) = 0$이면, $H^{m + 1}(K)$는 유한 표시
$\mathcal{O}$-가군이다.
\end{enumerate}
\end{lemma}

\begin{proof}
(1)의 증명. $U$를 $\mathcal{C}$의 대상이라 하자.
$H^m(K)$가 $U$의 어떤 덮개의 원소들 위에서 유한 개의 단면으로
생성될 수 있음을 보여야 한다(Modules on Sites, Definition
\ref{sites-modules-definition-site-local}을 보라).
따라서 증명하는 동안 (유한 번) 덮개 $\{U_i \to U\}$를 택하여
$\mathcal{C}$를 $\mathcal{C}/U_i$로, $U$를 $U_i$로 바꾸어도 된다.
특히 정의에 의해 강한 완전 복합체 $\mathcal{E}^\bullet$와
사상 $\alpha : \mathcal{E}^\bullet \to K$가 존재하여,
$> m$인 차수의 코호몰로지에서는 동형을, 차수 $m$에서는 전사를
유도한다고 가정해도 된다. $\mathcal{E}^\bullet$에 대해 결론을
증명하면 충분하다. $\mathcal{E}^n \not = 0$인 가장 큰 정수를
$n$이라 하자. $n = m$이면 $H^m(\mathcal{E}^\bullet)$는
$\mathcal{E}^n$의 몫이므로 결론은 명백하다.
$n > m$이면 $H^n(E^\bullet) = 0$이므로
$\mathcal{E}^{n - 1} \to \mathcal{E}^n$은 전사이다.
Lemma \ref{sites-cohomology-lemma-local-lift-map}에 의해 ($U$를 어떤 덮개의
원소들로 바꾼 뒤) 이 전사의 단면을 찾아
$\mathcal{E}^{n - 1} = \mathcal{E}' \oplus \mathcal{E}^n$으로
쓸 수 있다. 따라서 차수 $n - 1$에서 $\mathcal{E}'$을,
차수 $n$에서 $0$을 갖는다는 점을 제외하면
$\mathcal{E}^\bullet$와 같은 복합체 $(\mathcal{E}')^\bullet$에
대해 결론을 증명하면 충분하다. $n$에 대한 귀납법으로 결론이 따른다.

\medskip\noindent
(2)의 증명. $\mathcal{C}$의 대상 $U$를 택하자.
(1)의 증명에서와 같이 $U$ 위에서 국소적으로 작업해도 된다.
따라서 강한 완전 복합체 $\mathcal{E}^\bullet$와 사상
$\alpha : \mathcal{E}^\bullet \to K$가 존재하여,
$> m$인 차수의 코호몰로지에서 동형을, 차수 $m$에서 전사를
유도한다고 가정해도 된다. (1)의 증명에서와 같이
$i > m + 1$에 대해 $\mathcal{E}^i = 0$인 경우로 귀착할 수 있다.
그러면
$H^{m + 1}(K) \cong H^{m + 1}(\mathcal{E}^\bullet) =
\Coker(\mathcal{E}^m \to \mathcal{E}^{m + 1})$
이고, 이는 유한 표시이다.
\end{proof}








\section{Tor 차원}
\label{sites-cohomology-section-tor}

\noindent
이 절에서는 평탄 가군에 의한 분해를 더 자세히 살펴본다.

\begin{definition}
\label{sites-cohomology-definition-tor-amplitude}
$(\mathcal{C}, \mathcal{O})$를 환 달린 사이트라 하자.
$E$를 $D(\mathcal{O})$의 대상이라 하자.
$a \leq b$인 $a, b \in \mathbf{Z}$를 택하자.
\begin{enumerate}
\item 모든 $\mathcal{O}$-가군 $\mathcal{F}$와 모든
$i \not \in [a, b]$에 대하여
$H^i(E \otimes_\mathcal{O}^\mathbf{L} \mathcal{F}) = 0$이면,
$E$가 {\it $[a, b]$ 안에서 Tor-진폭을 가진다}고 한다.
\item 어떤 $a, b$에 대하여 $E$가 $[a, b]$ 안에서
Tor-진폭을 가지면 {\it 유한 Tor 차원을 가진다}고 한다.
\item $\mathcal{C}$의 임의의 대상 $U$에 대하여 덮개
$\{U_i \to U\}$가 존재하여 모든 $i$에 대해 $E|_{U_i}$가
유한 Tor 차원을 가지면, $E$가 {\it 국소적으로 유한 Tor 차원을
가진다}고 한다.
\end{enumerate}
$\mathcal{O}$-가군 $\mathcal{F}$에 대하여, $D(\mathcal{O})$의
대상으로 본 $\mathcal{F}[0]$가 $[-d, 0]$ 안에서 Tor-진폭을 가지면
그 가군이 {\it Tor 차원 $\leq d$}를 가진다고 한다.
\end{definition}

\noindent
정의에 나오는 $E$가 유한 Tor 차원을 가지면 $E$는
$D^b(\mathcal{O})$의 대상이다. 실제로 위 정의에서
$\mathcal{F} = \mathcal{O}$를 택하면 이를 알 수 있다.

\begin{lemma}
\label{sites-cohomology-lemma-last-one-flat}
$(\mathcal{C}, \mathcal{O})$를 환 달린 사이트라 하자.
$\mathcal{E}^\bullet$을 $[a, b]$ 안에서 Tor-진폭을 가지는
평탄한 $\mathcal{O}$-가군들의 위로 유계인 복합체라 하자.
그러면 $\Coker(d_{\mathcal{E}^\bullet}^{a - 1})$은 평탄한
$\mathcal{O}$-가군이다.
\end{lemma}

\begin{proof}
평탄 가군들의 위로 유계인 복합체 $\mathcal{E}^\bullet$에 대해
임의의 $\mathcal{O}$-가군 $\mathcal{F}$마다
$\mathcal{E}^\bullet \otimes_\mathcal{O} \mathcal{F} =
\mathcal{E}^\bullet \otimes_\mathcal{O}^{\mathbf{L}} \mathcal{F}$
임을 알 수 있다. 따라서 모든 $\mathcal{O}$-가군 $\mathcal{F}$에
대하여 수열
$$
\mathcal{E}^{a - 2} \otimes_\mathcal{O} \mathcal{F} \to
\mathcal{E}^{a - 1} \otimes_\mathcal{O} \mathcal{F} \to
\mathcal{E}^a \otimes_\mathcal{O} \mathcal{F}
$$
은 가운데에서 완전하다. 한편
$\mathcal{E}^{a - 2} \to \mathcal{E}^{a - 1} \to \mathcal{E}^a \to
\Coker(d^{a - 1}) \to 0$
은 평탄 분해이므로, 모든 $\mathcal{O}$-가군 $\mathcal{F}$에 대하여
$\text{Tor}_1^\mathcal{O}(\Coker(d^{a - 1}), \mathcal{F}) = 0$
이다. 이는 $\Coker(d^{a - 1})$이 평탄함을 뜻한다.
Lemma \ref{sites-cohomology-lemma-flat-tor-zero}를 보라.
\end{proof}

\begin{lemma}
\label{sites-cohomology-lemma-tor-amplitude}
$(\mathcal{C}, \mathcal{O})$를 환 달린 사이트라 하고,
$E$를 $D(\mathcal{O})$의 대상이라 하자.
$a \leq b$인 $a, b \in \mathbf{Z}$를 택하자. 다음은 동치이다.
\begin{enumerate}
\item $E$는 $[a, b]$ 안에서 Tor-진폭을 가진다.
\item $E$는 평탄한 $\mathcal{O}$-가군들의 복합체
$\mathcal{E}^\bullet$로 표시되며,
$i \not \in [a, b]$이면 $\mathcal{E}^i = 0$이다.
\end{enumerate}
\end{lemma}

\begin{proof}
(2)가 성립하면
$E \otimes_\mathcal{O}^\mathbf{L} \mathcal{F} =
\mathcal{E}^\bullet \otimes_\mathcal{O} \mathcal{F}$
로 계산할 수 있으므로 (1)이 성립함은 분명하다.

\medskip\noindent
(1)이 성립한다고 하자. Section \ref{sites-cohomology-section-flat}에 따라 $E$를
평탄한 $\mathcal{O}$-가군들의 위로 유계인 복합체
$\mathcal{K}^\bullet$로 표시할 수 있다.
$\mathcal{K}^n \not = 0$인 가장 큰 정수를 $n$이라 하자.
$n > b$이면 $H^n(\mathcal{K}^\bullet) = 0$이므로
$\mathcal{K}^{n - 1} \to \mathcal{K}^n$은 전사이다.
$\mathcal{K}^n$이 평탄하므로
$\Ker(\mathcal{K}^{n - 1} \to \mathcal{K}^n)$도 평탄하다
(Modules on Sites, Lemma \ref{sites-modules-lemma-flat-ses}).
따라서 $\mathcal{K}^\bullet$을
$\tau_{\leq n - 1}\mathcal{K}^\bullet$로 바꿀 수 있다.
그러므로 $n$에 관한 귀납법에 의해 $K^\bullet$이
$i > b$일 때 $\mathcal{K}^i = 0$인 평탄한
$\mathcal{O}$-가군들의 복합체인 경우로 귀착된다.

\medskip\noindent
$\mathcal{E}^\bullet = \tau_{\geq a}\mathcal{K}^\bullet$로 놓자.
$\mathcal{E}^a$가 평탄하다는 점을 제외하면 모두 분명하며,
그 점은 Lemma \ref{sites-cohomology-lemma-last-one-flat}와 정의에서 곧바로 따른다.
\end{proof}

\begin{lemma}
\label{sites-cohomology-lemma-bounded-below-tor-amplitude}
$(\mathcal{C}, \mathcal{O})$를 환 달린 사이트라 하고,
$E$를 $D(\mathcal{O})$의 대상이라 하자.
$a \in \mathbf{Z}$라 하자. 다음은 동치이다.
\begin{enumerate}
\item $E$는 $[a, \infty]$ 안에서 Tor-진폭을 가진다.
\item $E$는 평탄한 $\mathcal{O}$-가군들의 K-평탄 복합체
$\mathcal{E}^\bullet$로 표시될 수 있으며,
$i \not \in [a, \infty]$이면 $\mathcal{E}^i = 0$이다.
\end{enumerate}
더 나아가 환 달린 사이트의 사상에 의한 임의의 당김이
평탄한 항들을 갖는 K-평탄 복합체가 되도록
$\mathcal{E}^\bullet$을 택할 수 있다.
\end{lemma}

\begin{proof}
(2) $\Rightarrow$ (1)은 즉시 따른다. (1)이 성립한다고 하자.
먼저 Lemma \ref{sites-cohomology-lemma-K-flat-resolution}에 따라 $E$를 표시하는,
평탄한 항들을 갖는 K-평탄 복합체 $\mathcal{K}^\bullet$을 택한다.
임의의 $\mathcal{O}$-가군 $\mathcal{M}$에 대하여
$$
\mathcal{K}^{n - 1} \otimes_\mathcal{O} \mathcal{M} \to
\mathcal{K}^n \otimes_\mathcal{O} \mathcal{M} \to
\mathcal{K}^{n + 1} \otimes_\mathcal{O} \mathcal{M}
$$
의 코호몰로지는 $H^n(E \otimes_\mathcal{O}^\mathbf{L} \mathcal{M})$을
계산한다. 이는 $n < a$일 때 언제나 영이다. 따라서 복합체
$\ldots \to \mathcal{K}^{a - 1} \to \mathcal{K}^a \to \mathcal{K}^{a + 1}$
에 Lemma \ref{sites-cohomology-lemma-last-one-flat}를 적용하면
$\mathcal{N} = \Coker(\mathcal{K}^{a - 1} \to \mathcal{K}^a)$가
평탄한 $\mathcal{O}$-가군임을 얻는다. 다음과 같이 놓자.
$$
\mathcal{E}^\bullet = \tau_{\geq a}\mathcal{K}^\bullet =
(\ldots \to 0 \to \mathcal{N} \to \mathcal{K}^{a + 1} \to \ldots )
$$
$\mathcal{K}^\bullet \to \mathcal{E}^\bullet$의 핵
$\mathcal{L}^\bullet$은 복합체
$$
\mathcal{L}^\bullet = (\ldots \to \mathcal{K}^{a - 1} \to
\mathcal{I} \to 0 \to \ldots)
$$
이다. 여기서 $\mathcal{I} \subset \mathcal{K}^a$는
$\mathcal{K}^{a - 1} \to \mathcal{K}^a$의 상이다.
짧은 완전열
$0 \to \mathcal{I} \to \mathcal{K}^a \to \mathcal{N} \to 0$
이 있고 뒤의 두 항이 평탄하므로
$\mathcal{I}$도 평탄한 $\mathcal{O}$-가군이다.
따라서 $\mathcal{L}^\bullet$은 평탄 가군들의 위로 유계인
복합체이므로 Lemma \ref{sites-cohomology-lemma-bounded-flat-K-flat}에 의해 K-평탄이다.
그러므로 Lemma \ref{sites-cohomology-lemma-K-flat-two-out-of-three-ses}에 의해
$\mathcal{E}^\bullet$도 K-평탄이다.

\medskip\noindent
마지막 주장을 증명하자. 환 달린 사이트의 사상
$f : (\mathcal{C}', \mathcal{O}') \to (\mathcal{C}, \mathcal{O})$를
택하자. Lemma \ref{sites-cohomology-lemma-pullback-K-flat}에 의해
복합체 $f^*\mathcal{K}^\bullet$은 평탄한 항들을 가지며 K-평탄이다.
$f^*\mathcal{L}^\bullet$은 평탄한 $\mathcal{O}'$-가군들의
위로 유계인 복합체이므로 K-평탄이다. $\mathcal{O}'$-가군
복합체들의 짧은 완전열
$$
0 \to f^*\mathcal{L}^\bullet \to f^*\mathcal{K}^\bullet \to
f^*\mathcal{E}^\bullet \to 0
$$
을 얻는다. 평탄 가군들의 짧은 완전열
$0 \to \mathcal{I} \to \mathcal{K}^a \to \mathcal{N} \to 0$
의 당김이 짧은 완전열이기 때문이다.
Lemma \ref{sites-cohomology-lemma-K-flat-two-out-of-three-ses}에 의해
$f^*\mathcal{E}^\bullet$은 K-평탄이고, 이로써 증명이 끝난다.
\end{proof}

\begin{lemma}
\label{sites-cohomology-lemma-tor-amplitude-pullback}
환 달린 사이트의 사상
$(f, f^\sharp) : (\mathcal{C}, \mathcal{O}_\mathcal{C}) \to
(\mathcal{D}, \mathcal{O}_\mathcal{D})$
이라 하고, $E$를 $D(\mathcal{O}_\mathcal{D})$의 대상이라 하자.
$E$가 $[a, b]$ 안에서 Tor-진폭을 가지면 $Lf^*E$도
$[a, b]$ 안에서 Tor-진폭을 가진다.
\end{lemma}

\begin{proof}
$E$가 $[a, b]$ 안에서 Tor-진폭을 가진다고 하자.
Lemma \ref{sites-cohomology-lemma-tor-amplitude}에 따라 $E$를,
$i \not \in [a, b]$이면 $\mathcal{E}^i = 0$인 평탄한
$\mathcal{O}$-가군들의 복합체 $\mathcal{E}^\bullet$로 표시할 수 있다.
그러면 $Lf^*E$는 $f^*\mathcal{E}^\bullet$로 표시된다.
Modules on Sites, Lemma \ref{sites-modules-lemma-pullback-flat}에 의해
가군 $f^*\mathcal{E}^i$는 평탄하다.
따라서 Lemma \ref{sites-cohomology-lemma-tor-amplitude}에 의해
$Lf^*E$가 $[a, b]$ 안에서 Tor-진폭을 가짐을 얻는다.
\end{proof}

\begin{lemma}
\label{sites-cohomology-lemma-cone-tor-amplitude}
$(\mathcal{C}, \mathcal{O})$를 환 달린 사이트라 하자.
$(K, L, M, f, g, h)$를 $D(\mathcal{O})$의 특수 삼각형이라 하고,
$a, b \in \mathbf{Z}$라 하자.
\begin{enumerate}
\item $K$가 $[a + 1, b + 1]$ 안에서 Tor-진폭을 가지고
$L$이 $[a, b]$ 안에서 Tor-진폭을 가지면,
$M$은 $[a, b]$ 안에서 Tor-진폭을 가진다.
\item $K$와 $M$이 $[a, b]$ 안에서 Tor-진폭을 가지면,
$L$도 $[a, b]$ 안에서 Tor-진폭을 가진다.
\item $L$이 $[a + 1, b + 1]$ 안에서 Tor-진폭을 가지고
$M$이 $[a, b]$ 안에서 Tor-진폭을 가지면,
$K$는 $[a + 1, b + 1]$ 안에서 Tor-진폭을 가진다.
\end{enumerate}
\end{lemma}

\begin{proof}
생략한다. 힌트: 특수 삼각형에 딸린 긴 완전 코호몰로지 수열과
$- \otimes_\mathcal{O}^{\mathbf{L}} \mathcal{F}$
가 특수 삼각형을 보존한다는 사실에서 곧바로 따른다.
(2)를 증명하는 것이 가장 쉽고, 나머지는 평행이동을 통해
그로부터 따른다.
\end{proof}

\begin{lemma}
\label{sites-cohomology-lemma-tensor-tor-amplitude}
$(\mathcal{C}, \mathcal{O})$를 환 달린 사이트라 하고,
$K, L$을 $D(\mathcal{O})$의 대상이라 하자.
$K$가 $[a, b]$ 안에서 Tor-진폭을 가지고
$L$이 $[c, d]$ 안에서 Tor-진폭을 가지면,
$K \otimes_\mathcal{O}^\mathbf{L} L$은
$[a + c, b + d]$ 안에서 Tor-진폭을 가진다.
\end{lemma}

\begin{proof}
생략한다. 힌트: Tor의 스펙트럼열을 사용하라.
\end{proof}

\begin{lemma}
\label{sites-cohomology-lemma-summands-tor-amplitude}
$(\mathcal{C}, \mathcal{O})$를 환 달린 사이트라 하고,
$a, b \in \mathbf{Z}$라 하자. $D(\mathcal{O})$의 대상
$K$, $L$에 대하여 $K \oplus L$이 $[a, b]$ 안에서
Tor-진폭을 가지면 $K$와 $L$도 그러하다.
\end{lemma}

\begin{proof}
Tor 함자들이 가법적이라는 사실에서 명백하다.
\end{proof}

\begin{lemma}
\label{sites-cohomology-lemma-bounded}
$(\mathcal{C}, \mathcal{O})$를 환 달린 사이트라 하자.
$\mathcal{I} \subset \mathcal{O}$를 아이디얼 층이라 하고,
$K$를 $D(\mathcal{O})$의 대상이라 하자.
\begin{enumerate}
\item $K \otimes_\mathcal{O}^\mathbf{L} \mathcal{O}/\mathcal{I}$가
위로 유계이면,
$K \otimes_\mathcal{O}^\mathbf{L} \mathcal{O}/\mathcal{I}^n$
은 모든 $n$에 대해 균일하게 위로 유계이다.
\item $D(\mathcal{O}/\mathcal{I})$의 대상으로서
$K \otimes_\mathcal{O}^\mathbf{L} \mathcal{O}/\mathcal{I}$가
$[a, b]$ 안에서 Tor-진폭을 가지면, 모든 $n$에 대해
$D(\mathcal{O}/\mathcal{I}^n)$의 대상으로서
$K \otimes_\mathcal{O}^\mathbf{L} \mathcal{O}/\mathcal{I}^n$도
$[a, b]$ 안에서 Tor-진폭을 가진다.
\end{enumerate}
\end{lemma}

\begin{proof}
(1)을 증명하자.
$K \otimes_\mathcal{O}^\mathbf{L} \mathcal{O}/\mathcal{I}$
가 위로 유계이고, 예를 들어
$H^i(K \otimes_\mathcal{O}^\mathbf{L} \mathcal{O}/\mathcal{I}) = 0$
이 $i > b$에 대해 성립한다고 하자. 특수 삼각형
$$
K \otimes_\mathcal{O}^\mathbf{L}
\mathcal{I}^n/\mathcal{I}^{n + 1} \to
K \otimes_\mathcal{O}^\mathbf{L}
\mathcal{O}/\mathcal{I}^{n + 1} \to
K \otimes_\mathcal{O}^\mathbf{L}
\mathcal{O}/\mathcal{I}^n \to
K \otimes_\mathcal{O}^\mathbf{L}
\mathcal{I}^n/\mathcal{I}^{n + 1}[1]
$$
들과 등식
$$
K \otimes_\mathcal{O}^\mathbf{L}
\mathcal{I}^n/\mathcal{I}^{n + 1} =
\left(
K \otimes_\mathcal{O}^\mathbf{L}
\mathcal{O}/\mathcal{I}\right)
\otimes_{\mathcal{O}/\mathcal{I}}^\mathbf{L}
\mathcal{I}^n/\mathcal{I}^{n + 1}
$$
이 있다. 귀납법에 의해 모든 $n$에 대해
$H^i(K \otimes_\mathcal{O}^\mathbf{L} \mathcal{O}/\mathcal{I}^n) = 0$
이 $i > b$에 대해 성립함을 얻는다.

\medskip\noindent
(2)를 증명하자. $D(\mathcal{O}/\mathcal{I})$의 대상으로서
$K \otimes_\mathcal{O}^\mathbf{L} \mathcal{O}/\mathcal{I}$가
$[a, b]$ 안에서 Tor-진폭을 가진다고 하자.
$\mathcal{F}$를 $\mathcal{O}/\mathcal{I}^n$-가군의 층이라 하자.
그러면 유한 여과
$$
0 \subset \mathcal{I}^{n - 1}\mathcal{F} \subset \ldots
\subset \mathcal{I}\mathcal{F} \subset \mathcal{F}
$$
가 있고, 그 연속 몫들은 $\mathcal{O}/\mathcal{I}$-가군의 층이다.
따라서 $K \otimes_\mathcal{O}^\mathbf{L} \mathcal{O}/\mathcal{I}^n$이
$[a, b]$ 안에서 Tor-진폭을 가짐을 보이려면, 모든
$\mathcal{O}/\mathcal{I}$-가군 $\mathcal{G}$에 대하여
$H^i(K \otimes_\mathcal{O}^\mathbf{L} \mathcal{O}/\mathcal{I}^n
\otimes_{\mathcal{O}/\mathcal{I}^n}^\mathbf{L} \mathcal{G})$
가 $i \not \in [a, b]$일 때 영임을 보이면 충분하다. 모든
$\mathcal{O}/\mathcal{I}$-가군의 층 $\mathcal{G}$에 대하여
$$
\left(K \otimes_\mathcal{O}^\mathbf{L} \mathcal{O}/\mathcal{I}^n\right)
\otimes_{\mathcal{O}/\mathcal{I}^n}^\mathbf{L} \mathcal{G}
=
\left(K \otimes_\mathcal{O}^\mathbf{L} \mathcal{O}/\mathcal{I}\right)
\otimes_{\mathcal{O}/\mathcal{I}}^\mathbf{L} \mathcal{G}
$$
이므로 결론이 따른다.
\end{proof}

\begin{lemma}
\label{sites-cohomology-lemma-tor-amplitude-stalk}
$(\mathcal{C}, \mathcal{O})$를 환 달린 사이트라 하고,
$E$를 $D(\mathcal{O})$의 대상이라 하자.
$a, b \in \mathbf{Z}$라 하자.
\begin{enumerate}
\item $E$가 $[a, b]$ 안에서 Tor-진폭을 가지면, 사이트
$\mathcal{C}$의 모든 점 $p$에 대하여 $D(\mathcal{O}_p)$의 대상
$E_p$도 $[a, b]$ 안에서 Tor-진폭을 가진다.
\item $\mathcal{C}$가 충분히 많은 점을 가지면 그 역도 참이다.
\end{enumerate}
\end{lemma}

\begin{proof}
(1)을 증명하자. $p$에서 줄기를 취하는 것은 환 달린 사이트의 사상
$(p, \mathcal{O}_p) \to (\mathcal{C}, \mathcal{O})$에 의한
당김과 같으므로 Lemma \ref{sites-cohomology-lemma-tor-amplitude-pullback}를 적용할 수 있다.

\medskip\noindent
(2)를 증명하자. $\mathcal{C}$가 충분히 많은 점을 가지면
$H^i(E \otimes_\mathcal{O}^\mathbf{L} \mathcal{F})$
의 소멸을 줄기에서 확인할 수 있다. Modules on Sites,
Lemma \ref{sites-modules-lemma-check-exactness-stalks}를 보라.
$H^i(E \otimes_\mathcal{O}^\mathbf{L} \mathcal{F})_p =
H^i(E_p \otimes_{\mathcal{O}_p}^\mathbf{L} \mathcal{F}_p)$이므로
결론이 따른다.
\end{proof}










\section{완전 복합체}
\label{sites-cohomology-section-perfect}

\noindent
이 절에서는 환 달린 사이트 위의 완전 복합체가 가지는 성질을 논한다.

\begin{definition}
\label{sites-cohomology-definition-perfect}
$(\mathcal{C}, \mathcal{O})$를 환 달린 사이트라 하자.
$\mathcal{E}^\bullet$을 $\mathcal{O}$-가군의 복합체라 하자.
$\mathcal{C}$의 모든 대상 $U$에 대하여 덮개 $\{U_i \to U\}$가 존재하고,
각 $i$에 대하여 복합체의 사상
$\mathcal{E}_i^\bullet \to \mathcal{E}^\bullet|_{U_i}$
가 존재하여 이 사상이 준동형이고 $\mathcal{E}_i^\bullet$이
강한 완전 복합체이면 $\mathcal{E}^\bullet$을 {\it 완전}이라고 한다.
$D(\mathcal{O})$의 대상 $E$가 완전한 $\mathcal{O}$-가군 복합체로
표시될 수 있으면 이 대상을 {\it 완전}이라고 한다.
\end{definition}

\noindent
$\Sh(\mathcal{C})$가 준콤팩트이면
(Sites, Section \ref{sites-section-quasi-compact})
$D(\mathcal{O})$의 완전 대상은 $D^b(\mathcal{O})$에 속한다.
그러나 그렇지 않은 경우에는 반드시 그러하지는 않다.

\begin{lemma}
\label{sites-cohomology-lemma-perfect-independent-representative}
$(\mathcal{C}, \mathcal{O})$를 환 달린 사이트라 하자.
$E$를 $D(\mathcal{O})$의 대상이라 하자.
\begin{enumerate}
\item $\mathcal{C}$가 끝 대상 $X$를 가지고, 덮개 $\{U_i \to X\}$,
$\mathcal{O}_{U_i}$-가군의 강한 완전 복합체
$\mathcal{E}_i^\bullet$, 그리고 $D(\mathcal{O}_{U_i})$에서의 동형
 $\alpha_i : \mathcal{E}_i^\bullet \to E|_{U_i}$가 존재하면
$E$는 완전하다.
\item $E$가 완전하면 $E$를 표시하는 임의의 복합체는 완전하다.
\end{enumerate}
\end{lemma}

\begin{proof}
보조정리 \ref{sites-cohomology-lemma-pseudo-coherent-independent-representative}의
증명과 동일하다.
\end{proof}

\begin{lemma}
\label{sites-cohomology-lemma-perfect-precise}
$(\mathcal{C}, \mathcal{O})$를 환 달린 사이트라 하자.
$E$를 $D(\mathcal{O})$의 대상이라 하자.
$a \leq b$를 정수라 하자. $E$가 $[a, b]$에서 Tor-진폭을 가지고
$(a - 1)$-유사코히런트이면 $E$는 완전하다.
\end{lemma}

\begin{proof}
$U$를 $\mathcal{C}$의 대상이라 하자. $U$를 어떤 덮개의 원소들로,
$\mathcal{C}$를 국소화 $\mathcal{C}/U$로 바꾸어, 강한 완전 복합체
$\mathcal{E}^\bullet$과 사상 $\alpha : \mathcal{E}^\bullet \to E$가
존재하고 $i \geq a$에 대하여 $H^i(\alpha)$가 동형이라고
가정할 수 있다. $\mathcal{E}^\bullet$을
$\sigma_{\geq a - 1}\mathcal{E}^\bullet$로 바꾸어 두자.
특수 삼각형
$$
\mathcal{E}^\bullet \to E \to C \to \mathcal{E}^\bullet[1]
$$
을 택하자. $E$와 $\mathcal{E}^\bullet$의 코호몰로지 층들이
소멸한다는 사실과 $\alpha$에 관한 가정으로부터
$C \cong \mathcal{K}[2 - a]$를 얻는다. 여기서
$\mathcal{K} = \Ker(\mathcal{E}^{a - 1} \to \mathcal{E}^a)$이다.
$\mathcal{F}$를 $\mathcal{O}$-가군이라 하자.
$- \otimes_\mathcal{O}^\mathbf{L} \mathcal{F}$를 적용하면,
$E$가 $[a, b]$에서 Tor-진폭을 가진다는 가정에 의해
$\mathcal{K} \otimes_\mathcal{O} \mathcal{F} \to
\mathcal{E}^{a - 1} \otimes_\mathcal{O} \mathcal{F}$의 상은
$\Ker(\mathcal{E}^{a - 1} \otimes_\mathcal{O} \mathcal{F}
\to \mathcal{E}^a \otimes_\mathcal{O} \mathcal{F})$이다.
따라서 $\mathcal{E}' = \Coker(\mathcal{E}^{a - 1} \to \mathcal{E}^a)$라
놓으면 $\text{Tor}_1^\mathcal{O}(\mathcal{E}', \mathcal{F}) = 0$이다.
그러므로 $\mathcal{E}'$은 평탄하다
(보조정리 \ref{sites-cohomology-lemma-flat-tor-zero}).
따라서 Modules on Sites, Lemma
\ref{sites-modules-lemma-flat-locally-finite-presentation}에 의해
$\mathcal{E}'|_{U_i}$가 유한 자유 가군의 직합인자가 되는 덮개
$\{U_i \to U\}$가 존재한다. 이에 따라 복합체
$$
\mathcal{E}'|_{U_i} \to \mathcal{E}^{a + 1}|_{U_i} \to \ldots \to
\mathcal{E}^b|_{U_i}
$$
는 $E|_{U_i}$와 준동형이고 $E$는 완전하다.
\end{proof}

\begin{lemma}
\label{sites-cohomology-lemma-perfect}
$(\mathcal{C}, \mathcal{O})$를 환 달린 사이트라 하자.
$E$를 $D(\mathcal{O})$의 대상이라 하자.
다음은 동치이다.
\begin{enumerate}
\item $E$는 완전하다.
\item $E$는 유사코히런트이고 국소적으로 유한 Tor 차원을 가진다.
\end{enumerate}
\end{lemma}

\begin{proof}
(1)을 가정하자. $U$를 $\mathcal{C}$의 대상이라 하자.
정의에 의해 덮개 $\{U_i \to U\}$가 존재하여
$E|_{U_i}$가 강한 완전 복합체로 표시된다.
따라서 보조정리
\ref{sites-cohomology-lemma-pseudo-coherent-independent-representative}에 의해
$E$는 유사코히런트이다(즉, 모든 $m$에 대하여
$m$-유사코히런트이다). 또한 유한 자유 가군의 직합인자는 평탄하므로
보조정리 \ref{sites-cohomology-lemma-tor-amplitude}에 의해
$E|_{U_i}$는 유한 Tor 차원을 가진다. 따라서 (2)가 성립한다.

\medskip\noindent
(2)를 가정하자. $U$를 $\mathcal{C}$의 대상이라 하자.
$U$를 어떤 덮개의 원소들로 바꾸어, $E|_U$가
$[a, b]$에서 Tor-진폭을 가지게 하는 정수 $a \leq b$가
존재한다고 가정할 수 있다. $E|_U$는 모든 $m$에 대하여
$m$-유사코히런트이므로 보조정리
\ref{sites-cohomology-lemma-perfect-precise}를 사용하여 결론을 얻는다.
\end{proof}

\begin{lemma}
\label{sites-cohomology-lemma-perfect-pullback}
$(f, f^\sharp) : (\mathcal{C}, \mathcal{O}_\mathcal{C}) \to
(\mathcal{D}, \mathcal{O}_\mathcal{D})$
를 환 달린 사이트의 사상이라 하자.
$E$를 $D(\mathcal{O}_\mathcal{D})$의 대상이라 하자.
$E$가 $D(\mathcal{O}_\mathcal{D})$에서 완전하면,
$Lf^*E$는 $D(\mathcal{O}_\mathcal{C})$에서 완전하다.
\end{lemma}

\begin{proof}
이는 보조정리 \ref{sites-cohomology-lemma-perfect},
\ref{sites-cohomology-lemma-tor-amplitude-pullback}, 그리고
\ref{sites-cohomology-lemma-pseudo-coherent-pullback}에서 따른다.
\end{proof}

\begin{lemma}
\label{sites-cohomology-lemma-two-out-of-three-perfect}
$(\mathcal{C}, \mathcal{O})$를 환 달린 사이트라 하자.
$(K, L, M, f, g, h)$를 $D(\mathcal{O})$의 특수 삼각형이라 하자.
$K, L, M$ 가운데 두 대상이 완전하면 나머지 하나도 완전하다.
\end{lemma}

\begin{proof}
첫 번째 증명: 보조정리 \ref{sites-cohomology-lemma-perfect},
\ref{sites-cohomology-lemma-cone-pseudo-coherent}, 그리고
\ref{sites-cohomology-lemma-cone-tor-amplitude}를 결합하면 된다.
두 번째 증명(개요): $K$와 $L$이 완전하다고 하자.
$U$를 $\mathcal{C}$의 대상이라 하자. $U$를 어떤 덮개의 원소들로
바꾸어 $K|_U$와 $L|_U$가 강한 완전 복합체
$\mathcal{K}^\bullet$과 $\mathcal{L}^\bullet$로 표시된다고
가정할 수 있다. 다시 $U$를 어떤 덮개의 원소들로 바꾸어 사상
$K|_U \to L|_U$가 복합체의 사상
$\alpha : \mathcal{K}^\bullet \to \mathcal{L}^\bullet$로 주어진다고
가정할 수 있다. 보조정리 \ref{sites-cohomology-lemma-local-actual}을 보라.
그러면 $M|_U$는 $\alpha$의 원뿔과 동형이고,
이는 보조정리 \ref{sites-cohomology-lemma-cone}에 의해 강한 완전 복합체이다.
\end{proof}

\begin{lemma}
\label{sites-cohomology-lemma-tensor-perfect}
$(\mathcal{C}, \mathcal{O})$를 환 달린 사이트라 하자.
$K, L$이 $D(\mathcal{O})$의 완전 대상이면
$K \otimes_\mathcal{O}^\mathbf{L} L$도 완전하다.
\end{lemma}

\begin{proof}
보조정리 \ref{sites-cohomology-lemma-perfect}, \ref{sites-cohomology-lemma-tensor-pseudo-coherent},
그리고 \ref{sites-cohomology-lemma-tensor-tor-amplitude}에서 따른다.
\end{proof}

\begin{lemma}
\label{sites-cohomology-lemma-summands-perfect}
$(\mathcal{C}, \mathcal{O})$를 환 달린 사이트라 하자.
$K \oplus L$이 $D(\mathcal{O})$의 완전 대상이면
$K$와 $L$도 완전하다.
\end{lemma}

\begin{proof}
보조정리 \ref{sites-cohomology-lemma-perfect}, \ref{sites-cohomology-lemma-summands-pseudo-coherent},
그리고 \ref{sites-cohomology-lemma-summands-tor-amplitude}에서 따른다.
\end{proof}














\section{쌍대}
\label{sites-cohomology-section-duals}

\noindent
이 절에서는 복합체 범주와 유도 범주의 쌍대화 가능한 대상을
특징짓는다. 특히 $D(\mathcal{O})$의 대상이 쌍대를 가질
필요충분조건은 그 대상이 완전한 것임을 보일 것이다(이는 예
\ref{sites-cohomology-example-dual-derived}와 보조정리
\ref{sites-cohomology-lemma-left-dual-derived}에서 따른다).

\begin{lemma}
\label{sites-cohomology-lemma-symmetric-monoidal-cat-complexes}
$(\mathcal{C}, \mathcal{O})$를 환 달린 사이트라 하자.
$\mathcal{O}$-가군들의 복합체 범주에 텐서곱을
$\mathcal{F}^\bullet \otimes \mathcal{G}^\bullet =
\text{Tot}(\mathcal{F}^\bullet \otimes_\mathcal{O} \mathcal{G}^\bullet)$
로 정의하면 이 범주는 대칭 모노이드 범주이다.
\end{lemma}

\begin{proof}
생략한다. 힌트: 단위원 $\mathbf{1}$로는 차수 $0$에서
$\mathcal{O}$이고 다른 차수에서는 영인 복합체를 택하고,
자명한 동형사상
$\text{Tot}(\mathbf{1} \otimes_\mathcal{O} \mathcal{G}^\bullet) =
\mathcal{G}^\bullet$과
$\text{Tot}(\mathcal{F}^\bullet \otimes_\mathcal{O} \mathbf{1}) =
\mathcal{F}^\bullet$을 사용한다.
보조정리를 증명하려면 여러 도식의 가환성을 확인해야 한다.
Categories, 정의
\ref{categories-definition-monoidal-category}와
\ref{categories-definition-symmetric-monoidal-category}.
각각의 경우 확인은 곧바르다.
\end{proof}

\begin{example}
\label{sites-cohomology-example-dual}
$(\mathcal{C}, \mathcal{O})$를 환 달린 사이트라 하자.
$\mathcal{F}^\bullet$을 $\mathcal{O}$-가군들의 복합체라 하되,
각 $U \in \Ob(\mathcal{C})$에 대해 덮개 $\{U_i \to U\}$가 존재하여
$\mathcal{F}^\bullet|_{U_i}$가 강한 완전 복합체라고 하자.
복합체
$$
\mathcal{G}^\bullet = \SheafHom^\bullet(\mathcal{F}^\bullet, \mathcal{O})
$$
를 절 \ref{sites-cohomology-section-hom-complexes}에서와 같이 생각하자.
$$
\eta :
\mathcal{O}
\to
\text{Tot}(\mathcal{F}^\bullet \otimes_\mathcal{O} \mathcal{G}^\bullet)
\quad\text{and}\quad
\epsilon :
\text{Tot}(\mathcal{G}^\bullet \otimes_\mathcal{O} \mathcal{F}^\bullet)
\to
\mathcal{O}
$$
를 $\eta = \sum \eta_n$과 $\epsilon = \sum \epsilon_n$으로 정의하자.
여기서 $\eta_n : \mathcal{O} \to
\mathcal{F}^n \otimes_\mathcal{O} \mathcal{G}^{-n}$
이고
$\epsilon_n : \mathcal{G}^{-n} \otimes_\mathcal{O} \mathcal{F}^n
\to \mathcal{O}$은 Modules on Sites, 예
\ref{sites-modules-example-dual}에서와 같다.
그러면 $\mathcal{G}^\bullet, \eta, \epsilon$은 Categories, 정의
\ref{categories-definition-dual}의 의미에서
$\mathcal{F}^\bullet$의 왼쪽 쌍대이다.
등식
$(1 \otimes \epsilon) \circ (\eta \otimes 1) = \text{id}_{\mathcal{F}^\bullet}$
과
$(\epsilon \otimes 1) \circ (1 \otimes \eta) =
\text{id}_{\mathcal{G}^\bullet}$의 확인은 생략한다.
More on Algebra, 보조정리
\ref{more-algebra-lemma-left-dual-complex}와 비교하라.
\end{example}

\begin{lemma}
\label{sites-cohomology-lemma-left-dual-complex}
$(\mathcal{C}, \mathcal{O})$를 환 달린 사이트라 하고,
$\mathcal{F}^\bullet$을 $\mathcal{O}$-가군들의 복합체라 하자.
$\mathcal{F}^\bullet$이 $\mathcal{O}$-가군 복합체들의 모노이드 범주에서
왼쪽 쌍대를 가지면
(Categories, 정의 \ref{categories-definition-dual}),
$\mathcal{C}$의 각 대상 $U$에 대해 덮개 $\{U_i \to U\}$가 존재하여
$\mathcal{F}^\bullet|_{U_i}$는 강한 완전 복합체이고,
그 왼쪽 쌍대는 예 \ref{sites-cohomology-example-dual}에서 구성한 것이다.
\end{lemma}

\begin{proof}
왼쪽 쌍대의 유일성
(Categories, 주석 \ref{categories-remark-left-dual-adjoint})에 의해,
$\mathcal{F}^\bullet$이 명시된 성질을 가짐을 보이면 마지막 명제가 따른다.
$\mathcal{G}^\bullet, \eta, \epsilon$을 왼쪽 쌍대라 하자.
$\eta = \sum \eta_n$과 $\epsilon = \sum \epsilon_n$으로 쓰자.
여기서 $\eta_n : \mathcal{O} \to
\mathcal{F}^n \otimes_\mathcal{O} \mathcal{G}^{-n}$
이고
$\epsilon_n : \mathcal{G}^{-n} \otimes_\mathcal{O} \mathcal{F}^n
\to \mathcal{O}$이다.
Categories, 정의 \ref{categories-definition-dual}에 의해
$(1 \otimes \epsilon) \circ (\eta \otimes 1) = \text{id}_{\mathcal{F}^\bullet}$
과
$(\epsilon \otimes 1) \circ (1 \otimes \eta) = \text{id}_{\mathcal{G}^\bullet}$
이므로 곧바로
$(1 \otimes \epsilon_n) \circ (\eta_n \otimes 1) = \text{id}_{\mathcal{F}^n}$
과
$(\epsilon_n \otimes 1) \circ (1 \otimes \eta_n) =
\text{id}_{\mathcal{G}^{-n}}$을 얻는다.
다시 말해 $\mathcal{G}^{-n}$은 $\mathcal{F}^n$의 왼쪽 쌍대이며,
Modules on Sites, 보조정리
\ref{sites-modules-lemma-left-dual-module}를 각 $\mathcal{F}^n$에 적용할 수 있다.
$U$를 $\mathcal{C}$의 대상이라 하자. 각 $i$에 대해 유한 개의
$\eta_n|_{U_i}$만 영이 아니게 하는 덮개 $\{U_i \to U\}$가 존재한다.
따라서 $U$를 $U_i$로 바꾸어 유한 개의 $\eta_n|_U$만 영이 아니라고
가정할 수 있고, 인용한 보조정리에 의해 이는 유한 개의
$\mathcal{F}^n|_U$만 영이 아님을 뜻한다. 보조정리를 다시 사용하면
각 $\mathcal{F}^n|_{U_i}$가 유한 자유 $\mathcal{O}$-가군의
직합인자가 되게 하는 덮개 $\{U_i \to U\}$를 찾을 수 있으므로
증명이 끝난다.
\end{proof}

\begin{lemma}
\label{sites-cohomology-lemma-dual-perfect-complex}
$(\mathcal{C}, \mathcal{O})$를 환 달린 사이트라 하고,
$K$를 $D(\mathcal{O})$의 완전 대상이라 하자.
그러면 $K^\vee = R\SheafHom(K, \mathcal{O})$도 완전 대상이고
$(K^\vee)^\vee \cong K$이다. 함자적 동형사상
$$
M \otimes^\mathbf{L}_\mathcal{O} K^\vee = R\SheafHom_\mathcal{O}(K, M)
$$
및
$$
H^0(\mathcal{C}, M \otimes^\mathbf{L}_\mathcal{O} K^\vee) =
\Hom_{D(\mathcal{O})}(K, M)
$$
이 $D(\mathcal{O})$의 $M$에 대해 성립한다.
\end{lemma}

\begin{proof}
내부 Hom의 구성이 제한과 가환한다는 사실
(보조정리 \ref{sites-cohomology-lemma-restriction-RHom-to-U})을 더 언급하지 않고 사용하겠다.
$U$를 $\mathcal{C}$의 임의의 대상이라 하자. $K^\vee$가 완전함을
확인하려면 모든 $i$에 대해 $K^\vee|_{U_i}$가 완전하게 하는
덮개 $\{U_i \to U\}$가 존재함을 보이면 충분하다.
표준 사상
$$
K = R\SheafHom(\mathcal{O}_X, \mathcal{O}_X)
\otimes_{\mathcal{O}_X}^\mathbf{L} K \longrightarrow
R\SheafHom(R\SheafHom(K, \mathcal{O}_X), \mathcal{O}_X) =
(K^\vee)^\vee
$$
이 있다. 보조정리 \ref{sites-cohomology-lemma-internal-hom-evaluate}를 보라.
모든 $i$에 대해 이 사상을 $\mathcal{C}/U_i$로 제한한 것이
동형사상이 되게 하는 덮개 $\{U_i \to U\}$가 존재함을 보이면 충분하다.
보조정리 \ref{sites-cohomology-lemma-dual}에 의해 마지막 명제를 보이려면 사상
(\ref{sites-cohomology-equation-eval})
$$
M \otimes^\mathbf{L}_\mathcal{O} K^\vee
\longrightarrow
R\SheafHom(K, M)
$$
이 동형사상임을 확인하면 충분하다. 이 또한 (위의 의미에서) 국소적인
문제이다. 따라서 $K$가 강한 완전 복합체로 표시되는 경우에
보조정리를 증명하면 충분하다.

\medskip\noindent
$K$가 강한 완전 복합체
$\mathcal{E}^\bullet$로 표시된다고 하자. 보조정리
\ref{sites-cohomology-lemma-Rhom-strictly-perfect}에 의해 $K^\vee$는 차수 $-n$의 항이
$(\mathcal{E}^n)^\vee =
\SheafHom_\mathcal{O}(\mathcal{E}^n, \mathcal{O})$
인 복합체로 표시된다. $\mathcal{E}^n$이 유한 자유 $\mathcal{O}$-가군의
직합인자이므로 $(\mathcal{E}^n)^\vee$도 그러하다.
따라서 $K^\vee$도 강한 완전 복합체로 표시되며,
그 결과 $K^\vee$가 완전함을 알 수 있다.
사상 $K \to (K^\vee)^\vee$는 부호를 제외하면 동형사상인 평가 사상
$\mathcal{E}^n \to ((\mathcal{E}^n)^\vee)^\vee$들로 주어지므로
동형사상이다. (\ref{sites-cohomology-equation-eval})이 동형사상임을 보이기 위해
$M$을 K-평탄 복합체 $\mathcal{F}^\bullet$로 표시하자.
보조정리 \ref{sites-cohomology-lemma-Rhom-strictly-perfect}에 의해
$R\SheafHom(K, M)$은 그 항들이
$$
\bigoplus\nolimits_{n = p + q}
\SheafHom_\mathcal{O}(\mathcal{E}^{-q}, \mathcal{F}^p)
$$
인 복합체로 표시된다. 한편 대상
$M \otimes^\mathbf{L}_\mathcal{O} K^\vee$는 그 항들이
$$
\bigoplus\nolimits_{n = p + q}
\mathcal{F}^p \otimes_\mathcal{O} (\mathcal{E}^{-q})^\vee
$$
인 복합체로 표시된다. 따라서 (\ref{sites-cohomology-equation-eval})이
동형사상이라는 명제는 표준 사상
$$
\mathcal{F}
\otimes_\mathcal{O}
\SheafHom_\mathcal{O}(\mathcal{E}, \mathcal{O})
\longrightarrow
\SheafHom_\mathcal{O}(\mathcal{E}, \mathcal{F})
$$
이, $\mathcal{E}$가 유한 자유 $\mathcal{O}$-가군의 직합인자이고
$\mathcal{F}$가 임의의 $\mathcal{O}$-가군일 때 동형사상이라는
명제로 귀착된다. 이는 $\mathcal{E}$가 유한 자유인 경우의
대응하는 명제에서 곧바로 따른다.
\end{proof}

\begin{lemma}
\label{sites-cohomology-lemma-symmetric-monoidal-derived}
$(\mathcal{C}, \mathcal{O})$를 환 달린 사이트라 하자.
유도 텐서곱을 텐서곱으로 하고 통상적인 결합성과 가환성 제약을
갖춘 유도 범주 $D(\mathcal{O})$는 대칭 모노이드 범주이다
(부호 규칙은 More on Algebra, 절
\ref{more-algebra-section-sign-rules}를 보라).
\end{lemma}

\begin{proof}
생략한다. 보조정리
\ref{sites-cohomology-lemma-symmetric-monoidal-cat-complexes}와 비교하라.
\end{proof}

\begin{example}
\label{sites-cohomology-example-dual-derived}
$(\mathcal{C}, \mathcal{O})$를 환 달린 사이트라 하고,
$K$를 $D(\mathcal{O})$의 완전 대상이라 하자. 보조정리
\ref{sites-cohomology-lemma-dual-perfect-complex}에서와 같이
$K^\vee = R\SheafHom(K, \mathcal{O})$로 놓자. 그러면 사상
$$
K \otimes_\mathcal{O}^\mathbf{L} K^\vee \longrightarrow R\SheafHom(K, K)
$$
은 (그 보조정리에 의해) 동형사상이다.
$$
\eta :
\mathcal{O}
\longrightarrow
K \otimes_\mathcal{O}^\mathbf{L} K^\vee
$$
를 위 동형사상 아래에서 $\text{id}_K$에 대응하는 절단으로
$1$을 보내는 사상이라 하자.
$$
\epsilon : 
K^\vee
\otimes_\mathcal{O}^\mathbf{L} K
\longrightarrow
\mathcal{O}
$$
를 평가 사상이라 하자(예를 들어 이를 구성할 때 보조정리
\ref{sites-cohomology-lemma-internal-hom-composition}을 사용할 수 있다).
그러면 $K^\vee, \eta, \epsilon$은 Categories, 정의
\ref{categories-definition-dual}의 의미에서 $K$의 왼쪽 쌍대이다.
등식
$(1 \otimes \epsilon) \circ (\eta \otimes 1) = \text{id}_K$
과
$(\epsilon \otimes 1) \circ (1 \otimes \eta) =
\text{id}_{K^\vee}$의 확인은 생략한다.
\end{example}

\begin{lemma}
\label{sites-cohomology-lemma-left-dual-derived}
$(\mathcal{C}, \mathcal{O})$를 환 달린 사이트라 하고,
$M$을 $D(\mathcal{O})$의 대상이라 하자.
$M$이 모노이드 범주 $D(\mathcal{O})$에서 왼쪽 쌍대를 가지면
(Categories, 정의 \ref{categories-definition-dual}),
$M$은 완전하고 그 왼쪽 쌍대는 예
\ref{sites-cohomology-example-dual-derived}에서 구성한 것이다.
\end{lemma}

\begin{proof}
$N, \eta, \epsilon$을 왼쪽 쌍대라 하자. $\mathcal{C}$의 임의의 대상
$U$에 대해 제한 $N|_U, \eta|_U, \epsilon|_U$가
$M|_U$의 왼쪽 쌍대임에 유의하자.

\medskip\noindent
$U$를 $\mathcal{C}$의 대상이라 하자. $M|_{U_i}$가
$D(\mathcal{O}_{U_i})$의 완전 대상이 되게 하는 $\mathcal{C}$의 덮개
$\{U_i \to U\}_{i \in I}$를 찾으면 충분하다. 따라서
$\mathcal{C}, \mathcal{O}, M, N, \eta, \epsilon$을
$\mathcal{C}/U, \mathcal{O}_U, M|_U, N|_U, \eta|_U, \epsilon|_U$로
바꾸어 $\mathcal{C}$가 끝 대상 $X$를 가진다고 가정할 수 있다.
더 나아가 증명하는 동안 유한 번 $X$를 $X$의 덮개
$\{U_i \to X\}$의 구성원들로 바꿀 수 있다.

\medskip\noindent
다음 논증을 여러 번 사용하겠다. $M$을 표시하는 임의의
$\mathcal{O}$-가군 복합체 $\mathcal{M}^\bullet$을 택하자.
$N$을 표시하며 그 항들이 평탄한 $\mathcal{O}$-가군인
K-평탄 복합체 $\mathcal{N}^\bullet$을 택하자.
보조정리 \ref{sites-cohomology-lemma-K-flat-resolution}를 보라. 사상
$$
\eta : \mathcal{O} \to
\text{Tot}(\mathcal{M}^\bullet \otimes_\mathcal{O} \mathcal{N}^\bullet)
$$
$X$를 덮개의 구성원들로 바꾼 뒤, 정수 $N$과
$i = 1, \ldots, N$에 대한 정수 $n_i \in \mathbf{Z}$ 및
$\mathcal{M}^{n_i}$와 $\mathcal{N}^{-n_i}$의 절단 $f_i$, $g_i$를 택하여
$$
\eta(1) = \sum\nolimits_i f_i \otimes g_i
$$
$\mathcal{K}^\bullet \subset \mathcal{M}^\bullet$을
모든 $i = 1, \ldots, N$에 대한 절단 $f_i$들을 포함하는
임의의 $\mathcal{O}$-가군 부분복합체라 하자.
$\mathcal{N}^n$들이 평탄하므로
$\text{Tot}(\mathcal{K}^\bullet \otimes_\mathcal{O} \mathcal{N}^\bullet)
\subset
\text{Tot}(\mathcal{M}^\bullet \otimes_\mathcal{O} \mathcal{N}^\bullet)$
이고, 따라서 $\eta$는
$$
\tilde \eta :
\mathcal{O} \to
\text{Tot}(\mathcal{K}^\bullet \otimes_\mathcal{O} \mathcal{N}^\bullet)
$$
$\mathcal{K}^\bullet$로 표시되는 $D(\mathcal{O})$의 대상을
$K$라 하면 다음 가환 도식을 얻는다.
$$
\xymatrix{
M \ar[rr]_-{\eta \otimes 1} \ar[rrd]_{\tilde \eta \otimes 1} & &
M \otimes^\mathbf{L} N \otimes^\mathbf{L} M
\ar[r]_-{1 \otimes \epsilon} &
M \\
& &
K \otimes^\mathbf{L} N \otimes^\mathbf{L} M
\ar[u] \ar[r]^-{1 \otimes \epsilon} &
K \ar[u]
}
$$
위 행의 합성이 $M$ 위의 항등사상이므로
$M$은 $D(\mathcal{O})$에서 $K$의 직합인자이다.

\medskip\noindent
위 논증의 첫 번째 응용으로, 부분복합체
$\mathcal{K}^\bullet = \sigma_{\geq a} \tau_{\leq b}\mathcal{M}^\bullet$
을 모든 $i = 1, \ldots, N$에 대하여 $a < n_i < b$가 되도록
택할 수 있다. 따라서 $M$은 $D(\mathcal{O})$에서 유계 복합체의
직합인자이며, $M$이 $D^b(\mathcal{O})$에 속한다고 가정할 수 있음을
얻는다. (위 과정에서 $X$를 덮개의 구성원들로 바꾼다는 점을 상기하라.)

\medskip\noindent
$M$이 $D^b(\mathcal{O})$에 속하므로 $\mathcal{M}^\bullet$을
평탄 가군들의 위로 유계인 복합체로 택할 수 있다(Modules, 보조정리
\ref{modules-lemma-module-quotient-flat}과 Derived Categories, 보조정리
\ref{derived-lemma-subcategory-left-resolution}에 의한다).
그러면 위 논증에서 모든 $i = 1, \ldots, N$에 대하여 $a < n_i$가 되도록
$\mathcal{K}^\bullet = \sigma_{\geq a}\mathcal{M}^\bullet$을 택할 수 있다.
따라서 $M$이 $D(\mathcal{O})$에서 평탄 가군들의 유계 복합체의
직합인자라고 가정할 수 있다. 특히 $M$은 유한 Tor-진폭을 가진다.

\medskip\noindent
$M$이 $[a, b]$ 안에서 Tor-진폭을 가진다고 하자.
$M$이 $m$-유사코히런트라고 가정하고, $X$를 덮개의 구성원들로 바꾼 뒤
$M$이 $(m - 1)$-유사코히런트라고 가정할 수 있음을 보이겠다.
그러면 보조정리 \ref{sites-cohomology-lemma-perfect-precise}와
어쨌든 $M$이 $(b + 1)$-유사코히런트라는 사실에 의해 증명이 끝난다.
$X$를 덮개의 구성원들로 바꾸어 강한 완전 복합체
$\mathcal{E}^\bullet$과 $D(\mathcal{O})$의 사상
$\alpha : \mathcal{E}^\bullet \to M$이 존재하고,
$i > m$에 대하여 $H^i(\alpha)$가 동형이며 $i = m$에 대하여
전사라고 가정할 수 있다. $i < m$이면 $\mathcal{E}^i = 0$이라고
가정해도 된다. 특수 삼각형
$$
\mathcal{E}^\bullet \to M \to L \to \mathcal{E}^\bullet[1]
$$
을 택하자. $i \geq m$이면 $H^i(L) = 0$임에 유의하자.
따라서 $L$을, $i \geq m$이면 $\mathcal{L}^i = 0$인 복합체
$\mathcal{L}^\bullet$로 표시할 수 있다. 사상
$L \to \mathcal{E}^\bullet[1]$은 복합체의 사상
$\mathcal{L}^\bullet \to \mathcal{E}^\bullet[1]$로 주어지며,
이는 차수 $m - 1$을 제외한 모든 차수에서 영이고 그 차수에서는 사상
$\mathcal{L}^{m - 1} \to \mathcal{E}^m$을 준다.
Derived Categories, 보조정리 \ref{derived-lemma-negative-exts}를 보라.
그러면 $M$은 복합체
$$
\mathcal{M}^\bullet :
\ldots \to
\mathcal{L}^{m - 2} \to
\mathcal{L}^{m - 1} \to
\mathcal{E}^m \to
\mathcal{E}^{m + 1} \to \ldots
$$
로 표시된다. 이 복합체에 둘째 문단의 논의를 적용하여
모든 $i = 1, \ldots, N$에 대해 $\mathcal{M}^{n_i}$의 절단
$f_i$를 얻는다. $n < m$에 대해서는
$\mathcal{K}^n \subset \mathcal{L}^n$을 $n_i = n$인 절단
$f_i$들과 $n_i = n - 1$인 $d(f_i)$들로 생성되는
$\mathcal{O}$-부분가군으로 정의하자. $n \geq m$에 대해서는
$\mathcal{K}^n = \mathcal{E}^n$으로 놓자.
그러면 명백히 다음 특수 삼각형들의 사상이 있다.
$$
\xymatrix{
\mathcal{E}^\bullet \ar[r] &
\mathcal{M}^\bullet \ar[r] &
\mathcal{L}^\bullet \ar[r] &
\mathcal{E}^\bullet[1] \\
\mathcal{E}^\bullet \ar[r] \ar[u] &
\mathcal{K}^\bullet \ar[r] \ar[u] &
\sigma_{\leq m - 1}\mathcal{K}^\bullet \ar[r] \ar[u] &
\mathcal{E}^\bullet[1] \ar[u]
}
$$
여기서 모든 사상은 위에서 지시한 것들이다.
복합체 $\mathcal{K}^\bullet$에 대응하는 $D(\mathcal{O})$의 대상을
$K$라 하자. 증명의 둘째 문단에 나온 논증으로 $D(\mathcal{O})$의
사상 $s : M \to K$를 얻을 수 있으며, 합성 $M \to K \to M$은
$M$ 위의 항등사상이다. 도식
$$
\xymatrix{
\mathcal{E}^\bullet \ar[r] &
\mathcal{K}^\bullet \ar@{=}[r] &
K \\
\mathcal{E}^\bullet \ar[u]^{\text{id}} \ar[r]^i &
\mathcal{M}^\bullet \ar@{=}[r] &
M \ar[u]_s
}
$$
이 가환하는지는 알 수 없지만, 사상 $K \to M$과 합성한 뒤에는
가환함을 알고 있다. 보조정리 \ref{sites-cohomology-lemma-local-actual}에 의해
$X$를 덮개의 구성원들로 바꾸어 $s \circ i$가 복합체의 사상
$\sigma : \mathcal{E}^\bullet \to \mathcal{K}^\bullet$로 주어진다고
가정할 수 있다. 같은 보조정리에 의해 $\sigma$와 포함사상
$\mathcal{K}^\bullet \subset \mathcal{M}^\bullet$의 합성은 어떤 호모토피
$\{h^i : \mathcal{E}^i \to \mathcal{M}^{i - 1}\}$에 의해 영과
호모토픽이라고 가정할 수 있다. 따라서 $\mathcal{K}^{m - 1}$을
$\mathcal{K}^{m - 1} + \Im(h^m)$으로 바꾼 뒤에도
($\mathcal{K}^{m - 1}$이 유한 개의 전역 절단으로 생성된다는
성질은 여전히 성립한다) $\sigma$ 자체가 영과 호모토픽임을 알 수 있다!
이는 다음 실선 도식이 가환함을 뜻한다.
$$
\xymatrix{
\mathcal{E}^\bullet \ar[r] &
M \ar[r] &
\mathcal{L}^\bullet \ar[r] &
\mathcal{E}^\bullet[1] \\
\mathcal{E}^\bullet \ar[r] \ar[u] &
K \ar[r] \ar[u] &
\sigma_{\leq m - 1}\mathcal{K}^\bullet \ar[r] \ar[u] &
\mathcal{E}^\bullet[1] \ar[u] \\
\mathcal{E}^\bullet \ar[r] \ar[u] &
M \ar[r] \ar[u]^s &
\mathcal{L}^\bullet \ar[r] \ar@{..>}[u] &
\mathcal{E}^\bullet[1] \ar[u]
}
$$
삼각 범주의 공리들에 의해 이 도식에 맞는 점선 화살표를 얻는다.
차수 $m - 1$의 코호몰로지 층을 살펴보면 다음을 얻는다.
$$
\xymatrix{
H^{m - 1}(M) \ar[r] &
H^{m - 1}(\mathcal{L}^\bullet) \ar[r] &
H^m(\mathcal{E}^\bullet) \\
H^{m - 1}(K) \ar[r] \ar[u] &
H^{m - 1}(\sigma_{\leq m - 1}\mathcal{K}^\bullet) \ar[r] \ar[u] &
H^m(\mathcal{E}^\bullet) \ar[u] \\
H^{m - 1}(M) \ar[r] \ar[u] &
H^{m - 1}(\mathcal{L}^\bullet) \ar[r] \ar[u] &
H^m(\mathcal{E}^\bullet) \ar[u]
}
$$
왼쪽 열과 오른쪽 열에서 수직 합성이 모두 항등사상이므로,
가운데의 수직 합성
$H^{m - 1}(\mathcal{L}^\bullet) \to
H^{m - 1}(\sigma_{\leq m - 1}\mathcal{K}^\bullet) \to
H^{m - 1}(\mathcal{L}^\bullet)$은 전사이다!
특히 $H^{m - 1}(\sigma_{\leq m - 1}\mathcal{K}^\bullet) \to
H^{m - 1}(\mathcal{L}^\bullet)$은 전사이다.
위 삼각형 사상에서 유도되는 코호몰로지 층의 긴 완전 수열의
사상을 사용하면, 도식 추적을 통해 이는 $i \geq m$에 대하여
$H^i(K) \to H^i(M)$이 동형이고 $i = m - 1$에 대하여 전사임을
함의한다. 구성에 의해 어떤 $r \geq 0$과 전사
$\mathcal{O}^{\oplus r} \to \mathcal{K}^{m - 1}$을 택할 수 있다.
그러면 합성
$$
(\mathcal{O}^{\oplus r} \to \mathcal{E}^m \to
\mathcal{E}^{m + 1} \to \ldots ) \longrightarrow
K \longrightarrow M
$$
은 차수 $\geq m$에서 코호몰로지 층에 동형사상을 유도하고
차수 $m - 1$에서 전사를 유도하므로 증명이 끝난다.
\end{proof}

\begin{lemma}
\label{sites-cohomology-lemma-colim-and-lim-of-duals}
\begin{slogan}
완전 대상들의 계에 대한 자명한 쌍대성.
\end{slogan}
$(\mathcal{C}, \mathcal{O})$를 환 달린 사이트라 하자.
$(K_n)_{n \in \mathbf{N}}$을 $D(\mathcal{O})$의 완전 대상들의 계라 하자.
$K = \text{hocolim} K_n$를 유도 여극한이라 하자
(Derived Categories, 정의 \ref{derived-definition-derived-colimit}).
그러면 $D(\mathcal{O})$의 임의의 대상 $E$에 대해
$$
R\SheafHom(K, E) = R\lim E \otimes^\mathbf{L}_\mathcal{O} K_n^\vee
$$
이다. 여기서 $(K_n^\vee)$는 쌍대 완전 복합체들의 역계이다.
\end{lemma}

\begin{proof}
보조정리 \ref{sites-cohomology-lemma-dual-perfect-complex}에 의해
$R\lim E \otimes^\mathbf{L}_\mathcal{O} K_n^\vee =
R\lim R\SheafHom(K_n, E)$
이고, 이는 특수 삼각형
$$
R\lim R\SheafHom(K_n, E) \to
\prod R\SheafHom(K_n, E) \to
\prod R\SheafHom(K_n, E)
$$
에 들어간다. $K$도 마찬가지로 특수 삼각형
$\bigoplus K_n \to \bigoplus K_n \to K$에 들어가므로
$\prod R\SheafHom(K_n, E) = R\SheafHom(\bigoplus K_n, E)$임을 보이면
충분하다. 이는 (\ref{sites-cohomology-equation-internal-hom})과 유도 텐서곱이
직합과 가환한다는 사실의 형식적 귀결이다.
\end{proof}






\section{유도 범주의 가역 대상}
\label{sites-cohomology-section-invertible-D-or-R}

\noindent
환 달린 사이트의 유도 범주에서 가역인 대상들을 특징짓는다
(국소적으로 환 달린 토포스인 경우와 일반적인 경우를 모두 다룬다).

\begin{lemma}
\label{sites-cohomology-lemma-category-summands-finite-free}
$(\mathcal{C}, \mathcal{O})$를 환 달린 사이트라 하자.
$R = \Gamma(\mathcal{C}, \mathcal{O})$로 놓자.
유한 자유 $\mathcal{O}$-가군의 직합인자인 $\mathcal{O}$-가군들의
범주는 유한 사영 $R$-가군들의 범주와 동치이다.
\end{lemma}

\begin{proof}
유한 사영 $R$-가군은 유한 자유 $R$-가군의 직합인자와 같은 것임에
유의하자. 동치는 함자 $\mathcal{E} \mapsto
\Gamma(\mathcal{C}, \mathcal{E})$로 주어진다.
역함자는 다음 구성으로 주어진다. 전역 절단을 취하는 함자 $f_*$와,
$\Sh(\text{pt})$의 대상인 집합 $S$를 값이 $S$인 상수층으로 보내는
함자 $f^{-1}$을 갖는 토포스의 사상
$f : \Sh(\mathcal{C}) \to \Sh(\text{pt})$을 생각하자.
항등사상 $R \to f_*\mathcal{O}$을 사용하여 환 달린 토포스의 사상
$(f, f^\sharp) : (\Sh(\mathcal{C}), \mathcal{O}) \to (\Sh(\text{pt}), R)$
을 얻는다. 그러면 역함자는 $f^*$로 주어진다.
\end{proof}

\begin{lemma}
\label{sites-cohomology-lemma-invertible-derived}
$(\mathcal{C}, \mathcal{O})$를 환 달린 사이트라 하고,
$M$을 $D(\mathcal{O})$의 대상이라 하자. 다음은 동치이다.
\begin{enumerate}
\item $M$은 $D(\mathcal{O})$에서 가역이다.
Categories, 정의 \ref{categories-definition-invertible}를 보라.
\item 국소 유한\footnote{이는 $\mathcal{C}$의 각 대상 $U$에 대해
덮개 $\{U_i \to U\}$가 존재하여, 각 $i$마다 유한 개의 $n$에 대해서만
층 $\mathcal{O}_n|_{U_i}$가 영이 아님을 뜻한다.} 직접곱 분해
$$
\mathcal{O} = \prod\nolimits_{n \in \mathbf{Z}} \mathcal{O}_n
$$
가 존재하고, 각 $n$에 대해 가역 $\mathcal{O}_n$-가군
$\mathcal{H}^n$이 존재하며
(Modules on Sites, 정의 \ref{sites-modules-definition-invertible-sheaf}),
$D(\mathcal{O})$에서 $M = \bigoplus \mathcal{H}^n[-n]$이다.
\end{enumerate}
(1)과 (2)가 성립하면 $M$은 $D(\mathcal{O})$의 완전 대상이다.
$(\mathcal{C}, \mathcal{O})$가 국소적으로 환 달린 사이트이면
이 조건들은 다음 조건과도 동치이다.
\begin{enumerate}
\item[(3)] $\mathcal{C}$의 각 대상 $U$에 대해 덮개
$\{U_i \to U\}$가 존재하고 각 $i$에 대해 정수 $n_i$가 존재하여,
$M|_{U_i}$는 차수 $n_i$에 놓인 가역 $\mathcal{O}_{U_i}$-가군으로
표시된다.
\end{enumerate}
\end{lemma}

\begin{proof}
(2)를 가정하자. 대상 $R\SheafHom(M, \mathcal{O})$과 합성 사상
$$
R\SheafHom(M, \mathcal{O}) \otimes_\mathcal{O}^\mathbf{L} M \to \mathcal{O}
$$
을 생각하자. 이것이 동형사상임을 보이기 위해 국소적으로 작업해도 된다.
따라서 $\mathcal{O} = \prod_{a \leq n \leq b} \mathcal{O}_n$이고
$M = \bigoplus_{a \leq n \leq b} \mathcal{H}^n[-n]$이라고
가정할 수 있다. 그러면
$$
R\SheafHom(\mathcal{H}^m, \mathcal{O})
\otimes_\mathcal{O}^\mathbf{L} \mathcal{H}^n
$$
이 $n \not = m$이면 영이고 $n = m$이면 $\mathcal{O}_n$과 같음을
보이면 충분하다. $n \not = m$인 경우는 $\mathcal{O}_n$과
$\mathcal{O}_m$이 평탄한 $\mathcal{O}$-대수이고
$\mathcal{O}_n \otimes_\mathcal{O} \mathcal{O}_m = 0$이라는
사실에서 따른다. 가역 $\mathcal{O}$-가군의 국소 구조
(Modules on Sites, 보조정리 \ref{sites-modules-lemma-invertible})를 사용하고
국소적으로 작업하면 $n = m$인 경우의 동형사상도 곧바로 따른다.
세부 사항은 생략한다. $D(\mathcal{O})$가 대칭 모노이드 범주이므로
$M$이 가역임을 얻는다.

\medskip\noindent
(1)을 가정하자. (2)의 서술은 $\mathcal{O}_n$의 후보가
$\SheafHom_\mathcal{O}(H^n(M), H^n(M))$임을 보여 준다.
이것이 국소 유한한 환의 층들의 족이고
$\mathcal{O} = \prod \mathcal{O}_n$이면, 직접곱 분해에서 오는
$\mathcal{O}$의 멱등원들을 사용하여 곧바로 직접합 분해
$M = \bigoplus H^n(M)[-n]$을 얻는다. 따라서 (2)를 증명할 때
다음 의미에서 국소적으로 작업해도 된다. $U$를 $\mathcal{C}$의
대상이라 하자. 위와 같은 $\mathcal{O}_n$에 대하여
$\mathcal{C}/U_i$로 제한한 뒤 위의 명제들과 (2)가 성립하게 하는
$U$의 덮개 $\{U_i \to U\}$가 존재함을 보여야 한다.
따라서 $\mathcal{C}$가 끝 대상 $X$를 가진다고 가정할 수 있고,
(2)의 증명 중 유한 번 $X$를 $X$의 덮개 구성원들로 바꿀 수 있다.

\medskip\noindent
$D(\mathcal{O})$의 대상 $N$과 동형사상
$M \otimes_\mathcal{O}^\mathbf{L} N \cong \mathcal{O}$를 택하자.
그러면 $N$은 모노이드 범주 $D(\mathcal{O})$에서 $M$의 왼쪽 쌍대이고,
보조정리 \ref{sites-cohomology-lemma-left-dual-derived}에 의해 $M$은 완전하다.
대칭성에 의해 $N$도 완전하다. $X$를 덮개의 구성원들로 바꾸어
$M$과 $N$이 강한 완전 복합체 $\mathcal{E}^\bullet$과
$\mathcal{F}^\bullet$로 표시된다고 가정할 수 있다. 그러면
$M \otimes_\mathcal{O}^\mathbf{L} N$은
$\text{Tot}(\mathcal{E}^\bullet \otimes_\mathcal{O} \mathcal{F}^\bullet)$로
표시된다. $X$를 $X$의 덮개 구성원들로 다시 바꾸어 서로 역인 동형사상
$\mathcal{O} \to M \otimes_\mathcal{O}^\mathbf{L} N$과
$M \otimes_\mathcal{O}^\mathbf{L} N \to \mathcal{O}$가 복합체의 사상
$$
\alpha : \mathcal{O} \to
\text{Tot}(\mathcal{E}^\bullet \otimes_\mathcal{O} \mathcal{F}^\bullet)
\quad\text{and}\quad
\beta :
\text{Tot}(\mathcal{E}^\bullet \otimes_\mathcal{O} \mathcal{F}^\bullet)
\to \mathcal{O}
$$
으로 주어진다고 가정할 수 있다. 보조정리
\ref{sites-cohomology-lemma-local-actual}을 보라. 그러면 복합체의 사상으로서
$\beta \circ \alpha = 1$이고, $D(\mathcal{O})$의 사상으로서
$\alpha \circ \beta = 1$이다. $X$를 $X$의 덮개 구성원들로 바꾸어
합성 $\alpha \circ \beta$가 성분
$$
\theta^n :
\text{Tot}^n(\mathcal{E}^\bullet \otimes_\mathcal{O} \mathcal{F}^\bullet)
\to
\text{Tot}^{n - 1}(
\mathcal{E}^\bullet \otimes_\mathcal{O} \mathcal{F}^\bullet)
$$
을 가지는 어떤 호모토피 $\theta$에 의해 $1$과 호모토픽이라고
가정할 수 있다. 앞과 같은 보조정리를 사용했다.
$R = \Gamma(\mathcal{C}, \mathcal{O})$로 놓자. 보조정리
\ref{sites-cohomology-lemma-category-summands-finite-free}에 의해 다음을 얻는다.
\begin{enumerate}
\item $M^\bullet = \Gamma(X, \mathcal{E}^\bullet)$은
유한 사영 $R$-가군들의 유계 복합체이다.
\item $N^\bullet = \Gamma(X, \mathcal{F}^\bullet)$은
유한 사영 $R$-가군들의 유계 복합체이다.
\item $\alpha$와 $\beta$는 각각 복합체의 사상
$a : R \to \text{Tot}(M^\bullet \otimes_R N^\bullet)$와
$b : \text{Tot}(M^\bullet \otimes_R N^\bullet) \to R$에 대응한다.
\item $\theta^n$은 사상
$h^n : \text{Tot}^n(M^\bullet \otimes_R N^\bullet) \to
\text{Tot}^{n - 1}(M^\bullet \otimes_R N^\bullet)$에 대응한다.
\item $b \circ a = 1$이고 $b \circ a - 1 = dh + hd$이다.
\end{enumerate}
따라서 $M^\bullet$과 $N^\bullet$은 $D(R)$의 서로 역인 대상들을
정의한다. More on Algebra, 보조정리
\ref{more-algebra-lemma-invertible-derived}에 의해 직접곱 분해
$R = \prod_{a \leq n \leq b} R_n$과 가역 $R_n$-가군 $H^n$들이
존재하여 $M^\bullet \cong \bigoplus_{a \leq n \leq b} H^n[-n]$이다.
$D(R)$에서의 이 동형사상은 복합체의 사상
$$
\bigoplus H^n[-n] \longrightarrow M^\bullet
$$
으로 올릴 수 있다. 각 $H^n$이 사영 $R$-가군이기 때문이다.
이에 대응하여 보조정리 \ref{sites-cohomology-lemma-category-summands-finite-free}를
다시 사용하면 복합체의 사상
$$
\bigoplus H^n \otimes_R \mathcal{O}[-n] \to \mathcal{E}^\bullet
$$
을 얻으며, 이는 $D(\mathcal{O})$에서 동형사상이다. 여기서
$M \otimes_R \mathcal{O}$는 보조정리
\ref{sites-cohomology-lemma-category-summands-finite-free}의 증명에서 구성한
유한 사영 $R$-가군에서 $\mathcal{O}$-가군으로 가는 함자를 뜻한다.
$\mathcal{O}_n = R_n \otimes_R \mathcal{O}$로 놓으면 (2)가 성립한다.

\medskip\noindent
$(\mathcal{C}, \mathcal{O})$가 국소적으로 환 달린 사이트라고 하자.
대상 $U$와 유한 직접곱 분해
$\mathcal{O}|_U = \prod_{a \leq n \leq b} \mathcal{O}_n|_U$가 주어지면,
각 $i$마다 $\mathcal{O}_n|_{U_i}$가 영이 아니게 하는 $n$이 많아야
하나뿐인 덮개 $\{U_i \to U\}$를 찾을 수 있다. 이는 Modules on Sites,
보조정리 \ref{sites-modules-lemma-locally-ringed}의 (2)와 Modules on Sites,
정의 \ref{sites-modules-definition-locally-ringed}에 주어진 국소적으로
환 달린 사이트의 정의에서 곧바로 따른다. 이로부터 함의
(2) $\Rightarrow$ (3)을 쉽게 얻는다. 함의 (3) $\Rightarrow$ (2)는
환 달린 사이트에 아무 가정이 없어도 성립한다. 세부 사항은 생략한다.
\end{proof}










\section{사영 공식}
\label{sites-cohomology-section-projection-formula}

\noindent
$f : (\Sh(\mathcal{C}), \mathcal{O}_\mathcal{C}) \to
(\Sh(\mathcal{D}), \mathcal{O}_\mathcal{D})$를 환 달린 토포스의
사상이라 하자. $E \in D(\mathcal{O}_\mathcal{C})$와
$K \in D(\mathcal{O}_\mathcal{D})$라 하자.
아무런 추가 가정 없이도 사상
\begin{equation}
\label{sites-cohomology-equation-projection-formula-map}
Rf_*E \otimes^\mathbf{L}_{\mathcal{O}_\mathcal{D}} K
\longrightarrow
Rf_*(E \otimes^\mathbf{L}_{\mathcal{O}_\mathcal{C}} Lf^*K)
\end{equation}
이 있다. 실제로 이는 표준 사상
$$
Lf^*(Rf_*E \otimes^\mathbf{L}_{\mathcal{O}_\mathcal{D}} K) =
Lf^*Rf_*E \otimes^\mathbf{L}_{\mathcal{O}_\mathcal{C}} Lf^*K
\longrightarrow
E \otimes^\mathbf{L}_{\mathcal{O}_\mathcal{C}} Lf^*K
$$
의 수반 사상이다. 이 표준 사상은 $Lf^*Rf_*E \to E$와 보조정리
\ref{sites-cohomology-lemma-pullback-tensor-product}, \ref{sites-cohomology-lemma-adjoint}에서 온다.
사영 공식의 상당히 일반적인 판은 다음과 같다.

\begin{lemma}
\label{sites-cohomology-lemma-projection-formula}
$f : (\Sh(\mathcal{C}), \mathcal{O}_\mathcal{C}) \to
(\Sh(\mathcal{D}), \mathcal{O}_\mathcal{D})$를 환 달린 토포스의
사상이라 하자. $E \in D(\mathcal{O}_\mathcal{C})$와
$K \in D(\mathcal{O}_\mathcal{D})$라 하자. $K$가 완전하면
$$
Rf_*E \otimes^\mathbf{L}_{\mathcal{O}_\mathcal{D}} K =
Rf_*(E \otimes^\mathbf{L}_{\mathcal{O}_\mathcal{C}} Lf^*K)
$$
이 $D(\mathcal{O}_\mathcal{D})$에서 성립한다.
\end{lemma}

\begin{proof}
(\ref{sites-cohomology-equation-projection-formula-map})이 동형사상임을 확인할 때
$\mathcal{D}$ 위에서 국소적으로 작업해도 된다. 즉,
$\mathcal{D}$의 임의의 대상 $V$에 대해 사상을 $V_j$로 제한하면
동형사상이 되게 하는 덮개 $\{V_j \to V\}$를 찾아야 한다.
완전 대상의 정의에 의해 이는 $K$가 $\mathcal{O}_\mathcal{D}$-가군들의
강한 완전 복합체로 표시된다고 가정해도 됨을 뜻한다.
완전히 일반적으로 명제가 $K = K_1 \oplus K_2$에 대해 참일
필요충분조건은 $K_1$과 $K_2$ 각각에 대해 참인 것이다.
따라서 $K$가 유한 자유 $\mathcal{O}_\mathcal{D}$-가군들의
유한 복합체라고 가정할 수 있다. 이 경우 단순 절단을 사용하는
간단한 논증으로 $K$가 유한 자유 $\mathcal{O}_\mathcal{D}$-가군으로
표시되는 경우로 귀착된다. 명제는 $K$ 변수의 유한 직합인자 아래
불변이므로, $K = \mathcal{O}_\mathcal{D}[n]$인 경우만 증명하면
충분하며 이 경우는 자명하다.
\end{proof}

\begin{remark}
\label{sites-cohomology-remark-compatible-with-diagram}
사상 (\ref{sites-cohomology-equation-projection-formula-map})은 다음 의미에서
주석 \ref{sites-cohomology-remark-base-change}의 밑변환 사상과 호환된다.
즉,
$$
\xymatrix{
(\Sh(\mathcal{C}'), \mathcal{O}_{\mathcal{C}'})
\ar[r]_{g'} \ar[d]_{f'} &
(\Sh(\mathcal{C}), \mathcal{O}_\mathcal{C}) \ar[d]^f \\
(\Sh(\mathcal{D}'), \mathcal{O}_{\mathcal{D}'})
\ar[r]^g &
(\Sh(\mathcal{D}), \mathcal{O}_\mathcal{D})
}
$$
가 환 달린 토포스의 가환 도식이라고 하자.
$E \in D(\mathcal{O}_\mathcal{C})$와
$K \in D(\mathcal{O}_\mathcal{D})$라 하자. 그러면 도식
$$
\xymatrix{
Lg^*(Rf_*E \otimes^\mathbf{L}_{\mathcal{O}_\mathcal{D}} K)
\ar[r]_p \ar[d]_t &
Lg^*Rf_*(E \otimes^\mathbf{L}_{\mathcal{O}_\mathcal{C}} Lf^*K)
\ar[d]_b \\
Lg^*Rf_*E \otimes^\mathbf{L}_{\mathcal{O}_{\mathcal{D}'}}
Lg^*K \ar[d]_b &
Rf'_*L(g')^*(E \otimes^\mathbf{L}_{\mathcal{O}_\mathcal{C}}
Lf^*K) \ar[d]_t \\
Rf'_*L(g')^*E \otimes^\mathbf{L}_{\mathcal{O}_{\mathcal{D}'}}
Lg^*K \ar[rd]_p &
Rf'_*(L(g')^*E \otimes^\mathbf{L}_{\mathcal{O}_{\mathcal{D}'}}
L(g')^*Lf^*K) \ar[d]_c \\
& Rf'_*(L(g')^*E \otimes^\mathbf{L}_{\mathcal{O}_{\mathcal{D}'}}
L(f')^*Lg^*K)
}
$$
은 가환한다. 여기서 $t$로 표시된 화살표는 보조정리
\ref{sites-cohomology-lemma-pullback-tensor-product}를 적용하여 얻고,
$b$로 표시된 화살표는 주석 \ref{sites-cohomology-remark-base-change}를 적용하여 얻으며,
$p$로 표시된 화살표는 (\ref{sites-cohomology-equation-projection-formula-map})을
적용하여 얻는다. 또한 $c$는
$L(g')^* \circ Lf^* = L(f')^* \circ Lg^*$에서 온다.
확인은 생략한다.
\end{remark}






\section{약하게 축약 가능한 대상}
\label{sites-cohomology-section-w-contractible}

\noindent
사이트의 대상 $U$가 {\it 약하게 축약 가능하다}고 함은 집합의 층의 모든 전사
$\mathcal{F} \to \mathcal{G}$가 전사
$\mathcal{F}(U) \to \mathcal{G}(U)$를 유도하는 것을 말한다. Sites, Definition
\ref{sites-definition-w-contractible}을 참조하라.

\begin{lemma}
\label{sites-cohomology-lemma-w-contractible}
$\mathcal{C}$를 사이트라 하고 $U$를 $\mathcal{C}$의 약하게 축약 가능한
대상이라 하자. 그러면
\begin{enumerate}
\item 함자 $\mathcal{F} \mapsto \mathcal{F}(U)$는 완전 함자
$\textit{Ab}(\mathcal{C}) \to \textit{Ab}$이고,
\item $H^p(U, \mathcal{F}) = 0$
는 모든 아벨 층 $\mathcal{F}$와 모든 $p \geq 1$에 대해 성립하며,
\item 임의의 군의 층 $\mathcal{G}$에 대해 모든 $\mathcal{G}$-토서는
$U$ 위의 단면을 갖는다.
\end{enumerate}
\end{lemma}

\begin{proof}
첫째 명제는 정의에서 곧바로 따른다
(Homology, Section \ref{homology-section-functors}도 참조하라).
Derived Categories, Lemma \ref{derived-lemma-right-derived-exact-functor}에
의해 고차 유도 함자들은 영이다. $\mathcal{F}$를 $\mathcal{G}$-토서라 하자.
그러면 $\mathcal{F} \to *$는 층의 전사이다. 따라서 (3)도 정의에서 따른다.
\end{proof}

\noindent
약하게 축약 가능한 대상이 충분히 많을 때의 몇 가지 귀결을 여기에
나열해 두면 편리하다.

\begin{proposition}
\label{sites-cohomology-proposition-enough-weakly-contractibles}
$\mathcal{C}$를 사이트라 하자. 모든 $U \in \mathcal{B}$가 약하게 축약
가능하고 $\mathcal{C}$의 모든 대상이 $\mathcal{B}$의 원소들로 이루어진
피복을 갖도록 $\mathcal{B} \subset \Ob(\mathcal{C})$를 택하자.
$\mathcal{O}$를 $\mathcal{C}$ 위의 환의 층이라 하자. 그러면
\begin{enumerate}
\item $\mathcal{O}$-가군의 복합체
$\mathcal{F}_1 \to \mathcal{F}_2 \to \mathcal{F}_3$가 완전일 필요충분조건은
모든 $U \in \mathcal{B}$에 대해
$\mathcal{F}_1(U) \to \mathcal{F}_2(U) \to \mathcal{F}_3(U)$가
완전인 것이다.
\item $D(\mathcal{O})$의 모든 대상 $K$는 그 표준 절단들의 유도 극한이다:
$K = R\lim \tau_{\geq -n} K$.
\item 전이 사상이 전사인 역계
$\ldots \to \mathcal{F}_3 \to \mathcal{F}_2 \to \mathcal{F}_1$
가 주어지면 사영 $\lim \mathcal{F}_n \to \mathcal{F}_1$은 전사이다.
\item $\textit{Mod}(\mathcal{O})$에서 직접곱은 완전하다.
\item $D(\mathcal{O})$의 직접곱은 임의의 대표 복합체들의 직접곱을
취하여 계산할 수 있다.
\item $(\mathcal{F}_n)$이 $\mathcal{O}$-가군의 역계이면 모든 $p > 1$에 대해
$R^p\lim \mathcal{F}_n = 0$이고
$$
R^1\lim \mathcal{F}_n  =
\Coker(\prod \mathcal{F}_n \to \prod \mathcal{F}_n)
$$
이다. 여기서 사상은 $(x_n) \mapsto (x_n - f(x_{n + 1}))$이다.
\item $(K_n)$이 $D(\mathcal{O})$의 대상들의 역계이면 짧은 완전열
$$
0 \to R^1\lim H^{p - 1}(K_n) \to H^p(R\lim K_n) \to
\lim H^p(K_n) \to 0
$$
이 존재한다.
\end{enumerate}
\end{proposition}

\begin{proof}
(1)의 증명. 열이 완전하면 Lemma \ref{sites-cohomology-lemma-w-contractible}에 의해
$\mathcal{C}$의 임의의 약하게 축약 가능한 원소에서 값을 취하여 얻은 열도
완전하다. 역으로, 모든 $U \in \mathcal{B}$에 대해
$\mathcal{F}_1(U) \to \mathcal{F}_2(U) \to \mathcal{F}_3(U)$
가 완전하다고 가정하자. $V$를 $\mathcal{C}$의 대상이라 하고
$s \in \mathcal{F}_2(V)$를 $\mathcal{F}_2 \to \mathcal{F}_3$의 핵에
속하는 원소라 하자. 가정에 의해 $U_i \in \mathcal{B}$인 피복
$\{U_i \to V\}$가 존재한다. 그러면 $s|_{U_i}$는 단면
$s_i \in \mathcal{F}_1(U_i)$로 올라간다. 따라서 $s$는 상층
$\Im(\mathcal{F}_1 \to \mathcal{F}_2)$의 단면이다. 다시 말해 열
$\mathcal{F}_1 \to \mathcal{F}_2 \to \mathcal{F}_3$
은 완전하다.

\medskip\noindent
(2)의 증명. 이는 $d = 0$인 Lemma \ref{sites-cohomology-lemma-is-limit-dimension}에서 따른다.

\medskip\noindent
(3)의 증명. $(\mathcal{F}_n)$을 (2)에서와 같은 계라 하고
$\mathcal{F} = \lim \mathcal{F}_n$으로 놓자. $U \in \mathcal{B}$이면
$\mathcal{F}(U) = \lim \mathcal{F}_n(U)$
은 $\mathcal{F}_1(U)$로 전사이다. 모든 전이 사상
$\mathcal{F}_{n + 1}(U) \to \mathcal{F}_n(U)$이 전사이기 때문이다.
따라서 Sites, Definition \ref{sites-definition-sheaves-injective-surjective}과
모든 대상이 $\mathcal{B}$의 원소들로 이루어진 피복을 갖는다는 가정에 의해
$\mathcal{F} \to \mathcal{F}_1$은 전사이다.

\medskip\noindent
(4)의 증명. $\mathcal{O}$-가군의 완전열의 족
$\mathcal{F}_{i, 1} \to \mathcal{F}_{i, 2} \to \mathcal{F}_{i, 3}$
을 택하자. 다음 열이 완전함을 보이고자 한다.
$\prod \mathcal{F}_{i, 1} \to \prod \mathcal{F}_{i, 2} \to
\prod \mathcal{F}_{i, 3}$. (1)의 판정법을 사용한다.
$U \in \mathcal{B}$라 하자. 그러면
$$
(\prod \mathcal{F}_{i, 1})(U) \to
(\prod \mathcal{F}_{i, 2})(U) \to
(\prod \mathcal{F}_{i, 3})(U)
$$
는
$$
\prod \mathcal{F}_{i, 1}(U) \to
\prod \mathcal{F}_{i, 2}(U) \to
\prod \mathcal{F}_{i, 3}(U)
$$
와 같다. 각각의 열
$\mathcal{F}_{i, 1}(U) \to \mathcal{F}_{i, 2}(U) \to \mathcal{F}_{i, 3}(U)$
은 (1)에 의해 완전하다. 따라서 Homology, Lemma
\ref{homology-lemma-product-abelian-groups-exact}에 의해 위에 표시한 열들은
완전하다. 다시 (1)을 적용하여 결론을 얻는다.

\medskip\noindent
(5)의 증명. (4)와 (조금 일반화한) Derived Categories, Lemma
\ref{derived-lemma-products}에서 따른다.

\medskip\noindent
(6)과 (7)의 증명. 유도 극한과 호모토피 극한 및 그 관계에 관한 논의는
Section \ref{sites-cohomology-section-derived-limits}을 참조하라. Derived Categories,
Definition \ref{derived-definition-derived-limit}에 의해 구별 삼각형
$$
R\lim K_n \to \prod K_n \to \prod K_n \to R\lim K_n[1]
$$
을 얻는다. 코호몰로지 층들의 긴 완전열을 취하면
$$
H^{p - 1}(\prod K_n) \to H^{p - 1}(\prod K_n) \to
H^p(R\lim K_n) \to H^p(\prod K_n) \to H^p(\prod K_n)
$$
를 얻는다. (4)에 의해 직접곱이 완전하므로 이는
$$
\prod H^{p - 1}(K_n) \to \prod H^{p - 1}(K_n) \to
H^p(R\lim K_n) \to \prod H^p(K_n) \to \prod H^p(K_n)
$$
가 된다. 먼저 $(\mathcal{F}_n)$이 (6)에서와 같은 경우의
$K_n = \mathcal{F}_n[0]$에 이를 적용하면 (6)을 얻는다. 다음으로
(7)에서와 같은 $(K_n)$에 적용하면 (7)을 얻는다.
\end{proof}






\section{콤팩트 대상}
\label{sites-cohomology-section-compact}

\noindent
이 절에서는 환 달린 사이트 위의 가군의 유도범주에서 콤팩트 대상을
연구한다. 콤팩트 대상은 Derived Categories, Definition
\ref{derived-definition-compact-object}에서 정의됨을 상기하자.

\begin{lemma}
\label{sites-cohomology-lemma-compact-in-terms-of-generators}
$\mathcal{A}$를 Grothendieck 아벨 범주라 하자.
$S \subset \Ob(\mathcal{A})$를 다음을 만족하는 대상들의 집합이라 하자.
\begin{enumerate}
\item $\mathcal{A}$의 임의의 대상은 $S$의 원소들의 직합의 몫이고,
\item 임의의 $E \in S$에 대해 함자 $\Hom_\mathcal{A}(E, -)$는
직합과 가환한다.
\end{enumerate}
그러면 $D(\mathcal{A})$의 모든 콤팩트 대상은 $S$의 원소들의 유한 직합들로
이루어진 유한 복합체의 $D(\mathcal{A})$ 안의 직합인자이다.
\end{lemma}

\begin{proof}
$K \in D(\mathcal{A})$가 콤팩트 대상이라고 가정하자. $K$를 복합체
$K^\bullet$로 나타내고 사상
$$
K^\bullet
\longrightarrow
\bigoplus\nolimits_{n \geq 0} \tau_{\geq n} K^\bullet
$$
을 생각하자. 여기서는 표준 절단을 사용했다. Homology, Section
\ref{homology-section-truncations}을 참조하라. 각 차수에서 오른쪽의 직합은
유한하므로 이 사상은 잘 정의된다. 가정에 의해 이 사상은 유한 직합을
경유한다. 따라서 적어도 하나의 $n$에 대해 $K \to \tau_{\geq n} K$는
영이다. 즉, $K$는 $D^{-}(R)$에 속한다.

\medskip\noindent
$K$를 각 항이 $S$의 대상들의 직합인 위로 유계인 복합체 $K^\bullet$로
나타낼 수 있다. Derived Categories, Lemma
\ref{derived-lemma-subcategory-left-resolution}을 참조하라. 다음에 유의하자.
$$
K^\bullet = \bigcup\nolimits_{n \leq 0} \sigma_{\geq n}K^\bullet
$$
여기서는 단순 절단을 사용했다. Homology, Section
\ref{homology-section-truncations}을 참조하라. 따라서 Derived Categories,
Lemmas \ref{derived-lemma-colim-hocolim} 및
\ref{derived-lemma-commutes-with-countable-sums}에 의해
$1 : K^\bullet \to K^\bullet$은 $D(R)$에서
$\sigma_{\geq n}K^\bullet \to K^\bullet$을 경유한다. 그러므로
$1 : K^\bullet \to K^\bullet$은 다음과 같이 분해된다.
$$
K^\bullet \xrightarrow{\varphi} L^\bullet \xrightarrow{\psi} K^\bullet
$$
여기서 이 분해는 $D(\mathcal{A})$ 안의 것이고, $L^\bullet$은 유계이며
각 항이 $S$의 원소들의 직합인 복합체이다. $i \not \in [a, b]$이면
$L^i$가 영이라고 하자. $i < c$일 때마다 $L^i$가 $S$의 원소들의
유한 직합이 되며 값이 $\leq b + 1$인 가장 큰 정수를 $c$라 하자.
주장: $c < b + 1$이면 $c$가 증가하도록 $L^\bullet$을 바꿀 수 있다.
이 주장과 귀납법에 의해 $1_K$의 분해
$$
K \xrightarrow{\varphi} L \xrightarrow{\psi} K
$$
를 얻는다. 여기서 이는 $D(\mathcal{A})$ 안의 분해이고, $L$은 $S$의
원소들의 유한 직합들로 이루어진 유한 복합체로 나타낼 수 있다.
$e = \varphi \circ \psi \in \text{End}_{D(\mathcal{A})}(L)$가 멱등원임에
유의하자. Derived Categories, Lemma
\ref{derived-lemma-projectors-have-images-triangulated}에 의해
$L = \Ker(e) \oplus \Ker(1 - e)$이다. 사상 $\varphi : K \to L$은
$D(R)$에서 $\Ker(1 - e)$와의 동형사상을 유도하므로 결론을 얻는다.

\medskip\noindent
주장의 증명. $L^c = \bigoplus_{\lambda \in \Lambda} E_\lambda$로 쓰자.
$L^{c - 1}$은 $S$의 원소들의 유한 직합이므로 가정 (2)에 의해
$L^{c - 1} \to L^c$가
$\bigoplus_{\lambda \in \Lambda'} E_\lambda \subset L^c$를 경유하게 하는
유한 부분집합 $\Lambda' \subset \Lambda$를 찾을 수 있다. 복합체의 사상
$$
\pi :
L^\bullet
\longrightarrow
(\bigoplus\nolimits_{\lambda \in \Lambda \setminus \Lambda'} E_\lambda)[-i]
$$
을 차수 $i$에서 $\Lambda \setminus \Lambda'$에 대응하는 인자들로의
사영으로 정의하자. $K$에 관한 가정에 의해, 필요하면 $\Lambda'$를 더 큰
유한 부분집합으로 바꾼 뒤 $D(\mathcal{A})$에서
$\pi \circ \varphi = 0$이라고 가정할 수 있다.
$(L')^\bullet \subset L^\bullet$을 $\pi$의 핵이라 하자. $\pi$는 전사이므로
복합체의 짧은 완전열을 얻고, 이는 $D(\mathcal{A})$에서 구별 삼각형을 준다
(Derived Categories, Lemma
\ref{derived-lemma-derived-canonical-delta-functor} 참조).
$\Hom_{D(\mathcal{A})}(K, -)$는 호몰로지적이고(Derived Categories,
Lemma \ref{derived-lemma-representable-homological} 참조)
$\pi \circ \varphi = 0$이므로, $(L')^\bullet \to L^\bullet$과 합성하면
$\varphi$가 되는 $D(\mathcal{A})$의 사상
$\varphi' : K^\bullet \to (L')^\bullet$을 찾을 수 있다. $\psi'$를 $\psi$와
$(L')^\bullet \to L^\bullet$의 합성으로 놓으면 새로운 분해를 얻는다.
$(L')^\bullet$은 차수 $c$를 제외하면 $L^\bullet$과 일치하고
$(L')^c = \bigoplus_{\lambda \in \Lambda'} E_\lambda$이므로 주장이 증명된다.
\end{proof}

\begin{lemma}
\label{sites-cohomology-lemma-compact-objects-if-enough-qc}
$(\mathcal{C}, \mathcal{O})$를 환 달린 사이트라 하자. $\mathcal{C}$의
모든 대상이 준콤팩트 대상들로 이루어진 피복을 갖는다고 가정하자. 그러면
$D(\mathcal{O})$의 모든 콤팩트 대상은, 각 항이 $j_!\mathcal{O}_U$ 꼴의
$\mathcal{O}$-가군들의 유한 직합인 유한 복합체의 $D(\mathcal{O})$ 안의
직합인자이다. 여기서 $U$는 $\mathcal{C}$의 준콤팩트 대상이다.
\end{lemma}

\begin{proof}
Lemma \ref{sites-cohomology-lemma-compact-in-terms-of-generators}를 적용한다. 여기서
$S \subset \Ob(\textit{Mod}(\mathcal{O}))$는 준콤팩트한
$U \in \Ob(\mathcal{C})$에 대한 $j_!\mathcal{O}_U$ 꼴의 가군들의
집합이다. Modules on Sites, Lemma
\ref{sites-modules-lemma-module-quotient-flat}과 모든 $U$가 준콤팩트
대상들로 피복된다는 가정에 의해 가정 (1)이 성립한다. 가정 (2)는
$$
\Hom_\mathcal{O}(j_!\mathcal{O}_U, \mathcal{F}) = \mathcal{F}(U)
$$
이고 이것이 Sites, Lemma \ref{sites-lemma-directed-colimits-sections}에 의해
직합과 가환한다는 사실에서 따른다.
\end{proof}

\noindent
위 보조정리의 상황에서도 가군 $j_!\mathcal{O}_U$가 항상
$D(\mathcal{O})$의 콤팩트 대상인 것은 아니다($U$가 $\mathcal{C}$의
준콤팩트 대상인 경우에도 그렇다). 다음 두 보조정리는 이 문제를 다룬다.


\begin{lemma}
\label{sites-cohomology-lemma-when-jshriek-lower-compact}
$(\mathcal{C}, \mathcal{O})$를 환 달린 사이트라 하고 $U$를
$\mathcal{C}$의 대상이라 하자. 함자
$\mathcal{F} \mapsto H^p(U, \mathcal{F})$가 직합과 가환한다고 가정하자.
그러면 $\mathcal{O}$-가군 $j_!\mathcal{O}_U$는 다음 의미에서
$D^+(\mathcal{O})$의 콤팩트 대상이다. 즉, $D(\mathcal{O})$에서
$M = \bigoplus_{i \in I} M_i$가 아래로 유계이면
$\Hom(j_{U!}\mathcal{O}_U, M) =
\bigoplus_{i \in I} \Hom(j_{U!}\mathcal{O}_U, M_i)$이다.
\end{lemma}

\begin{proof}
Modules on Sites, Equation
(\ref{sites-modules-equation-map-lower-shriek-OU-into-module})에 의해
$\Hom(j_{U!}\mathcal{O}_U, -)$는 함자
$\mathcal{F} \mapsto \mathcal{F}(U)$와 같으므로
$H^p(U, M) = \bigoplus H^p(U, M_i)$임을 보이면 충분하다.
$\mathcal{I}_i$, $i \in I$를 단사 $\mathcal{O}$-가군들의 족이라 하자.
가정에 의해
$$
H^p(U, \bigoplus\nolimits_{i \in I} \mathcal{I}_i) =
\bigoplus\nolimits_{i \in I} H^p(U, \mathcal{I}_i) = 0
$$
가 모든 $p$에 대해 성립한다. $M = \bigoplus M_i$는 아래로 유계이므로
$n < a$일 때 $H^n(M_i) = 0$이 되는 $a \in \mathbf{Z}$가 존재한다.
따라서 $M_i$를 나타내고 $n < a$일 때 $\mathcal{I}_i^n = 0$인 단사
$\mathcal{O}$-가군의 복합체 $\mathcal{I}_i^\bullet$을 택할 수 있다.
Derived Categories, Lemma \ref{derived-lemma-injective-resolutions-exist}을
참조하라. Injectives, Lemma \ref{injectives-lemma-derived-products}에 의해
직합 복합체 $\bigoplus \mathcal{I}_i^\bullet$은 $M$을 나타낸다.
Leray 비순환성(Derived Categories, Lemma
\ref{derived-lemma-leray-acyclicity})에 의해
$$
R\Gamma(U, M) = \Gamma(U, \bigoplus \mathcal{I}_i^\bullet) =
\bigoplus \Gamma(U, \bigoplus \mathcal{I}_i^\bullet) =
\bigoplus R\Gamma(U, M_i)
$$
를 얻는다. 이것이 원하는 바이다.
\end{proof}

\begin{lemma}
\label{sites-cohomology-lemma-when-jshriek-lower-compact-worked-out}
$(\mathcal{C}, \mathcal{O})$를 피복들의 집합 $\text{Cov}_\mathcal{C}$가
주어진 환 달린 사이트라 하자. $\mathcal{B} \subset \Ob(\mathcal{C})$와
$\text{Cov} \subset \text{Cov}_\mathcal{C}$를 부분집합들이라 하자.
다음을 가정한다.
\begin{enumerate}
\item 모든 $\mathcal{U} \in \text{Cov}$에 대해
$\mathcal{U} = \{U_i \to U\}_{i \in I}$이고, $I$는 유한하며,
$U, U_i \in \mathcal{B}$이고, 모든
$U_{i_0} \times_U \ldots \times_U U_{i_p} \in \mathcal{B}$이다.
\item 모든 $U \in \mathcal{B}$에 대해 $\text{Cov}$에 나타나는 $U$의
피복들은 $U$의 피복들의 공종계를 이룬다.
\end{enumerate}
그러면 $U \in \mathcal{B}$에 대해 대상 $j_{U!}\mathcal{O}_U$는 다음
의미에서 $D^+(\mathcal{O})$의 콤팩트 대상이다. 즉, $D(\mathcal{O})$에서
$M = \bigoplus_{i \in I} M_i$가 아래로 유계이면
$\Hom(j_{U!}\mathcal{O}_U, M) =
\bigoplus_{i \in I} \Hom(j_{U!}\mathcal{O}_U, M_i)$이다.
\end{lemma}

\begin{proof}
이는 Lemma \ref{sites-cohomology-lemma-when-jshriek-lower-compact}과
Lemma \ref{sites-cohomology-lemma-colim-works-over-collection}에서 따른다.
\end{proof}

\begin{lemma}
\label{sites-cohomology-lemma-when-jshriek-compact}
$(\mathcal{C}, \mathcal{O})$를 환 달린 사이트라 하고 $U$를
$\mathcal{C}$의 대상이라 하자. 다음을 만족하는 정수 $d$가 존재하면
$\mathcal{O}$-가군 $j_!\mathcal{O}_U$는 $D(\mathcal{O})$의 콤팩트 대상이다.
\begin{enumerate}
\item 모든 $p > d$에 대해 $H^p(U, \mathcal{F}) = 0$이고,
\item 함자 $\mathcal{F} \mapsto H^p(U, \mathcal{F})$는 직합과 가환한다.
\end{enumerate}
\end{lemma}

\begin{proof}
(1)과 (2)를 가정하자. $D(\mathcal{O})$의 $K$에 대해
$\Hom(j_!\mathcal{O}_U, K) = R\Gamma(U, K)$임을 상기하자. 따라서
$R\Gamma(U, -)$가 직합과 가환함을 보여야 한다. 첫째 가정은 함자
$F = H^0(U, -)$가 유한 코호몰로지 차원을 갖는다는 뜻이다. 또한 둘째
가정에 의해 단사 가군들의 임의의 직합은 $F$-비순환이다. $K_i$를
$D(\mathcal{O})$의 대상들의 족이라 하자. $K_i$를 나타내고 각 항이
단사인 K-단사 대표 $I_i^\bullet$을 택하자. Injectives, Theorem
\ref{injectives-theorem-K-injective-embedding-grothendieck}을 참조하라.
비순환 대상들로 이루어진 임의의 복합체에 $F$를 적용하여 $RF$를
계산할 수 있고(Derived Categories, Lemma
\ref{derived-lemma-unbounded-right-derived}), $\bigoplus K_i$는
$\bigoplus I_i^\bullet$로 나타나므로(Injectives, Lemma
\ref{injectives-lemma-derived-products}),
$R\Gamma(U, \bigoplus K_i)$는 $\bigoplus H^0(U, I_i^\bullet)$로 나타난다.
따라서 원하는 대로 $R\Gamma(U, -)$는 직합과 가환한다.
\end{proof}

\begin{lemma}
\label{sites-cohomology-lemma-quasi-compact-weakly-contractible-compact}
$(\mathcal{C}, \mathcal{O})$를 환 달린 사이트라 하자. $U$를
$\mathcal{C}$의 준콤팩트이고 약하게 축약 가능한 대상이라 하자. 그러면
$j_!\mathcal{O}_U$는 $D(\mathcal{O})$의 콤팩트 대상이다.
\end{lemma}

\begin{proof}
Lemmas \ref{sites-cohomology-lemma-when-jshriek-compact} 및
\ref{sites-cohomology-lemma-w-contractible}과 Modules on Sites, Lemma
\ref{sites-modules-lemma-sections-over-quasi-compact}을 결합하면 된다.
\end{proof}

\begin{lemma}
\label{sites-cohomology-lemma-perfect-is-compact}
$(\mathcal{C}, \mathcal{O})$를 환 달린 사이트라 하자. $\mathcal{C}$가
다음 성질들을 갖는다고 가정하자.
\begin{enumerate}
\item $\mathcal{C}$는 준콤팩트 끝 대상 $X$를 갖는다.
\item $\mathcal{C}$의 모든 준콤팩트 대상은 준콤팩트 대상들로 이루어진
유한 피복들의 공종계를 갖는다.
\item $U$와 $U_i$가 준콤팩트인 유한 피복
$\{U_i \to U\}_{i \in I}$에 대해 섬유곱 $U_i \times_U U_j$는
준콤팩트이다.
\end{enumerate}
$K$를 $D(\mathcal{O})$의 완전 대상이라 하자. 그러면
\begin{enumerate}
\item[(a)] $K$는 다음 의미에서 $D^+(\mathcal{O})$의 콤팩트 대상이다.
$M = \bigoplus_{i \in I} M_i$가 아래로 유계이면
$\Hom(K, M) = \bigoplus_{i \in I} \Hom(K, M_i)$이다.
\item[(b)] $(\mathcal{C}, \mathcal{O})$가 유한 코호몰로지 차원을 갖는다고
하자. 즉, 임의의 $\mathcal{O}$-가군 $\mathcal{F}$에 대해 $i > d$이면
$H^i(X, \mathcal{F}) = 0$이 되게 하는 $d$가 존재한다고 하자. 그러면
$K$는 $D(\mathcal{O})$의 콤팩트 대상이다.
\end{enumerate}
\end{lemma}

\begin{proof}
$K^\vee$를 $K$의 쌍대라 하자. Lemma
\ref{sites-cohomology-lemma-dual-perfect-complex}을 참조하라. 그러면
$$
\Hom_{D(\mathcal{O})}(K, M) =
H^0(X, K^\vee \otimes_\mathcal{O}^\mathbf{L} M)
$$
이고, 이는 $D(\mathcal{O})$의 $M$에 대해 함자적이다.
$K^\vee \otimes_\mathcal{O}^\mathbf{L} -$는 직합과 가환하므로
$R\Gamma(X, -)$가 해당 직합들과 가환함을 보이면 충분하다.

\medskip\noindent
(a)의 증명. 위와 같이 다시 쓰면 이는 $U = X$인 Lemma
\ref{sites-cohomology-lemma-when-jshriek-lower-compact-worked-out}의 특수한 경우이다.

\medskip\noindent
(b)의 증명. $R\Gamma(X, K) = R\Hom(\mathcal{O}, K)$이고 Lemma
\ref{sites-cohomology-lemma-colim-works-over-collection}에 의해 $H^p(X, -)$가 직합과
가환하므로 이는 Lemma \ref{sites-cohomology-lemma-when-jshriek-compact}의 특수한 경우이다.
\end{proof}






\section{코호몰로지 층이 국소 상수인 복합체}
\label{sites-cohomology-section-locally-constant}

\noindent
국소 상수층은 Modules on Sites, Section
\ref{sites-modules-section-locally-constant}에서 도입하였다.
$\mathcal{C}$를 사이트, $\Lambda$를 환이라 하자.
$D(\mathcal{C}, \Lambda)$로 $\mathcal{C}$ 위의 $\underline{\Lambda}$-가군들의
아벨 범주의 유도 범주를 나타낸다.

\begin{lemma}
\label{sites-cohomology-lemma-locally-constant}
$\mathcal{C}$를 끝 대상 $X$를 갖는 사이트라 하고,
$\Lambda$를 뇌터환이라 하자.
$K \in D^b(\mathcal{C}, \Lambda)$의 각 $H^i(K)$가
유한형 $\Lambda$-가군의 국소 상수층이라고 하자. 그러면 각
$K|_{U_i}$를 유한형 $\Lambda$-가군의 국소 상수층들로 이루어진
복합체로 나타낼 수 있는 덮개 $\{U_i \to X\}$가 존재한다.
\end{lemma}

\begin{proof}
$i \not \in [a, b]$이면 $H^i(K) = 0$이 되도록 $a \leq b$를 택하자.
$b - a$에 관한 귀납법으로 다음을 보이겠다. 덮개 $\{U_i \to X\}$가
존재하여 $K|_{U_i}$를 복합체 $\underline{M^\bullet}_{U_i}$로 나타낼 수
있고, 여기서 $M^p$는 유한형 $\Lambda$-가군이며
$p \not \in [a, b]$이면 $M^p = 0$이다. $b = a$인 경우는 명백하다.
일반적으로 $X$를 한 덮개의 원소들로 바꾸어 $H^b(K)$가 상수라고,
이를테면 $H^b(K) = \underline{M}$이라고 가정해도 된다. Modules on Sites, Lemma
\ref{sites-modules-lemma-locally-constant-finite-type}
에 의해 가군 $M$은 유한 $\Lambda$-가군이다. $M$의 생성원
$x_1, \ldots, x_r$가 주는 전사 $\Lambda^{\oplus r} \to M$을 택하자.

\medskip\noindent
Lemma \ref{sites-cohomology-lemma-kill-cohomology-class-on-covering}를 조금 일반화하면
(세부사항은 생략한다), $x_i \in H^0(X, H^b(K))$가 $H^b(U_i, K)$의
원소로 올라가게 하는 덮개 $\{U_i \to X\}$가 존재한다. 따라서
$X$를 $U_i$들로 바꾸면 사상
$\underline{\Lambda^{\oplus r}}[-b] \to K$
이 차수 $b$의 코호몰로지 층 위에서 전사를 유도하는 상황에 이른다.
특수 삼각형
$$
\underline{\Lambda^{\oplus r}}[-b] \to K \to L \to
\underline{\Lambda^{\oplus r}}[-b + 1]
$$
을 택하자. 이제 $L$의 코호몰로지 층은 구간 $[a, b - 1]$에서만
0이 아닐 수 있고, 구간 $[a, b - 2]$에서는 $K$의 코호몰로지 층과
일치하며, 차수 $b - 1$에서는 짧은 완전열
$$
0 \to H^{b - 1}(K) \to H^{b - 1}(L) \to
\underline{\Ker(\Lambda^{\oplus r} \to M)} \to 0
$$
이 있다. Modules on Sites, Lemma
\ref{sites-modules-lemma-kernel-finite-locally-constant}
에 의해 $H^{b - 1}(L)$은 유한형 국소 상수층이다. 귀납 가정으로부터
$D(\mathcal{C}, \underline{\Lambda})$ 안의 동형
$\underline{M^\bullet} \to L$을 얻는다. 여기서 $M^p$는 유한
$\Lambda$-가군이고 $p \not \in [a, b - 1]$이면 $M^p = 0$이다.
사상 $L \to \Lambda^{\oplus r}[-b + 1]$은 국소적으로 상수인 사상
$\underline{M^{b - 1}} \to \underline{\Lambda^{\oplus r}}$을 준다
(Modules on Sites, Lemma
\ref{sites-modules-lemma-morphism-locally-constant}).
따라서 이는 사상 $M^{b - 1} \to \Lambda^{\oplus r}$로 주어진다고
가정해도 된다. 특수 삼각형은 합성
$M^{b - 2} \to M^{b - 1} \to \Lambda^{\oplus r}$이 0임을 보이며,
삼각 범주의 공리들은 동형
$$
\underline{M^a \to \ldots \to M^{b - 1} \to \Lambda^{\oplus r}}
\longrightarrow K
$$
을 $D(\mathcal{C}, \Lambda)$ 안에서 준다.
\end{proof}

\noindent
$\mathcal{C}$를 사이트, $\Lambda$를 환이라 하자. 사상
$\Sh(\mathcal{C}) \to \Sh(pt)$를 사용하면 함자
$D(\Lambda) \to D(\mathcal{C}, \Lambda)$, $K \mapsto \underline{K}$가 있음을 알 수 있다.

\begin{lemma}
\label{sites-cohomology-lemma-map-out-of-locally-constant}
$\mathcal{C}$를 끝 대상 $X$를 갖는 사이트라 하고, $\Lambda$를 환이라 하자.
다음 데이터를 택하자.
\begin{enumerate}
\item $D(\Lambda)$의 완전 대상 $K$,
\item $K$를 나타내는 유한 사영 $\Lambda$-가군들의 유한 복합체 $K^\bullet$,
\item $\Lambda$-가군층들의 복합체 $\mathcal{L}^\bullet$, 그리고
\item $D(\mathcal{C}, \Lambda)$ 안의 사상
$\varphi : \underline{K} \to \mathcal{L}^\bullet$.
\end{enumerate}
그러면 덮개 $\{U_i \to X\}$와 $\varphi|_{U_i}$를 나타내는 복합체 사상
$\alpha_i : \underline{K}^\bullet|_{U_i} \to \mathcal{L}^\bullet|_{U_i}$가 존재한다.
\end{lemma}

\begin{proof}
Lemma \ref{sites-cohomology-lemma-local-actual}에서 곧바로 따른다.
\end{proof}

\begin{lemma}
\label{sites-cohomology-lemma-locally-constant-map}
$\mathcal{C}$를 끝 대상 $X$를 갖는 사이트라 하고, $\Lambda$를 환이라 하자.
$K, L$을 $D(\Lambda)$의 대상들이라 하되 $K$는 완전하다고 하자.
$\varphi : \underline{K} \to \underline{L}$를 $D(\mathcal{C}, \Lambda)$ 안의
사상이라 하자. 그러면 덮개 $\{U_i \to X\}$가 존재하여
$\varphi|_{U_i}$는 $D(\Lambda)$ 안의 어떤 사상 $\alpha_i : K \to L$에
대하여 $\underline{\alpha_i}$와 같다.
\end{lemma}

\begin{proof}
Lemma \ref{sites-cohomology-lemma-map-out-of-locally-constant}와 Modules on Sites, Lemma
\ref{sites-modules-lemma-morphism-locally-constant}에서 따른다.
\end{proof}

\begin{lemma}
\label{sites-cohomology-lemma-locally-constant-tensor-product}
$\mathcal{C}$를 사이트, $\Lambda$를 뇌터환이라 하자.
$K, L \in D^-(\mathcal{C}, \Lambda)$라 하자. $K$와 $L$의 코호몰로지 층들이
유한형 $\Lambda$-가군의 국소 상수층이면,
$K \otimes_\Lambda^\mathbf{L} L$
의 코호몰로지 층들도 유한형 $\Lambda$-가군의 국소 상수층이다.
\end{lemma}

\begin{proof}
이를 Lemma \ref{sites-cohomology-lemma-locally-constant}의 응용으로 증명하겠다.
$H^i(K \otimes_\Lambda^\mathbf{L} L)$은
$H^i(\tau_{\geq i - 1}K \otimes_\Lambda^\mathbf{L} \tau_{\geq i - 1}L)$과
같음에 유의하자. 따라서 $K$와 $L$은 유계라고 가정해도 된다.
Lemma \ref{sites-cohomology-lemma-locally-constant}에 의해 $K$와 $L$은 유한형
$\Lambda$-가군의 국소 상수층들로 이루어진 복합체로 나타난다고 가정해도 된다.
그런 다음 이 복합체들을 유한 자유 $\Lambda$-가군들의 위로 유계인 복합체로
바꿀 수 있다. 이 경우 결과는 명백하다.
\end{proof}

\begin{lemma}
\label{sites-cohomology-lemma-locally-constant-bounded}
$\mathcal{C}$를 사이트, $\Lambda$를 뇌터환이라 하자.
$I \subset \Lambda$를 아이디얼이라 하고, $K \in D^-(\mathcal{C}, \Lambda)$라 하자.
$K \otimes_\Lambda^\mathbf{L} \underline{\Lambda/I}$의 코호몰로지 층들이
유한형 $\Lambda/I$-가군의 국소 상수층이면,
$K \otimes_\Lambda^\mathbf{L} \underline{\Lambda/I^n}$
의 코호몰로지 층들은 모든 $n \geq 1$에 대해 유한형
$\Lambda/I^n$-가군의 국소 상수층이다.
\end{lemma}

\begin{proof}
유한형 $\Lambda$-가군의 국소 상수층들은 모든 $\underline{\Lambda}$-가군의
범주의 약한 세르 부분범주를 이룸을 상기하자. Modules on Sites, Lemma
\ref{sites-modules-lemma-kernel-finite-locally-constant}.
따라서 코호몰로지 층들이 유한형 $\Lambda$-가군의 국소 상수층인 복합체들로
이루어진 $D(\mathcal{C}, \Lambda)$의 부분범주는
$D(\mathcal{C}, \Lambda)$의 엄밀히 충만하고 포화된 삼각부분범주이다.
Derived Categories, Lemma \ref{derived-lemma-cohomology-in-serre-subcategory}를 보라.
이제 특수 삼각형들
$$
K \otimes_\Lambda^\mathbf{L} \underline{I^n/I^{n + 1}} \to
K \otimes_\Lambda^\mathbf{L} \underline{\Lambda/I^{n + 1}} \to
K \otimes_\Lambda^\mathbf{L} \underline{\Lambda/I^n} \to
K \otimes_\Lambda^\mathbf{L} \underline{I^n/I^{n + 1}}[1]
$$
과 동형들
$$
K \otimes_\Lambda^\mathbf{L} \underline{I^n/I^{n + 1}}
=
\left(K \otimes_\Lambda^\mathbf{L} \underline{\Lambda/I}\right)
\otimes_{\Lambda/I}^\mathbf{L} \underline{I^n/I^{n + 1}}
$$
을 생각하자. 이를 Lemma \ref{sites-cohomology-lemma-locally-constant-tensor-product}와 결합하면
결론을 얻는다.
\end{proof}
